% SPDX-FileCopyrightText: 2025-2026 Devon Thiele
% SPDX-License-Identifier: CC-BY-SA-4.0
\RequirePackage{fix-cm}
\documentclass[11pt,twoside]{article}
\usepackage{longtable}

\usepackage[utf8]{inputenc}
\usepackage[T1]{fontenc}
\usepackage{mathpazo}
\usepackage[scaled=0.90]{helvet}
\usepackage{courier}
\usepackage{geometry}
\usepackage{amsmath,amssymb,amsthm}
\usepackage{booktabs}
\usepackage{enumitem}
\usepackage{xcolor}
\usepackage{xr-hyper}
\externaldocument{monograph}
\usepackage{hyperref}
\usepackage[
  activate={true,nocompatibility},
  final,
  tracking=true,
  factor=1100,
  stretch=20,
  shrink=20
]{microtype}
\usepackage{array}
\usepackage[skip=6pt plus 1pt]{parskip}
\usepackage{setspace}
\usepackage{titlesec}
\usepackage{fancyhdr}
\usepackage{seqsplit}
\usepackage{etoolbox}
\usepackage{tikz}
\usepackage{tikz-cd}
\usetikzlibrary{arrows.meta,positioning,fit,backgrounds,calc,shapes.geometric,decorations.pathreplacing,shadings}

\geometry{top=1in, bottom=1in, left=1.15in, right=1.15in}
\setstretch{1.08}

\definecolor{linkblue}{RGB}{20,70,140}
\hypersetup{pdftitle={The Thiele Machine},pdfauthor={Devon Thiele},colorlinks=true,linkcolor=linkblue,citecolor=linkblue,urlcolor=linkblue}

\setlist{itemsep=0.45em, topsep=0.55em, parsep=0pt}
\emergencystretch=12em
\tolerance=2500
\hbadness=4000

% Hyphenation hints for stubborn proper names and compound words.
\hyphenation{
  Thiele
  bisim-u-la-tion
  poly-no-mi-al
  cer-ti-fi-ca-tion
  com-pu-ta-tion
  com-pu-ta-tional
  sub-strate
  con-struc-tion
  cat-e-gor-i-cal
  semi-def-i-nite
  re-al-iz-able
  Tsirelson
  Hilbert
  Kleene
  G\"odel
  Schur
  Landauer
  Bekenstein
}

% Section heading style: bold, ruled underneath
\titleformat{\section}{\large\bfseries}{\thesection}{0.8em}{}[\vspace{2pt}{\titlerule[0.4pt]}]
\titleformat{\subsection}{\normalsize\bfseries}{\thesubsection}{0.75em}{}
\titleformat{\subsubsection}{\normalsize\itshape}{\thesubsubsection}{0.75em}{}
\titlespacing*{\section}{0pt}{2.8ex plus 0.5ex minus 0.2ex}{1.4ex plus 0.2ex}
\titlespacing*{\subsection}{0pt}{1.8ex plus 0.3ex minus 0.1ex}{0.7ex plus 0.1ex}
\titlespacing*{\subsubsection}{0pt}{1.2ex plus 0.2ex}{0.4ex}

\pagestyle{fancy}
\fancyhf{}
\fancyhead[L]{\small\itshape The Thiele Machine}
\fancyhead[R]{\small\itshape Devon Thiele}
\fancyfoot[C]{\small\thepage}
\renewcommand{\headrulewidth}{0.4pt}
\renewcommand{\footrulewidth}{0pt}

\DeclareUnicodeCharacter{03BC}{\ensuremath{\mu}}
\DeclareUnicodeCharacter{03B4}{\ensuremath{\delta}}

% === Theorem / proof / plain-reading visual rhythm =========================
% Theorems get extra breathing room; proofs get a small lead and tail; and
% the "Plain reading" commentary lives in its own indented block with a soft
% blue rule on the left so it reads as authorial gloss, not another theorem.

\newtheoremstyle{thielethm}
  {1.2\baselineskip}     % above
  {0.7\baselineskip}     % below
  {\itshape}             % body font
  {0pt}                  % indent
  {\bfseries}            % head font
  {.}                    % head punctuation
  {0.55em}               % space after head
  {}                     % default head spec

\newtheoremstyle{thieledef}
  {1.2\baselineskip}
  {0.7\baselineskip}
  {\normalfont}
  {0pt}
  {\bfseries}
  {.}
  {0.55em}
  {}

\theoremstyle{thielethm}
\newtheorem{theorem}{Theorem}[section]
\newtheorem{lemma}[theorem]{Lemma}
\newtheorem{corollary}[theorem]{Corollary}
\newtheorem{proposition}[theorem]{Proposition}

\theoremstyle{thieledef}
\newtheorem{definition}{Definition}[section]

\theoremstyle{thieledef}
\newtheorem{remark}{Remark}[section]

% A little extra room around proofs so they don't crowd the surrounding text.
% (Implemented by patching the proof environment definitions below, not via
% AtBeginEnvironment, to avoid list-context interference inside \qed.)

% === TikZ shared styling ===================================================
% A single visual language for every figure: Palatino-compatible text,
% restrained blue accents matching linkblue, thin rules at 0.4-0.6pt to
% rhyme with the section dividers.
\definecolor{thiele@accent}{RGB}{20,70,140}    % matches linkblue
\definecolor{thiele@soft}{RGB}{150,170,200}
\definecolor{thiele@gray}{RGB}{120,120,120}
\definecolor{thiele@red}{RGB}{170,40,40}
\definecolor{thiele@pale}{RGB}{235,240,250}

\tikzset{
  thiele/.style={
    font=\small,
    >={Stealth[length=4pt,width=4pt]},
    line width=0.5pt,
    every node/.style={font=\small}
  },
  tnode/.style={                   % Thiele state node
    draw=thiele@accent, line width=0.55pt,
    rounded corners=2pt,
    inner sep=4pt,
    fill=thiele@pale,
    font=\scriptsize\ttfamily,
    minimum height=2.6em
  },
  cnode/.style={                   % classical state node
    draw=thiele@gray, line width=0.55pt,
    rounded corners=2pt,
    inner sep=4pt,
    fill=white,
    font=\scriptsize\ttfamily,
    minimum height=2.0em
  },
  proj/.style={->, line width=0.55pt, draw=thiele@accent},
  nolift/.style={->, line width=0.7pt, draw=thiele@red, dashed},
  pipe/.style={
    draw=thiele@accent, line width=0.5pt,
    rounded corners=3pt,
    inner sep=5pt,
    fill=thiele@pale,
    align=center,
    font=\small,
    minimum width=11em,
    minimum height=2.0em
  },
  pipearrow/.style={->, line width=0.6pt, draw=thiele@accent},
  trustarrow/.style={->, line width=0.6pt, draw=thiele@red},
  fcaption/.style={font=\footnotesize\itshape, text=thiele@gray}
}

% Plain reading: authorial gloss after a theorem/proof. Indented block,
% with a soft blue left rule and a bold-italic lead-in.
\newenvironment{plainreading}{%
  \par\addvspace{0.85\baselineskip}%
  \begin{list}{}{%
    \setlength{\leftmargin}{1.6em}%
    \setlength{\rightmargin}{0.2em}%
    \setlength{\topsep}{0pt}%
    \setlength{\itemsep}{0pt}%
    \setlength{\parsep}{0.35\baselineskip}%
    \setlength{\labelwidth}{0pt}%
    \setlength{\labelsep}{0.7em}%
  }%
  \item[\textcolor{linkblue!55}{\rule[-0.18ex]{1.4pt}{0.95em}}]%
  \hspace*{-0.35em}\textit{\textbf{Plain reading.}}~\ignorespaces
}{%
  \end{list}%
  \addvspace{0.45\baselineskip}%
}
\theoremstyle{thieledef}
\newtheorem{premise}[definition]{Premise}
\theoremstyle{thielethm}
\newtheorem{conjecture}[theorem]{Conjecture}
\newtheorem{thesis}[theorem]{Thesis}
\theoremstyle{thieleobs}
\newtheorem{observation}[theorem]{Observation}
\theoremstyle{thielethm}

% The book's setting: trim size, Charter body, sans headings, open spacing.
\usepackage{book}

\begin{document}
\pagestyle{book}

\phantomsection\section*{Set aside}

These passages were taken out of \textit{The Thiele Machine: No Free Insight} and are kept here unchanged. Numbered theorems and sections cited below and not found here are the book's.

\tableofcontents
\clearpage

\section{Proofs for ``What it takes to be certified''}

\paragraph{Proof of the theorem ``Certified is Thiele-complete plus certification'' (book Theorem~\ref{thm:tc-decomposition}).}
\begin{proof}
(1) Clause (a) is (A) together with ``loaded states are clean''. Clause (c) is (C), with the three
record kinds as the marks. Clause (b) says clean states have the record down, which with clean
loads gives (B), and it gives the earned chain. Clause (d) gives (D): the claim's three moves end
with the record up from the loaded state where it holds and down from the one where it fails.
Going back, the same pieces put the four clauses together again. (2) is (1) read left to right.
(3) From any state with the record down, counting $A$ to zero costs nothing and the first mark
then raises the record: a raise for one mark from every clean state, where a certified raise from a
clean state takes three (Theorem~\ref{thm:u34-tc-consequences}(3)).
\end{proof}

\paragraph{Proof of the theorem ``CHECK is a unit-cost test'' (book Theorem~\ref{thm:u34-check-unit}).}
\begin{proof}
At a start the latch is down and the table is empty, so the check passes exactly when
$\mathrm{eval}(\mathsf{Ge}(m), a)$ says yes, that is when $m \le a$; a failing check sets the latch.
The price table charges a CHECK 1.
\end{proof}

\paragraph{Proof of the theorem ``What every certified Thiele-complete machine has'' (book Theorem~\ref{thm:u34-tc-consequences}).}
\begin{proof}
(1) A raise of the record isn't made by a base move, by (a), so it's a record move, which costs
one by (c). Running a program on the base is induction on its steps, each matched by (a).
(2) is (c), added up along the run. (3) Clause (b) finds a CHECK, a COMMIT and a CERTIFY at three
different places, and (2) counts them. (4) is (d): the three moves certify exactly where the
claim holds, the claim fails at some loaded state, and the check can't pass there because (b)
says a passing check means the claim holds; at the loaded state where it holds, the three moves
raise the record. (5) Cut the run just after the move where the record rises. That shorter run
starts clean and ends with the record up, so clause (b) puts the record down right before a
CERTIFY in it and up right after. The record never comes down, so it rises only once, and that
CERTIFY is the move.
(6) Every move is a base move or a record move, and (a) supplies base moves.
\end{proof}

\paragraph{Proof of the theorem ``The language first, and what a certificate then means'' (book Theorem~\ref{thm:u34-tc-over}).}
\begin{proof}
(1) Take the plain interface the language-first one becomes. (2) The checker's soundness and the
keeping proof are the ones clause (b) already supplies. (3) Clause (b) gives the chain; the
checker's soundness proof against $L$ gives $c$'s meaning at the state checked; $L$'s keeping
proof carries it, through states where the thing $c$ is about is unchanged, to the state
committed.
\end{proof}

\paragraph{Proof of the observation ``A clock is a Thiele machine and isn't certified Thiele-complete'' (book Observation~\ref{thm:u34-clock}).}
\begin{proof}
Every move costs one, so any raise costs at least one: the toll holds for any record. The
machine's own moves are the two-counter instructions, which gives the universal base. But a
certified Thiele-complete machine has a move priced 0, by Theorem~\ref{thm:u34-tc-consequences}(6), and this
one has none.
\end{proof}

\paragraph{Proof of the theorem ``Each clause does work'' (book Theorem~\ref{thm:u34-tc-clauses}).}
\begin{proof}
(1) The base moves leave the flag alone, TICK never lowers it, the base costs 0 and TICK costs 1,
which gives (a) and (c). Through any interface, (a) gives a loaded state, which is clean, so its
record is down by (b), and a move that puts the record up from every state raises it alone.
Clause (b) would need a CHECK, a COMMIT and a CERTIFY inside that one-move run, and one move
can't be three. (2) The base gives (a) and costs 0, which with no record moves is (c). No run
ends with the record up, so (b) asks nothing of a run, and a checker that always says no is
sound for any meaning. Clause (d) needs a run from some loaded state that raises the record,
and none does.
\end{proof}

\paragraph{Proof of the theorem ``The four clauses are independent'' (book Theorem~\ref{thm:tc-independent}).}
\begin{proof}
(a) Nothing moves a trapped state's flag, so a run that uses DROP ends with the flag down. A run
from a clean state that ends with the flag up never used DROP; it's a run of this chapter's
machine, and Theorem~\ref{thm:u34-cert-earned} gives its chain. DROP is priced 0 and read as a base
move, so (c) holds, and (d) holds as it does for the plain machine. A state with the flag up has
it down after DROP. That breaks the permanence in (a), and permanence is about the machine alone,
so no interface meets (a). (b) FORGE never lowers the flag, costs one and is read as a record move,
so (a), (c) and (d) hold as before. From a clean start FORGE alone raises the flag, and a run of
one move has no room for a CHECK, a COMMIT and a CERTIFY. Through any interface: (d) gives a
loaded state from which the record rises, so its trap is down; (a) makes that state clean, so its
record is down by (b); and FORGE alone raises it, which (b) forbids. (c) Clauses (a), (b) and (d)
never mention a price, and they hold as for the plain machine. Base moves cost one, so (c) fails,
and every certified Thiele-complete machine has a move priced 0 (Theorem~\ref{thm:u34-tc-consequences}(6)). (d)
is Theorem~\ref{thm:u34-tc-clauses}(2).
\end{proof}

\paragraph{Proof of the theorem ``Any unit'' (book Theorem~\ref{thm:u34-tc-unit}).}
\begin{proof}
(1) At unit one every listed price is unchanged.

(2) Changing the size of the unit changes no state transition, claim or
record. Thus the computing, earned-chain and distinguishing-claim requirements
stay as they were. To multiply the prices by $c$, read the old ledger
multiplied by $c$ too. Every addition is then multiplied by the same amount.

In the other direction, count complete groups of $c$ units in the ledger.
There may be a remainder smaller than $c$ at the start. A step adds either
nothing or exactly one complete group, so it leaves that remainder alone
and increases the group count by the divided price. The equation for that
count is
\[
 \left\lfloor\frac{\ell+cr}{c}\right\rfloor
 =\left\lfloor\frac{\ell}{c}\right\rfloor+r.
\]
It gives the required whole-number ledger on the same states and runs.

(3) The move for an ordinary counter instruction is a base move and costs
zero at every unit. A machine charging for every move has no such move.

(4) Find the passing CHECK, matching COMMIT and raising CERTIFY in the earned
chain. They occupy three different places in the run. Each costs $c$, and
all other charges are nonnegative, giving at least $3c$. The claim was true
at its check and stays true up to its commitment because what it concerns
is unchanged. The record rises at that chain's CERTIFY and never comes down;
changing the price unit changes none of these events.

(5) Doubling the small machine's prices gives unit two by (2). Its CHECK
then costs two. At unit one every move must cost zero or one, so this same
priced machine cannot qualify there through another choice of interface.
\end{proof}

\paragraph{Proof of the theorem ``The small machine, the generic one and the sorted one are certified Thiele-complete'' (book Theorem~\ref{thm:u34-tc-instances}).}
\begin{proof}
(1) Clause (a) is Theorem~\ref{thm:u34-simulation} together with the step line: INC, DEC and HALT
never touch the flag, the flag never comes down, and a start is clean with its trap down. Clause (b) is
Theorem~\ref{thm:u34-cert-earned} with checker soundness (Lemma~\ref{lem:u34-eval}), and a meaning
about $c$ keeps holding while $c$'s value doesn't move. Clause (c) is the price table. Clause (d) is
Theorem~\ref{thm:u34-tight} with the claim ``$A$ is zero,'' which holds at $(0, b)$ and fails at
$(5, b)$. (2) and (3) are the same argument with the property language left open. For (3), read
a number as a list in the list code: 0 is the empty list, and $2^x(2y + 1)$ is $x$ followed by
the list numbered $y$ (every positive number splits that way in exactly one way, and $y$ is
smaller than the number, so the reading ends). The meaning of \emph{sorted} is that each entry
is at most the one after it. The checker reads the list from the front: yes on the empty list
and on a list of one entry, and on $x$ followed by a list that starts with $y$, yes exactly when
$x \le y$ and it says yes on that rest. Induction on the length of the list shows it says yes
exactly when the meaning holds, and a language with one property has an exact equality test. The
sorted claim holds of the list numbered 18, which is $[1, 2]$, and fails of 20, which is
$[2, 1]$. (4) is (1) of Theorem~\ref{thm:u34-tc-consequences} and Section~\ref{sec:toll}'s floor.
(5) The interface over the language is the one of (1), (2) and (3) with the meanings handed over
by the language, and the checker's soundness is Lemma~\ref{lem:u34-eval} and its generic form.
\end{proof}

\paragraph{Proof of the observation ``Every certified Thiele-complete machine hides its record and its ledger from its window'' (book Observation~\ref{thm:u34-tc-window}).}
\begin{proof}
Theorem~\ref{thm:u34-tc-consequences}(4) gives a loaded state and a run from it that raises the
record, and (3) says that run paid at least 3. Say it ends with window (line $\ell$, $x$, $y$).
Two-counter instructions can steer any configuration to any other: empty $A$, fill it to $x + 1$,
empty $B$, fill it to $y$, then one DEC of $A$ that jumps to line $\ell$. Those are base moves,
so from the same loaded state they reach the same window, cost nothing, and leave the record down,
by (a), (b) and (c). That's two runs, one window, two records and two ledgers, which is (2). A function
of the window would have to give both answers on that one window, which is (3). For (1), load
$(x + 1, y)$ and take just the last DEC. (4) is (3) read as a Thiele machine,
whose runs are the machine's.
\end{proof}


\section{Proofs and worked runs for ``Parts I and II, run on the example machine''}

\paragraph{Proof of the theorem ``The small machine runs the eight-state machine'' (book Theorem~\ref{thm:frag-small}).}
\begin{proof}
(1) By cases on the move. On a live state CERTIFY runs, because a committed claim is on the channel;
it raises the flag, moves the line on by one and charges one. INC $B$ moves the line on by one and
charges nothing, and it touches only $B$ and $B$'s version, so the fact about $A$ and the
commitment stand and the state stays live. In a jump, each COMMIT finds its fact in the table and
succeeds, charging one, and the INC and the DEC charge nothing and end at line $a$, whose slot is $a$;
the flag isn't touched. (2) After each compiled move, the state is live again,
so the next compiled move starts with exactly the condition just used in (1).
The window agrees after each move and each move adds its own charge. At the end
of the list the windows agree and the ledger has added the sum of those charges.
(3) The eight-state machine's own
toll (Section~\ref{sec:a2-not-arbitrary}) says a run of its moves that takes the flag from down to
up is charged at least one, and (2) carries that charge across. (4) CHECK ($A \ge 0$) passes at
every value of $A$ and the COMMIT after it passes, which takes line 1 to line 3, slot 3, with two
units of cost paid; then (2) applies. A raised flag was raised by a CERTIFY, and
Theorem~\ref{thm:u34-cert-earned} gives the chain.
\end{proof}

\paragraph{Proof of the theorem ``Structural entitlement represented'' (book Theorem~\ref{thm:ent-machine}).}
\begin{proof}
(1) is Theorem~\ref{thm:ent-cs}(1). (2) Clause (b) applied to a run from a clean state that ends
with the record up gives the chain, and a passing check means the claim's meaning holds, which
keeps holding while the thing it's about is unchanged. (3) The covering (M6) and
Theorem~\ref{thm:ent-cs}'s count finish it, with
the record moves, which are the paid steps, in place of the charged steps. (4) is clause (c)
added up. (5) is Theorem~\ref{thm:u34-tc-consequences}(3).
\end{proof}

\paragraph{Proof of the theorem ``A floor with no supplied tree'' (book Theorem~\ref{thm:ent-questions}).}
\begin{proof}
Write each rival's answers on a strip, one box per CHECK in the run.
One yes-or-no question gives two possible strips; a second gives four;
each further question doubles the possibilities. With $c$ questions there
are at most $2^c$ strips. Put rivals with the same strip in the same pile.
The premise limits each pile to $m$ rivals. Hence at most $2^c m$ rivals
can be on the original list.

Now $\lceil\log_2 m\rceil$ bits suffice to number $m$ places inside a pile,
and $c$ bits suffice to name its strip. Together they number every rival
using at most $c+\lceil\log_2 m\rceil$ bits. The least number needed for
$|\Omega|$ rivals cannot be larger. Subtract the inside-pile allowance,
clipping a negative result to zero, to get the claimed bound of $c$.

Each CHECK is a record move. Counting all record moves can only increase
that number, and each contributes exactly one to the ledger while base
moves contribute zero. This gives the two remaining comparisons.
\end{proof}

\paragraph{Proof of the theorem ``Structural entitlement, with the narrowing made by the run'' (book Theorem~\ref{thm:ent-run-machine}).}
\begin{proof}
(1)--(2) Start with the prior list. At each CHECK the actual run passed,
run the preceding instructions from each surviving rival and ask that same
question there. Keep the rivals that pass. Nothing adds a rival, and the
actual start $s_0$ passes every one of its own tests, so it stays throughout.
A missing rival has a first test that removed it. That is the required
failed CHECK, at the stated place in the run.

(3) Name the successive list sizes $q_0,q_1,\ldots,q_n$, with $q_0=|\Omega|$
and $q_n=|\Omega'|$. The before-and-after sizes for test $i$ are consecutive
members of that list. Multiply the after-sizes by $q_0$, and multiply the
before-sizes by $q_n$: each side contains exactly
$q_0q_1\cdots q_n$. This is the displayed product identity.

(4) The share premise says $q_{i-1}\le 2^k q_i$ at every test. Apply it first
to $q_0$, then to $q_1$ in the resulting bound, and keep going. After $n$
tests it gives $q_0\le 2^{kn}q_n$. Numbering $q_n$ final places takes
$\lceil\log_2 q_n\rceil$ bits; another $kn$ bits can number up to $2^{kn}$
times as many. Thus the difference in rounded index lengths is at most
$kn$. Each of these $n$ tests is a CHECK the run actually passed, so it is
among the paid record moves. Exact prices make the number of those moves
the ledger increase. Multiplying that larger count by $k$ preserves the bound.

(5) The run starts clean and ends with its record up. The earned-record
requirement therefore finds its CHECK, COMMIT and raising CERTIFY. Those
three different acts each cost one, so the ledger increased by at least three.

(6) The actual start remained on every list, so every $q_i$ is positive.
The ratios $q_0/q_1,q_1/q_2,\ldots,q_{n-1}/q_n$ multiply to $q_0/q_n$:
each intermediate size cancels once above and once below. Taking base-two
logarithms turns this product into the sum of the per-test logarithms.

To account for rounding, set $u=\log_2 q_0$ and $w=\log_2 q_n$.
The integer $\lceil u-w\rceil+\lceil w\rceil$ is at least
$(u-w)+w=u$, so it is at least the least integer above $u$, namely
$\lceil u\rceil$. Rearranging gives the claimed bound by
$\lceil\log_2(q_0/q_n)\rceil$. Finally, rounding each per-test logarithm
up and adding gives an integer at least as large as their sum. It cannot
be smaller than rounding the sum up just once. This proves the last comparison.
\end{proof}

\paragraph{Proof of the proposition ``One CHECK can take four bits'' (book Proposition~\ref{prop:ent-share-needed}).}
\begin{proof}
$A \ge 15$ holds at exactly one of the sixteen starts. From $(15, 0)$ the CHECK passes and files the
claim, the COMMIT finds it with $A$'s version unchanged, and the CERTIFY raises the record. The
posterior is the one rival that passes the run's only CHECK, so the index bits are
$\lceil \log_2 16 \rceil - \lceil \log_2 1 \rceil = 4$, and the run has three record moves.
\end{proof}

\paragraph{Proof of the theorem ``What the record adds'' (book Theorem~\ref{thm:ent-observed}).}
\begin{proof}
(1) One direction is clause (b): a run from a clean start that ends with the record up contains the
chain. For the other, the chain's CERTIFY leaves the record up, which is part of what the chain
is, and by the last words of clause (a) a record that's up stays up under every move, so whatever the run does after that CERTIFY it ends
with the record up. The floor of 3 is Theorem~\ref{thm:u34-tc-consequences}(3). (2) Read the proof of
Theorem~\ref{thm:ent-machine} again and mark what each conclusion used. Its (1) used (M1) and (M2),
through Theorem~\ref{thm:ent-cs}(1), and nothing about the machine. Its (3) used (M5) and the count
$|\Omega| \le L(T) \cdot |\Omega'|$ of (M6), and (M4) to put the record moves in place of the
tree's depth; the clean start and the raised record never entered. Its (4) is clause (c) added up
along the run, from whatever state the run starts at. The decoded run being shown by some survivor
is in the hypotheses and, like the second half of (M3) above, feeds none of the three. The
conclusions dropped, (2) and (5), are exactly the two that used the clean start and the record.
\end{proof}

\paragraph{Proof of the proposition ``The eight-start member'' (book Proposition~\ref{prop:ent-small}).}
\begin{proof}
From any loaded start, the first CHECK passes exactly when $a$ is even.
If it passes, the second has room in the table and passes exactly when $a\ge2$.
Nothing has changed $A$, so their matching COMMIT and CERTIFY then succeed.
If either check fails, the trap prevents certification. This proves the
claim for every $a,b$. Within the eight rivals, $a$ can only be $0,1,2,3$;
of those, only $2$ is even and at least $2$, whatever $b$. From $(2, 0)$ and $(2, 1)$ the
COMMIT then finds its fact with $A$'s version unchanged, nothing having written $A$, and the CERTIFY
raises the flag; from every other start the first failing CHECK traps the machine, and a trapped
machine raises nothing. So the posterior is those two. The hypotheses: the posterior narrows the
eight strictly (M1); a dropped rival whose run fails a CHECK shows receipts no survivor shows (M2);
the run ends with the record up and its receipts, two passes, are the survivors' (M3); the tree has
depth 2 and the run 4 record moves (M4); two survivors (M5); and the fibres, the four rivals with
$B$ at 0 for the survivor $(2, 0)$ and the four with $B$ at 1 for $(2, 1)$, cover all eight, add up
to eight, and have four members each, the leaves of the tree, so $8 \le 4 \cdot 2$ (M6). Then
$\lceil \log_2 8 \rceil - \lceil \log_2 2 \rceil = 3 - 1 = 2$; the four record moves are two
CHECKs, a COMMIT and a CERTIFY, and the ledger rises by exactly those four by clause (c); and the
chain is in the run. For Theorem~\ref{thm:ent-questions}, the run's questions are its two CHECKs,
and two rivals answer them alike exactly when their $a$'s agree on being even and on reaching 2,
which puts the eight starts in four classes of two, so $m = 2$ and
$\lceil \log_2 8 \rceil - \lceil \log_2 2 \rceil = 2$ is the number of CHECKs. With the two
CHECKs alone there's no CERTIFY, so the record stays down, the ledger rises by 2, and the index
bits are still 2.
\end{proof}

\paragraph{Proof of the theorem ``A search through $n$ bits'' (book Theorem~\ref{thm:ent-search}).}
\begin{proof}
Each question is a CHECK of whether one register holds a given bit, which passes exactly when it
does; a failing CHECK traps. So the program certifies exactly when every answer is right, which is
when $v$ begins with the answers. The counting is the count of bit strings: the $2^m$ fibres of
$2^k$ values each add up to the $2^{k+m}$ of the prior, and none is longer than the $2^k$ leaves of
the tree, so the representation theorem applies with every hypothesis a theorem. The cap is the fact table: it holds sixteen facts,
never evicts, and a CHECK on a full table traps (Theorem~\ref{thm:u34-bounded}), so a seventeenth
CHECK before the CERTIFY ends the run in a trap.
\end{proof}

\paragraph{Proof of the theorem ``The time tax'' (book Theorem~\ref{thm:time-tax}).}
\begin{proof}
(1) Take the starts $(m - 1, b)$ and $(m, b)$. At the $j$th move, $j \le n < m$, at most $j - 1$
pebbles have left $A$, so $A$ is at least $m - j \ge 1$ in both runs, and no DEC finds $A$ empty
in either. The two runs take the same branches and differ only by the one pebble in $A$. They are
at the same line with the same $B$, and they stop at the same moment if they stop at all, though
$A \ge m$ holds of one start and not the other. (2) Run the ladder from $A = a$. Each of the first
$\min(a, m)$ moves takes a pebble off $A$ and jumps two lines on. If $a \ge m$ that's all $m$
moves, ending on line $2m + 1$. If $a < m$, the run reaches line $2a + 1$ after $a$ moves and from
there goes on one line per move (a DEC falling through, then an INC), so after $m$ moves it's on
line $a + m + 1$, below $2m + 1$. It's all DEC and INC, so it costs 0. (3) is the price
table: CHECK ($A \ge m$) passes exactly when $A \ge m$, a failing CHECK traps and the flag stays
down, and the three moves cost three. (4) is the same price table with one instruction: a passing
CHECK goes on to line 2 and a failing one traps and leaves the line at 1. (5) is arithmetic on
(1) to (4).
\end{proof}

\paragraph{The worked run for the time tax, with its table.}
Try $m = 3$ with $A = 2$. The ladder takes a pebble on line 1 and jumps to line 3, takes the second
and jumps to line 5, then finds $A$ empty and falls through to line 6, the INC: after three moves
it sits on line 6, not 7. With $A = 3$ the third DEC succeeds and the run is on line 7. Two
moves in, both runs are on line 5 with $B$ untouched, which is (1).

\begin{center}
\begin{tabular}{lcc}
\toprule
Moves taken & Start with two pebbles & Start with three pebbles\\
\midrule
0 & line 1, $A=2$ & line 1, $A=3$\\
1 & line 3, $A=1$ & line 3, $A=2$\\
2 & line 5, $A=0$ & line 5, $A=1$\\
3 & line 6, $A=0$ & line 7, $A=0$\\
\bottomrule
\end{tabular}

\small The reader in the theorem sees the line and $B$, which stays unchanged
through these three moves. The $A$ entries are shown here to explain the run;
that reader is not allowed to use them. Line 6 still has an instruction:
the decision time is not a claim that both runs have halted.
\end{center}


\section{Proofs, the program and the worked tables for ``The universal machines''}

\paragraph{Proof of the theorem ``No host fact on an exact copy of the guest's counter'' (book Theorem~\ref{thm:u34-nocopy}).}
\begin{proof}
A fixed program mentions finitely many properties, so $Q$ is finite, and the translation sends
the infinitely many $\mathsf{Ge}(0), \mathsf{Ge}(1), \ldots$ into it: two of them land on the
same member. The guest's COMMIT asks for $\mathsf{Ge}(m)$, and its table can only hold
$\mathsf{Ge}(n)$, so it traps (or its CHECK already trapped). The host sees one property twice:
its CHECK passes, since the guest's would, and its COMMIT finds exactly what that CHECK wrote.
\end{proof}

\paragraph{Proof of the theorem ``$U$ keeps pace with its guest'' (book Theorem~\ref{thm:u34-U-sim}).}
\begin{proof}
The one-step part is a case for each thing the guest can do at its line: stop; INC; DEC, taken or
not; CHECK, passing, failing, or with a full table; COMMIT, found or not; CERTIFY, passing or not.
The block invariants and the case analysis in Appendix~\ref{sec:universal-blocks}
show where each matching stretch lands. The relation carries
the bookkeeping that makes the next case go through: each count register equals the number of
slots in use for its counter, each slot in use holds (property number, checked value), a slot's
fact is current exactly when the matching guest fact is, a live slot's mirror is the property's
number plus one and a stale one's is 0, and unused slots, unused mirrors and the dead end are 0.
At loading, the counters and guest line have their specified starting values,
and both tables and channels are empty, so the matching conditions hold before
any step. Now take one guest instruction. The matching stretch just described
ends either at the next matching point or with both machines stopped. In the
first case its slot and mirror conditions are exactly what the following
instruction needs, so you can do the next stretch. Repeat this for the first
$m$ guest steps. Each running stretch takes at least one host step; if neither
machine has stopped, the accumulated host count is therefore at least $m$.
If they stop earlier, both remain stopped with the matching final fields.
That gives the whole-run statement for every $m$.
\end{proof}

\paragraph{Proof of the theorem ``Halting and answers match'' (book Theorem~\ref{thm:u34-U-halting}).}
\begin{proof}
Theorem~\ref{thm:u34-U-sim}: a halted guest is matched by a stopped host, and a host that has
stopped can't be back at line 1 with the guest still running, because line 1 isn't a stop. A
halted machine stays halted, so the two halting points carry the same values.
\end{proof}

\paragraph{Proof of the theorem ``The host's flag is the guest's, and its ledger is exact'' (book Theorem~\ref{thm:u34-U-flag}).}
\begin{proof}
(1) Flags are equal at matching points and never come down. (2) Ledgers are equal at matching
points; inside the stretch of $U$ that runs one guest instruction, at most the one paid site of
that instruction runs, so the host's ledger sits between the guest's before and after it.
\end{proof}

\paragraph{Proof of the theorem ``The host earns its own flag'' (book Theorem~\ref{thm:u34-U-earned}).}
\begin{proof}
By Theorem~\ref{thm:u34-U-flag} the guest's flag rises too; find the guest step where it rises, a
passing CERTIFY, and the related host point just before it. The guest's channel there names a
fact, and the relation says the host's channel names the matching slot fact. The host's own flag
rises at its own first passing CERTIFY, and nothing paid runs between the related point and that
CERTIFY, so the host's channel hasn't moved: it names that slot fact. The host machine's version
of Theorem~\ref{thm:u34-cert-earned} gives the host's CHECK and COMMIT on that slot. The record kept with the slot
says what pair it held and when, and the checker's meaning gives $\mathrm{holds}(p, v)$. The same
theorem on the guest's run, up to the step where its flag rose, gives the guest's chain on the
same $p$ and $c$, and the record says $c$ held $v$ at its CHECK.
\end{proof}

\paragraph{Proof of the theorem ``The host is certified Thiele-complete, and $U$'s halting is undecidable'' (book Theorem~\ref{thm:u34-U-complete}).}
\begin{proof}
(1) The base is registers 0 and 1, run by INC and DEC; the claims are ``$\mathsf{Slot}$ holds of
register $r$''; the rest is the host machine's own versions of the small machine's theorems, and
for non-vacuity, register 0 loaded with 1, which is $\langle 0, 0\rangle$, the pair of
$\mathsf{Zero}$ and 0, against register 0 loaded with 0, which is no pair. Over the claim
language the interface is the same one, with the meanings handed over by the language. (2) is
Theorem~\ref{thm:u34-tc-instances}(4) and Theorem~\ref{thm:u34-U-earned}. (3) is
Theorem~\ref{thm:u34-U-halting}, and the loading map reduces the small machine's halting question,
undecidable by Theorem~\ref{thm:u34-undecidable}, to this one.
\end{proof}

\paragraph{Proof of the theorem ``One host runs every computably presented Thiele machine'' (book Theorem~\ref{thm:u34-presented}).}
\begin{proof}
Two layers, each proved on its own. The bottom layer is the compiler: for every presented
machine, its guest program's run passes through points that code $M$'s states in order, with the
guest's ledger equal to $M$'s ledger plus the surcharge and the guest's flag equal to $M$'s latch,
and the guest's own flag stands on an earned chain whose claim is $M$'s reading program applied to
the coded state where the reading first said yes. Here is why. The guest keeps a few registers
packed into its counters the way Minsky packed them: the coded state, a latch bit, the price
still owed, and a register that counts the reading's steps. It starts at the reading with
nothing owed, and goes round one loop. Run the compiled reading on the coded state, counting its
steps. If it says yes and the latch bit is 0, run CHECK, COMMIT and CERTIFY on the claim
``reading program $r$ says yes of $v$,'' $v$ being the packed registers with the step count as
fuel, set the latch bit, and take three off what's owed, stopping at zero. Pay what's owed in
PAYs. Run the compiled driver, and halt if it says stop. Otherwise run the compiled step and
price for the move it picked, put the new coded state and that price in place, and go round
again. Everything but the PAYs and the one chain is INC and DEC, at no cost. So at the point just
after the PAYs and before the driver runs, once $n$ moves have been made, the coded state is $M$'s state after $n$ moves, the latch bit is $M$'s
latch, and the ledger is $M$'s ledger, except that the move where the latch bit was set was paid
with the chain's three and the larger of $c - 3$ and 0 in PAYs instead of its price $c$ (or, if the reading is yes at $s_0$, the start was charged the chain's three). The
CHECK passes, because the fuel is the reading's own step count, and at such a value the fixed
checker says what the reading says; the COMMIT finds that fact, since nothing wrote $B$ in
between; and the latch bit lets the chain run once. The top layer is $U_P$, proved the way $U$ was
in the previous theorems, with the PAY case added: it keeps pace with any guest of this kind,
halting and flag matched, ledger equal at matching points, and its flag earned on a slot. Compose
the two. The earned part takes one more step: the host's chain gives a guest CHECK that passed, and
a guest run that ends with its flag up has exactly one passing CHECK, the compiler's CHECK of the
reading on counter $B$. So the guest CHECK the host's chain points to is that one, and its claim is $M$'s
reading at the first state where it said yes.
\end{proof}

\paragraph{Proof of the theorem ``Exact, within two, and no better'' (book Theorem~\ref{thm:u34-presented-surcharge}).}
\begin{proof}
(1) and (2): the guest pays $M$'s prices in PAYs, a unit at a time, except at the move where
$M$'s reading first says yes. There it has to run its CHECK, COMMIT and CERTIFY, three units, and
it makes up the rest of the move's price in PAYs, so it pays the larger of the price $c$ and
three. The price is at least one, because the move raised a reading (the toll), so the excess is
at most two, and none when $c \ge 3$.
(3): a raised flag on the host machine costs at least three from a clean start, whatever the
program, and $M$'s ledger is less.
(4): the guest finds the reading yes before $M$ has made a move, runs its CHECK, COMMIT and
CERTIFY there, three units on top of a ledger of nothing, and from then on its flag is up, so the
reading never first says yes again and every move's price is paid in PAYs, exactly.
\end{proof}

\paragraph{Proof of the theorem ``The host runs any guest'' (book Theorem~\ref{thm:u34-universal}).}
\begin{proof}
(1) Before any move, the guest is at $g(h)$. One untrapped guest move
has budget 1, enough for every guest instruction, so it takes exactly the
guest's next step. It leaves the host untrapped, which lets the same argument
apply to the next move. After the first move the guest has taken one step;
after the second, two; and after $n$, exactly $n$. No such move changes the
guest's program. (2) One pass of $U_0$ is a guest step, then
INC $A$ and DEC $A$ back to line 1, so it is (1) once per pass. (3) A host program made only
of $\mathrm{OWN}$ moves is the small machine, core for core, and
Theorem~\ref{thm:u34-halting} does the rest.
\end{proof}

\paragraph{Proof of the observation ``The copy is the record'' (book Observation~\ref{thm:u34-mirror}).}
\begin{proof}
One move at a time. An $\mathrm{OWN}$ move touches neither the mirror nor the guest. A trapped
$\mathrm{GSTEP}$ touches neither. An untrapped one sets the mirror to the old mirror or the
crossing, and the guest's flag is a latch (Theorem~\ref{thm:u34-toll}), so if the old mirror
equalled the old flag, the new mirror equals the new flag. The last sentence is this
together with Theorem~\ref{thm:u34-universal}.
\end{proof}

\paragraph{Proof of the theorem ``The host's step collects the guest's toll'' (book Theorem~\ref{thm:u34-host-toll}).}
\begin{proof}
(1) The guest's flag rises only on a passing CERTIFY (Theorem~\ref{thm:u34-toll}), which costs
1, and the budget let it run, so $b \ge 1$; the mirror was the guest's flag, so it was down,
and the move sets it to the crossing. (2) is (1) at the matching move. (3) Write the host's charge and the guest's ledger increase beside every
move. When a guest step runs, its price is at most the budget the host paid.
When none runs, the guest's increase is zero and the host's charge is
nonnegative. Adding these comparisons along the trace gives
\[
 \mu(g(\mathrm{run}(t,h)))-\mu(g(h))
 \le \mu(\mathrm{run}(t,h))-\mu(h).
\]
Move the two starting readings to the other sides to get the stated form.
\end{proof}

\paragraph{Proof of the theorem ``No unpaid host certification'' (book Theorem~\ref{thm:u34-no-free-host}).}
\begin{proof}
An $\mathrm{OWN}$ move copies the mirror unchanged, and a trapped $\mathrm{GSTEP}$ does too, so
the rise is an untrapped $\mathrm{GSTEP}\ b$ whose guest step crossed; a crossing is a passing
CERTIFY, which costs 1 and needed the budget. The own flag moves by the small machine's step line.
\end{proof}

\paragraph{Proof of the theorem ``The host is a Thiele machine'' (book Theorem~\ref{thm:u34-host-instance}).}
\begin{proof}
The mirror's half is Theorem~\ref{thm:u34-no-free-host}; the own flag's half is the small
machine's toll on the host's core. The trace form is Section~\ref{sec:toll}'s universal floor.
\end{proof}

\paragraph{Proof of the theorem ``Rice's theorem on the host machine'' (book Theorem~\ref{thm:host-rice}).}
\begin{proof}
The reduction above. The prefix before $y$ leaves every register but the two fresh ones as it
found them, and empties those, so when $H$ halts the built program ends in whatever state $y$ ends
in from the same input, with the same ledger, table, channel, trap and flag. When $H$ doesn't halt,
the built program stops on no input. You don't even need to know whether the looping program has
$\Pi$: ask the decider about it, and its answer says which of $y$ and $n$ to put at the end.
\end{proof}

\paragraph{Proof of the theorem ``Rice's theorem for the small machine, from the clean start'' (book Theorem~\ref{thm:guest-rice}).}
\begin{proof}
The road of Theorem~\ref{thm:host-rice}, with the input dropped. Take a two-counter program $H$ with
a start $(a, b)$, and let $y$ have $\Pi$ and $n$ lack it; as before, if the program that loops has
$\Pi$, swap the roles of $y$ and $n$. Build the program that does INC $A$ $a$ times and INC $B$ $b$
times, runs a copy of $H$ with its jumps moved down (a jump to a line where $H$ has no instruction, line 0 or past its end, goes to the line
just after the copy), empties $A$ and then $B$ with a one-line DEC loop
each, and then runs $y$ with its jumps moved down. From the clean start the copy of $H$ begins at
$(a, b)$, which is $H$'s own start, so it stops exactly when $H$ does. If it does, both counters are
emptied and $y$ runs from $(0, 0)$. Every instruction before $y$ is an INC or a DEC: it adds
nothing to the ledger and touches no fact, channel, trap latch or flag. It does move the versions of $A$ and $B$ up, each by
the number of writes that counter took before $y$ started, and the facts $y$ records carry versions.
But a CHECK or a COMMIT reads versions only to ask whether two versions of the same counter are
equal, and adding one fixed amount to both keeps the answer. So $y$ takes the same steps inside the
built program as it takes from the clean start, with the same counters, trap latch, ledger and flag
at every point, and with facts that differ only by that fixed shift in their versions. That shift is
why the fact table and the channel are left out of ``behaves the same from the clean start.'' If $H$ doesn't stop,
neither does the built program: the copy of $H$ never reaches a HALT, and nothing before $y$ can
trap. So the built program behaves the same as $y$ when $H$ stops from $(a, b)$, and the same as the
two-line program that loops, INC $A$ then DEC $A$ back to line 1, when it doesn't. A decider for
$\Pi$ would then decide whether $H$ stops from $(a, b)$, with the consequence stated in
Theorem~\ref{thm:u34-undecidable}; which of $y$ and $n$ goes at the end is the decider's own
answer about the looping program. ``Halts
from the clean start'' is carried by behaving the same from the clean start, held by the empty
program and lacked by the looping one, so it's covered.
\end{proof}

\paragraph{Proof of the theorem ``s-m-n, and a universal program for host programs'' (book Theorem~\ref{thm:host-smn}).}
\begin{proof}
(1) The prefix writes register 2 $c$ times and nothing else, so it moves that register's version from
0 to $c$. A fact records a version, so after the prefix every later version of register 2, and the
version in every fact the run records about register 2, is $c$ more than in a clean run. A CHECK or a
COMMIT reads versions only to ask whether two versions of the same register are equal, and adding $c$
to both keeps the answer, so the two runs take the same steps. The run is $V$'s run from the stated
start, moved down by $c$ lines. (2) is the evaluator. A finite algorithm decodes a number into a program, runs it step by
step on given registers, record layer included, and, searching over the number of steps,
returns the output of the first halting run. It is a finite tape algorithm, and the translation
of Appendix~\ref{sec:counter-constructions} turns any finite tape algorithm into a finite program
of INC and DEC on finitely many counters. The host runs those instructions on registers of its
own, with HALT marking the returned output. That's $W$.
\end{proof}

\paragraph{Proof of the theorem ``Kleene's recursion theorem on the host machine'' (book Theorem~\ref{thm:host-kleene}).}
\begin{proof}
(1) Let $E$ be the following program of INC and DEC, built the same way as $W$ from one finite
evaluator, with $T$'s number $t$ written in. On inputs $x$ and $c$ it computes the number of $\mathrm{spec}(c, P_c)$, where $P_c$ is the
program numbered $c$, runs $T$ on that number to get $d$, and then runs the program numbered $d$ on
$x$. Let $c_0$ be the number of $E$, and let $e = \mathrm{spec}(c_0, E)$. By (1) of
Theorem~\ref{thm:host-smn}, $e$ on input $x$ does what $E$ does on $x$ and $c_0$. And $E$ on
$c_0$ builds the number of $\mathrm{spec}(c_0, P_{c_0})$, which is the number of $e$, since $P_{c_0}$
is $E$. So $e$ hands its own number to $T$, gets $d$, and runs $d$ on $x$. (2) and (3) are (1) with $T$
carrying out $F$ or $f$.
\end{proof}

\paragraph{Proof of the theorem ``The recursion theorem, up to the record'' (book Theorem~\ref{thm:host-kleene-record}).}
\begin{proof}
The construction is Theorem~\ref{thm:host-kleene}'s, with a different evaluator. The first step is to
see that a record is small. From a clean start, a host program that stops has a record I can read off
five numbers, in this order: what it left in register 0; how many facts the table holds; what is
left of the ledger once the facts, the trap and the flag are paid for, which I'll call the replay allowance;
1 if the flag is up and 0 if not; and 1 if the trap latch is up and 0 if not. The record obeys a
short rulebook, and you can read each rule off the price table. A run stops at its first trap, so
its ledger is its passing CHECKs, plus its passing COMMITs, plus its passing CERTIFYs, plus one if
it trapped (only those three can fail, and the failing one is charged). Its passing CHECKs are its
facts, at most sixteen. Its flag is up exactly when some CERTIFY passed. So the allowance is the
passing COMMITs plus the passing CERTIFYs less one for the flag, which is never negative, and the
ledger is the facts, plus that allowance, plus the flag, plus the trap. A COMMIT needs a fact
behind it, so a positive allowance means at least one fact. A CERTIFY needs a commitment, so a
raised flag means a passing COMMIT, and so a positive allowance. And the channel is empty exactly
when no COMMIT passed, which is exactly when the allowance is zero, since a CERTIFY can't pass
before a COMMIT has. The allowance counts repeated CERTIFYs as well as COMMITs; it counts the cost the
replay can spend on commitments, not the original commitment events. Those are the only properties
of a record the replay will use.

The second step is five blocks, one for each number. Each block reads the code of $T$, runs $T$ on the
number of the specialised program, runs the program that comes out on $x$, and returns its number. All of
that happens on plain data: the blocks are programs of INC and DEC, built the way $W$ was, from a
finite evaluator translated into INC and DEC. So if $T$ has marks
of its own, they are simulated and never paid.

The third step is the payment. The evaluator puts 1 in a fresh register, which is
$\langle 0, 0 \rangle$, the pair of $\mathsf{Zero}$ and 0, so it satisfies the host's one property,
and CHECKs it once for every fact; there are at most sixteen, so the table has room. It COMMITs to that claim once for every unit in the replay allowance. It CERTIFYs
once if the flag is up. It copies the first number into register 0. And if the trap latch is up, it
CHECKs a register that holds 0, which fails, since 0 is no pair, and the host traps. Everything
before the payments is INC and DEC and costs nothing, and the payments cost what the record costs, so the ledger is exact. Specialise the
evaluator to its own number, as before.
\end{proof}

\paragraph{Proof of the theorem ``No fixed point of full behaviour'' (book Theorem~\ref{thm:host-no-exact}).}
\begin{proof}
Let $F(p)$ be the one-line program INC $r$, with $r$ one above the number of $p$. To carry this out on numbers, add one to the input code, encode INC with that register,
and encode the one-element instruction list. These are finite arithmetic operations of the
counter construction in Appendix~\ref{sec:counter-constructions}. A program's number is at
least every register it names: the number of an instruction is above its register, and the code of a
list is above every instruction in it. So $e$ can't name $r$, the register $F(e)$ writes. On input 0 the
run of $F(e)$ ends with 1 in register $r$, and if $e$ stops at all, it ends with 0 there. If $e$ doesn't
stop, then it already differs from $F(e)$, which does. Either way they don't behave the same.
\end{proof}

\paragraph{Proof of the theorem ``The diagonal on the host machine'' (book Theorem~\ref{thm:host-wall}).}
\begin{proof}
(1) The flip, send programs that $d$ calls good to the program lacking $\Pi$ and the others to the
one having it, is carried out by a host program by hypothesis. Theorem~\ref{thm:host-kleene}(2) gives
a program $e$ that computes what its image computes; $\Pi$ holds of $e$ exactly when it holds of the
image, and the image lacks it exactly when $d$ calls $e$ good, so $d$ is wrong at $e$. (2) From
a deciding program $D$ build a host program that computes $D$'s flip on numbers, by the finite-simulation search of Appendix~\ref{sec:counter-constructions}, reading the
Boolean from the first halting run of $D$; then (1) applies. This one is a plain no: the flip (1) needs is built from $D$ itself, so no premise about what can be represented is left over.
\end{proof}

\paragraph{Proof of the theorem ``Rice's theorem for the two-counter machine, on any inputs'' (book Theorem~\ref{thm:tc-rice}).}
\begin{proof}
For every input: given a two-counter instance $H$ starting from counters $a$ and $b$, build a program
that turns the input of $n$ pebbles in bin A into $(6n + 5) \cdot 2^a \cdot 3^b$, runs $H$ inside the
exponents of 2 and 3, divides out every 2 and every 3 when $H$ stops, takes away 5, divides by 6, and
then runs the program that has the property. Because $6n + 5$ shares no factor with 2 or 3, the coding
never mistakes it for an exponent, and the input comes back whole. If $H$ halts, the built program
behaves the same as the one with the property on every input; if not, it never stops on any input,
like the program that loops. As before, if the looping program has the property, put a program
that lacks it at the end instead, and the decider's own answer about the looping program says
which. So a decider for the property would decide two-counter halting.
\end{proof}

\paragraph{Proof of the theorem ``Kleene for the two-counter machine, inputs written as powers of two'' (book Theorem~\ref{thm:tc-kleene}).}
\begin{proof}
(1) The host's three pieces, with two counters in place of registers. First, the universal program.
Load a program's number $e$ and its input $x$ together as $2^x \cdot 3^e$ in bin A, with bin B empty.
The evaluator of Appendix~\ref{sec:counter-constructions}, the one that searches budgets and runs the
program numbered $e$ on $x$, is a finite counter program on finitely many counters, and packing those
counters into two by the prime-exponent construction of Appendix~\ref{sub:prime-packing} gives
$U_2$: started on $2^x \cdot 3^e$ it stops with $2^y$ in bin A exactly when the program numbered $e$,
started on $2^x$, stops with $2^y$. Second, specialising. For a number $c$ and a program $V$, let
$\mathrm{spec}(c, V)$ be $c$ copies of a block that multiplies bin A by three, followed by $V$ with
its jumps moved down. The block is INC and DEC only, so it adds nothing to the ledger and records nothing, and
started on $2^x$ the $c$ copies leave $2^x \cdot 3^c$; so $\mathrm{spec}(c, V)$ on $2^x$ does what $V$
does on $2^x \cdot 3^c$. Third, the fixed point. Let $T$ be a program that carries out $F$: on $2^n$,
$n$ the number of a program $p$, it stops with $2^m$, $m$ the number of $F(p)$. Build $E$ from the
same finite evaluator, with $T$'s number written in, and pack it into two counters as before: on
$2^x \cdot 3^c$ it computes the number of $\mathrm{spec}(c, P_c)$, where $P_c$ is the program numbered
$c$, runs $T$ on that number to get $d$, and runs the program numbered $d$ on $2^x$. Let $c_0$ be the
number of $E$ and $e = \mathrm{spec}(c_0, E)$. On $2^x$, $e$ does what $E$ does on
$2^x \cdot 3^{c_0}$: it builds the number of $\mathrm{spec}(c_0, P_{c_0})$, which is the number of $e$
itself since $P_{c_0}$ is $E$, hands it to $T$, gets the number of $F(e)$, and runs $F(e)$ on $2^x$.
The machine is deterministic, so a packed output is unique, and for every $x$ and $y$, $e$ started
on $2^x$ stops with $2^y$ exactly when $F(e)$ does. (2) is the counterexample that follows.
\end{proof}

\paragraph{Proof of the theorem ``No plain recursion on two counters'' (book Theorem~\ref{thm:no-plain-recursion}).}
\begin{proof}
(2) from (1): let $F(p)$ be a program that multiplies bin A by the constant
$c$ determined by the length of $p$. Multiplication by a fixed constant is the
transfer loop of Appendix~\ref{sec:counter-constructions}. Decoding a finite
program, counting its lines, finding the least prime above a number and writing that loop are
computable operations, so a program carries out $F$ on encoded program numbers.
A fixed point would multiply every plain input by the constant forbidden by (1).

For (1), suppose an $n$-line program did it. Split each state into the two bin counts and a
finite \emph{control}: the line (one of the $n$ lines, or ``stopped'' for any line with no
instruction), the trap, whether the channel is empty, and the table read entry by entry as a
fact's property, its counter, and whether it is \emph{current} (its version is its counter's
version now) or stale. Only properties the program names can enter the table, at most $n$ of
them, so an entry is one of at most $4n$ kinds, and a table of at most sixteen entries is one of
at most $\sum_{k \le 16} (4n)^k \le (4n+1)^{16}$ lists. So there are at most
$Q = 4(n+1)(4n+1)^{16}$ controls. The control and the two counts decide the next control and
counts: a CHECK writes a current fact, a write to a counter makes that counter's facts stale, a
COMMIT passes exactly when a current fact with its property and counter is in the table, a
CERTIFY passes exactly when the channel isn't empty, and the ledger, the flag, the channel's
contents and the versions themselves never decide a branch. Call the control with the two counts
a \emph{configuration}.

Pick a threshold $h \ge 1$ above every number $k$ in a property $\mathsf{Ge}(k)$ the program
names. At $h$ or above, a count answers every question the program can put to it (is it zero, is
it even, is it at least $k$, can a DEC take one off) the same way as any other count of the same
parity at $h$ or above. Call a control together with the parities of both counts a
\emph{pattern}; there are at most $4Q = N$ of them.

\emph{The forced run.} A \emph{stage start} is a configuration with one count below $h$, in the
\emph{exact} bin, and the other at $h$ or above, in the \emph{large} bin, which holds $v$. From
it, run the program while pretending the large bin can never fall below $h$: follow the exact bin
exactly, the large bin only by its parity, and keep a tally of the large bin's net change. This
forced run depends on the control, the exact count and the parity of $v$, and not on $v$ itself.
The real run does exactly what the forced run does, with the large bin at $v$ plus the tally, up
to the first step where $v$ plus the tally drops below $h$, the \emph{crossing}. A move changes a
count by at most one, so at the crossing the large bin holds exactly $h - 1$.

\emph{What a forced run does.} One of three things, fixed by the stage start. (a) It stops. (b)
It comes back to a step with the same control, the same exact count below $h$ and the same
parity of the large bin as an earlier step; it's deterministic in those, so it repeats that
stretch forever, and the exact count stays bounded. (c) Neither. Then it has at most $2Qh$ steps
with the exact count below $h$, one per combination, so from some step on the exact count stays
at $h$ or above for good. Among any $N + 1$ consecutive steps of that tail, two have the same
pattern, at steps $t_1 < t_2 \le t_1 + N$, and between them the exact count changes by some
$\delta$, even because the parities match. If $\delta < 0$, the run from $t_2$ repeats the run
from $t_1$ with the exact count lowered by $|\delta|$ each time round, so after finitely many
rounds it goes below $h$, which the tail never does. So $\delta \ge 0$, and the stretch repeats
forever, lifted by $\delta$ each time round, never dipping below $h$. In every case but (a),
then, the forced run from some step on repeats one stretch forever, and each time round the exact
bin changes by a fixed $g$ and the large bin by a fixed amount; in (b) $g = 0$, and in (c) the
stretch is at most $N$ moves long.

\emph{What the real run does,} once $v$ is large enough for that stage start, meaning the large
bin can't reach the crossing before the forced run stops or its repeating stretch has gone round once in full. If the
forced run stops, the real run stops the same way, with the exact bin at a count fixed by the
stage start and the large one at $v$ plus a fixed amount. If the large bin doesn't lose each time
round, the real run never crosses and never stops. If it loses and $g = 0$, the real run crosses
with both bins below a bound fixed by the stage start. The remaining case is a \emph{transfer}:
each time round the large bin loses $m > 0$ and the exact one gains $g > 0$. That's case (c), so
the stretch has at most $N$ moves, each moving one bin by at most one, and $m, g \le N$; both are
even. Now the stage rule. Start the same stage with $v + m$ in place of $v$. The parity is the
same, so the forced run is the same, and since the crossing comes after the stretch has gone
round once, the bin needs exactly one more time round to get down to $h - 1$. So the real run
crosses one time round later, with the same control, the formerly large bin at $h - 1$, and the
other bin larger by $g$. With $v + km$ it crosses $k$ times round later, the other bin larger by
$kg$. After a transfer, the crossing is a new stage start with the bins' jobs swapped. For
example, a loop that takes two pebbles from the full bin and puts four in the other has $m = 2$
and $g = 4$: two more pebbles going in buy one more trip round the loop, and four more coming
out.

\emph{Choosing the input.} Apart from $v$, a stage start is one of finitely many things: a
control, an exact count below $h$, a parity, and which bin is which. Take one bound $V \ge h$
that's large enough, in the sense above, for every one of them, and above every count that was
called fixed or bounded above: the exact bin of a stage that stops, and both bins at a crossing
with $g = 0$. The configurations with one count below $h$ and the other below $V$ are finitely
many. If the program multiplied every input by $c$, no two different inputs could pass through
the same configuration, because everything after it, the output included, would be the same,
while $cx \ne cx'$. So only finitely many inputs pass through that finite set. Pick an input $x$
that passes through none of it and has $cx$ above $V$. Since $(x, 0)$ is itself a configuration,
$x \ge V$.

\emph{Following that input.} The run from $(x, 0)$ starts a stage: bin B exact at 0, bin A large
at $x$. Every later stage start along the run has its large bin at $V$ or above, or the run would
pass through the finite set. The run stops, since the program multiplies, so no stage runs
forever. No stage crosses with $g = 0$, since that crossing is in the finite set. So each stage
before the last is a transfer, with its own pair $(m_i, g_i)$, and the last one stops. It stops
with its exact bin below $V$ and its large bin at its $v$ plus a fixed amount, and the output $cx$
is above $V$, so the output bin A is the large one. There are finitely many stages, because each
takes at least one step of a run that stops.

\emph{The comparison.} Put $D = 2\prod_i m_i$ and run the input $x + D$. It is $x$ plus a
multiple of $m_1$, so by the stage rule the first stage crosses at the same configuration except
that the other bin is larger by $(D / m_1) g_1 = 2 g_1 \prod_{i \ge 2} m_i$, which is a multiple of
$m_2$. So the second stage crosses the same way, with the next bin larger by
$2 g_1 g_2 \prod_{i \ge 3} m_i$, and so on through every stage. The last stage starts with its
large bin larger by $2\prod_i g_i$, an even amount, so it stops the same way with bin A larger by
exactly that. So the outputs from $x + D$ and from $x$ differ by $2\prod_i g_i$. But
multiplication by $c$ makes them differ by $cD = 2c\prod_i m_i$. Cancel two:
\[
c\prod_i m_i=\prod_i g_i.
\]
Empty products are 1. Every factor on the right is between 1 and $N$, and
$c$ is a prime above $N$, so it divides none of them, and a prime that divides a
product divides one of its factors. So $c$ can't divide the right side. That's the
contradiction.
\end{proof}

\paragraph{One slot, done by hand.}
Do one slot by hand. The property ``even'' has number 1. Checking it on value 4
puts $\langle1,4\rangle=(2\cdot4+1)2^1=18$ in a separate slot. To read that
pair, halve 18 until what's left is odd: one halving gives 9, so the
property number is 1; subtract one from 9 and halve to recover value 4. The host
checks that property of that value. A second register, the mirror, holds 2: the property
number plus one. COMMIT can read this second register while leaving the checked slot
untouched.

When the guest changes the counter that fact concerned, nudge the slot
$18\to19\to18$ and set the mirror to zero. The number is back where it was, but the version
went up twice, so the host's old fact about that slot is stale. Putting a number back
doesn't put an old check back. The mirror is what tells the search to leave that slot alone,
and if the search tried it anyway, the version rule would refuse it.

\begin{center}
\small
\begin{tabular}{@{}lccc@{}}
Act & slot value & slot edition & searchable mirror \\ \hline
Check evenness of encoded 4 & 18 & $v$ & 2 \\
Guest writes its counter: INC slot & 19 & $v+1$ & \\
DEC slot; clear mirror & 18 & $v+2$ & 0 \\ \hline
\end{tabular}
\end{center}

\paragraph{The program, register by register and job by job.}
\textbf{The program.} I'll describe $U$ the way Turing would, one register at a time and then
one job at a time. Registers 0 and 1 hold the guest's $A$ and $B$. Register 2 holds the guest's
whole program as one number: each instruction gets a number (an opcode paired with its operands,
INC as 0, DEC as 1, HALT as 2, CHECK as 3, COMMIT as 4, CERTIFY as 5), and the list of them is
packed with $\langle \cdot, \cdot \rangle$, the empty list as 0. Register 3 is the guest's line.
Registers 4 to 13 are scratch. Registers 14 and 15 count how many CHECKs on $A$ and on $B$ have
passed. Registers 16 to 47 are \emph{mirrors}, sixteen for each guest counter, and registers 48 to
79 are the slots, sixteen for each guest counter. The guest's table holds sixteen facts in
total, across both counters, and all sixteen can be about one counter, so the host gives each
counter room for all of them. The reservation is the host's, and the guest's own capacity stays
sixteen in all. Register
80 is a dead end that's never written.

Each guest step goes like this. At line 1 the program copies the program number and the guest's
line into scratch, putting the originals back as it goes, because copying a counter on this
machine means emptying it into two others and refilling it. If the guest's line is 0, or past
the end of the program, it stops. Otherwise it peels instructions off the front of the copy, one
per line, until it holds the one at the guest's line, and splits that into opcode and operands.
Then a chain of DECs on the opcode sends it to one of six handlers.

\begin{itemize}
\item INC $c$: INC the guest's counter $c$, nudge every slot about $c$ and zero its mirrors, add
one to the guest's line, back to line 1.
\item DEC $c$ $j$: DEC the guest's counter $c$. If it was 0, add one to the guest's line. If it
wasn't, nudge every slot about $c$, zero its mirrors, and set the guest's line to $j$. Back to
line 1.
\item HALT: stop.
\item CHECK $p$ $c$: let $k$ be the count of passed CHECKs on $c$. If $k$ is below 16, build
$\langle$number of $p$, value of $c\rangle$ in slot $k$ of $c$, run CHECK $\mathsf{Slot}$ on that
slot, set its mirror to the number of $p$ plus one, add one to the count and to the guest's line,
and go back to line 1. If $k$ is 16, run CHECK $\mathsf{Slot}$ on the dead end, which holds 0,
which is no pair, so it traps, just as the guest's seventeenth CHECK does.
\item COMMIT $p$ $c$: look through the sixteen mirrors of $c$ for one equal to the number of $p$
plus one. At the first match, run COMMIT $\mathsf{Slot}$ on that slot. With no match, run COMMIT
$\mathsf{Slot}$ on the dead end, which traps.
\item CERTIFY: run CERTIFY, add one to the guest's line, back to line 1.
\end{itemize}

Why the mirrors? COMMIT has to find the slot holding the claim, and it can't read the slot to
find out, because reading would make the fact stale. So beside every slot sits a mirror the
program is allowed to read: the property's number plus one while the slot's fact is live, 0 once
the guest has written the counter. If a guest CHECK fails its property, the host's CHECK on the slot fails too,
since the slot holds exactly the pair the guest checked. If the guest table is
full across both counters, the host's own sixteen-entry table is full as well,
even when this counter's slot count is below sixteen; its CHECK then traps for
lack of room. So both kinds of refusal come through.

\paragraph{The two tally lines for the surcharge.}
Keep two tally lines for a guest whose first three moves each cost one and whose
reading first says yes after the third. The guest's totals are 1, 2, 3. This
compiler pays one PAY on each of the first two moves, then its three-act chain:
the host totals are 1, 2, 5. The excess is two, added once. If the third move
instead costs four, the host runs the chain and one PAY there; both totals then
read 6. These are counts of charged instructions, not counts of all the work needed to decode the guest. Stopping at zero throws away what the chain paid beyond the move's price. A compiler that kept it as credit against the moves after the raise would bring the surcharge back to zero once the guest had paid three from the raising move on, and a guest whose reading starts at yes would be matched exactly from its third unit of cost on. I keep the simpler rule, and every surcharge stated below is this compiler's.


\paragraph{A worked replay, with its table.}
Try CHECK, COMMIT, CERTIFY, CERTIFY on a true claim. The ledger is four, the table holds one
fact, the flag is up and the trap is down, so the replay allowance is two. The replay pays
CHECK, COMMIT, COMMIT, CERTIFY: the same ledger of four and the same summary, in a different order.
The theorem claims that agreement and no more.

\begin{center}
\small
\begin{tabular}{@{}lccccc@{}}
Run & checks & commits & certificates & total paid & final flag \\ \hline
Original & 1 & 1 & 2 & 4 & up \\
Replay & 1 & 2 & 1 & 4 & up \\ \hline
\end{tabular}
\end{center}

Each table has one fact and each channel is nonempty. The counts of the
individual acts differ; the selected final summary agrees. In the replay, the trap latch is up exactly when the fifth number is 1, and the flag
exactly when the fourth is positive.

\paragraph{A worked packing, with its table.}
Try input 2 while the simulated counters hold 1 and 2. The pile is
$17\cdot2\cdot9=306$. Raising the first simulated counter doubles it to 612; lowering
the second divides it by three, to 204. Strip off every factor of two and three and 17
remains; subtract five and divide by six, and the original input 2 is back.

\begin{center}
\small
\begin{tabular}{@{}lccc@{}}
Act & packed pile & simulated counters & input factor \\ \hline
Load input 2; counters 1, 2 & $306=17\cdot2\cdot3^2$ & $(1,2)$ & 17 \\
Increment first counter & $612=17\cdot2^2\cdot3^2$ & $(2,2)$ & 17 \\
Decrement second counter & $204=17\cdot2^2\cdot3$ & $(2,1)$ & 17 \\
Strip factors 2 and 3 & 17 & discarded & 17 \\
Subtract 5; divide by 6 & 2 & & input recovered \\ \hline
\end{tabular}
\end{center}


\section{Running it}

Everything so far is definitions and proofs about them. Does any of it run? This appendix is
me finding out. I built the small machine twice, once straight out of the definitions and once
by hand from the same description, ran both, and compared the whole state after every step.
Then I ran the universal programs on guests small enough to finish, and real programs through
the part of the compiler that executes.

One rule holds for the whole appendix: the runs test the builds, and the proofs carry the theorems. A
run shows what a build did on the inputs it was given, and no theorem in the book leans on one: every theorem in it stands on its argument, on the page.

There are two counts to keep apart as you read. The ledger counts what the price list charges. The
number of interpreter steps counts the work this particular build does to carry the run out.
One guest step can take millions of host steps and add one to the ledger.

\begin{center}
\small
\begin{tabular}{@{}p{.27\linewidth}p{.63\linewidth}@{}}
Stage & What the next stage must preserve \\ \hline
Source orders & Variable values and completed procedure calls \\
Assignments and jumps & The same branches, loops and final answer \\
Counter blocks & Values at block boundaries; spare counters restored \\
Host instructions & The same result, with the opposite DEC convention translated \\
Packed two-counter guest & The result as an exponent in the packed counter \\
Fixed interpreter $U_P$ & Guest counters at matching points; decode the exponent to read the answer \\ \hline
\end{tabular}
\end{center}
This is a path through the compiler, not a claim that one step at each stage
takes the same time. The source compiler inserts no certification chain. Running the executable uses physical hardware and energy. These runs did not measure
heat per host step, and the ledger is not a heat measurement. The times only
tell you which runs are practical on this build. The table shows what ran and what didn't.

\begin{center}
\begin{tabular}{@{}p{0.42\linewidth}p{0.5\linewidth}@{}}
\textbf{What} & \textbf{Status} \\ \hline
The small machine & run: the executable versions agree on every tested step \\
The host machines, with and without PAY & run, compared the same way \\
$U$ and $U_P$ on tiny guests & run, against a separate guest implementation \\
$U$ on a guest that certifies & proved and counted, not run \\ $U_P$ on a guest that certifies & proved, not run \\
Source programs through the compiler & run as far as the host program \\
The two-counter guest and $U_P$ running it & proved, not run \\
The presented machine carried by $U_P$ & proved, not run \\
\end{tabular}
\end{center}

\subsection{Taking the coat off}

The first executable is generated straight from the formal definitions, the way you'd
strip the coat off a proof and keep the program underneath. The second was written
separately, by hand, from the same description. Both implement the small machine, the host with a register for every whole number, the host with PAY and the loaders, and both run the universal programs; their instruction lists the second build reads as data printed from the definitions, so it checks them by running them, not by writing them again. The compiler below has an executable runner too.

Whole numbers matter here. A machine word in one build is 63 bits wide, and the guest
codes $U$ reads go up past $2^{100}$, so ordinary fixed-width arithmetic would wrap around
and lie. The mathematical machine never overflows, so its executable uses integers of any
size, with the arithmetic chosen to match the definitions: subtraction stops at zero and
division by zero gives zero. What sits between the definitions and the running code, the
mappings, the compilers and the routines that read input and write output, is code I had to
trust.

\subsection{A second machine, written by hand, to check the first}

Comparing two builds step by step catches one of them disagreeing with the description, on the
runs I tried. They're compared on the \emph{whole} state after \emph{every} step: the program line, every counter
and register and its version, the fact table in order, the channel, the trap latch, the ledger and
the flag. The comparison covers exhaustive sets of short programs, random programs and both hosts, and
they agree on every tested step. $U$ and $U_P$ on tiny guests are compared at fifteen checkpoints two hundred steps apart, over the first three thousand host steps of each run, instead of after every step, and the longer runs below against a separate guest
implementation, once per pass of the loop and at the stop. I also broke the machine on purpose, giving it a fact table one too small, a
CHECK, COMMIT and CERTIFY that cost nothing, a commitment that ignores versions, and a write that doesn't bump the version, and each break
gets caught. A second set pins the boundaries this book draws in words, and those runs come out as the words say: two runs that end in the same state with different first certifications, a check that
can't fail certifying from every start, a CHECK and six commitments costing seven without a flag, a write
after the commitment and a flag over a claim that's now false, sixteen copies of one CHECK filling
the table, a CERTIFY costing one at a ledger already at two, and a merge between two yes-states
that leaves the entering state alone.

A mistake both builds share gets past the comparison, and so does a mistake in the definitions
themselves, since both builds follow them.

\subsection{\texorpdfstring{$U$}{U} on a real guest}

Now the universal program, which is where I expected trouble. The guest I wanted to run first was
the obvious one, a guest that goes the whole way: CHECK, COMMIT, CERTIFY, flag up, run inside $U$ on
the built machine. That isn't what ran, and the reason is the way the guest is stored. $U$ keeps the whole guest program as one number,
a number that grows like two raised to the sum of the codes for the guest's instructions (for a
guest of one instruction it's exactly a power of two), and it reads that number
the only way a counter can, one unit at a time. So its speed depends on the code, and the code grows
fast: a guest with one CHECK in it has code $2^{24}$, and a guest with the whole chain has a code
near $2^{106}$.

So I came down to what could finish. The table below is what ran, on the executable host, compared against a separate
guest implementation at every pass of the loop and at the stop: the host stopped exactly when the guest
did, and at the stop showed the guest's counters, trap latch, ledger and flag.

\begin{center}
\begin{tabular}{@{}llrr@{}}
guest & host & host steps & time \\ \hline
count A down from 2 & $U$ & 25,953,317 & \\
one passing CHECK & $U$ & 402,654,059 & 618 s \\
one failing CHECK & $U$ & 201,327,155 & 305 s \\
one passing CHECK & $U_P$ & 402,654,059 & 617 s \\
one failing CHECK & $U_P$ & 201,327,155 & 313 s \\
\end{tabular}
\end{center}

The first row's count is pinned exactly: if the build changed so that it moved by one step, the
comparison would fail. The times are from one ordinary computer and are only there to give you the
scale. The counts work out to about twelve host steps for every unit of the guest's code each time $U$ goes round
its loop, which it does once for every guest step and once more to see that the guest has
stopped: the failing CHECK traps on its first pass, $12 \times 2^{24}$ and a little, and the
passing one takes two passes, $24 \times 2^{24}$ and a little. The growth is exponential in the size of
the guest. A guest with one CERTIFY in it would have code $2^{32}$, about fifty billion host
steps, and at the measured rate of about 650,000 steps a second it would take something like
twenty-two hours, nearly a day. The guest I wanted, with the whole chain, has a code near
$2^{106}$, which at the same twelve steps a unit is about $10^{33}$ host steps for each pass; at the
measured rate, with this encoding on that computer, a run of its three steps comes to around
$10^{20}$ years. So $U$ on a guest that certifies is proved, in Appendix~\ref{sec:universal-host},
and has never been run. Runs this long only fit in memory at all because the build can repack the
host's register store in the middle of a run; repacking changes how the registers are stored, not
what they hold.

What a run that long can be is counted. Four loop shapes do most of $U$'s work.
Count from the DEC instruction that tests the loop counter; any instructions before
that test are still counted separately. You can derive each shortcut from the arrows:
\begin{itemize}
\item In the emptying loop, a positive counter $v$ takes one DEC back to its own test
and becomes $v-1$. At zero, one DEC exits. Thus $E(0)=1$ and
$E(v+1)=1+E(v)$, giving $E(v)=v+1$.
\item In the copying loop, a positive DEC jumps back one line to INC on a distinct
destination, which returns to the test. Each unit moved therefore takes two steps,
and the zero test takes one: $C(0)=1$, $C(v+1)=2+C(v)$, so $C(v)=2v+1$.
\item With two destination INCs before the test, each positive pass takes three steps.
Thus $D(0)=1$, $D(v+1)=3+D(v)$ and $D(v)=3v+1$. If the two destinations are the same
register, it receives two increments per pass; the step count stays three.
\item In the halving loop, two successful DECs on the source and one INC on the
quotient return to the same test. After $k=\lfloor v/2\rfloor$ such passes,
$3k$ steps have removed $2k$ from the source and added $k$ to the quotient.
The remaining source is zero or one. Follow that final branch instruction by instruction.
\end{itemize}

Each successful decrement or increment also raises its register's version once.
The zero test does not write. So the same counts determine the version changes as well
as the counter values and next line. These loops neither charge nor change the table,
channel, trap or flag. Replacing one completed loop stretch by its derived effect and
count therefore preserves the whole endpoint and the number of steps. Repeating that
replacement throughout a run preserves both by addition.

The numerical tally below applies those four rules and steps the other instructions
individually. It agrees with the three runs above that were also stepped all the way.
For the guest CHECK, COMMIT, CERTIFY, using the zero predicate on $A$, the passing start
$(0,0)$ ends with the flag up and ledger three. From $(5,0)$ the first check fails and
the guest stops with ledger one. Appendix~\ref{sec:universal-blocks} supplies the program
blocks and their guest correspondence; the displayed totals count that particular encoding.

\begin{center}
\small
\begin{tabular}{@{}lr@{}}
start & host steps \\ \hline
$(0,0)$ & 3,894,222,629,847,358,066,688,933,266,198,944 \\
$(5,0)$ & 973,555,646,581,507,138,873,939,791,446,643 \\
\end{tabular}
\end{center}

What that table shows is a count, not a run. Nothing stepped those $10^{33}$ instructions, and the row in the table
at the top of this section says so.

Those big counts come from this encoding and from reading it one pebble at a time. They
measure the work this interpreter does, and that's all they measure. Whether some other
interpreter, or some other base, could do the same job fast is a separate construction with its
own running-time analysis. Nothing measured here puts a floor under every way of doing it.

\subsection{From source programs to the host}

Running guests one CHECK at a time doesn't tell you whether the built machine can do anything you'd
recognise as a program. You'd want to hand it gcd, or a sort, written the way people write them, and
get the answer out, and writing gcd by hand in INC and DEC is no way to find out. So the last piece
of the build is a compiler. The source is a small language you could write by hand: whole-number
variables, addition, subtraction that stops at zero, comparisons, \emph{if}, \emph{while}, and
procedures that return several results, with the rule that a procedure can only call ones defined
before it, so there's no recursion. Each translation stage has a proof that it preserves the program's result: the procedures are
inlined, the result is flattened to assignments and jumps, that is compiled to Minsky's counter
machine through the counter-machine compiler, that to a program for the host machine, that to a
two-counter guest, and the guest to a run of $U_P$.

\begin{theorem}[The compiler is correct]
For every source program with no recursion, and every input: the program computes $y$ exactly when
the runner (the bounded evaluator of the compiled host program) halts with $y$ for some amount of fuel; exactly when the compiled host program
halts with $y$; exactly when the two-counter guest halts with $y$; and exactly when $U_P$ running that
guest halts with $y$, with its first counter at 0, its ledger at 0 and its flag down. So a terminating compiled program delivers its answer and certifies nothing. To certify a
claim about that answer, a target program must add a CHECK, a COMMIT and a CERTIFY;
this compiler does not insert that chain. The runner also
takes exactly as many steps as the host program does, with the rest of its fuel left over, and a run
that doesn't halt has used all its fuel.
\end{theorem}

\begin{proof}
Six arrows and a runner, each argued below, and the statement is their composition. Four of the
arrows are between machines that take steps, and for those one argument does both directions, so
I give it first.

\emph{The simulation argument.} Take two machines $S$ and $T$, each deterministic, where $S$ has a
next state from every state inside its program, and a relation $\sim$ between their states such
that: a state of $S$ inside its program is related only to states of $T$ inside its program, and
a state of $S$ outside it only to states of $T$ outside it; and from any related pair with $S$
inside its program, one step of $S$ is matched by one or more steps of $T$ that land related to
$S$'s next state. Then from related starts, $S$ halts exactly when $T$ does, and they halt
related. Forward: start with the related pair. Follow the stretch of $T$ that
matches $S$'s first step; its endpoint is the related start needed for the
second stretch. Keep joining stretches until $S$ stops. There are only
finitely many steps to match, so this gives a finite run of $T$ whose final
state is related to $S$'s final state. The latter is outside its program,
so the former is outside its own. Backward: suppose $T$ halts after $k$ steps and $S$ doesn't halt. Then $S$ takes at least
$k + 1$ steps without leaving its program, the matched stretches give $T$ a run of at least
$k + 1$ steps all of whose stretch ends are inside its program, and $T$, being deterministic, has
only one run from its start, which halted at step $k$. That's a contradiction, so $S$ halts, and
by the forward direction it halts related to where $T$ did.

\emph{Unfolding the calls.} Imagine a program written on cards. A call says
``carry out the orders on that other card, then come back.'' Here a card may
call only a lower-numbered card. Card zero therefore calls nothing. Replace
every call on card one by the orders from card zero; then do card two, using
the already expanded lower cards, and continue. There are finitely many cards,
and every replacement has already been built. The process ends with calls
replaced by actual orders.

The variables need their own places too. Give each active procedure a block
of numbered boxes, its \emph{frame}. If the caller's block starts at box $b$
and uses $n$ boxes, put the called procedure's block immediately after it,
starting at $b+n$. A call becomes four pieces: copy the arguments into that
new block; set its other boxes to zero; carry out the expanded orders using
those box numbers; copy the answers back into the caller's named boxes.
A local variable numbered $x$ in the called procedure is thus the global
box numbered $b+n+x$. No assignment reaches below the procedure's own block.

Why does this preserve a terminating run? Begin with the lowest-numbered
procedure. It has no calls, so its expanded orders are its original orders
with the boxes renamed. Assignments change the same values, tests choose
the same branches, and loops repeat the same bodies. Now take a procedure
whose lower-numbered calls have already been settled. Each replacement call
builds exactly the arguments and zeros that the original call would receive.
The expanded lower procedure returns the same answers, and the copy-back
puts them in the same caller boxes. The caller's surrounding orders then
continue from the same values. Working upward through the numbered procedures
settles the main program.

The comparison works backward as well. Whenever an expanded call finishes,
its four pieces have run in their written order. The middle piece is the
expanded lower procedure, whose terminating runs I have already matched to
the original lower procedure. Its preparation and copy-back are just the
original call's input and output operations. Replace that finished four-piece
execution by the call it implements. The surrounding assignments, branches
and completed loop repetitions then give the same original execution.
A finite execution in either description therefore gives a finite execution
in the other, with the same answer. Neither can terminate when the other does
not.

In symbols, matching a source arrangement $e$ to the global boxes $G$ in a
frame starting at $b$ means $G(b+x)=e(x)$ for each local box $x$. The comparison
above preserves that equality at the beginning and end of every translated
statement.

\emph{Flattening.} The flat code of a statement placed at address $a$: an assignment is one
assignment; a sequence is the first statement's code followed by the second's; \emph{if} $b$ is a
jump on false to the else part, the then part, a jump past the else part, and the else part; and
\emph{while} $b$ with a body of $l$ instructions is a jump on false to $a + l + 2$, the body at
$a + 1$, and a jump back to $a$. Every jump in the code of a statement lands inside that code or
at its end, the address just past it. Now follow a completed structured execution.
An assignment performs the same change; a sequence joins its two runs; a branch
follows the same test answer. For each loop iteration, the test falls through,
the body's code takes the run to $a+l+1$, and the jump returns it to $a$.
Repeat for the remaining iterations; the final false test jumps to the end.
Thus an execution from variables $e$ to $e_1$ gives a flat run from address $a$
to the end of the code with those same final variables.

Read a completed flat run in the other direction. Because no jump escapes past
the code's end, a run that leaves it must visit that end first. Assignments and
branches recover their structured operations immediately; a sequence splits at
the second statement's first address. For a
loop, the run's first step is the test. If it is false, the run jumps to the end, and the loop's
derivation is the one with a false test. If it is true, the run enters the body's code, which it
can leave only through its end at $a + l + 1$ (the body's own backward statement), where the jump
returns it to $a$ with a strictly shorter run left. Read the next iteration the same way.
This cannot continue beyond the finite run I started with: its final test must
be false, giving the last piece of the structured loop execution. So the
flat program at address 1 halts exactly when the structured statement has a derivation, with the
same final variables.

\emph{Counter machine.} Here a DEC jumps when its register is zero and goes on otherwise. Variable
$x$ lives in register $x + 1$, register 0 is kept at 0, so DEC 0 $j$ is a plain jump to $j$, and
registers above the variables are working space. The relation: register 0 is 0, registers 1 to
$n$ hold the flat variables, and every register above is 0. Each flat instruction becomes a block built from small loops. Follow one
pass and then count how many pebbles remain to be moved:
Emptying a register: each pass takes one off, and the loop exits at zero. Moving $s$ into $d$:
each pass takes one off $s$ and adds one to $d$, so $s + d$ never changes, and at exit $s$ is 0
and $d$ holds the old sum. Copying: move $s$ into two registers at once, then move one of them
back. Cut-off subtraction $t - u$: take one off each while both are positive; at exit one of them
is 0, and the other holds the difference, or 0. Comparison: take one off each in step; whichever
reaches zero first, or both together, decides $<$ and $=$, and the block empties both. Build a larger expression from these smaller jobs. A constant is a
fixed number of increments into an empty working register. A variable is
copied there with its original restored. For a sum or cut-off subtraction,
compute the two smaller expressions in separate working registers, combine
them with the loops just described, and empty the temporary registers.
Starting with the innermost pieces and working outward, the final value
lands in register $t$, the registers above $t$ return to zero, and those
below $t$ are left alone. A test computes its comparison in the same spare
registers, chooses the true or false exit, and clears the spare registers
on either exit. Combining tests follows the selected exits of their smaller
tests, each of which has already restored the registers. Thus a test leaves
the registers as it found them and changes only which address comes next. So an
assignment block computes the value into the first working register, empties register $x + 1$,
moves the value in, and lands related; a jump on false runs the test's code with the true exit at
the next block and the false exit at the block of $j$; a jump is DEC 0 $j$. The linker places the
blocks one after another and sends every jump to the start of its target's block. That's the
one-step matching. The flat machine has a next state from every address in its program, and an
instruction that names a variable outside the frame just advances, so the step is total. The
counter machine is deterministic. The simulation argument gives both directions: the counter
program from the loaded inputs leaves its code with register $\mathit{out} + 1$ equal to $y$
exactly when the flat program halts with $y$ in $\mathit{out}$.

\emph{Host.} The host's DEC is the other way round: it jumps when its register is positive, taking
one off, and goes on at zero. An INC $x$ stays INC $x$. A DEC $x$ $j$ at host address $a$ becomes
three instructions: DEC $x$ $(a + 3)$, then INC 0, then DEC 0 $j'$, where $j'$ is the host address
of $j$'s block. If register $x$ is positive, the first takes one off and jumps to $a + 3$, the
start of the next block, which is what the counter machine does. If it is 0, the run goes on to
INC 0, which raises register 0 to one more than it held, and DEC 0 $j'$ finds it positive, takes
the one back off and jumps to $j'$, which is the counter machine's jump at zero, with register 0
as it was. The relation: the same registers, at corresponding addresses. Each counter-machine step
is matched by one or three host steps. The host program has no CHECK, COMMIT, CERTIFY or PAY, and
INC and DEC change only a register and the line, at no cost, so every host state in the run has
the start's facts, channel, trap latch, ledger and flag. The counter machine's step is total, the
host is deterministic, and the simulation argument gives: the host program halts with register
$\mathit{out} + 1$ equal to $y$ exactly when the counter program leaves its code with that
register equal to $y$.

\emph{Guest.} The host program's registers $r_0, \ldots, r_{k-1}$ are packed into the guest's
counter $B$ as $\prod_i p_i^{r_i}$, one prime $p_i$ per register, with counter $A$ the spare. That
is the packing of Appendix~\ref{sub:prime-packing}, with the two counters' jobs swapped. The
relation: the guest is at the start of the block for the matching instruction, $B$ holds the code
of the registers, and $A$ is 0. An INC on register $i$ is the stretch that multiplies $B$ by
$p_i$; a DEC on register $i$ is the stretch that divides $B$ by $p_i$ with remainder and takes the
positive branch when the remainder is 0, the zero branch otherwise, restoring $B$. The appendix
argues each stretch: it ends, because each of its loops empties the counter it counts down; it
lands with $A$ at 0; and, by unique factorisation, it changes the exponent of $p_i$ alone, by the
right amount. The host's step is total on these programs and the guest is deterministic, so the
simulation argument gives: the guest halts with $B$ coding registers whose register
$\mathit{out} + 1$ is $y$, that is, with $y$ the exponent of that register's prime in $B$, exactly
when the host program halts with that register at $y$. The guest has no record instruction, so its
ledger stays 0 and its flag down.

\emph{$U_P$.} $U_P$ keeps pace with its guest the way $U$ does, with the PAY case added (the proof of Theorem~\ref{thm:u34-presented}). This guest has no PAY and no record instruction, so on it $U_P$ runs only branches $U$ has, and the argument of Theorems~\ref{thm:u34-U-sim} and~\ref{thm:u34-U-halting} goes through: from a freshly loaded host $U_P$ halts exactly when the guest does, with its registers 0 and 1 holding the guest's $A$ and $B$, and its
ledger and flag equal to the guest's. So it halts with register 0 at 0, register 1 the code with
$y$ in it, ledger 0 and flag down.

\emph{The runner.} The runner keeps the host program and the registers in binary tries, a tree
indexed by the bits of a number. Read the address one bit at a time: a zero chooses one branch and a one
the other. Looking up an address follows that route to its stored value,
or returns the default zero if no value was stored there. Storing at an
address follows the same route and changes its final value, creating missing
branches as needed. Every other address either branches off before that end
or names a different node, and keeps its value. Thus the tree performs the
same lookup and single-address update as the register map it represents. So each round of the runner's loop does
exactly what one host step does to the same registers, and the runner stops when the line leaves
the program or the fuel is spent. If it halts after $k$ rounds, the host run of $k$ steps is
halted at the end of the program with the runner's registers, and the fuel left is the fuel given
less $k$. If the source program has a derivation, the arrows above give a halting host run, of
some length $k$, and the host, being deterministic, has only that run; any fuel of at least $k$
makes the runner take those $k$ steps and stop. If the runner hasn't halted, it used every unit of
fuel on host steps that stayed inside the program, so the host is still running after that many
steps, and a derivation, which would give a halting host run of some length $k$, would need $k$
larger than the fuel.

Composing: the program computes $y$ exactly when the unfolded program does, exactly when the flat
program halts with $y$, exactly when the counter program does, exactly when the host program does,
exactly when the guest halts with $y$ coded in $B$, exactly when $U_P$ halts as stated, and exactly
when the runner halts with $y$ on some fuel.
\end{proof}

The theorem says nothing about how many steps the compiled program takes next to the source; it's
about answers.

The runner and the compiler are implemented, and the build has been run on a set of programs I picked
to be annoying: gcd, factorial, a primality program written with procedures, an Ackermann tower, a sorting
network on a list packed into one number, and a tiny interpreter for another machine written in the
source language. Then come 220 random programs that halt, 160 random programs with loops that might not,
programs that never halt run for a fixed budget, and programs mutated on purpose to make sure the
comparison notices. Answers are compared with direct arithmetic and separate interpreters of the source and
compiled programs. Ackermann's function at $(3, 5)$ is 253; the
source program takes 64,768 operations to get there, and the compiled host program takes about 119
million steps. The executable runner does a few million host steps a second, so that's under a minute.

That works out to about 1,800 host steps per source operation for Ackermann, and up to about a hundred thousand
for multiplication by repeated addition, because every addition copies a running total one unit at a
time. On this build the work grows with the size of the numbers being handled, because every
operation is one pebble at a time. How it grows with the size of the program I haven't
measured. The last two stages, the two-counter guest and $U_P$ running it, are proved and have
never been run, and $U_P$ on a compiled guest would be slow for the same reason $U$ on a
certifying guest is. Still, I like that 253. It came out of nothing but INC and DEC, and a proof
that they mean what the source meant.

\subsection{What wasn't built}

There's no chip. The model is not a circuit, and the lack of one says nothing about it. The example was never cut into metal: there's no circuit for the small machine, no
simulation of one at the level of gates, and no board ever ran any of this. The small machine exists
as mathematical definitions and executable implementations that
agree on every tested step. That's the whole of what ``it runs'' means in this book.

I'd love to see it in metal some day. But builds that can be held against each other on every step are worth a lot to
me.
A board with nothing to hold it up against isn't, and that's the trade I made.

%===========================================================================


\section{Constructing the quoting operation}
\label{sec:quoting-construction}

The proof in Section~\ref{sec:diagonal} uses a program that takes the code of a term
and returns the code of that code. Here's how to build it.

Draw the instruction ``add one'' on a card. Now make a second card describing the
first card's marks, rather than carrying out the addition. The count has not risen:
you have described an instruction. Describe the second card in turn and you have
a description of a description. Quoting does that second job mechanically. In the
language below, take the variable of index 0. Its code is a tree with three lambda
nodes, then an application of variable 2 to the numeral for 0. Quoting that code
builds a description of those three lambdas and that application; it does not
look up variable 0 or execute the term it describes.

Nothing in the operation lets a program inspect itself by magic: its input is an
ordinary tree encoded as a closed lambda term.

The construction follows Forster and Smolka (2017), whose construction for L has the same
three syntax builders on Scott-encoded children, an auxiliary procedure for numerals, and a quote
procedure $Q$ with the same specification, defined by the same recursion on codes; the
self-application that uses $Q$ in Section~\ref{sec:diagonal} is their fixed-point lemma, after
Barendregt's (1984) proof for the full lambda calculus. What follows is my rebuilding of it, with
the reductions written out.

\subsection{Numerals and syntax as data}

For readability, the formulas below give names to bound variables. Renaming a bound
variable changes nothing; the underlying syntax uses the de Bruijn indices of
Definition~\ref{def:u12-L}. Applications associate to the left. Write $X\Downarrow Y$
here for a finite sequence of the weak call-by-value steps ending in the value $Y$.

Scott numerals are the following closed values:
\[
\overline{0}=\lambda z.\lambda s.z,
\qquad
\overline{n+1}=\lambda z.\lambda s.s\,\overline n.
\]
So a numeral selects either the supplied zero case or the supplied successor case,
passing its predecessor to the latter. The cases below are values, so choosing one
does not prematurely run the other.

The three syntax constructors, operating on already encoded children, are
\[
\begin{aligned}
\mathrm{mkvar}&=\lambda n.\lambda v.\lambda a.\lambda l.(v\,n),\\
\mathrm{mkapp}&=\lambda p.\lambda q.\lambda v.\lambda a.\lambda l.a\,p\,q,\\
\mathrm{mklam}&=\lambda p.\lambda v.\lambda a.\lambda l.l\,p.
\end{aligned}
\]
Each builder is a term of the language $L$ of Definition~\ref{def:u12-L}.
For a variable of index $n$, an application and an abstraction, respectively, define
their codes by the following resulting values:
\[
\begin{aligned}
\ulcorner\#n\urcorner&=\lambda v.\lambda a.\lambda l.(v\,\overline n),\\
\ulcorner p\,q\urcorner&=\lambda v.\lambda a.\lambda l.a\,\ulcorner p\urcorner\,\ulcorner q\urcorner,\\
\ulcorner\lambda p\urcorner&=\lambda v.\lambda a.\lambda l.l\,\ulcorner p\urcorner.
\end{aligned}
\]
In particular, $\mathrm{mkvar}\,\overline n$, $\mathrm{mkapp}\,\ulcorner p\urcorner\,\ulcorner q\urcorner$ and
$\mathrm{mklam}\,\ulcorner p\urcorner$ reduce to these three codes. A code is closed even when
the term it describes has a free variable: the index of that variable is stored
as a numeral. It is data inside the code, and the code itself has no free variable.

\subsection{Recursion without running under a lambda}

The quoting procedure will have to call itself on the smaller pieces of a
description. There is a catch: this language works out an argument before
handing it to a function. Hand over the recursive calculation itself and it
starts calculating at once, before the receiving function has decided whether
it needs it. That can send you straight back into the same calculation forever.

Instead, hand over a recipe waiting for an ingredient. In this language a
lambda is such a recipe, and it counts as a finished value even though its
body has not run. The receiving function can keep the recipe and call it when
it reaches a smaller piece. In the expression below, the little recipe
$\lambda a.x\,x\,f\,a$ is the part doing that job. The letter $a$ names the
argument it is waiting for; $f$ names the work to do once the recipe and the
current input have been handed over.
\[
W=\lambda x.\lambda f.\lambda v.\ f\,(\lambda a.x\,x\,f\,a)\,v,
\qquad
\operatorname{rec}(F)=\lambda v.W\,W\,F\,v.
\]
When $F$ and $u$ are closed values, beta reduction gives
\[
\operatorname{rec}(F)\,u
\longrightarrow W\,W\,F\,u
\longrightarrow^{*} F\,(\lambda a.W\,W\,F\,a)\,u
=F\,(\operatorname{rec}(F))\,u.
\]
The recursive argument is a lambda, so it's already a value. It starts work only when
$F$ applies it to an argument. That is the whole reason this construction survives
the call-by-value rule.

\subsection{Quoting numerals}

Let $E_j=\ulcorner\#j\urcorner$ for $j=0,1,2$, and let
$Z=\ulcorner\overline0\urcorner$, obtained directly from the finite syntax of
$\overline0=\lambda\lambda.1$. Write $\mathrm{mklam}^k(X)$ for $k$ nested applications of the
syntax builder $\mathrm{mklam}$ to $X$; it is an abbreviation for constructing $k$ lambda nodes, and the
language gains no new operation from it.

Define
\[
F_N=\lambda r.\lambda n.\ n\,Z\,
       (\lambda k.\mathrm{mklam}^2(\mathrm{mkapp}\,E_0\,(r\,k))),
\qquad Q_N=\operatorname{rec}(F_N).
\]
Follow the two branches. Zero chooses the already built description $Z$.
A number one larger than $n$ chooses the other branch and hands it the
predecessor $n$. That branch first asks the same procedure to describe the
predecessor. Then it puts the description inside an application of variable
zero and wraps two lambdas around it. Those are precisely the marks in the
numeral for $n+1$.

Zero therefore works. If the procedure describes $n$ correctly, the wrapping
just described makes it describe $n+1$ correctly. Starting at zero and taking
one more step proves the claim for every natural number. Every call on a
positive number asks about a smaller one, so it reaches zero and finishes.
In the notation for termination and returned value, the result is
$Q_N\,\overline n\Downarrow\ulcorner\overline n\urcorner$.

Do the smallest case by hand. Call the successor case $G$. Then $Q_N\,\overline1$ steps to
$F_N\,Q_N\,\overline1$ and on to $\overline1\,Z\,G$; the numeral hands its predecessor
$\overline0$ to $G$, the recursive call $Q_N\,\overline0$ comes back as $Z$, and the builders
return $\mathrm{mklam}(\mathrm{mklam}(\mathrm{mkapp}\,E_0\,Z))$. That is the code of
$\lambda\lambda.(\#0\,\overline0)$, which is $\overline1$.

\subsection{Quoting terms}

For a proposed recursive call $r$, define three handlers:
\[
\begin{aligned}
H_V&=\lambda n.\mathrm{mklam}^3(\mathrm{mkapp}\,E_2\,(Q_N\,n)),\\
H_A(r)&=\lambda p.\lambda q.\mathrm{mklam}^3(\mathrm{mkapp}\,(\mathrm{mkapp}\,E_1\,(r\,p))\,(r\,q)),\\
H_L(r)&=\lambda p.\mathrm{mklam}^3(\mathrm{mkapp}\,E_0\,(r\,p)).
\end{aligned}
\]
Now set
\[
F_Q=\lambda r.\lambda t.t\,H_V\,(H_A(r))\,(H_L(r)),
\qquad Q=\operatorname{rec}(F_Q).
\]
The handlers are lambdas and so are values. The Scott code of the input chooses
one of them and hands it the encoded children. Start at the leaves and build
outward. At each node the handler builds the syntax of that node's encoding:
$Q\,\ulcorner t\urcorner\Downarrow\ulcorner\ulcorner t\urcorner\urcorner$. There are
three kinds of node to consider:

\begin{itemize}
\item For $t=\#n$, the first handler constructs three lambda nodes over an application
of variable $2$ to the quoted numeral. That is the syntax of
$\lambda\lambda\lambda.(\#2\,\overline n)$, the code of $\#n$.
\item For $t=p\,q$, quote the two children first. The second handler
constructs three lambdas over variable $1$ applied to those two codes. That is
the syntax of the application code.
\item For $t=\lambda p$, the last handler constructs three lambdas over variable $0$
applied to the quoted child code. That is the syntax of the abstraction code.
\end{itemize}

Each recursive call works on a proper child of the finite input tree, or on a
smaller numeral. No chain of calls can keep descending through a finite tree
or decreasing a natural number forever. Each builder takes finitely many beta
steps, so the construction finishes and returns the code just described.

The $\mathrm{mkapp}$ used in Section~\ref{sec:diagonal} is this builder. With $Q$ and that
builder written out, every piece of the self-referential term there is on the page.

%===========================================================================

\section{Computing with counters}
\label{sec:counter-constructions}

Put a few beans in one cup and leave a second empty. Write a short list of orders:
add one, take one away if there is one, go to this line, stop. A finger on the page
keeps your place. Carry out the orders literally. The question in this appendix
is how to give the same job to another person with different equipment (a strip
of paper, stacks of cards, or a different set of cups) and get the same answer
whenever the first person stops.

Agree on how the input is laid out and how to read the output before either starts.
One person must not quietly solve the problem while preparing the other's cups.
The reductions in this appendix concern functions on natural numbers. To compute a
relation $R(\vec x,y)$ means to halt with output $y$ exactly when $R(\vec x,y)$ holds,
and to produce only numerical outputs on halting runs. The machines below are
deterministic, so any relation they compute is the graph of a partial function.
A translation must preserve both the output and whether the run terminates.
Its input and output encodings are part of the translation.

\subsection{Finite programs made from two instructions}

A counter holds a natural number. INC adds one and advances to the next line.
DEC subtracts one and jumps to its named line when the counter is positive; at zero,
it leaves the counter alone and advances. A finite program uses finitely many counters.
An unconditional jump to $j$ is INC $r$ followed by DEC $r$ $j$: it restores $r$'s
value and always jumps. A zero test is DEC $r$ to a positive branch that first restores
$r$; its fall-through is the zero branch. These operations preserve values, not versions.
They are used here for computation, where there are no live facts to preserve.

Try copying a cup with three beans. Put two empty cups beside it and keep a supply
of spare beans. Each time you remove a bean from the first cup, put one in each of
the other two. When the first is empty, move the third cup's beans back into it.
Now the first and second each have three, and the third is empty. Try zero beans,
then five. At every pass one more of the original beans has been counted into
both other cups; that is why it works for any starting count, not just the three
you tried. These are the moves behind the copy $C(x,y,t)$ in the next appendix.

Emptying a counter is a DEC loop back to its own line. Moving a pile decrements
its source and increments its destination until the source is empty. Copying
uses two such destinations and then restores the source from one of them.
Addition, multiplication by a fixed natural number, and bounded loops are finite
combinations of these moves. Division by a fixed positive number $b$ needs only
$b$ control states: decrement the input, cycle through remainders $0,\ldots,b-1$,
and increment a quotient counter whenever the cycle returns to zero. On exit the
quotient is in that counter and the remainder is in the control state.

\subsection{A tape in finitely many counters}

Cut a strip into numbered squares and put a finger on one square. Lift the squares
on its left into a stack, nearest on top, and do the same on its right. Moving your
finger one place right is now two acts: put the square you were reading on top of
the left stack, then take the top of the right stack as the new square. Write a
new mark before moving and the stack keeps the changed square. Move left and the
two stacks exchange jobs. You can walk the strip either way without losing its
order. An exhausted stack supplies a fresh blank square.

Now replace each stack by a count. Use digit 1 for a blank and digit 2 for a mark,
with base 3. Push a mark onto an empty stack: its code is 2. Push a blank on top:
its code is $3\cdot2+1=7$. Count seven beans in groups of three. The one left over
says ``blank on top''; the two groups say the remaining stack has code 2. Pop that
and the marked square comes back. Multiplication pushes and division pops.

In general, number a finite tape alphabet by $1,\ldots,q$, including its blank,
and put $b=q+1$.
Encode a finite stack whose top symbol is $d$ and whose remaining stack has code $n$
by $bn+d$; the empty stack is 0. A push is multiplication by $b$ followed by $d$
increments. A pop is division by $b$: the remainder is the top symbol and the quotient
is the remaining stack. The digits are nonzero, so no nonempty stack has code 0.
Popping an empty stack gives remainder 0, which is no digit; the control reads it as the blank and leaves the stack empty.

A tape configuration is a control state, the scanned symbol, and two stacks: the
symbols to the left of the head, nearest first, and the symbols to its right,
nearest first. The control state and scanned symbol have finitely many possibilities,
so they fit in the program's line number. To write $d$ and move right, push $d$ onto
the left stack and pop the right stack for the new scanned symbol. A left move does
the reverse. The arithmetic just described implements each transition in a finite
terminating stretch of counter instructions. Extra stored blanks do not change the
represented tape. A finite number of tapes uses two stacks per tape and finitely
many scratch counters. Halting states translate to halting lines. Induction over tape
steps proves the simulation; termination of every arithmetic stretch proves its converse.

In the other direction a tape can hold the finitely many counter values as unary
strings separated by delimiters. INC appends a mark in the selected string; DEC erases
one when present and otherwise detects the empty string. The current instruction is
finite control. Each scan and rewrite terminates, so halting and the output are preserved.
Multiple tapes may be represented on one tape by a list of tracks with marked head
positions: a complete scan reads the marked symbols, a second scan writes their changes
and shifts their markers. Extend the represented interval by a blank cell when needed.
A fixed-length binary code for the finite alphabet implements each scan with binary
symbols and finitely many intermediate control states. So this argument also applies
to the single-tape binary halting question.

\subsection{Packing finitely many counters into two}
\label{sub:prime-packing}


Try holding two counts in one pile of eighteen beans. Count how many times you can
halve eighteen exactly: once. Count how many times you can divide it exactly by
three: twice. Reassemble the eighteen beans after each reading. Those are the two stored counts,
1 and 2. Double the pile to 36 and
only the first stored count rises. Divide 36 by three to get 12 and only the second
falls. Divide by three again: four beans, storing 2 and 0.

Try taking one more off that second count. Four beans do not divide evenly into
threes. The attempted division must restore the four and take the zero branch;
keeping only the quotient would corrupt the first count. This is why the procedure
below remembers the remainder. Two and three share no factor, so changing the
number of one kind of factor leaves the other alone. Assign another prime to
each extra count. An all-zero collection is represented by one bean: the empty
product is 1.

Let $p_0,\ldots,p_{k-1}$ be distinct primes. At the boundary between simulated
instructions, keep
\[
A=\prod_{i=0}^{k-1}p_i^{r_i},\qquad B=0,
\]
where $r_i$ are the simulated counter values. The simulated line is finite control.
The representation is unique by prime factorization. There are enough primes for
any finite program: a prime divisor of $1+\prod_i p_i$ is not among the primes already
chosen, so the list can always be extended.

To simulate INC on $r_i$, empty $A$ into $B$, adding $p_i$ to $B$ per decrement,
then move $B$ back to $A$. The result is $p_iA_{\rm old}$ and $B=0$.
To simulate DEC, divide $A$ by $p_i$ using the finite remainder cycle described above.
On exit $A=0$, $B=q$ and the control holds $d$, where $A_{\rm old}=p_iq+d$ and
$0\le d<p_i$. If $d=0$, the original number was divisible by $p_i$, so $r_i>0$.
Move $B$ to $A$ and take the positive branch. Its exponent of $p_i$ has decreased by
one and every other exponent is unchanged. If $d>0$, then $r_i=0$. Restore $A$ by
moving $B$ back with $p_i$ increments per unit, add $d$, and take the zero branch.
Both branches restore $B=0$. All loops terminate because their source counter decreases.
Induction gives the represented run at every simulated instruction boundary, and
no simulated instruction can conceal an infinite internal run.

This is universality with an encoding. With input $n$ in the first simulated counter
and the others zero, the packed input is $2^n$ when $p_0=2$. It is supplied to the
two-counter machine as its input; the construction does not claim that two counters
starting at $(n,0)$ can first manufacture $2^n$. That is precisely the distinction
Schroeppel's theorem in Section~\ref{sec:smallest} turns on, and the undecidability used by the
one-counter theorem in Appendix~\ref{sec:every-model} ranges over packed starts too.
For a tuple of inputs, supply its product of prime powers. If a numerical output is
needed in ordinary notation, the external decoding reads the designated prime's
exponent. A host with finitely many additional counters can perform this decoding itself
by repeated fixed-prime division.

\subsection{Why halting has no decider}

Write two responses on a card: ``If the predictor says STOP, keep going forever.
If it says KEEP GOING, stop.'' Suppose somebody offers a predictor that always
finishes and always gets a program's eventual behaviour right. Give it this very
card, with these instructions for using its answer. Work through both possible
answers. STOP makes the card keep going; KEEP GOING makes it stop. Waiting longer
for the predictor does not create a third correct answer: its promise was to
finish with one of the two.

The card can be a program because each response is a finite instruction, and
waiting for and branching on the predictor's answer are computable operations.
Here is how to include the input explicitly. Finite tape programs can be written
as finite binary words, and a program can interpret such a word by looking up
the next transition and maintaining the simulated tape.
Invalid descriptions may be assigned one fixed halting program. Suppose a terminating
algorithm $H(p,x)$ decided whether the program described by $p$ halts on $x$.
Build a program $D$ which, on input $p$, computes $H(p,p)$, loops forever if the answer
is yes, and halts if it is no. Let $d$ describe $D$. Either answer to $H(d,d)$ contradicts
what $D(d)$ does. So no such deciding algorithm exists.

Translate a tape program and its input by the stack construction and then prime packing.
Both translations are finite algorithms, and the translated two-counter run halts exactly
when the tape run halts. A decider for two-counter halting would then give $H$.
The same reduction, followed by Theorem~\ref{thm:u34-halting}, applies to the small machine.
Halting instances can be enumerated by simulating every description and input for successively
larger finite budgets. If the nonhalting instances could also be enumerated, interleaving
those two enumerations would decide each instance. This proves the impossibility underlying
the reduction formulation of Theorem~\ref{thm:u34-undecidable}.

\subsection{Seven descriptions of the same computation}
\label{sub:seven-constructions}

Here are the translations behind Theorem~\ref{thm:lift-models}. Each one keeps the
partial function, including where it runs forever. The constructions
are the familiar ones behind the results of Minsky, Shepherdson and Sturgis, and
Forster, Larchey-Wendling and their coauthors; the encodings and the reason each
translation works are given here.

\emph{Binary stacks.} A binary stack is a finite word with push and pop operations,
with a distinct empty case. Encode the empty word as 1 and pushing $d\in\{0,1\}$
as $n\mapsto2n+d$. Division by two pops it, with 1 tested before division.
Finite counters implement these operations. Conversely two stacks hold a tape as
above; encode every symbol of a larger alphabet by a fixed-length binary block.
A push or pop of that block takes a fixed finite number of stack operations, and the
control stores a partially read block. A tape simulates counters, so the translations
connect binary stack machines and counter machines in both directions.

\emph{The two decrement conventions.} Some counter machines jump on zero and
fall through after a positive decrement; the convention of this book does the reverse.
To translate a conditional decrement, send its two outcomes to separate short blocks
that unconditionally jump to the intended destinations. In this book's convention an unconditional jump is the restoring INC/DEC pair above. Under jump-on-zero that pair falls through, because the DEC finds the counter positive; there the jump is a DEC on one extra counter that nothing ever increments, so it is always zero. So either convention implements the other's instruction, with the same final values on the original counters and one finite stretch per step, and the translation into jump-on-zero spends one counter more. Relabel each target
to the beginning of its translated block. This proves the equivalence of the two Minsky
instruction conventions without identifying their individual steps.

\emph{FRACTRAN.} A program is an ordered finite list of positive rational numbers.
At a positive integer $N$, multiply by the first fraction whose product with $N$ is
an integer; halt if none applies. Counter machines simulate this by fixed-denominator
division, multiplication and an ordered scan of the finite list.
For the reverse translation, assign distinct primes $p_r$ to counters and $q_i$ to
control states. Represent a counter configuration by
$q_i\prod_r p_r^{n_r}$. Every translated instruction uses a fresh intermediate state
prime, so its state marker never cancels out of a fraction.
For INC $r$ going to $j$, use
$t_i p_r/q_i$, followed by the transfer $q_j/t_i$.
Try it with a counter represented by factors of 2, a starting line marked by 5,
a destination marked by 7 and an intermediate marker 11. Ten beans, $5\cdot2$,
mean line 5 with counter value 1. The fraction $22/5$ takes ten to 44; $7/11$
takes 44 to 28. Divide the result by its line marker 7: four remains, two factors
of 2. The line has changed and the counter has risen by one. The intermediate 11
keeps the two parts of the instruction separate.
For DEC $r$ with positive target $j$ and zero target $k$, use, in this priority order,
\[
\frac{t_i}{q_i p_r},\qquad \frac{u_i}{q_i},
\]
with transfers $q_j/t_i$ and $q_k/u_i$. All state and intermediate primes are distinct
from the counter primes. Exactly one state marker is present. Fractions belonging
to other states cannot apply. At state $i$, the first decrement fraction applies exactly
when $n_r>0$; otherwise the second applies. The next fraction transfers the intermediate
marker to the destination marker. A halting state has no outgoing fraction. Each
counter step then takes two fraction steps, and the prime exponents preserve all
counter values. Read the output as the exponent of the output counter's prime.

\emph{Partial recursive functions.} Compare two orders. ``Try the first ten tickets''
gives you a stopping point before you begin. ``Try tickets in order until one
passes'' does not. If each test finishes, the second order can still keep you
working forever because every ticket fails. If one test itself never finishes,
you never reach the next ticket. Those are the differences between a bounded
loop and the search used here.
Start with zero, successor and projections;
close under composition, primitive recursion and minimization. In minimization,
search $0,1,2,\ldots$ for the first argument where the tested function returns zero;
every earlier test must terminate with a positive result. A counter program implements
the initial functions directly. Disjoint scratch counters preserve arguments across
composition. For primitive recursion it keeps the recursion index and accumulated result,
increments the index, and runs the step program until it reaches the supplied bound.
For minimization it runs the tests successively; an undefined test leaves the run there
forever, as the definition requires. Induction on these constructions supplies a finite
counter program for every partial recursive function.

Conversely encode a counter configuration by nested pairs
$\langle m,n\rangle=(2n+1)2^m$, with a separate finite code for the line.
Addition, multiplication, powers, bounded division and bounded search are primitive
recursive: their loops have a supplied natural bound and can be written as recursion
on that bound. Pair decoding is bounded by the encoded positive number. The next
configuration and the halt test are then primitive recursive for any fixed finite
program. The configuration after $t$ steps is primitive recursion on $t$, keeping halted
configurations fixed. Minimize the total halt test over $t$, then read the output counter.
This is the program's partial function, and it is undefined exactly when the program
never halts.

\emph{The language $L$.} Fold two instruction cards shut. One says ``put three beans
in the output cup and stop''; the other says ``keep adding beans forever.'' A zero
input is to choose the first card. If your rule makes you carry out both cards
before choosing, you never deliver the three beans. You have to keep the
instructions folded until the choice is made. An abstraction does that job in
weak call-by-value: its body waits until the abstraction is applied.
The Scott numeral definitions and the delayed recursion
construction in Appendix~\ref{sec:quoting-construction} give zero, successor, predecessor,
case distinction and guarded recursion as terms. A case on a numeral $n$ can be delayed
explicitly as
\[
n\,(\lambda d.u)\,(\lambda p.\lambda d.(v\,p))\,(\lambda a.a):
\]
at zero it evaluates $u$, and at a successor it evaluates $v$ on the predecessor.
Both branches are values until selected, so call-by-value does not evaluate the unused
branch. Composition is application; primitive recursion repeats a step on the predecessor;
minimization is a guarded search that evaluates each test before moving to its successor.
Build the term in the same order as the function: begin with its initial
functions, then use the corresponding term operation each time composition,
recursion or minimization is used. Each operation consumes the results of the
already built pieces exactly as the function definition requires. On numerical
inputs, a computation that finishes therefore returns the Scott numeral of its
answer. If a required subcomputation never finishes, evaluation waits for it;
if minimization never finds a successful test, its search keeps going. In either
case there is no resulting value.

In the other direction, encode a finite term as a finite string or as the nested-pair
code of its syntax tree. Reading the outer constructor, finding a variable index,
renaming bound indices and substituting a finite term are terminating traversals of
finite trees. A tape algorithm performs them with an explicit stack of unvisited children.
For weak call-by-value evaluation it evaluates an application's function and argument,
then substitutes the argument value in the abstraction's body; it never reduces inside
an unapplied abstraction. Each operation is a finite tape computation. Induction on
a terminating evaluation gives a terminating tape run with its value; reading a terminating
run backwards gives the evaluation tree. Scott numerals are decoded by counting successor
constructors. So an $L$ computation of a numerical relation is a tape computation of it.

The tape and counter translations, the two counter conventions, binary stacks, FRACTRAN,
partial recursive functions and $L$ now lie in one connected family of translations.
Composing the translations preserves their input/output relation. None of them keeps time,
space, or a ledger of what was paid just because it keeps a partial function. Where I need a ledger
kept, I build that in on purpose, as the universal program of
Appendix~\ref{sec:universal-blocks} does; time and space I don't claim at all.

\subsection{Searching for a finite computation}

Give yourself one token and run one instruction of the program you want to test.
If it has not stopped, reset its starting arrangement, give yourself two tokens
and try again. Then three. A program that stops after seven instructions is found
on the seven-token attempt. Seeing six attempts fail says only that it has not
stopped within six instructions; it does not tell you it runs forever.

The evaluators used in Appendix~\ref{sec:small-diagonal} work by that search. A program
and a natural budget determine a finite simulation. Enumerate budgets
$0,1,2,\ldots$, and stop at the first one that exhibits a halting run, returning its output.
If no such run exists, the search continues forever. Encoding a finite program, decoding
an instruction, and performing one step use the finite operations above, so this search
is itself a partial recursive function and has a finite counter program.
The host can execute that program instruction for instruction, using as many registers
as the construction requires. Its computation uses only INC, DEC and HALT, so it adds nothing to the ledger and leaves the fact table, channel and flag alone. This supplies the evaluators and
program transformations used in the host's s-m-n and fixed-point arguments.


\section{The universal program in blocks}
\label{sec:universal-blocks}

Put the guest's instructions on one sheet and the host's on another, with a finger
marking the current line on each. The guest takes one instruction. The host takes
as many of its smaller instructions as it needs, then returns to the head of its
sheet. Compare the two arrangements there: counts, live facts, commitment, flag
and ledger. The host is allowed to do more work between comparisons. It is not
allowed to lose a fact, revive a dead one, or charge another unit for its bookkeeping.
The lists below let you carry out that comparison instruction by instruction.

Here is the finite list of instructions that is $U$. A block is a list, not an extra
instruction: replace every block name by its list, place the lists consecutively, and
resolve each label to its instruction number. The first instruction is at address 1.
Within a block, $o$ is its first address and $o+k$ its offset $k$. Unless a branch says
otherwise, execution continues into the next block. The only final instructions are
INC, DEC, HALT, CHECK, COMMIT and CERTIFY, with the meanings and prices of
Definition~\ref{def:u34-step} on the host's registers (Definition~\ref{def:u34-multi}).

\subsection{Arithmetic blocks}

Write $J(a,p)$ for the two-instruction list INC $a$; DEC $a$ $p$.
It jumps to $p$ and restores the value of $a$.
The following lists specify the other arithmetic blocks. Semicolons concatenate lists;
a named block in a list is placed at the next available address. A target such as
$o+3$ refers to the enclosing block's start, not the nested block's start.

\begingroup\small
\begin{longtable}{@{}>{\raggedright\arraybackslash}p{0.19\linewidth}>{\raggedright\arraybackslash}p{0.72\linewidth}@{}}
\toprule
Block & Expansion at address $o$ \\
\midrule\endhead
$Z(x)$ & DEC $x$ $o$ \\
$J_0(x,p)$ & DEC $x$ $(o+3)$; $J(x,p)$; INC $x$ \\
$M(x,y)$ & $J(x,o+3)$; INC $y$; DEC $x$ $(o+2)$ \\
$M_2(x,y,z)$ & $J(x,o+4)$; INC $y$; INC $z$; DEC $x$ $(o+2)$ \\
$C(x,y,t)$ & $Z(y)$; $Z(t)$; $M_2(x,y,t)$; $M(t,x)$ \\
$H(a,x,p)$ & $Z(a)$; $M(x,a)$; $J(a,o+8)$; INC $x$; DEC $a$ $(o+11)$; $J(a,p)$; DEC $a$ $(o+7)$ \\
$P(a,x,y)$ & $Z(a)$; $M_2(x,a,a)$; INC $a$; $M(a,x)$; $J(a,o+22)$; $M_2(x,a,a)$; $M(a,x)$; DEC $y$ $(o+13)$ \\
$N(a,x,y)$ & $J_0(x,0)$; $Z(y)$; $J(a,o+8)$; INC $y$; $H(a,x,o+7)$ \\
\bottomrule
\end{longtable}
\endgroup

Their lengths, in the order $Z,J_0,M,M_2,C,H,P,N$, are
$1,4,4,5,11,12,23,20$. The source and destination of a move are distinct;
$M_2$ may have equal destinations, as in $P$. In $C$ the three registers are distinct;
in $H$, $a\ne x$; in $P$ and $N$, all three are distinct.

Each claim about these blocks follows from a decreasing counter. $Z$ decrements to
zero. $M$ adds the old source to the destination and empties the source; $M_2$ does it
to both destinations. So $C$ leaves $x$ unchanged, writes its old value to
$y$, and leaves $t=0$. $J_0$ tests zero without changing the value: a positive
counter is decremented and then restored, while zero takes the jump.

In $H$, first move the old $x$ into $a$. Each complete pass removes two from $a$ and
adds one to $x$. So it leaves $a=0$ and $x=\lfloor x_{\rm old}/2\rfloor$;
an even old value jumps to $p$, an odd one falls through. In $P$, the first four
blocks make $x=2x_{\rm old}+1$ and $a=0$. Each subsequent pass doubles $x$ while
decrementing $y$, so it ends at $x=(2x_{\rm old}+1)2^{y_{\rm old}}$, with $a=y=0$.
In $N$, zero jumps to address 0, which holds no instruction, so the run stops. A positive number is repeatedly halved
until odd, counting the even divisions in $y$. So
$N$ sends $x=\langle m,n\rangle$ to $x=n$, $y=m$, $a=0$.
These invariants prove both the returned values and termination on the stated inputs.

\subsection{Selection, comparison and fetching}

The selector $D(x,t;p_0,\ldots,p_{k-1})$ has length $3k$.
For $i=0,\ldots,k-1$, place at offset $3i$ the three instructions
DEC $x$ $(o+3i+3)$; $J(t,p_i)$.
If the entering value is $i<k$, it takes target $p_i$ with $x=0$.
If it is at least $k$, it falls through with $x$ reduced by $k$.
This follows by counting the decrements before the first zero.

The equality block $E(x,y;t_1,t_2,u;p)$ first makes copies
$C(x,t_1,u)$ and $C(y,t_2,u)$, occupying offsets 0 to 21. Its remaining list is
\[
\begin{array}{r l}
22:&\mathrm{DEC}\ t_1\ (o+29),\\
23:&J_0(t_2,p),\\
27:&J(u,o+32),\\
29:&\mathrm{DEC}\ t_2\ (o+22),\\
30:&J(u,o+32),\\
32:&Z(t_1);\ Z(t_2).
\end{array}
\]
The listed offsets count expanded instructions, so the block has length 34.
The temporary registers are distinct from each other and from $x,y$.
The two copies count down together. If they reach zero together, the block jumps to
$p$; if one reaches zero first, it clears both and falls through. So it preserves
$x,y$, leaves its temporary counters zero at either exit, and branches exactly on equality.

For fetching, define $F(q,l;w,k,h,a,t;p_0,p_e)$ by the list
\[
\begin{array}{r l}
0:&C(q,w,t),\\
11:&C(l,k,t),\\
22:&J_0(k,p_0),\\
26:&\mathrm{DEC}\ k\ (o+27),\\
27:&J_0(w,p_e),\\
31:&N(a,w,h),\\
51:&J_0(k,o+56),\\
55:&\mathrm{DEC}\ k\ (o+27).
\end{array}
\]
Its length is 56. Register $q$ holds a list encoded by 0 for empty and
$\langle\text{head},\text{tail}\rangle$ for a nonempty list, and $l$ is a line number.
The scratch registers are distinct and separate from $q,l$. A zero line takes $p_0$.
Otherwise decrementing the line copy once makes $k=l-1$. Each pass removes the next
list head into $h$ and its tail into $w$. If $k=0$, that head is the selected
instruction and the block exits; otherwise it decrements $k$ and repeats.
An exhausted list takes $p_e$. The original program and line values are restored by
the copies. Every call terminates: each complete pass decreases $k$, or exits.

\subsection{Registers and instruction codes}

Put $T_i=i+4$ for $0\le i<10$. For bank $c\in\{0,1\}$ and slot $0\le k<16$, put
\[
R_c=c,\quad Q_c=14+c,\quad V_{c,k}=16+16c+k,\quad S_{c,k}=48+16c+k.
\]
Register 2 is the program, register 3 its line, and register 80 is the dead end.
$V$ are the mirrors, $S$ the slots, and $Q$ the counts of successful checks in each bank.
Property codes are $0$ for zero, $1$ for even, and $n+2$ for at least $n$.
The code of an instruction is the pair of its opcode and operand:
INC has opcode 0 and operand $c$; DEC has opcode 1 and operand $\langle c,j\rangle$;
HALT has opcode 2 and operand 0; CHECK has opcode 3 and operand $\langle c,p\rangle$;
COMMIT has opcode 4 and that same operand form; CERTIFY has opcode 5 and operand 0.
A program is their list code. These rules specify the loaded number in register 2.

A bump of slot $S$ and mirror $V$ is INC $S$; DEC $S$ to the following instruction;
$Z(V)$. It has length 3. It leaves the slot's value alone, increases its version by
two, and empties the mirror. Let $B_c$ concatenate the sixteen bumps of bank $c$,
in increasing slot order; its length is 48.

\subsection{The fifteen blocks of \texorpdfstring{$U$}{U}}

The following table fixes the order, start address and expanded length of every major
block. Labels in the definitions below mean these addresses; $H$ without arguments denotes the head at 1
and $K$ the halt block at 98.

\begingroup\small
\begin{longtable}{@{}>{\raggedright\arraybackslash}p{0.46\linewidth}rr@{}}
\toprule
Block & Start & Length \\
\midrule\endhead
Head & 1 & 97 \\
Halt & 98 & 11 \\
INC bank selection & 109 & 8 \\
INC bank 0 & 117 & 52 \\
INC bank 1 & 169 & 52 \\
DEC bank selection & 221 & 28 \\
DEC bank 0 & 249 & 60 \\
DEC bank 1 & 309 & 60 \\
CHECK bank selection & 369 & 28 \\
CHECK bank 0 & 397 & 1022 \\
CHECK bank 1 & 1419 & 1022 \\
COMMIT bank selection & 2441 & 28 \\
COMMIT bank 0 & 2469 & 628 \\
COMMIT bank 1 & 3097 & 628 \\
CERTIFY & 3725 & 4 \\
\bottomrule
\end{longtable}
\endgroup

The head is
$F(2,3;T_0,T_1,T_2,T_3,T_4;K,K)$;
$Z(T_0)$; $N(T_3,T_2,T_5)$;
$D(T_5,T_4;109,\allowbreak 221,\allowbreak 98,\allowbreak 369,\allowbreak 2441,\allowbreak 3725)$; $J(T_4,K)$.
Fetching supplies an instruction code in $T_2$. Unpacking leaves its operand in $T_2$
and its opcode in $T_5$. The six-way selector chooses its handler. The final jump
handles an invalid opcode, though a loaded program has none.
The halt block is $Z(T_0);\cdots;Z(T_9)$ followed by HALT.

INC's bank selector is $D(T_2,T_4;117,169);J(T_4,K)$.
Each other bank selector first runs $N(T_3,T_2,T_6)$, leaving the property code or
jump target in $T_2$ and the bank in $T_6$, then selects its two bank blocks with
$D(T_6,T_4;\text{bank 0},\text{bank 1});J(T_4,K)$.

The INC block for bank $c$ is INC $R_c$; $B_c$; INC 3; $J(T_4,H)$.
The DEC block, starting at $o$, is
DEC $R_c$ $(o+5)$; INC 3; $Z(T_2)$; $J(T_4,H)$;
$B_c$; $Z(3)$; $M(T_2,3)$; $J(T_4,H)$.
A zero counter follows the first return, advancing the guest line without bumping
slots. A positive counter jumps over it, bumps the bank and installs the target line.

A CHECK slot block for $(c,k)$ is the following 60 instructions after expansion:
\[
\begin{aligned}
&C(R_c,T_8,T_4);\ C(T_2,T_9,T_4);\ P(T_3,T_8,T_9);\ M(T_8,S_{c,k});\\
&\mathrm{CHECK}\ \mathsf{Slot}\ S_{c,k};\ Z(V_{c,k});\ M(T_2,V_{c,k});\\
&\mathrm{INC}\ V_{c,k};\ \mathrm{INC}\ Q_c;\ \mathrm{INC}\ 3;\ J(T_4,H).
\end{aligned}
\]
The CHECK bank block at $o$ is $C(Q_c,T_7,T_4)$, then a selector on $T_7$ whose
sixteen targets are $o+62+60k$, then CHECK $\mathsf{Slot}$ 80 and $J(T_4,K)$,
then the sixteen slot blocks. Its length is
$11+48+1+2+16\cdot60=1022$.
The selected unused slot starts at zero, so the move into it installs exactly the
pair of property code and guest value. The host's own table has capacity sixteen:
if it is full, its CHECK traps even when this bank's count is below sixteen.
So the capacity is shared between the banks, just as it is in the guest. A count of sixteen in one bank already fills the table, so the dead-end check and the full-table trap are the same case.

A COMMIT slot block is COMMIT $\mathsf{Slot}$ $S_{c,k}$;
$Z(T_2)$; INC 3; $J(T_4,H)$, of length 5.
The COMMIT bank block at $o$ first increments $T_2$, making the wanted mirror value
one more than the property code. It then runs, for $k=0,\ldots,15$,
$E(V_{c,k},T_2;T_7,T_8,T_4;o+548+5k)$.
After those comparisons come COMMIT $\mathsf{Slot}$ 80 and $J(T_4,K)$, followed
by the sixteen slot blocks. Its length is $1+16\cdot34+1+2+16\cdot5=628$.
The comparison reads the mirror, leaving the slot untouched. Its first match
commits that slot's current fact; absence of a match reaches the dead end and traps.
The CERTIFY block is CERTIFY; INC 3; $J(T_4,H)$, of length 4.

These expansions determine every instruction and target. Adding the fifteen lengths
gives 3728. There are sixteen ordinary CHECK sites and one dead-end CHECK in each
bank, sixteen ordinary COMMIT sites and one dead-end COMMIT in each bank, and one
CERTIFY site: $2\cdot17+2\cdot17+1=69$ paid sites. Every other instruction costs nothing.

\subsection{Why the blocks simulate a guest}

Run a case that catches more than a wrong sum. Start the guest with four beans
in $A$ and run CHECK Even $A$. Put its fact on a card saying ``even, version 0.''
The host copies the value into an unused slot as the pair (Even, 4), whose code
is $(2\cdot4+1)2^1=18$, and checks that slot. Put one unit on each ledger.

Now let the guest add one bean and take it away. Its pile is four again, but its
version is 2. The old fact still says version 0. On the host, each of those guest
writes bumps the slot by one and back: 18 to 19 to 18. The slot's value ends where
it began, while its version rises and its mirror is cleared. Try COMMIT Even $A$
using the old check. Both machines must refuse, even though four is even. Both
pay one more unit for that failed attempt. A translation that merely compares
the piles at the end would miss the failure this experiment is designed to catch.

Start again and leave out the two bean moves. CHECK, COMMIT, CERTIFY must succeed
on both sheets, raising both flags for exactly three units. The longer list of host instructions adds nothing to the ledger for copying or looking up a line. These
two runs test the reason for keeping both values and versions. The argument for
all runs has to keep the following correspondence at every return to the head.

At the head, registers 0, 1 and 3 give the guest counters and line. Each successful
guest CHECK has one host slot fact holding its property and checked value. Its mirror
is the property code plus one while the guest fact is current, and zero after the
guest counter has been written. The matching host slot fact is current in exactly
the first case. Count registers say how many slots have been used in their bank;
unused slots and mirrors are zero. Tables have corresponding entries in the same
order, and channels name corresponding entries. The ledgers and flags agree.
Loading a guest establishes every part of this relation.

Fetching and dispatching terminate and preserve that relation's stored values.
For INC or a positive DEC, the guest writes its counter; the selected bank's bumps
invalidate exactly its former slot facts and clear their mirrors. A zero DEC writes
nothing and bumps nothing. These cases add nothing to either ledger.
For CHECK, an unused slot receives the exact pair tested by the guest. Both checkers
answer the same, and the tables have the same length. A passing check appends the
corresponding fact and updates its mirror and count. A false check or a full table
traps both machines, after the same single paid act. For COMMIT, a nonzero mirror
matches exactly a current guest fact with the requested property and counter. Its
host slot has not been read or changed, so the corresponding COMMIT succeeds.
If none matches, both machines trap. Again exactly one unit is paid. For CERTIFY,
channels are either both empty or both hold corresponding commitments, so the
instruction either traps both or raises both flags, for one unit.

All arithmetic between these paid acts terminates by the block invariants. A normal
case returns to the head; a halting or trapped case stops. No arithmetic block has a
paid instruction. This proves the one-step correspondence and bounds every
intermediate host ledger between the guest's old and new ledgers. Induction proves
the run correspondence of Theorems~\ref{thm:u34-U-sim} and~\ref{thm:u34-U-flag}.
It also proves the reverse halting implication: an infinite guest run entails
infinitely many nonempty host stretches, and a running guest instruction cannot
produce a stopped host except in its own halting or trapping case.

\subsection{Adding PAY}

For $U_P$, use opcode 6 with operand 0 for PAY, append a PAY handler to the six handlers,
and append its block after CERTIFY. The block is PAY; INC 3; $J(T_4,H)$, of length 4.
The head's selector now has seven entries and is three instructions longer. Every
block after the head then starts three addresses later. Resolve targets to those
new addresses; the head itself remains at address 1.
The new PAY block starts at 3732. The total length is $3728+3+4=3735$ and there
are $69+1=70$ paid sites. The interpretation of $\mathsf{Slot}$ uses the priced
host's property language from Appendix~\ref{sec:presented-universal}; the arithmetic
blocks and their invariant do not depend on the particular meaning of a property code.
A PAY guest step is one PAY host step surrounded by arithmetic that costs nothing, so the
same ledger argument applies.


\section{The CHSH derivation, in full}

Appendix~\ref{sec:chsh} tells the story; this appendix does the algebra. Each item gives
the statement with its hypotheses and the steps you need to check it. Nothing here says
that a matrix is a physical state. Where quantum strategies come in, I say exactly how they
stand to the matrices.

\subsection{Setting and notation}


Start with four piles of tally sheets, one for each pair of settings Alice
and Bob can be given. On each round their answers either agree or disagree.
Give agreement a score of $+1$ and disagreement a score of $-1$. In each pile,
add the scores and divide by the number of rounds. That pile's average is its
\emph{correlator}. It says how strongly the two answers agree for that pair of
settings.

For example, three agreements and one disagreement give a total of two over
four rounds, hence an average of one half. The four averages make a little
square table: Alice's setting chooses the row, Bob's chooses the column.
The score I will study adds three cells of that table and subtracts the
bottom-right one. Here are the names for those quantities.

\begin{definition}[Settings, outcomes, correlators]
Alice receives a setting $x \in \{0,1\}$ and Bob a setting $y \in \{0,1\}$. Each answers a
bit, recoded as an outcome $a, b \in \{+1, -1\}$ (bit 0 becomes $+1$, bit 1 becomes $-1$).
The correlator $E_{xy}$ is the average of the product $ab$ over the rounds with settings
$(x, y)$. The CHSH quantity is
\[
S = E_{00} + E_{01} + E_{10} - E_{11}.
\]
I collect the four correlators into the $2 \times 2$ matrix
\[
E = \begin{pmatrix} E_{00} & E_{01} \\ E_{10} & E_{11} \end{pmatrix},
\]
whose rows are indexed by Alice's setting and whose columns by Bob's. The rows of $E$ are
$(E_{00}, E_{01})$ and $(E_{10}, E_{11})$; the columns are $(E_{00}, E_{10})$ and
$(E_{01}, E_{11})$.
\end{definition}

\begin{definition}[Counters]
There are eight counts: for each setting pair, $\mathrm{same}_{xy}$ counts rounds
whose two answer bits matched and $\mathrm{diff}_{xy}$ counts rounds whose bits differed. A
round adds one to exactly one of the eight. Write
\[
N_{xy} = \mathrm{same}_{xy} + \mathrm{diff}_{xy}, \qquad
D_{xy} = \mathrm{same}_{xy} - \mathrm{diff}_{xy}, \qquad
E_{xy} = \frac{D_{xy}}{N_{xy}}.
\]
\end{definition}

A matched round contributes $ab = +1$ and a differing round $ab = -1$, so $D_{xy}$ is the sum
of the products and $D_{xy}/N_{xy}$ is their average. Since $|D_{xy}| \le N_{xy}$, every
correlator from counters lies in $[-1,1]$. When a bucket is empty, its correlator is defined
as 0. The score is $E_{00}+E_{01}+E_{10}-E_{11}$.

\subsection{The classical bound}


Write a fixed plan on a card: Alice's answer to setting zero, her answer to
setting one, and Bob's two answers. Each answer is a sign, plus or minus.
Before multiplying anything, look at Bob's pair. Either the signs agree or
they differ. If they agree, adding them gives plus or minus two and subtracting
them gives zero. If they differ, the sum is zero and the difference is plus
or minus two. One of those two quantities vanishes; the other cannot exceed
two in size. The next calculation uses exactly that split.

\begin{lemma}[Deterministic strategies]\label{lem:chsh-deterministic}
Let $a_0, a_1, b_0, b_1 \in \{+1, -1\}$. Then $a_0 b_0 + a_0 b_1 + a_1 b_0 - a_1 b_1$ is $2$ or $-2$.
\end{lemma}

\begin{proof}
Factor: $a_0 b_0 + a_0 b_1 + a_1 b_0 - a_1 b_1 = a_0(b_0 + b_1) + a_1(b_0 - b_1)$. If
$b_0 = b_1$ then $b_0 + b_1 = \pm 2$ and $b_0 - b_1 = 0$; if $b_0 \ne b_1$ the other way
round. In both cases exactly one of the two brackets is $\pm 2$ and the other is $0$, so the
expression is $\pm 2$ times a sign, which is $+2$ or $-2$. There are sixteen sign choices,
and every one of them gives exactly $\pm 2$. None gives $0$ or $\pm 4$.
\end{proof}

You can also just list the sixteen sign choices and add each one up. Nothing in the argument needs real
numbers: the signs are whole numbers, so the bound holds the same way whether you read the
correlators as rationals (what a counter gives) or as reals.

A fixed card is only the first case. You can shuffle among several cards,
perhaps using some more often than others. Every card still tells each player
an answer using that player's setting alone. The shuffle changes how often
each card contributes; it does not change the score that a particular card
would produce. Since each card's score lies between minus two and two, their
average must stay between those same ends. The following definition writes
out that weighted shuffle.

\begin{definition}[Finite factorizable model]
Four real correlators $E_{00}, E_{01}, E_{10}, E_{11}$ are factorizable when there are a
natural number $\lambda_{\max}$, weights $p(\lambda)$ and response signs $A(x, \lambda)$,
$B(y, \lambda)$ for $\lambda = 0, \ldots, \lambda_{\max}$ such that:
$0 \le p(\lambda) \le 1$;
$\sum_{\lambda = 0}^{\lambda_{\max}} p(\lambda) = 1$; every $A(x, \lambda)$ and every
$B(y, \lambda)$ is $-1$ or $+1$; and for all $x, y$,
$E_{xy} = \sum_{\lambda=0}^{\lambda_{\max}} p(\lambda) A(x, \lambda) B(y, \lambda)$.
In words: a hidden card $\lambda$ is drawn with probability $p(\lambda)$, and each player reads
an answer off the card using only their own setting. Only the correlators are modelled.
\end{definition}

\begin{theorem}[The classical bound for a finite local model]
If $E$ is factorizable, then $-2 \le E_{00} + E_{01} + E_{10} - E_{11} \le 2$.
\end{theorem}

\begin{proof}
Substitute the factorization: $S = \sum_\lambda p(\lambda) S_\lambda$, where $S_\lambda$ is
the deterministic expression of Lemma~\ref{lem:chsh-deterministic} at the signs
$A(\cdot, \lambda), B(\cdot, \lambda)$. Each $S_\lambda \in [-2, 2]$, the weights are
nonnegative and sum to 1, so $-2 = \sum_\lambda p(\lambda)(-2) \le S \le
\sum_\lambda p(\lambda) \cdot 2 = 2$.
\end{proof}

\begin{remark}[Scope of the classical bound]\label{rem:chsh-fine-scope}
Fine's 1982 theorem is an equivalence between several descriptions of local models and the
Bell inequalities. The preceding theorem proves one direction: a finite factorizable model
obeys the bound. The converse of this single inequality is false: $(1, 1, -1, 1)$ has $S = 0$,
but it breaks another of the sign variants ($E_{00} + E_{01} - E_{10} + E_{11} = 4$), so no
local model gives it. The correlation-form equivalence needs each correlator between $-1$ and $1$
and all eight inequalities; with those, Section~\ref{sub:finite-local-construction} builds the
local model. The preceding theorem is stated for finitely many hidden cards ($\lambda\le\lambda_{\max}$), and that loses nothing. Whatever the hidden variable is and however it is drawn, the correlators see it only through its four answer signs, so lumping together the cards that carry the same four signs gives at most sixteen cards and the same four correlators. All the lumping needs is a probability for each of the sixteen sign patterns. \end{remark}

The finite-tally statement concerns one fixed strategy replayed on every recorded round.

\begin{definition}[A tally consistent with one local plan]
For bits $a_0, a_1, b_0, b_1 \in \{0, 1\}$ and a counter record $w$, $w$ is locally consistent with
$(a_0, a_1, b_0, b_1)$ when, for each setting pair $(x, y)$: if $a_x = b_y$ then
$\mathrm{diff}_{xy} = 0$, otherwise $\mathrm{same}_{xy} = 0$; and every bucket is sampled,
$N_{00}, N_{01}, N_{10}, N_{11} > 0$.
\end{definition}

\begin{theorem}[A tally above 2 is not one replayed strategy]
Let $w$ be a counter record whose rational score $S(w)$ is greater than 2.
Then for all bits $a_0, a_1, b_0, b_1$, $w$ is not locally consistent with
$(a_0, a_1, b_0, b_1)$. Every locally consistent record has $S(w) = 2$ or $S(w) = -2$; the first statement is that fact turned around.
\end{theorem}

\begin{proof}
Under local consistency each sampled bucket is pure: if $a_x = b_y$ every round in the
bucket matched, so $E_{xy} = N_{xy}/N_{xy} = +1$; otherwise every round differed and
$E_{xy} = -1$. In signs, $E_{xy} = \alpha_x \beta_y$ with $\alpha_x = (-1)^{a_x}$ and
$\beta_y = (-1)^{b_y}$. The score is then the deterministic expression of
Lemma~\ref{lem:chsh-deterministic}, which is 2 or $-2$. A score above 2 contradicts this.
\end{proof}

The theorem rules out one explanation, a single four-bit plan used on every round. It says
nothing about a local model that redraws its hidden state every round; Appendix~\ref{sec:chsh} gives a four-round coin-flip tally with score 4. The clause that every bucket is sampled does work here. An empty bucket reads as 0, and the plan that answers 0 to everything, never asked the pair 11, would score $1 + 1 + 1 - 0 = 3$.

\subsection{Tsirelson's bound from algebra}


Now put the cards aside and look just at four numbers in the square. I want
to bound their three-plus-one-minus score from their squared sizes. The useful
trick is to write the proposed allowance minus the squared score as a pile
of squares. A square cannot be negative. If the difference really is that
pile, the score cannot outrun the allowance.

The identity below is the receipt for that trick. The letters $a,b,c,d$ are
the four entries read across the first row, then across the second. It holds
for any four real numbers. After proving it, I will put in the two bounds
on row length that this part of the checker supplies.

\begin{lemma}[The six-square identity]
For all real $a, b, c, d$,
\[
\begin{aligned}
&4(a^2 + b^2 + c^2 + d^2) - (a + b + c - d)^2 \\
&\quad = (a-b)^2 + (a-c)^2 + (a+d)^2 + (b-c)^2 + (b+d)^2 + (c+d)^2.
\end{aligned}
\]

\end{lemma}

\begin{proof}
Expand the square on the left:
$(a+b+c-d)^2 = a^2 + b^2 + c^2 + d^2 + 2ab + 2ac - 2ad + 2bc - 2bd - 2cd$. So the left side is
\[
3(a^2 + b^2 + c^2 + d^2) - 2ab - 2ac + 2ad - 2bc + 2bd + 2cd.
\]
On the right, each variable appears in exactly three of the six brackets ($a$ in
$a-b$, $a-c$, $a+d$; $b$ in $a-b$, $b-c$, $b+d$; $c$ in $a-c$, $b-c$, $c+d$; $d$ in $a+d$,
$b+d$, $c+d$), so the squared terms add to $3(a^2+b^2+c^2+d^2)$. The cross terms are
$-2ab$ from $(a-b)^2$, $-2ac$ from $(a-c)^2$, $+2ad$ from $(a+d)^2$, $-2bc$ from $(b-c)^2$,
$+2bd$ from $(b+d)^2$ and $+2cd$ from $(c+d)^2$. Both sides agree term by term.
\end{proof}

Since the right side is a sum of squares, $S^2 \le 4 \sum_{xy} E_{xy}^2$ for any four reals.

\begin{theorem}[Row-bound form of Tsirelson's bound]
For all real $E_{00}, E_{01}, E_{10}, E_{11}$ with $E_{00}^2 + E_{01}^2 \le 1$ and
$E_{10}^2 + E_{11}^2 \le 1$, $S^2 \le 8$. The absolute form $|S| \le \sqrt 8 = 2\sqrt 2$ follows.
\end{theorem}

\begin{proof}
Adding the two row bounds gives $\sum_{xy} E_{xy}^2 \le 2$. By Lemma~\ref{lem:chsh-sos},
$S^2 \le 4 \sum_{xy} E_{xy}^2 \le 8$.
\end{proof}

\begin{theorem}[Tightness]
The point $E_{00} = E_{01} = E_{10} = 1/\sqrt 2$, $E_{11} = -1/\sqrt 2$ satisfies both row
bounds and has $S = \sqrt 8$.
\end{theorem}

\begin{proof}
Each row has squared length $1/2 + 1/2 = 1$, and $S = 4/\sqrt 2 = 2\sqrt 2 = \sqrt 8$.
\end{proof}

\begin{remark}[Where equality can happen]
Equality $S^2 = 8$ under the row bounds forces
equality in both steps of the proof above: the six squares must all vanish, which gives
$a = b = c = -d$, and the rows must be tight, which gives $2a^2 = 1$. So the point of
Theorem~\ref{thm:chsh-tight} and its negative are the only places the bound is reached. A tally's correlators are
ratios of whole numbers and $2a^2 = 1$ has no such solution, so no tally reaches the bound; the integer
check's passing tallies stay strictly below it and come as near as you like. (Take rational
$a, b$ with $a^2 + b^2 = 1$, from a Pythagorean triple, and the correlators $(a, b, b, -a)$: they pass
the check of Section~\ref{sub:chsh-integer} with no slack, and $S = 2(a + b)$, which gets as close to
$2\sqrt 2$ as you like as $a/b$ approaches 1.)
\end{remark}

\begin{lemma}[A rational constant above the bound]\label{lem:chsh-rational}
$\sqrt 8 < 5657/2000$.
\end{lemma}

\begin{proof}
Both numbers are positive, so it suffices to compare squares. $5657^2 = 32{,}001{,}649$ and
$8 \cdot 2000^2 = 32{,}000{,}000$, so $(5657/2000)^2 = 32{,}001{,}649/4{,}000{,}000 > 8 =
(\sqrt 8)^2$. If $\sqrt 8 \ge 5657/2000$ held, squaring both nonnegative sides would give
$8 \ge (5657/2000)^2$, a contradiction.
\end{proof}

The constant is the four-decimal round-up of $2\sqrt 2 = 2.8284271\ldots$; the gap is about
$7.3 \times 10^{-5}$. It lets rational code test ``above the Tsirelson bound'' without a square
root. It plays no part in the integer check of Section~\ref{sub:chsh-integer}.

\subsection{The zero-marginal \texorpdfstring{$5\times5$}{5 by 5} moment matrix}


The four averages fill only the Alice--Bob part of a larger table. Give that
table five labels: a constant reading 1, Alice's two possible readings, and
Bob's two. Each cell records an average product of its row and column
readings. A reading multiplied by itself is 1, so the diagonal is all ones.
The four averages measured fill the cross-player cells and their reflected
partners.

Other cells are still empty. This particular check fills them with zero:
each player's average answer against the constant reading, and the cells
between that player's two settings. These are choices of the check; the
four cross-player averages alone did not supply them. I will later let the
two within-player cells vary and ask whether any filling works. For now,
write down this one fixed filling and test it.

\begin{definition}[The pinned level-1 matrix]
Take rows and columns $I,A_0,A_1,B_0,B_1$. Put 1 on the diagonal and the four correlators in
the Alice--Bob cells, and set the four one-party entries (each player's own average answer) and
the two within-party entries ($A_0$ against $A_1$, $B_0$ against $B_1$) to zero. The resulting
matrix is:
\end{definition}

\[
M(E) \;=\;
\begin{array}{c@{\quad}c}
 & \begin{array}{ccccc} I & A_0 & A_1 & B_0 & B_1 \end{array} \\
\begin{array}{c} I \\ A_0 \\ A_1 \\ B_0 \\ B_1 \end{array} &
\begin{pmatrix}
1 & 0 & 0 & 0 & 0 \\
0 & 1 & 0 & E_{00} & E_{01} \\
0 & 0 & 1 & E_{10} & E_{11} \\
0 & E_{00} & E_{10} & 1 & 0 \\
0 & E_{01} & E_{11} & 0 & 1
\end{pmatrix}
\end{array}
\]

What do I ask of this table? Choose five numbers, one beside each row and
one beside the matching column. For each cell, multiply its entry by its row's
number and its column's number, then add all twenty-five results. Do this
with any five numbers you like, positive, negative or zero. I require the
total never to be negative. That requirement is called \emph{positive
semidefiniteness}, abbreviated PSD.

Write the five chosen numbers as a column $v$. The superscript $\mathsf T$
turns a column into a row, so ordinary matrix multiplication writes that same
weighted total as $v^{\mathsf T}Mv$. The condition is therefore
$v^{\mathsf T}Mv\ge0$ for every choice of $v$. The table is symmetric too:
reflect it across its diagonal and the entries agree.

The first row and column have only one nonzero entry, the corner 1. Their
contribution is the square of the first chosen number. The remaining four
rows and columns can be tested on their own. Grouping them gives

\[
M(E) \;=\; \begin{pmatrix} 1 & 0 \\ 0 & G \end{pmatrix},
\qquad
G \;=\; \begin{pmatrix} I_2 & E \\ E^{\mathsf T} & I_2 \end{pmatrix}.
\]

Here $I_2$ is the two-by-two identity table: ones on its diagonal and zeros
elsewhere. The superscript on $E$ again swaps rows and columns. The weighted
total for $M(E)$ is the first number's square plus the weighted total for $G$.
If every total for $G$ is nonnegative, adding a square keeps it nonnegative.
If some total for $G$ is negative, choose the first number to be zero and the
same negative total appears for $M(E)$. Thus the two tables pass or fail
together.

The next operation is completing a square, with several numbers at once.
For one number, $u^2+2uw+rw^2$ becomes $(u+w)^2+(r-1)w^2$.
The first square cannot hurt, and choosing $u=-w$ makes it vanish.
So the whole expression stays nonnegative for every choice exactly when
the remaining part does. The matrix formula below does the same job for
two groups of numbers.

Here \emph{positive definite} means that the weighted total for $P$ is
strictly positive for every nonzero choice of weights. In particular, $P$
has an inverse, written $P^{-1}$: applying it undoes multiplication by $P$.
The matrix left after completing the square is called the \emph{Schur
complement}. Its name is less useful here than the operation that produces it.

\begin{lemma}[Schur complement, the standard form]
Let $P$ be symmetric positive definite. The symmetric block matrix
$\begin{pmatrix} P & Q \\ Q^{\mathsf T} & R \end{pmatrix}$ is PSD if and only if
$R - Q^{\mathsf T} P^{-1} Q$ is PSD.
\end{lemma}

\begin{proof}
For $u, w$ of the right sizes,
\[
\begin{pmatrix} u \\ w \end{pmatrix}^{\!\mathsf T}
\begin{pmatrix} P & Q \\ Q^{\mathsf T} & R \end{pmatrix}
\begin{pmatrix} u \\ w \end{pmatrix}
= (u + P^{-1} Q w)^{\mathsf T} P \,(u + P^{-1} Q w)
+ w^{\mathsf T} (R - Q^{\mathsf T} P^{-1} Q)\, w.
\]
The first term is never negative because $P$ is positive definite, and it is zero at
$u = -P^{-1} Q w$. So the form is nonnegative for all $(u, w)$ exactly when the second term
is nonnegative for all $w$.
\end{proof}

The identity itself holds for any invertible symmetric $P$. One direction (block PSD gives the
complement PSD) needs only that: put $u = -P^{-1}Qw$. The other direction is where $P$ has to be
positive definite. For the pinned matrix I don't need the lemma in the abstract: I carry out its one instance with
$P = I_2$, $Q = E$, $R = I_2$, where the complement is
\[
I_2 - E^{\mathsf T} E =
\begin{pmatrix}
1 - E_{00}^2 - E_{10}^2 & -(E_{00} E_{01} + E_{10} E_{11}) \\
-(E_{00} E_{01} + E_{10} E_{11}) & 1 - E_{01}^2 - E_{11}^2
\end{pmatrix}.
\]
The diagonal of $E^{\mathsf T} E$ holds the squared lengths of the columns of $E$, which is
why the conditions below are column conditions.

\begin{definition}[Column contractive]
Write
$p = 1 - E_{00}^2 - E_{10}^2$, $q = 1 - E_{01}^2 - E_{11}^2$ and
$s = E_{00} E_{01} + E_{10} E_{11}$. The correlators are column contractive when
\[
p \ge 0, \qquad q \ge 0, \qquad p\,q - s^2 \ge 0.
\]
\end{definition}

\begin{lemma}[The $2\times2$ test]\label{lem:chsh-psd2}
If $\alpha \ge 0$, $\delta \ge 0$ and
$\alpha\delta - \beta^2 \ge 0$, then $\alpha u^2 + 2\beta u w + \delta w^2 \ge 0$ for all real
$u, w$.
\end{lemma}

\begin{proof}
If $w = 0$ the form is $\alpha u^2 \ge 0$. Otherwise put $t = u/w$, so the form is
$w^2(\alpha t^2 + 2\beta t + \delta)$. If $\alpha = 0$, then $-\beta^2 \ge 0$ forces
$\beta = 0$ and the bracket is $\delta \ge 0$. If $\alpha > 0$,
$\alpha t^2 + 2\beta t + \delta = \alpha (t + \beta/\alpha)^2 + (\alpha\delta - \beta^2)/\alpha$,
a sum of two nonnegative terms.
\end{proof}

\begin{theorem}[Contractive implies PSD]\label{thm:chsh-cc-psd}
If $p \ge 0$, $q \ge 0$ and $pq - s^2 \ge 0$, then $M(E)$ is positive semidefinite.
\end{theorem}

\begin{proof}
Write $v = (v_0, x_0, x_1, y_0, y_1)$. Reading off the entries of $M(E)$,
\[
v^{\mathsf T} M(E)\, v = v_0^2 + x_0^2 + x_1^2 + y_0^2 + y_1^2
+ 2x_0(E_{00} y_0 + E_{01} y_1) + 2x_1(E_{10} y_0 + E_{11} y_1).
\]
Complete the squares in $x_0$ and $x_1$:
\[
\begin{aligned}
v^{\mathsf T} M(E)\, v ={}& v_0^2 + (x_0 + E_{00} y_0 + E_{01} y_1)^2
+ (x_1 + E_{10} y_0 + E_{11} y_1)^2 \\
&+ p\, y_0^2 - 2 s\, y_0 y_1 + q\, y_1^2.
\end{aligned}
\]
To check the last three terms, expand the two squares: they contribute
$(E_{00} y_0 + E_{01} y_1)^2 + (E_{10} y_0 + E_{11} y_1)^2$ beyond the cross terms, which is
$(E_{00}^2 + E_{10}^2) y_0^2 + 2 s\, y_0 y_1 + (E_{01}^2 + E_{11}^2) y_1^2$; subtracting it
from $y_0^2 + y_1^2$ leaves $p\, y_0^2 - 2 s\, y_0 y_1 + q\, y_1^2$. The first three terms are
squares. The last is the $2 \times 2$ form with $\alpha = p$, $\beta = -s$, $\delta = q$,
nonnegative by Lemma~\ref{lem:chsh-psd2}.
\end{proof}

\begin{lemma}[Discriminant]\label{lem:chsh-discriminant}
If $a + 2bt + ct^2 \ge 0$ for every real $t$, then
$b^2 \le ac$.
\end{lemma}

\begin{proof}
If $c < 0$ the quadratic tends to $-\infty$, so $c \ge 0$. If $c = 0$, the line
$a + 2bt$ is nonnegative for every $t$ only if $b = 0$, and then $b^2 = 0 = ac$. If $c > 0$,
evaluate at $t = -b/c$: $a - b^2/c \ge 0$, so $b^2 \le ac$.
\end{proof}

\begin{theorem}[PSD implies contractive]\label{thm:chsh-psd-cc}
If $M(E)$ is positive semidefinite, then $p \ge 0$, $q \ge 0$ and $pq - s^2 \ge 0$.
\end{theorem}

\begin{proof}
Feed the quadratic form three test vectors, written in the order
$(I, A_0, A_1, B_0, B_1)$.
Each is the choice $x = -E y$ that zeroes both completed squares in the proof of
Theorem~\ref{thm:chsh-cc-psd}, so only $y^{\mathsf T}(I_2 - E^{\mathsf T} E) y$ is left.

First, $v = (0, -E_{00}, -E_{10}, 1, 0)$ gives $v^{\mathsf T} M v = p$, so $p \ge 0$.

Second, $v = (0, -E_{01}, -E_{11}, 0, 1)$ gives $v^{\mathsf T} M v = q$, so $q \ge 0$.

Third, for each real $t$, $v(t) = (0, -(E_{00} t + E_{01}), -(E_{10} t + E_{11}), t, 1)$
gives $v(t)^{\mathsf T} M v(t) = p\,t^2 - 2 s\,t + q \ge 0$. Apply Lemma~\ref{lem:chsh-discriminant}
with $a = q$, $b = -s$, $c = p$: $s^2 \le pq$.
\end{proof}

\begin{corollary}[The equivalence]
For all real $E_{00},E_{01},E_{10},E_{11}$, the symmetric matrix $M(E)$ is positive
semidefinite if and only if the correlators are column contractive.
\end{corollary}

\begin{proof}
Theorem~\ref{thm:chsh-cc-psd} is one direction and Theorem~\ref{thm:chsh-psd-cc} the other.
\end{proof}

The equivalence above is about the columns. The bound can go through either the rows or the columns: the six-square identity only needs the four squares to add to at most 2, and the two row bounds give that, and so do $p\ge0$ and $q\ge0$.

\begin{theorem}[PSD implies Tsirelson]
If the pinned matrix $M(E)$ is symmetric and positive semidefinite at
$E_{00},E_{01},E_{10},E_{11}$, then $S^2\le8$ and $|S|\le\sqrt8$.
\end{theorem}

\begin{proof}
Feed the form the vector that is 1 at $A_0$, $-E_{00}$ at $B_0$, $-E_{01}$ at $B_1$ and 0
elsewhere. Reading off the entries, the form is
$1 + E_{00}^2 + E_{01}^2 - 2E_{00}^2 - 2E_{01}^2 = 1 - E_{00}^2 - E_{01}^2$ (the $B_0B_1$ cell is
0), so $E_{00}^2 + E_{01}^2 \le 1$. The vector that is 1 at $A_1$, $-E_{10}$ at $B_0$ and
$-E_{11}$ at $B_1$ gives $E_{10}^2 + E_{11}^2 \le 1$ the same way. Then
Theorem~\ref{thm:chsh-row-tsirelson} applies.
\end{proof}

Here are the three conditions and the score at a few points, all by direct arithmetic.
In this table, $a = 3/5$, $b = 4/5$ and $t = 1/\sqrt 2$.

\begingroup\small
\begin{longtable}{@{}>{\raggedright\arraybackslash}p{0.214\linewidth}>{\raggedright\arraybackslash}p{0.122\linewidth}>{\raggedright\arraybackslash}p{0.122\linewidth}>{\raggedright\arraybackslash}p{0.122\linewidth}>{\raggedright\arraybackslash}p{0.122\linewidth}>{\raggedright\arraybackslash}p{0.122\linewidth}@{}}
\toprule
\textbf{point $(E_{00}, E_{01}, E_{10}, E_{11})$} & \textbf{$S$} & \textbf{$p$} & \textbf{$q$} & \textbf{$pq - s^2$} & \textbf{contrac-\newline tive} \\
\midrule
\endhead
\bottomrule
\endlastfoot
$(0, 0, 0, 0)$ & 0 & 1 & 1 & 1 & yes \\
$(0.6, 0.6, 0.6, -0.6)$ & 2.4 & 0.28 & 0.28 & 0.0784 & yes \\
$(a, b, b, -a)$ & 2.8 & 0 & 0 & 0 & yes, with no slack \\
$(t, t, t, -t)$ & $2\sqrt 2$ & 0 & 0 & 0 & yes, with no slack \\
$(1, 1, 1, 1)$, all-ones plan & 2 & $-1$ & $-1$ &  & no \\
$(1, 0, 1, 0)$, a classical mixture & 2 & $-1$ & 1 &  & no \\
$(1, 1, 1, -1)$, PR box & 4 & $-1$ & $-1$ &  & no \\
\end{longtable}
\endgroup

Every deterministic plan has all four correlators equal to $\pm 1$, so $p = 1 - 1 - 1 = -1$
and every one of the sixteen fails.

\begin{remark}[What PSD of the pinned matrix means]\label{rem:chsh-scope-five}
$M(E)$ is the level-1 moment matrix with the two within-party entries $\langle A_0 A_1
\rangle$ and $\langle B_0 B_1 \rangle$ pinned to zero. Its PSD is sufficient for the
correlators to admit some PSD completion (Section~\ref{sub:chsh-elliptope}), and it
is not necessary: the all-ones plan has a completion but fails here. So a check built on
$M(E)$ is sound for the completion set and not complete for it, and it rejects some honest
classical tallies. That correlators with a PSD completion are exactly those reachable by
quantum measurements is Tsirelson's theorem, and both halves are proved here in finite dimension.
For real and for complex amplitudes, a strategy's correlators are inner products of unit vectors
(with complex amplitudes, use real vectors of twice the length) and its level-1 moment matrix is PSD, so they have a PSD completion. The other half, that every PSD completion comes from a quantum
strategy, comes from this construction: the strategy is built on two four-level systems from
the completion's Gram vectors (Section~\ref{sub:four-vectors}).
\end{remark}

\subsection{The integer check}


The machine has whole-number tallies. The three conditions above were written
using averages, which are fractions. You can test the same conditions without
ever dividing: first make sure each denominator is positive, then multiply
it away. Multiplying by a positive number cannot turn a passing inequality
into a failing one or a failing one into a passing one.

For the first column, the condition starts as one minus two squared averages.
Multiplying by the product of their squared denominators leaves three
whole-number products. I call that result $A$. The second column gives $B$
the same way. The cross-column term is a sum of two products of averages;
multiplying by all four denominators leaves a whole number I call $C$.
The last condition becomes $C^2\le AB$, because both sides have had the same
positive denominator removed. Four positive bucket sizes and these three
comparisons make seven checks. Here are the products to calculate.

\begin{definition}[The seven-integer check]\label{def:chsh-check}
From the counters of Definition~\ref{def:chsh-counters}, with all
arithmetic in $\mathbb{Z}$, let
\[
\begin{aligned}
A &= N_{00}^2 N_{10}^2 - D_{00}^2 N_{10}^2 - D_{10}^2 N_{00}^2, \\
B &= N_{01}^2 N_{11}^2 - D_{01}^2 N_{11}^2 - D_{11}^2 N_{01}^2, \\
C &= D_{00} D_{01} N_{10} N_{11} + D_{10} D_{11} N_{00} N_{01}.
\end{aligned}
\]
The check is the conjunction of seven facts, in this order:
$0 < N_{00}$, $0 < N_{01}$, $0 < N_{10}$, $0 < N_{11}$, $0 \le A$, $0 \le B$, and
$C \cdot C \le A \cdot B$.
\end{definition}

There is no division, no real number, no square root and no $S$ in it.

\begin{lemma}[The seven facts are the three conditions with denominators cleared]
Assume $N_{00}, N_{01}, N_{10}, N_{11} > 0$ and put $E_{xy} = D_{xy}/N_{xy}$. Then, with
$p, q, s$ as in Definition~\ref{def:chsh-cc},
\[
\begin{gathered}
p = \frac{A}{N_{00}^2 N_{10}^2}, \qquad q = \frac{B}{N_{01}^2 N_{11}^2}, \qquad
s = \frac{C}{N_{00} N_{01} N_{10} N_{11}}, \\
pq - s^2 = \frac{AB - C^2}{(N_{00} N_{01} N_{10} N_{11})^2}.
\end{gathered}
\]
Because every denominator is positive, $0 \le A \iff p \ge 0$, $0 \le B \iff q \ge 0$ and
$C^2 \le AB \iff s^2 \le pq$.
\end{lemma}

\begin{proof}
For $p$: $1 - D_{00}^2/N_{00}^2 - D_{10}^2/N_{10}^2$; multiply by $N_{00}^2 N_{10}^2$ to get
$N_{00}^2 N_{10}^2 - D_{00}^2 N_{10}^2 - D_{10}^2 N_{00}^2 = A$. For $q$ the same with
columns $01$ and $11$, multiplied by $N_{01}^2 N_{11}^2$. For $s$:
\[
\frac{D_{00}}{N_{00}} \frac{D_{01}}{N_{01}} + \frac{D_{10}}{N_{10}} \frac{D_{11}}{N_{11}}
= \frac{D_{00} D_{01} N_{10} N_{11} + D_{10} D_{11} N_{00} N_{01}}{N_{00} N_{01} N_{10} N_{11}}
= \frac{C}{N_{00} N_{01} N_{10} N_{11}}.
\]
Then
$pq = AB / (N_{00}^2 N_{10}^2 N_{01}^2 N_{11}^2) = AB/(N_{00} N_{01} N_{10} N_{11})^2$ and
$s^2 = C^2/(N_{00} N_{01} N_{10} N_{11})^2$; subtract. Multiplying an inequality by a
positive number preserves it in both directions.
\end{proof}

The check has the following guarantee:

\begin{theorem}[Soundness of the integer check]
If a tally $w$ passes the seven-integer check, then its four correlators
$E_{xy}=D_{xy}/N_{xy}$ are column contractive. Consequently its pinned matrix is positive
semidefinite.
\end{theorem}

\begin{proof}
The first four tests say that every bucket contains a round, so each average
is its signed total divided by a positive count. The next two tests say that
$A$ and $B$ are nonnegative. Dividing each by its positive denominator gives
$p\ge0$ and $q\ge0$. The last test says $AB-C^2\ge0$. Divide that difference
by the positive square of the product of all four counts; the result is
$pq-s^2\ge0$, as the preceding calculation shows. Those are exactly the
three column conditions, and Theorem~\ref{thm:chsh-cc-psd} shows that they
make the matrix pass the nonnegative-total test.
\end{proof}

\begin{corollary}[A passed check implies the Tsirelson bound]
If a tally $w$ passes the seven-integer check, the correlators of $w$ satisfy
$S^2 \le 8$. It is the composition of Theorem~\ref{thm:chsh-check-sound},
Corollary~\ref{cor:chsh-iff} and Theorem~\ref{thm:chsh-psd-tsirelson}.\proofclosed \end{corollary} Two of the seven facts carry the bound by themselves. With every pair sampled, $0\le A$ and $0\le B$ are $p\ge0$ and $q\ge0$, so the four squared correlators add to at most 2 and Lemma~\ref{lem:chsh-sos} gives $S^2\le8$. The fact $C\cdot C\le A\cdot B$ is what positive semidefiniteness adds, and the bound doesn't use it.
For the form this takes on the small machine, where a raised flag has to stand on a passing
check, see Appendix~\ref{sec:chsh}.

A worked tally. Take $(\mathrm{same}, \mathrm{diff}) = (4,1), (4,1), (4,1), (1,4)$ for the
pairs $00, 01, 10, 11$. Then $N_{xy} = 5$ for all four, $D_{00} = D_{01} = D_{10} = 3$,
$D_{11} = -3$, and the correlators are $0.6, 0.6, 0.6, -0.6$ with $S = 2.4$. Compute:
\[
\begin{aligned}
A &= 25 \cdot 25 - 9 \cdot 25 - 9 \cdot 25 = 625 - 225 - 225 = 175, \\
B &= 25 \cdot 25 - 9 \cdot 25 - 9 \cdot 25 = 175, \\
C &= 3 \cdot 3 \cdot 5 \cdot 5 + 3 \cdot (-3) \cdot 5 \cdot 5 = 225 - 225 = 0.
\end{aligned}
\]
All seven facts hold: four positive $N$, $0 \le 175$ twice, and $0 \le 175 \cdot 175$.
Dividing back, $p = 175/625 = 0.28$ and $s = 0$, the second row of the table in Section~\ref{sub:chsh-five}.

Two more, for contrast. One round per pair from the all-ones plan, $(1,0)$ four times, gives
$N = D = 1$ everywhere and $A = 1 - 1 - 1 = -1 < 0$: the check fails. The tally $(4,1), (9,1), (9,1),
(1,4)$ gives correlators $3/5, 4/5, 4/5, -3/5$, with $N = 5, 10, 10, 5$ and
$D = 3, 8, 8, -3$:
$A = 25 \cdot 100 - 9 \cdot 100 - 64 \cdot 25 = 0$,
$B = 100 \cdot 25 - 64 \cdot 25 - 9 \cdot 100 = 0$ and
$C = 3 \cdot 8 \cdot 10 \cdot 5 + 8 \cdot (-3) \cdot 5 \cdot 10 = 0$. It passes with no
slack, at $S = 2.8$.

\subsection{Level 1+AB}


The first table had one place for each player's answer. Enlarge it by giving
each Alice--Bob product a place too. There are four such products, so the
five places become nine. A cell still asks for the average product of its
row and column labels. Some of those products reduce to entries I already
have; five more numbers are needed for the particular filling below.
Call them $\gamma_1$ through $\gamma_5$. The longer labels say which
products they stand for.

The five additional parameters name the moments
$\gamma_1=\langle A_0 A_1 B_0\rangle$,
$\gamma_2=\langle A_0 A_1 B_1\rangle$,
$\gamma_3=\langle A_0 B_0 B_1\rangle$,
$\gamma_4=\langle A_1 B_0 B_1\rangle$ and
$\gamma_5=\langle A_0 A_1 B_0 B_1\rangle$.
These are labels for the entries; the definition below does not construct operators.

\begin{definition}[The level-1+AB matrix]
The $9 \times 9$ real matrix on rows and
columns $I, A_0, A_1, B_0, B_1, A_0B_0, A_0B_1, A_1B_0, A_1B_1$, with the following entries:
\end{definition}

\[
\setlength{\arraycolsep}{3.2pt}
\begin{array}{c|ccccccccc}
 & I & A_0 & A_1 & B_0 & B_1 & A_0B_0 & A_0B_1 & A_1B_0 & A_1B_1 \\ \hline
I      & 1 & 0 & 0 & 0 & 0 & E_{00} & E_{01} & E_{10} & E_{11} \\
A_0    & 0 & 1 & 0 & E_{00} & E_{01} & 0 & 0 & \gamma_1 & \gamma_2 \\
A_1    & 0 & 0 & 1 & E_{10} & E_{11} & \gamma_1 & \gamma_2 & 0 & 0 \\
B_0    & 0 & E_{00} & E_{10} & 1 & 0 & 0 & \gamma_3 & 0 & \gamma_4 \\
B_1    & 0 & E_{01} & E_{11} & 0 & 1 & \gamma_3 & 0 & \gamma_4 & 0 \\
A_0B_0 & E_{00} & 0 & \gamma_1 & 0 & \gamma_3 & 1 & 0 & 0 & \gamma_5 \\
A_0B_1 & E_{01} & 0 & \gamma_2 & \gamma_3 & 0 & 0 & 1 & -\gamma_5 & 0 \\
A_1B_0 & E_{10} & \gamma_1 & 0 & 0 & \gamma_4 & 0 & -\gamma_5 & 1 & 0 \\
A_1B_1 & E_{11} & \gamma_2 & 0 & \gamma_4 & 0 & \gamma_5 & 0 & 0 & 1
\end{array}
\]

Two constraints are fixed by this layout. First, the four-body cells carry opposite signs:
$+\gamma_5$ at $(A_0B_0, A_1B_1)$ and $-\gamma_5$ at $(A_0B_1, A_1B_0)$. In a real
moment matrix those two cells are $\langle A_0 A_1 B_0 B_1 \rangle$ and $\langle A_0 A_1 B_1 B_0
\rangle$, so the tie says $\langle A_0 A_1 (B_0 B_1 + B_1 B_0) \rangle = 0$, which for a real
matrix is $\langle \{A_0, A_1\} \otimes \{B_0, B_1\} \rangle = 0$. The level-1 pinnings don't
imply it. Take a classical strategy: Alice's two answers $a_0, a_1$ are independent fair signs,
and Bob copies them, $b_0 = a_0$ and $b_1 = a_1$. Every pinned cell of its moment matrix is 0 and
both four-body cells are 1, so no choice of the nine numbers gives its matrix. So the matrix
here is level 1+AB cut down to the slice $\langle \{A_0, A_1\} \otimes \{B_0, B_1\} \rangle = 0$,
and that is what ``level 1+AB'' names in this book. The slice holds whenever $A_0$ and $A_1$ anticommute.
The soundness theorems below are about this matrix, so they stand as stated. Second, the matrix
keeps the level-1 pinning: every cell whose label reduces to a marginal or to
$\langle A_0 A_1 \rangle$ or $\langle B_0 B_1 \rangle$ is zero, for instance
$(A_0, A_0B_0)$, $(A_0B_0, A_0B_1)$ and $(A_0B_0, A_1B_0)$. The displayed matrix is symmetric.

I want to test this larger table with the same operation as before:
put a weight beside each row and column and add the weighted cells. The
first three weights will each be collected into one completed square.
Everything involving the other six will be left over. That leftover is
what I call the residual; it is the part whose sign I still need to settle.

Write a vector as $(v_I, v_{A_0}, v_{A_1}, v_{B_0}, v_{B_1}, v_{00}, v_{01}, v_{10}, v_{11})$,
where $v_{xy}$ is the coordinate of $A_xB_y$. Put
\[
\begin{aligned}
L_1 &= E_{00} v_{00} + E_{01} v_{01} + E_{10} v_{10} + E_{11} v_{11}, \\
L_2 &= E_{00} v_{B_0} + E_{01} v_{B_1} + \gamma_1 v_{10} + \gamma_2 v_{11}, \\
L_3 &= E_{10} v_{B_0} + E_{11} v_{B_1} + \gamma_1 v_{00} + \gamma_2 v_{01}.
\end{aligned}
\]

\begin{definition}[The residual]
\[
\begin{aligned}
R ={}& v_{B_0}^2 + v_{B_1}^2 + v_{00}^2 + v_{01}^2 + v_{10}^2 + v_{11}^2 \\
&+ 2 v_{B_0} (\gamma_3 v_{01} + \gamma_4 v_{11}) + 2 v_{B_1} (\gamma_3 v_{00} + \gamma_4 v_{10}) \\
&+ 2 \gamma_5 (v_{00} v_{11} - v_{01} v_{10}) - L_1^2 - L_2^2 - L_3^2.
\end{aligned}
\]
The condition called column contractivity at level 1+AB is $R \ge 0$ for all six real inputs.
\end{definition}

\begin{lemma}[Three squares plus the residual]\label{lem:chsh-nine-sos}
For every vector,
\[
v^{\mathsf T} M_9\, v = (v_I + L_1)^2 + (v_{A_0} + L_2)^2 + (v_{A_1} + L_3)^2 + R.
\]
\end{lemma}

\begin{proof}
The coordinates $v_I$, $v_{A_0}$, $v_{A_1}$ appear in the form only through their diagonal
squares and the cross terms $2 v_I L_1$, $2 v_{A_0} L_2$, $2 v_{A_1} L_3$, read off rows
$I$, $A_0$, $A_1$ of the matrix. Completing each square adds $L_1^2 + L_2^2 + L_3^2$, which
the residual subtracts. What is left in $R$ is the diagonal of the remaining six coordinates
and the cross terms among them, read off the lower right $6 \times 6$ block.
\end{proof}

\begin{theorem}[PSD at level 1+AB]
For all real $E_{xy}$ and $\gamma_1, \ldots, \gamma_5$,
the matrix of Definition~\ref{def:chsh-nine} is positive semidefinite if and only if
the residual $R$ is nonnegative for every choice of its six real inputs.
\end{theorem}

\begin{proof}
If $R \ge 0$ everywhere, Lemma~\ref{lem:chsh-nine-sos} writes the form as three squares plus $R$.
Conversely, given the six residual inputs, choose $v_I = -L_1$, $v_{A_0} = -L_2$,
$v_{A_1} = -L_3$; the three squares vanish and PSD gives $R \ge 0$. This is the level-1
argument of Section~\ref{sub:chsh-five} one size up.
\end{proof}

As a quadratic form in $(v_{B_0},v_{B_1},v_{00},v_{01},v_{10},v_{11})$, numbered 1 to 6,
the residual has the symmetric coefficient matrix $H$ with these upper entries. Each entry
can be read back off $R$ by collecting the coefficient of $v_i v_j$.

\begin{align*}
H_{11} &= 1 - E_{00}^2 - E_{10}^2, & H_{22} &= 1 - E_{01}^2 - E_{11}^2, \\
H_{12} &= -(E_{00}E_{01} + E_{10}E_{11}), & H_{13} &= -E_{10}\gamma_1, \\
H_{14} &= \gamma_3 - E_{10}\gamma_2, & H_{15} &= -E_{00}\gamma_1, \\
H_{16} &= \gamma_4 - E_{00}\gamma_2, & H_{23} &= \gamma_3 - E_{11}\gamma_1, \\
H_{24} &= -E_{11}\gamma_2, & H_{25} &= \gamma_4 - E_{01}\gamma_1, \\
H_{26} &= -E_{01}\gamma_2, & H_{33} &= 1 - E_{00}^2 - \gamma_1^2, \\
H_{34} &= -(E_{00}E_{01} + \gamma_1\gamma_2), & H_{35} &= -E_{00}E_{10}, \\
H_{36} &= \gamma_5 - E_{00}E_{11}, & H_{44} &= 1 - E_{01}^2 - \gamma_2^2, \\
H_{45} &= -\gamma_5 - E_{01}E_{10}, & H_{46} &= -E_{01}E_{11}, \\
H_{55} &= 1 - E_{10}^2 - \gamma_1^2, & H_{56} &= -(E_{10}E_{11} + \gamma_1\gamma_2), \\
H_{66} &= 1 - E_{11}^2 - \gamma_2^2.
\end{align*}

The top left $2 \times 2$ block is $I_2 - E^{\mathsf T} E$ from Section~\ref{sub:chsh-five}.

The full check does not decide $R \ge 0$ directly. It tests the stronger fact that $H$ is
positive definite, by two Schur steps and a Sylvester test, all fraction-free.

\begin{lemma}[Fraction-free Schur step]\label{lem:chsh-schur-step}
For a symmetric $n \times n$ matrix $h$ and any vector $v$, with
$S_{ij} = h_{11} h_{ij} - h_{1i} h_{1j}$ for $i, j \ge 2$ (here $S$ is the Schur matrix;
the score plays no part
in this lemma),
\[
h_{11} \, v^{\mathsf T} h\, v = \Bigl(\sum_{j} h_{1j} v_j\Bigr)^2 + \sum_{i,j \ge 2} S_{ij} v_i v_j.
\]
So if $h_{11} > 0$ and the $(n-1) \times (n-1)$ matrix $S$ gives a nonnegative form, then $h$
does too. The cascade uses this identity at sizes six and five.
\end{lemma}

\begin{proof}
Expand the square: $h_{11}^2 v_1^2 + 2 h_{11} v_1 \sum_{j \ge 2} h_{1j} v_j +
(\sum_{j \ge 2} h_{1j} v_j)^2$. The left side is
$h_{11}^2 v_1^2 + 2 h_{11} v_1 \sum_{j \ge 2} h_{1j} v_j + h_{11} \sum_{i,j \ge 2} h_{ij} v_i v_j$.
The difference is $\sum_{i,j \ge 2} (h_{11} h_{ij} - h_{1i} h_{1j}) v_i v_j$. Dividing by
$h_{11} > 0$ keeps the sign.
\end{proof}

\begin{lemma}[Sylvester for $4\times4$]\label{lem:chsh-sylvester}
Let $h$ be symmetric $4 \times 4$ with leading principal
minors $d_1, d_2, d_3, d_4$. If $d_1, d_2, d_3$ are strictly positive and $d_4 \ge 0$, then $v^{\mathsf T} h\, v \ge 0$ for all $v$.
\end{lemma}

\begin{proof}
Use the fraction-free factorization
\[
d_1^2 d_2 d_3 \, v^{\mathsf T} h\, v = d_1 d_2 d_3 \, P_1^2 + d_1 d_3 \, Q_1^2 + Q_2^2 +
d_1^2 d_2 d_4 \, v_4^2,
\]
where, writing $e_{ij}=h_{11}h_{ij}-h_{1i}h_{1j}$,
\[
\begin{aligned}
P_1 &= \sum_j h_{1j}v_j,\\
Q_1 &= d_2v_2+e_{23}v_3+e_{24}v_4,\\
Q_2 &= d_1d_3v_3+(d_2e_{34}-e_{23}e_{24})v_4.
\end{aligned}
\]
The identity is Lemma~\ref{lem:chsh-schur-step} three times over, each time on the matrix the
last one left. With $f_{ij}=e_{22}e_{ij}-e_{2i}e_{2j}$ for $i,j\ge3$, the lemma gives
$d_1\,v^{\mathsf T}h\,v = P_1^2+\sum_{i,j\ge2}e_{ij}v_iv_j$, then
$e_{22}\sum_{i,j\ge2}e_{ij}v_iv_j = Q_1^2+\sum_{i,j\ge3}f_{ij}v_iv_j$, then
$f_{33}\sum_{i,j\ge3}f_{ij}v_iv_j = Q_2^2+(f_{33}f_{44}-f_{34}^2)\,v_4^2$. The corners are
minors of $h$, by one fact about determinants: the $k\times k$ leading minor of $e$ is
$h_{11}^{k-1}$ times the $(k+1)\times(k+1)$ leading minor of $h$. To see it, multiply rows 2 to
$k+1$ of that block of $h$ by $h_{11}$, which multiplies its determinant by $h_{11}^k$, then
subtract $h_{1i}$ times the first row from row $i$, which changes nothing; those rows now read
$0, e_{i2}, \ldots, e_{i,k+1}$, so the determinant is $h_{11}$ times the leading minor of $e$.
So $e_{22}=d_2$ and $f_{33}=e_{22}e_{33}-e_{23}^2=d_1d_3$, and the same fact applied to $e$,
whose corner is $d_2$, gives $f_{33}f_{44}-f_{34}^2=d_2\det e=d_1^2d_2d_4$. Multiply the first
equation by $d_1d_2d_3$ and substitute the other two. With $d_1, d_2, d_3 > 0$ and $d_4 \ge 0$ the right side is nonnegative and the factor on the left is positive.
\end{proof}

The cascade for the full check runs on $H$. Step one: $H_{11} > 0$ and the $5 \times 5$
matrix $S^{(6)}_{ij} = H_{11} H_{ij} - H_{1i} H_{1j}$, $2 \le i, j \le 6$, gives a nonnegative
form. Step two: $S^{(6)}_{22} > 0$ and the $4 \times 4$ matrix
$S^{(5)}_{ij} = S^{(6)}_{22} S^{(6)}_{ij} - S^{(6)}_{2i} S^{(6)}_{2j}$, $3 \le i, j \le 6$,
gives a nonnegative form. Step three: the four leading minors of $S^{(5)}$ are positive. Together these conditions make the residual nonnegative, and so the nine-by-nine matrix positive semidefinite.

\begin{definition}[What the integer check computes]\label{def:chsh-g12345}
The supplied moments arrive as bucket pairs: $\gamma_k = N_{g_k}/D_{g_k}$ with
$N_{g_k} = \mathrm{same}_{g_k} - \mathrm{diff}_{g_k}$ and
$D_{g_k} = \mathrm{same}_{g_k} + \mathrm{diff}_{g_k}$. Watch the letters: for these five
buckets $N$ names the numerator and $D$ the denominator: the difference and the count, the reverse
of Section~\ref{sub:chsh-notation}. Let
$K = (N_{00} N_{01} N_{10} N_{11} D_{g_1} D_{g_2} D_{g_3} D_{g_4} D_{g_5})^2$. The check computes the 21 integers
$K H_{ij}$, the 15
integers of $S^{(6)}$ at scale $K^2$ and the 10 integers of $S^{(5)}$ at scale $K^4$, each by
$h_{11} h_{ij} - h_{1i} h_{1j}$ on integers. The
decider is the conjunction of: $N_{xy} > 0$ for the four
pairs; $D_{g_k} > 0$ for $k = 1..5$; $-D_{g_k} < N_{g_k} < D_{g_k}$ for $k = 1..5$; and six
strict positivities, namely $K H_{11} > 0$, $K^2 S^{(6)}_{22} > 0$, and the four leading
minors of the integer $S^{(5)}$. The full check
combines the seven-integer check of Definition~\ref{def:chsh-check} with this decider.
\end{definition}

Each cleared integer, read as a real, equals the
corresponding real quantity times the stated positive power of $K$. A $k \times k$ leading
minor of a matrix scaled by $c > 0$ scales by $c^k$. So every integer positivity holds exactly
when the real one does.

Two parts of the full check are there to fail early, and no proof below needs them. A passing
decider makes $H$ positive definite (Remark~\ref{rem:chsh-interior}), and the top left two-by-two
block of $H$ has $p$ and $q$ on its diagonal and $-s$ off it, so $p > 0$, $q > 0$ and $pq > s^2$,
and the seven-integer check passes. And with $H$ positive definite, the nine-by-nine form is zero
only at zero: it is three squares plus $R$ (Lemma~\ref{lem:chsh-nine-sos}), so where it is zero the
six inputs of $R$ are zero, then so are $L_1$, $L_2$ and $L_3$, and then so are $v_I$, $v_{A_0}$
and $v_{A_1}$. So the nine-by-nine matrix is positive definite; with ones on its diagonal, that puts every $\gamma_k$ strictly between
$-1$ and 1, which is the bucket bound $-D_{g_k} < N_{g_k} < D_{g_k}$.

\begin{theorem}[The full 1+AB check implies PSD]\label{thm:chsh-g12345}
Let the counts make every $N_{xy}$ and every $D_{g_k}$ positive, and let
the decider just defined pass on them. Then the matrix of
Definition~\ref{def:chsh-nine} at the four correlators and $\gamma_k = N_{g_k}/D_{g_k}$ is
symmetric and positive semidefinite.
\end{theorem}

\begin{proof}
The bucket conjuncts make every denominator positive, the scaling remark just above carries the
six integer positivities over to the real cascade, and Lemma~\ref{lem:chsh-schur-step} twice plus
Lemma~\ref{lem:chsh-sylvester} give $R \ge 0$, and so positive semidefiniteness by Theorem~\ref{thm:chsh-nine-iff}.
\end{proof}

\begin{remark}[Interior only]\label{rem:chsh-interior}
Every test in the cascade is strict, so a passing check certifies that $H$ is positive
definite, which is stronger than semidefinite. A point where $H$ is PSD but singular cannot pass, and
neither can a moment with $|\gamma_k| = 1$. For example, the tally of Section~\ref{sub:chsh-integer} with correlators
$(3/5, 4/5, 4/5, -3/5)$ has $H_{11} = 1 - 9/25 - 16/25 = 0$, so the full 1+AB check fails on it
for every choice of $\gamma$, though the level-1 check of Section~\ref{sub:chsh-integer} accepts it. The argument establishes soundness only; no converse is stated.
\end{remark}

A worked example of the cascade in exact rational arithmetic. Take the correlators $(0.6, 0.6, 0.6, -0.6)$ and
$\gamma = (0, 0, 0, 0, -1/2)$; as buckets, $\gamma_1$ to $\gamma_4$ are $(1,1)$, giving
$0/2$, and $\gamma_5$ is $(1,3)$, giving $-2/4$, both inside the strict bounds. Then
$H_{11} = H_{22} = 7/25$, $H_{12} = 0$, and every entry $H_{1j}$, $H_{2j}$ with $j \ge 3$ is
zero. So $S^{(6)}_{22} = (7/25)^2 = 49/625$, and $S^{(5)}$ is a positive multiple of the lower
right block of $H$:
\[
S^{(5)} = \frac{343}{781250}
\begin{pmatrix}
32 & -18 & -18 & -7 \\
-18 & 32 & 7 & 18 \\
-18 & 7 & 32 & 18 \\
-7 & 18 & 18 & 32
\end{pmatrix},
\]
where $50 H_{33} = 32$, $50 H_{34} = 50 H_{35} = -18$, $50 H_{36} = 50(-1/2 + 0.36) = -7$,
$50 H_{45} = 50(1/2 - 0.36) = 7$, and $50 H_{46} = 50 H_{56} = 18$. The leading minors of the
integer matrix are $32$, $700$, $15000$ and $140625$, all positive. With
$H_{11} = 7/25 > 0$ and $S^{(6)}_{22} > 0$, all six tests pass, and with Section~\ref{sub:chsh-integer}'s $A = B = 175$,
$C = 0$ the full check passes. At $\gamma_5 = 0$ instead, the same block becomes
$I_4 - \mathbf e \mathbf e^{\mathsf T}$ with $\mathbf e = (0.6, 0.6, 0.6, -0.6)$, whose third
leading minor is negative ($\tfrac{1}{25^3}$ times $-1250$), and the check fails.

That last fact is general. The zero-moment check takes no $\gamma$ buckets; it combines
the level-1 check with the integer form of
$E_{00}^2 + E_{01}^2 + E_{10}^2 + E_{11}^2 \le 1$ with denominators cleared by
$(N_{00} N_{01} N_{10} N_{11})^2$.

\begin{theorem}[The zero-moment check sits inside the classical bound]\label{thm:chsh-gzero}
If the zero-moment check accepts a tally $w$, then the four correlators of $w$ satisfy
$E_{00}^2 + E_{01}^2 + E_{10}^2 + E_{11}^2 \le 1$
 and $|S| \le 2$.
\end{theorem}

\begin{proof}
The check's extra conjunct, divided back by the positive factor
$(N_{00} N_{01} N_{10} N_{11})^2$, is $r = E_{00}^2 + E_{01}^2 + E_{10}^2 + E_{11}^2 \le 1$.
Then Lemma~\ref{lem:chsh-sos} gives $S^2 \le 4 r \le 4$.
\end{proof}

The level-1 half of that check follows from $r \le 1$ on its own: $p$ and $q$ are each at least
$1 - r$, and $pq - s^2 = 1 - r + (E_{00}E_{11} - E_{01}E_{10})^2$.

This is exactly the level-1+AB condition at $\gamma = 0$. There the residual splits into the
level-1 form in $(v_{B_0}, v_{B_1})$ and
$|\mathbf v|^2 - \langle \mathbf e, \mathbf v \rangle^2$ in
$\mathbf v = (v_{00}, v_{01}, v_{10}, v_{11})$, where
$\mathbf e = (E_{00}, E_{01}, E_{10}, E_{11})$. The second piece is nonnegative for every
$\mathbf v$ exactly when $r \le 1$: if $r \le 1$, Cauchy--Schwarz gives
$\langle \mathbf e, \mathbf v \rangle^2 \le r\,|\mathbf v|^2 \le |\mathbf v|^2$; if $r > 1$, the
choice $\mathbf v = \mathbf e$ gives $r - r^2 < 0$.

So the zero-moment check certifies a strict part of the classical region: no point
above $|S| = 2$, and none of the sixteen deterministic points either, since each has
$|\mathbf e|^2 = 4$. A point above 2 needs supplied moments.
The matrix is positive semidefinite at correlators $(3/5, 3/5, 3/5, -3/5)$,
$\gamma = (0,0,0,0,-1/2)$, with $S = 12/5$. The cleanest way to see it is the $\gamma_5$ slice,
$\gamma_1 = \cdots = \gamma_4 = 0$. There the residual splits into the level-1 form in
$(v_{B_0}, v_{B_1})$ and the form of $J - \mathbf e\mathbf e^{\mathsf T}$ in the four $v_{xy}$,
where $J$ is $I_4$ with $\gamma_5$ at $(00, 11)$ and $-\gamma_5$ at $(01, 10)$. For
$-1 < \gamma_5 < 1$, $J$ is positive definite, and the second form is nonnegative exactly when
$\mathbf e^{\mathsf T} J^{-1} \mathbf e \le 1$: apply Lemma~\ref{lem:chsh-schur} to the block
matrix $\begin{pmatrix} 1 & \mathbf e^{\mathsf T} \\ \mathbf e & J \end{pmatrix}$ once with the corner
$1$ as $P$ and once, rows reordered, with $J$ as $P$. Inverting the two $2 \times 2$ blocks of $J$ and multiplying out, that is
\[
\begin{aligned}
&(1 - \gamma_5)\bigl[(E_{00} + E_{11})^2 + (E_{01} - E_{10})^2\bigr] \\
&\quad + (1 + \gamma_5)\bigl[(E_{00} - E_{11})^2 + (E_{01} + E_{10})^2\bigr]
\le 2(1 - \gamma_5^2).
\end{aligned}
\]
At this point the left side is $\tfrac32 \cdot 0 + \tfrac12 \cdot (\tfrac{36}{25} +
\tfrac{36}{25}) = \tfrac{36}{25} = 1.44$ and the right side is $2 \cdot \tfrac34 = 1.5$, and the
level-1 part holds ($p = q = 0.28$, $s = 0$). The slice test is an illustration, not a third
check: on the counters $(4,1), (4,1), (4,1), (1,4)$ and the bucket $(1,3)$, the same comparison
with the denominators multiplied out says yes, and the full check of
Definition~\ref{def:chsh-g12345} accepts the same counters with $\gamma_1$ to $\gamma_4$ given as
the buckets $(1,1)$, as the worked example above shows.

Both 1+AB checks, the full check and the zero-moment check, contain
the seven-integer check as a conjunct, so a passing 1+AB check also passes the
level-1 check and inherits $S^2 \le 8$ from Section~\ref{sub:chsh-integer}. This follows directly by composing the definitions and the bound.

\subsection{The elliptope completion}


Go back to the two within-player cells I filled with zero. Leave them blank
this time. Keep all four measured cross-player averages fixed, and ask whether
there is any way to fill those two blanks that makes the table pass. A failure
with zeros no longer settles the question; another filling may work.

Call the first missing number $x$ and the second $y$, and copy each into its
reflected cell so the table stays symmetric. The collection of averages that
admit a passing filling has a name, the \emph{elliptope} completion set.
Membership means that at least one choice of those two numbers works. It
does not require every choice to work, or require the zero choice to work.

\begin{definition}[Completion and membership]
Complete the matrix of Section~\ref{sub:chsh-five} by setting the
within-party entries to free reals: $x$ in the $(A_0, A_1)$ cells and $y$ in the
$(B_0, B_1)$ cells, marginals still zero. The correlators are elliptope realizable
 when some $x, y$ make the completed matrix symmetric and PSD.
\end{definition}

\[
M(E; x, y) \;=\;
\begin{pmatrix}
1 & 0 & 0 & 0 & 0 \\
0 & 1 & x & E_{00} & E_{01} \\
0 & x & 1 & E_{10} & E_{11} \\
0 & E_{00} & E_{10} & 1 & y \\
0 & E_{01} & E_{11} & y & 1
\end{pmatrix}
\]

Its quadratic form at $v = (v_0, \ldots, v_4)$ is
$\sum_k v_k^2 + 2x v_1 v_2 + 2y v_3 v_4 + 2E_{00} v_1 v_3 + 2E_{01} v_1 v_4 + 2E_{10} v_2 v_3
+ 2E_{11} v_2 v_4$. The pinned matrix is the case $x = y = 0$,
so PSD of $M(E)$ implies membership.

\begin{theorem}[Deterministic plans are inside]\label{thm:chsh-ell-det}
If $a_0^2 = a_1^2 = b_0^2 = b_1^2 = 1$, then
$(a_0 b_0, a_0 b_1, a_1 b_0, a_1 b_1)$ is elliptope realizable, with $x = a_0 a_1$ and
$y = b_0 b_1$.
\end{theorem}

\begin{proof}
With these entries the form equals
$v_0^2 + (a_0 v_1 + a_1 v_2 + b_0 v_3 + b_1 v_4)^2$. Expanding the second square gives
$a_0^2 v_1^2 + \cdots$ on the diagonal, which is $v_1^2 + \cdots$ by the unit squares, and the
cross terms $2 a_0 a_1 v_1 v_2 = 2x v_1 v_2$, $2 a_0 b_0 v_1 v_3 = 2E_{00} v_1 v_3$, and so on
for all six pairs. A sum of two squares is never negative.
\end{proof}

The all-ones plan is one instance. The classical mixture $(1,0,1,0)$ has a completion at
$x=1$, $y=0$, with form $v_0^2+(v_1+v_2+v_3)^2+v_4^2$.

\begin{theorem}[Convexity and finite mixtures]\label{thm:chsh-ell-convex}
For $0 \le \lambda \le 1$, if $E$ and $F$ are elliptope realizable, so is
$\lambda E + (1 - \lambda) F$. For finitely
many members with nonnegative weights summing to at most 1, the weighted sum is realizable. Every finite mixture of deterministic plans is realizable.
\end{theorem}

\begin{proof}
The form is affine in $(E, x, y)$ for each fixed $v$. Take the completion
$(\lambda x_1 + (1-\lambda) x_2, \lambda y_1 + (1-\lambda) y_2)$; the form at the mixture is
the same mixture of two nonnegative forms. For several members, weight each one's completed matrix by its assigned
share and add them. If the shares total less than one, give the missing
share to the zero-correlator point with both missing entries zero; its
completed matrix is the identity. The diagonal of the weighted matrix is
then one, its cross entries are the desired weighted correlators, and every
weighted total is a sum of nonnegative totals with nonnegative weights.
It therefore supplies the required completion.
The last statement combines this with Theorem~\ref{thm:chsh-ell-det}.
\end{proof}

\begin{theorem}[Tsirelson's bound holds on the completion set]\label{thm:chsh-ell-tsirelson}
If $E$ is elliptope realizable then $S^2 \le 8$.
\end{theorem}

\begin{proof}
For a symmetric PSD matrix the bilinear form obeys Cauchy-Schwarz,
$\mathcal B(u, w)^2 \le \mathcal Q(u)\, \mathcal Q(w)$, because
$\mathcal Q(u + t w) = \mathcal Q(u) + 2t\,\mathcal B(u,w) + t^2 \mathcal Q(w) \ge 0$ for all
$t$ and Lemma~\ref{lem:chsh-discriminant} applies. Take $u_1 = e_{A_0}$ and
$w_+ = e_{B_0} + e_{B_1}$: $\mathcal B(u_1, w_+) = E_{00} + E_{01}$, $\mathcal Q(u_1) = 1$,
$\mathcal Q(w_+) = 2 + 2y$. Take $u_2 = e_{A_1}$ and $w_- = e_{B_0} - e_{B_1}$:
$\mathcal B(u_2, w_-) = E_{10} - E_{11}$, $\mathcal Q(w_-) = 2 - 2y$. So
$c_1^2 \le 2 + 2y$ and $c_2^2 \le 2 - 2y$ with $c_1 = E_{00} + E_{01}$,
$c_2 = E_{10} - E_{11}$, and both budgets are nonnegative. Then
$c_1^2 c_2^2 \le (2 + 2y)(2 - 2y) = 4 - 4y^2 \le 4$, so $c_1 c_2 \le 2$, and
$S^2 = (c_1 + c_2)^2 = c_1^2 + 2 c_1 c_2 + c_2^2 \le (2 + 2y) + 4 + (2 - 2y) = 8$.
\end{proof}

The completion can make a row of $E$ longer than 1 (the all-ones plan has rows of squared
length 2), which is why this proof does not use the row bounds of Section~\ref{sub:chsh-five}.

\begin{corollary}[The PR box is outside]\label{cor:chsh-ell-pr}
 $(1, 1, 1, -1)$ is not elliptope realizable, since its
$S^2 = 16 > 8$.\proofclosed
\end{corollary}

\begin{theorem}[Strictly beyond the classical bound]\label{thm:chsh-ell-beyond}
 $(3/5, 3/5, 3/5, -3/5)$ is elliptope realizable and has
$S = 12/5 > 2$.
\end{theorem}

\begin{proof}
At $x = y = 0$ the form equals
$v_0^2 + (v_1 + \tfrac35 v_3 + \tfrac35 v_4)^2 + (v_2 + \tfrac35 v_3 - \tfrac35 v_4)^2 +
\tfrac{7}{25} v_3^2 + \tfrac{7}{25} v_4^2$. Check: the two middle squares give
$v_3^2$ coefficient $\tfrac{9}{25} + \tfrac{9}{25} = \tfrac{18}{25}$, plus $\tfrac{7}{25}$ is
1; the $v_3 v_4$ terms are $+\tfrac{18}{25}$ and $-\tfrac{18}{25}$ and cancel; the cross
terms are $2 \cdot \tfrac35 v_1 v_3$, $2 \cdot \tfrac35 v_1 v_4$, $2 \cdot \tfrac35 v_2 v_3$
and $-2 \cdot \tfrac35 v_2 v_4$, as the entries require.
\end{proof}

By Theorem~\ref{thm:chsh-fine} no finite factorizable model has $S > 2$, so this point is not a finite
mixture of plans, while Theorem~\ref{thm:chsh-ell-convex} puts every such mixture inside. The completion
set then strictly contains the finite classical set, and the PR box lies outside it. It sits
inside the no-signalling set, for the plain reason that every correlator in it is between $-1$
and $1$ (the vector that is 1 at $A_x$ and $\mp 1$ at $B_y$ gives $2 \mp 2E_{xy} \ge 0$), and any
four such numbers come from the no-signalling box that answers $(a, b)$ at settings $(x, y)$ with
probability $(1 + ab\,E_{xy})/4$, each player's own answer alone a fair coin.

The elliptope gate is an integer decider with a soundness proof. Its
completion entries arrive as bucket pairs, and an empty bucket is read as $0/1$,
so unlike Section~\ref{sub:chsh-integer} it does not require every pair to be sampled.

\begin{definition}[The two branches]
The positive-minor branch: with the completion supplied as buckets and all
entries scaled to integers by a positive common factor $K$, the four leading minors of
the $4 \times 4$ block on $A_0, A_1, B_0, B_1$ are strictly positive. The factorization branch: with the completion and a rational certificate supplied (six subdiagonal entries over
a common positive denominator, four pivots $d_k \ge 0$ over another), the ten entry equations
of $M = L D L^{\mathsf T}$ hold in cross-multiplied integer arithmetic. The full gate is their disjunction.
\end{definition}

\begin{theorem}[Soundness of the gate]
If the full gate returns yes on a counter record, any
completion buckets and any certificate, then the four correlators of the record are
elliptope realizable. On the PR box
tally the gate returns no for every choice of the remaining arguments.
\end{theorem}

\begin{proof}
In both branches it is enough to show the $4 \times 4$ block on $A_0, A_1, B_0, B_1$ gives a
nonnegative form, because the $I$ row and column are 1 on the diagonal and 0 elsewhere. The
first branch: dividing by the positive factor $K$ (a $k \times k$ minor scales by $K^k$) gives four
positive leading minors of the real block, and Lemma~\ref{lem:chsh-sylvester} applies. The
second: cross-multiplying by positive denominators changes no equation, so $M = LDL^{\mathsf T}$
holds over the rationals with $L$ unit lower triangular, and the form is
$v^{\mathsf T} L D L^{\mathsf T} v = \sum_k d_k P_k^2$ with $P = L^{\mathsf T} v$ and every
$d_k \ge 0$. The PR statement is soundness followed by Corollary~\ref{cor:chsh-ell-pr}.
\end{proof}

Examples by direct computation: the all-ones plan passes through the factorization branch
at $x=y=1$ with pivots $(1,0,0,0)$; the classical mixture $(1,0,1,0)$ passes through the same
branch at $x=1$, $y=0$ with pivots $(1,0,0,1)$; the $12/5$ point passes through the
positive-minor branch at $x=y=0$; and $(3/5,4/5,4/5,-3/5)$ passes through the factorization
branch at $x=y=0$ with pivots $(1,1,0,0)$, since $B_0 = \tfrac35 A_0 + \tfrac45 A_1$ and
$B_1 = \tfrac45 A_0 - \tfrac35 A_1$ reproduce every entry.

That's all the gate is: no property of the small machine is built from it.

\subsection{The four-vector construction}


Put two equal-length pencils down with their ends together, and draw a line along
one of them. Keep that pencil still and turn the other. Use the square corner of
a sheet of paper to bring the moving pencil's tip straight across to the line,
and mark where it lands. The mark moves from a full pencil-length forward,
through the joined ends when the pencils make a square corner, to a full
pencil-length backward. Call those readings 1, 0 and $-1$. That is the signed
shadow I want. Between those positions the reading is the cosine of the angle.
A dot product of unit vectors measures that shadow.

Now ask for four directions, two labelled Alice and two Bob, with a prescribed
shadow for each Alice--Bob pair. Try this request: Alice's first points exactly
along both of Bob's; Alice's second points along Bob's first and exactly against
his second. Each request alone is easy. Put them together and your hands cannot do
it: Bob's pencils would have to point both the same way and opposite ways. Four
numbers in range need not describe four directions that can coexist.

Write those four readings as $x=E_{00}$, $y=E_{01}$, $z=E_{10}$ and $w=E_{11}$.
I haven't specified the shadow between Alice's own two directions, or between
Bob's own two. Those are the two blanks, $t$ and $u$, in the following table.
Can they be filled so that the whole table comes from directions? The pencils let
you ask it in a plane; the construction will allow four coordinates. A completion is
a positive-semidefinite matrix
\[
G=\begin{pmatrix}
1&t&x&y\\t&1&z&w\\x&z&1&u\\y&w&u&1
\end{pmatrix}.
\]
The extra row and column labelled $I$ in Section~\ref{sub:chsh-elliptope} have diagonal 1 and all other entries
zero, so they do not change this condition.

Why positivity? Lay arrows end to end, allowing a backward arrow when its weight
is negative. Square the length from the start to the finish. Expanding that square
adds up the arrows' pairwise shadows, with their weights. It cannot be negative.
A proposed shadow table that gives a negative answer for some choice of weights
cannot come from arrows. That is the PSD test. A Gram matrix is simply the table
of all those pairwise dot products.

The converse has to build the arrows. Pick the first direction as a baseline.
Give every other arrow the component along it that its shadow demands, then
build the remaining components perpendicular to it. What remains to be fitted is
a smaller shadow table. Here is that construction, including its boundary. If its first diagonal entry is zero, positivity of
$h_{11}a^2+2h_{1j}ab+h_{jj}b^2$ for every $a,b$ forces $h_{1j}=0$:
with $b=1$, a nonzero $h_{1j}$ would let the linear term go negative without bound.
Delete that zero row and column, construct vectors for the smaller matrix, and put back
the zero vector. If $h_{11}>0$, set
\[
K_{ij}=h_{ij}-\frac{h_{1i}h_{1j}}{h_{11}}\quad(i,j>1).
\]
In the quadratic form of $h$, take the first coordinate to be
$-\sum_{j>1}h_{1j}v_j/h_{11}$. The remaining form is $v^{\mathsf T}Kv$, so $K$ is PSD.
Repeat the construction on this smaller matrix. Each repetition removes a row
and column, so it ends at a single nonnegative entry, represented by its square
root, or at the empty table. Working back up supplies vectors $r_i$ with
$r_i\cdot r_j=K_{ij}$. Then
$u_1=(\sqrt{h_{11}},0)$ and $u_i=(h_{1i}/\sqrt{h_{11}},r_i)$ have Gram matrix $h$.
The dimension is at most the number of rows. Applied to $G$, this gives four unit
vectors $a_0,a_1,b_0,b_1$ in $\mathbb R^4$, with the required cross inner products.

Three dimensions are always enough. The table fixes the cross entries and leaves two entries of
$G$ open, $a_0\cdot a_1$ and $b_0\cdot b_1$. Hold the first and let the second be $t$. The values
of $t$ that keep $G$ PSD form an interval: if two values do, so does every value between, because
the quadratic form at a mixture is the same mixture of the two forms; the set is closed, because a
form that is nonnegative at every value of a convergent sequence is nonnegative at its limit; and
it is bounded, because a PSD matrix with ones on its diagonal has every entry between $-1$ and 1.
At an end of that interval $G$ is singular: if it were positive definite there, it would stay
positive definite under every small enough change of $t$, on both sides. Take $G$ at an end. A
singular PSD $G$ has a nonzero $c$ with $c^{\mathsf T}Gc=0$, and $c^{\mathsf T}Gc$ is the squared
length of $\sum_i c_iv_i$ for the four vectors $v_i$, so that sum is zero and the four vectors lie
in a space of at most three dimensions. Coordinates in an orthonormal basis of that space keep every
inner product, so $a_0,a_1,b_0,b_1$ can be taken in $\mathbb R^3$.

The next question is harder: how do directions become a prescription for a
quantum measurement? Start with two signed cards in a row. One move swaps them;
another changes the sign of the card in the right-hand place. Start at $(+,+)$.
Swap then change the right sign: $(+,-)$. Start again, change the right sign then
swap: $(-,+)$. Every sign is reversed. Try unequal numbers on the cards: the two
orders still give negatives of each other. ``Anticommute'' means exactly that
relation between the two moves. Do either move twice and everything is back.

Write the swap as $X$, the sign change as $Z$, and doing nothing as $I$:
\[
X=\begin{pmatrix}0&1\\1&0\end{pmatrix},\qquad
Z=\begin{pmatrix}1&0\\0&-1\end{pmatrix},\qquad I=\begin{pmatrix}1&0\\0&1\end{pmatrix}.
\]
For the larger construction, label four places by the two-bit words 00 through
11 and put a signed number at each. An $X$ in one bit-position swaps each pair of
places differing at that bit; a $Z$ reverses the signs at places whose bit there is 1.
The two positions give the following three moves, written as real $4\times4$
matrices:
\[
\Gamma_1=X\otimes I,\quad
\Gamma_2=Z\otimes X,\quad
\Gamma_3=Z\otimes Z.
\]
The tensor product has entries $(A\otimes B)_{(i,k),(j,l)}=A_{ij}B_{kl}$.
Each $\Gamma_i$ is symmetric and squares to the identity. For $i<j$, the first
position where the two products differ has $X$ against $Z$, and $XZ=-ZX$; all
other positions commute. So $\Gamma_i\Gamma_j=-\Gamma_j\Gamma_i$.
For a unit vector $a$, the symmetric matrix $A(a)=\sum_i a_i\Gamma_i$ then
satisfies $A(a)^2=\sum_i a_i^2 I=I$. Its two-outcome measurement has projectors
$(I+A(a))/2$ and $(I-A(a))/2$: they are symmetric, square to themselves, are
orthogonal, and sum to the identity.

Take
$\lvert\Phi_4\rangle=4^{-1/2}\sum_{k=1}^{4}\lvert k\rangle\otimes\lvert k\rangle$.
Expanding this sum gives
\[
\langle\Phi_4\rvert A\otimes B\lvert\Phi_4\rangle
=\frac14\sum_{i,j}A_{ij}B_{ij}=\frac14\operatorname{tr}(A^{\mathsf T}B).
\]
For $i\ne j$, cyclicity of the trace and anticommutation give
$\operatorname{tr}(\Gamma_i\Gamma_j)=-\operatorname{tr}(\Gamma_i\Gamma_j)=0$;
for $i=j$ the trace is 4. So the correlator of $A(a)$ and $A(b)$ is
$a\cdot b$. Give Alice $A(a_0),A(a_1)$ and Bob $A(b_0),A(b_1)$.
Their strategy realizes $x,y,z,w$. All three $\Gamma_i$ have trace zero, so the
one-party expectations vanish too. The same real matrices and state are also a
complex strategy.

With complex amplitudes two levels on each side are enough. Add
$Y=\begin{pmatrix}0&-i\\ i&0\end{pmatrix}$ to $X$ and $Z$: the three are self-adjoint, square
to $I$ and anticommute in pairs, and they have trace zero. Give Alice
$a_1X+a_2Y+a_3Z$ for her vector $a$, and Bob $b_1X-b_2Y+b_3Z$ for his vector $b$. The expansion
above holds for complex entries too, with $\Phi_2$ and a factor $\frac12$, and the transpose of
Alice's matrix is $a_1X-a_2Y+a_3Z$, since $Y$'s transpose is $-Y$. So the correlator is
$\frac12\operatorname{tr}\bigl((a_1X-a_2Y+a_3Z)(b_1X-b_2Y+b_3Z)\bigr)=a\cdot b$.

What would that mean at a bench? Prepare the stated shared state, choose one of
the two specified measurements on each side, and write down the two signs returned.
For each pair of settings, count agreements as $+1$ and disagreements as $-1$.
Their average is the correlator. The model predicts the four averages $x,y,z,w$;
a finite tally fluctuates around its expectation. The signed-card moves let you
check the measurement algebra by hand. Realizing that state and those measurements
in a device, and comparing its tallies with the predictions, is the physical experiment.

Conversely, in a normalized quantum state $\psi$ with self-adjoint measurements
$A_x^2=B_y^2=I$ on separate systems, put
$a_x=(A_x\otimes I)\psi$ and $b_y=(I\otimes B_y)\psi$. They are unit vectors and
$\operatorname{Re}\langle a_x,b_y\rangle=E_{xy}$. Replace each complex coordinate by
its real and imaginary parts. This preserves those real inner products, and their
Gram matrix is a completion. A mixed state gives the same construction by taking
the direct sum of the vectors of its pure components, weighted by the square roots
of their probabilities. So the completion condition and finite-dimensional quantum
realizability are equivalent.

The necessary direction does not require finite-dimensional vectors. If a positive
normalized expectation $\omega$ on an algebra has self-adjoint $A_x,B_y$, with
$\omega(A_x^2),\omega(B_y^2)\le1$, then the real symmetric matrix with entries $\operatorname{Re}\,\omega(X_iX_j)$ is PSD: its quadratic form is $\omega(C^*C)\ge0$ for the real linear combination $C$
of the four observables. Increasing each diagonal entry to 1 preserves PSD. When
Alice's observables commute with Bob's, its cross entries are precisely the real
correlators. This supplies the same $G$. Including the unit as a fifth observable
also gives the level-1 moment matrix with its marginals. The sufficient direction
already used only two four-level systems.

\subsection{The arcsine boundary, in both directions}


Keep Bob's two directions hinged at one end. Alice's first direction has to make
a specified angle with each. How far can Bob's hinge open while that is possible?
Hold Alice's direction upright and sweep each of Bob's around it at its required
angle. Point them towards the same side to close the hinge as far as it will go;
point them towards opposite sides to open it as far as it will go. There is a
whole interval of openings between those limits. Alice's second direction asks
for another interval.

Draw both intervals on one strip marked from 0 to 180 degrees. If they miss, no
hinge setting satisfies both. If they overlap, choose an opening in the overlap
and construct each Alice direction at that opening. The algebra below permits
each its own perpendicular component in four dimensions. The pencils show why
the shared opening is the question; they need not fit the entire construction
on a tabletop. In the impossible request above, one interval is just 0 degrees
and the other just 180. There is nowhere to put the hinge.

Assume $x,y,z,w\in[-1,1]$, and put
$s_x=\arcsin x$, with analogous notation for the other three entries. A unit vector
making angles $\alpha,\beta\in[0,\pi]$ with two unit vectors can exist at an angle
$\delta$ between those two exactly when
\[
|\alpha-\beta|\le\delta\le
\min(\alpha+\beta,2\pi-\alpha-\beta).
\]
To see this, decompose the two vectors along and perpendicular to the first. Their
inner product ranges from $\cos\alpha\cos\beta-\sin\alpha\sin\beta$ to
$\cos\alpha\cos\beta+\sin\alpha\sin\beta$, by Cauchy--Schwarz; every value in that
interval is attained by choosing the angle between their perpendicular components.
Taking arccosines gives the displayed interval, including zero perpendicular components.

Apply this to Alice's first vector and Bob's two. Since
$\alpha=\pi/2-s_x$ and $\beta=\pi/2-s_y$, Bob's angle must lie in
\[
I_0=[\,|s_x-s_y|,\ \pi-|s_x+s_y|\,].
\]
Alice's second vector gives
$I_1=[\,|s_z-s_w|,\ \pi-|s_z+s_w|\,]$.
Each interval is nonempty because $|a-b|+|a+b|=2\max(|a|,|b|)\le\pi$
for $a,b\in[-\pi/2,\pi/2]$. They intersect exactly when
\[
|s_x-s_y|+|s_z+s_w|\le\pi,\qquad
|s_z-s_w|+|s_x+s_y|\le\pi.
\]
Using $|a|+|b|=\max(|a+b|,|a-b|)$, these are the four inequalities
\[
\begin{aligned}
|s_x+s_y+s_z-s_w|&\le\pi,&|s_x+s_y-s_z+s_w|&\le\pi,\\
|s_x-s_y+s_z+s_w|&\le\pi,&|-s_x+s_y+s_z+s_w|&\le\pi.
\end{aligned}
\]
So a completion implies all four.

For the converse, choose any $\delta\in I_0\cap I_1$ and set
$b_0=e_1$, $b_1=\cos\delta\,e_1+\sin\delta\,e_2$.
If $0<\delta<\pi$, put
\[
\begin{aligned}
c_0&=\frac{y-x\cos\delta}{\sin\delta},&
c_1&=\frac{w-z\cos\delta}{\sin\delta},\\
a_0&=xe_1+c_0e_2+\sqrt{1-x^2-c_0^2}\,e_3,\\
a_1&=ze_1+c_1e_2+\sqrt{1-z^2-c_1^2}\,e_4.
\end{aligned}
\]
The interval conditions say exactly that the two radicands are nonnegative:
after multiplication by $\sin^2\delta$, the first is
$1+2xy\cos\delta-x^2-y^2-\cos^2\delta\ge0$, and likewise for $z,w$.
These four vectors have unit length and the prescribed cross inner products.
If $\delta=0$, the intervals force $x=y$ and $z=w$; take $b_0=b_1=e_1$,
$a_0=xe_1+\sqrt{1-x^2}e_3$, $a_1=ze_1+\sqrt{1-z^2}e_4$.
If $\delta=\pi$, they force $y=-x$ and $w=-z$; use the same Alice vectors and
$b_1=-b_0$. Their Gram matrix supplies the completion in every case.
Together with the preceding construction, this proves the stated quantum characterization.

\subsection{The finite local model, in both directions}


Make sixteen cards, each with four boxes labelled $A_0,A_1,B_0,B_1$, and put every
possible arrangement of plus and minus signs on them, one arrangement per card.
Draw a card before choosing which setting to ask each player about. Alice reads
her chosen box; Bob reads his. Neither needs the other's setting. How often each
card is drawn is the only choice left in this model.

Try to get perfect agreement for the first three pairs and perfect disagreement
for the fourth. Any card that ever gets drawn would need $A_0=B_0$, $A_0=B_1$
and $A_1=B_0$, which makes $A_1=B_1$ too. Changing how often you draw the cards
cannot fix it. For less extreme requests the question is quantitative: can you
assign nonnegative weights to the cards, adding to one, that give all four
requested averages? The following construction supplies those weights exactly
when all eight CHSH inequalities hold and the correlators are in $[-1,1]$.

There is a useful way to assemble the deck. First allow only Bob's first question
and make a distribution of the three answers $A_0,A_1,B_0$. Then do the same for
his second question. Both smaller decks must draw Alice's answer pairs in the
same proportions. The adjustable number $t$ is Alice's agreement-minus-disagreement
average across her two boxes. Choose it in both intervals
\[
J_0=[\,-1+|x+z|,\ 1-|x-z|\,],\qquad
J_1=[\,-1+|y+w|,\ 1-|y-w|\,].
\]
Each interval is nonempty. They overlap exactly when
$|x+z|+|y-w|\le2$ and $|y+w|+|x-z|\le2$, which expand to the eight
CHSH inequalities. For $a_0,a_1,b\in\{-1,1\}$ set
\[
\begin{aligned}
P_0(a_0,a_1,b)&=\frac{1+t a_0a_1+x a_0b+z a_1b}{8},\\
P_1(a_0,a_1,b)&=\frac{1+t a_0a_1+y a_0b+w a_1b}{8}.
\end{aligned}
\]
When $a_0=a_1$, nonnegativity of $P_0$ is $1+t\ge|x+z|$;
when $a_0=-a_1$, it is $1-t\ge|x-z|$. These are precisely $t\in J_0$.
The same check for $P_1$ uses $J_1$. Summing signs shows that both distributions
have mass 1 and the desired pair correlators. Their common Alice marginal is
$P_A(a_0,a_1)=(1+t a_0a_1)/4$.

To join the two smaller decks, first draw Alice's pair with its common probability.
Among the first deck's cards bearing that pair, draw Bob's first answer; among
the second deck's cards bearing the same pair, draw his second. Those conditional
draws give a four-answer card. A pair that never occurs needs no card. Written
as probabilities, this procedure is
\[
P(a_0,a_1,b_0,b_1)=
\begin{cases}
P_0(a_0,a_1,b_0)P_1(a_0,a_1,b_1)/P_A(a_0,a_1),&P_A(a_0,a_1)>0,\\
0,&P_A(a_0,a_1)=0.
\end{cases}
\]
If $P_A=0$, nonnegativity forces both corresponding triple distributions to be zero.
Otherwise summing over either Bob answer cancels the denominator. So $P$ has the
two triple distributions as marginals and total mass 1. Sample one of its sixteen
answer quadruples and let each party answer with its assigned sign. This is a finite
mixture of deterministic local plans with exactly the four required correlators.

The other direction is one line per card. On a single card,
$A_0(B_0+B_1)+A_1(B_0-B_1)$ is $\pm2$, because one bracket is 0 and the other is $\pm2$.
Each of the other seven CHSH expressions is this one with the boxes relabelled or a sign
flipped, so it is $\pm2$ on every card too, and any weighting of the cards keeps all eight
averages between $-2$ and 2 (Theorem~\ref{thm:chsh-fine} is the first of them). So a
deck exists exactly when the correlators are in $[-1,1]$ and the eight inequalities hold.

\subsection{Scope of the constructions}


The constructions above pin down the correlators. What follows is where they stop: they
aren't a completeness claim for the integer check, they say nothing about full behaviours
with marginals, and they say nothing about data from a lab.

The arcsine boundary. For two settings and two outcomes per side, Tsirelson, Landau (1988) and
Masanes (2003) give theorems characterizing the correlators reachable by quantum measurements:
$|\arcsin E_{00} + \arcsin E_{01} + \arcsin E_{10} - \arcsin E_{11}| \le \pi$ and its sign
variants. The characterization holds as an if and only if with the
completion set and, through Tsirelson's representation, with the quantum
set. The point $(3/5, 4/5, 4/5, -3/5)$ is on the boundary:
$\arcsin \tfrac35 + \arcsin \tfrac45 = \tfrac{\pi}{2}$ by the 3-4-5 triangle, so the sum is $\pi$; Tsirelson's point is on it too.
The necessary direction holds in any dimension, for any expectation on an algebra of
observables in which Alice's commute with Bob's.

Equality with the quantum set. That the completion set of Section~\ref{sub:chsh-elliptope} is exactly the set of
quantum correlators is Tsirelson's representation theorem, and it is proved in finite dimension:
for both real and complex amplitudes, with two four-level systems enough for every point, and two
two-level ones with complex amplitudes. That PSD of the level-1 moment matrix is necessary for
quantum behaviour is proved for strategies in finite dimension, with real amplitudes
and with complex ones, and in any dimension, marginals included, for any expectation on
an algebra of observables.

Achievability. Theorem~\ref{thm:chsh-tight} is a point of $\mathbb{R}^4$. A state and
measurements that reach it are explicit: the two-qubit state
$(|00\rangle + |11\rangle)/\sqrt 2$ with $A_0 = Z$, $A_1 = X$, $B_0 = (Z + X)/\sqrt 2$ and
$B_1 = (Z - X)/\sqrt 2$ has correlators $1/\sqrt 2$, $1/\sqrt 2$, $1/\sqrt 2$, $-1/\sqrt 2$ and
$S = 2\sqrt 2$, and every strategy in finite dimension has
$S^2 \le 8$, with real amplitudes and with complex ones. Both are arithmetic about the model.
How close a real laboratory comes, and what that costs in loopholes closed, is the story told in
Appendix~\ref{sec:chsh}.

Completeness of the checks. The 1+AB cascade and the elliptope gate are proved sound. The cascade has no proved converse; the gate's factorization branch accepts every rational PSD completion, by the row-peeling of Section~\ref{sub:four-vectors} with each pivot kept. The level-1 integer check is the exception: on the small machine the same
seven integer facts hold exactly when every pair was sampled and the pinned matrix is positive
semidefinite (Theorem~\ref{thm:chsh-exact}), so that check is
complete for what it tests. The 1+AB cascade is strict and misses boundary points.
The two 1+AB checks don't sit inside one another. The full check accepts the $12/5$ tally above, which the zero-moment check refuses, and it refuses the tally with $E_{00}=1$ and the other three correlators 0, where $H_{11}=0$, which the zero-moment check accepts.

Finite levels. For full behaviours, marginals included, the NPA levels are relaxations.
Chaturvedi (arXiv 2607.14569, 2026) argues, in a preprint, that no finite NPA level characterizes the complete
quantum set even in this simplest Bell scenario. I haven't gone through that proof, and nothing
here rests on it.

Statistics and physics. Nothing here says how many rounds make a tally informative, and
nothing says the counters came from honest play; a program can INC whatever counts it likes.

\subsection{Prior work}


Tsirelson's 1980 paper gives the bound $2\sqrt 2$ for quantum correlations, and his theorem
says that quantum two-outcome correlators are inner products of real unit vectors. Landau (1988)
and Masanes (2003, arXiv quant-ph/0309137) each give a condition that, by their proofs,
characterizes the quantum correlators exactly for two settings and two outcomes. Wehner (2006,
arXiv quant-ph/0510076) gets Tsirelson bounds, CHSH among them, from a semidefinite program on
the $4 \times 4$ Gram matrix of those vectors. Navascués,
Pironio and Acín (2007, arXiv quant-ph/0607119; 2008, arXiv 0803.4290) built the hierarchy of
moment-matrix tests; their 2008 paper applies a Schur complement to the level-1 matrix of this
scenario, recovers the arcsine condition, and discusses the 1+AB level. I went back to each of
those papers for these points, and that is what I found in them. The results are theirs. Two of
them are proved again here, in finite dimension: Tsirelson's representation
and the arcsine characterization; the rest of
their proofs are not redone. Echenim and Mhalla give a detailed derivation of Tsirelson's
bound (arXiv 2306.12535), and exact rational
sum-of-squares certificates go back at least to Peyrl and Parrilo (2008). So the
moment-matrix route to $2\sqrt 2$ is standard. What is specific here is narrower: an integer
certificate computed from outcome counts, and the argument from ``the counts pass'' to PSD
of the pinned matrix and on to $S^2 \le 8$. I did not find that chain elsewhere, and I make no
claim beyond it.
On the host machine, the same chain is earned: a raised flag from a clean start
stands on a passing check of the packed counts, and the check passes exactly when the pinned matrix
is positive semidefinite and every pair was sampled (Appendix~\ref{sec:chsh},
Theorems~\ref{thm:chsh-exact} and~\ref{thm:chsh-earned}). The converse direction of the integer check,
integers from PSD, is proved there too, and there too I make no claim beyond the chain.

%===========================================================================


%===========================================================================
\section{From the front matter}
%===========================================================================
That sentence walks straight into one of the oldest fights in mathematics, so I'll take a short
trip into it first. G. H. Hardy, a Cambridge number theorist who wrote \emph{A Mathematician's Apology} in 1940 to
defend doing mathematics for its own sake, came down on
the found side about as hard as anyone ever has: mathematical reality, he wrote, lies outside
the people who study it, and the theorems they like to call their creations are really just
notes on what they saw there. The German mathematician Leopold Kronecker supposedly split the difference: God made the
whole numbers, and everything else is the work of people. Supposedly, because push on that
line and it gets thin fast. It comes down through Heinrich Weber, another German mathematician, in a memorial piece written
after Kronecker died, pointing back to a lecture from 1886, and that's the only source I could find anybody citing. It's a great line with one witness. (The record chapter of Part~I, Section~\ref{sec:toll}, stops
early on to ask how many witnesses a claim ought to need. The answer is more than one.) None of this
settles whether I found the Thiele Machine or made it up. A proof doesn't settle
it either: showing what follows from the premises doesn't decide whether you found the
theorem or made it. That's why it's labelled a conviction. What it changes is how I hold
the thing: like something that was already there, which means I can be wrong about its shape
and the shape won't care. And to be exact about what ``found'' covers, it covers the parts the counting
forces on you, the latch shape and the merge or the fresh states. The event to record, the price list, and putting
the count inside the step are choices, mine, and the three lists below say which is which.

\subsection*{The four steps of the argument}

The four steps above are the argument. The proofs establish the counts and constructions; the
claims section at the back lists them with what each stands on.

Argued only: that the structure a computation makes is kept distinctions, priced by the toll, while
a reading that undoable moves keep never rises and is never charged; that's a choice of what to
price, and on a finite machine the second half is proved.

Argued only: that an account of computation \emph{should} carry what a computation made the case, and
what it cost, in how it steps. The proofs show that it hangs together, what it follows from, and what
it costs; whether it ought to be adopted is a claim about what matters, which no theorem decides.

My proposed reason for choosing certification is that its records get copied and relied on
by parties who didn't produce them. That is the pointer thesis in Part~III. It connects the
machine to how people use records; the counting arguments work with whichever record is chosen.

The results also have a history. By my own sorting at the back (Appendix~\ref{sec:whats-mine}), each piece has its neighbours standing
right there. Start from the top. The small machine and its universal programs carry the reference
monitor's every-program move (Anderson in 1972) to a record that's priced and earned; the two-player
game check in the appendices (Appendix~\ref{sec:chsh}) is careful engineering on known mathematics.
The latch has the 1918 flip-flop and write-once memory for neighbours. The window fact is an old fact
about functions, from Edward Moore's 1956 thought experiments and Rudolf Kalman's observability, with a
price put on it. The toll derived from finiteness, permanence and a price on merging is a ledger form
of the core underneath. And that core, what step one says keeping takes, is old mathematics: the pigeonhole principle,
Dedekind's 1888 definition of infinite and the finite heart of Ernst Zermelo's 1896 recurrence
argument, with Landauer and Bennett on merging and Thomas Ouldridge and David Wolpert's 2023 price in
entropy on the merges a finite automaton can't avoid. For deterministic machines the arrangement
is old too: Ming Li and Paul Vitányi (1996) wrote that a bounded memory forces erasure in the long
run, and Demaine and his coauthors (2016) that extra space only postpones it. What I'd defend most
is there, at the bottom, in the form I didn't find: one law with three columns and no fourth, stated
alike for deterministic, noisy and quantum steps, with two small exact lemmas, the count of fresh states
for a bounded memory and every scale paying the same way, and a toll on a reading that stays, which charges for keeping whichever column
pays for it.

The route through the book is the route through the model: follow a step, see which
distinctions survive it, and count what the next step has to work with. Then put the record
and its price into the step itself. The example at the end shows those choices can all be made at once.


%===========================================================================
\section{From ``A distinction that stays''}
%===========================================================================
The walker is the case proved here. Other finite-time results use other,
explicitly specified dynamics. Drag a bead held in a spring-like trap through water from one place to another in a set
time and the least work you can get away with is the drag times the distance squared, divided by
the time plus twice the drag over the stiffness, which at long times is the drag times the distance
squared over the time (Schmiedl and Seifert,
2007, who worked out the best protocols for doing it in finite time). Within their overdamped diffusion or Markov jump models, change a probability
distribution into a different one in a set time and the entropy produced is at least a fixed fraction of the
square of how far apart the two spreads are, divided by the time and by how mobile the states are: the diffusion, for a bead that drifts (Aurell, Mej\'ia-Monasterio and Muratore-Ginanneschi, 2011), and how busy the jumping is, for states that jump (Shiraishi, Funo and Saito, 2018). Those are
their results, and I cite them as physics. These bounds require their stated dynamics, a nonzero change of distribution,
and finite activity or mobility. For the jump-process bound, local detailed
balance identifies the usual thermal entropy production; without it the cited
result uses a different entropy-production quantity. A thermal environment alone
does not supply these assumptions. The bounds concern entropy produced, not
positive heat on every step: a step that leaves the spread exactly as it was
owes nothing on these bounds, and the heat a single step gives off can be negative when the device's
own entropy goes up. An idealised machine cut off from everything, frictionless billiard balls
bouncing their way through a computation (Edward Fredkin and Tommaso Toffoli's, 1982), could in
principle run without dissipating at all; Bennett (2003) wrote that such ballistic machines compute
at constant velocity without dissipating any energy, and that they must be kept isolated from heat
baths and noise. Keeping one on course against noise means correcting its errors, and whether that
brings the dissipation back I don't settle here; in 2003 Bennett called the thermodynamics of error
correction a subject that appears to be in its infancy. The bead of Jun, Gavrilov and Bechhoefer
(2014) shows the other half. Put through the same moves as an erasure but with no tilt, so
nothing is squashed, its work extrapolated to the slow limit came out at $0.05 \pm 0.03\,kT$,
within two error bars of zero: no floor under it, which is what the premise says.


I also built a numerical bead to see the steps, using a quartic double well.
It follows the experiment's sequence: lower the barrier for one second, ramp
the tilt, remove the tilt, then raise the barrier for one second. This is a
chosen potential and a numerical approximation, not the measured trap. The
starting positions are sampled from a numerical approximation to the thermal
distribution in that potential, with equal chances of starting in either well.

Keep three accounts as it moves. Changing the trap while holding the position
fixed does work on the bead. Moving the bead in the fixed trap transfers heat
to the bath. Whatever is left changes the bead's potential energy. These are
the steps the script actually adds up:

\begin{center}
\begin{tikzpicture}[thiele]
\node[pipe,text width=5.8cm,minimum width=0pt] (a)
  {Old trap\\Bead at its old position};
\node[pipe,text width=5.8cm,minimum width=0pt,below=1.1cm of a] (b)
  {Changed trap\\Same position};
\node[pipe,text width=5.8cm,minimum width=0pt,below=1.1cm of b] (c)
  {Same changed trap\\New position};
\draw[pipearrow] (a.south) -- node[right=8pt]{work in} (b.north);
\draw[pipearrow] (b.south) -- node[right=8pt]{heat out} (c.north);
\end{tikzpicture}

\small Work in minus heat out equals the change in stored energy.
Either transfer can be negative on an individual step.
\end{center}

The work account includes removing the tilt and restoring the starting trap.
The heat account uses the potential difference during each displacement.
Those two accounts balance on every simulated path to rounding precision.
That checks the accounting; it does not show that the time steps are small
enough to reproduce continuous motion exactly.

Here is what came out, with 3,000 paths at each setting and a time step of
half a millisecond. Each setting uses a specified tilt, chosen earlier to
make roughly nine erasures in ten work. Success means ending in the right well.
The last two columns average only those successes; energies are in units of
$kT$. The time column is the tilt ramp, with the two one-second barrier ramps
in addition.

\begin{center}\small
\begin{tabular}{rrrr}
\toprule
Tilt time (s) & Success rate & Heat out on success & Work in on success\\
\midrule
2  & 0.909 & 2.517 & 2.501\\
5  & 0.916 & 1.466 & 1.490\\
10 & 0.913 & 1.073 & 1.095\\
20 & 0.899 & 0.830 & 0.907\\
40 & 0.893 & 0.773 & 0.782\\
80 & 0.869 & 0.693 & 0.719\\
\bottomrule
\end{tabular}
\end{center}

Heat and work are close in some rows and less close in others. The difference
is measured by the simulation, not supplied as a constant correction. At the
last setting the sampling standard errors are about $0.023\,kT$ for heat and
$0.011\,kT$ for work. They do not include error from the time step or the model.

Now try to make a floor out of those six points. Fit the successful heat to a
constant plus a number divided by the tilt time: the constant is
$0.674\,kT$, with sampling error about $0.015\,kT$. Halve the time step and
repeat the six settings: it becomes $0.732\,kT$, with sampling error about
$0.016\,kT$. The second run also samples different paths, so this comparison
mixes sampling and numerical error. It is enough to stop me treating a fitted
constant near $\ln 2$ as a confirmation. The barrier ramps are still one second
long, the tilt changes between settings, and this fit does not take the whole
protocol to a slow limit.

There is a useful theoretical comparison, but it is a comparison with work.
For the continuous thermal model, beginning in equilibrium and ending with the
same symmetric trap, the detailed Jarzynski relation gives the successful
paths a mean-work bound $kT\ln(2p)$, where $p$ is the success probability.
Averaging the successful and failed paths gives the bound
$kT[\ln 2+p\ln p+(1-p)\ln(1-p)]$. At nine successes in ten these are about
$0.588\,kT$ and $0.368\,kT$. B\'erut, Petrosyan and Ciliberto's 2013 paper,
\href{https://arxiv.org/abs/1302.4417}{Detailed Jarzynski Equality applied to a Logically Irreversible Procedure},
works through the success-conditioned relation. Neither formula is a heat
bound for this simulation without also accounting for its change in stored
energy. Nor does an exact bookkeeping identity make the discrete simulation
obey every continuum fluctuation relation exactly.

The script is \texttt{scripts/bead\_erasure\_sim.py}. It reports heat, work,
stored-energy change, sampling errors and the descriptive fits, and accepts a
smaller time step for checking them. The slower settings here use less mean
work and release less mean heat. That is what this calculation demonstrates.
It does not determine a universal floor or reproduce the laboratory's fitted
constant.

\begin{longtable}{p{0.28\textwidth} p{0.64\textwidth}}
A merge, the floor & Jun, Gavrilov and Bechhoefer (2014): one bead in a feedback trap.
Squash two wells to one: $0.71 \pm 0.03\,kT$ of work, in the slow limit. Run the same
moves without the tilt, so nothing is squashed: $0.05 \pm 0.03\,kT$, also in the slow limit, within two error bars of zero. $\ln 2$ is $0.693$. Same bead, same trap. \\[3pt]
The floor, three more ways & B\'erut and others (2012), a colloidal bead; Hong, Lambson, Dhuey
and Bokor (2016), nanomagnetic memory bits, within two standard deviations of Landauer;
Gaudenzi, Burzur\'i, Maegawa, van der Zant and Luis (2018), molecular spins at one kelvin. \\[3pt]
Lopsided wells & Gavrilov and Bechhoefer (2016), the feedback trap of Jun and others, with an asymmetric double well:
the average work to erase came out below $kT\ln 2$. The floor as stated is for a symmetric memory. \\[3pt]
Time against heat & B\'erut and others (2012; B\'erut, Petrosyan and Ciliberto 2015): heat
$= kT\ln 2 + B/\tau$ fits with $B = 8.15\,kT\cdot\mathrm{s}$; left free, the floor of the fit
comes out at $0.72\,kT$. The floor is reached only as the time goes to infinity. \\[3pt]
The distinction itself & The same bead: the paper estimates a characteristic
Kramers escape time of about 3,000 seconds at $8\,kT$, and about ten seconds
at $2.2\,kT$. Those are statistical time scales. The barrier height controls
retention; it is not itself a measured heat payment. \\[3pt]
The loop & Rauer, Erne, Schweigler, Cataldini, Tajik and Schmiedmayer (2018): recurrences of
coherence in an isolated gas of thousands of atoms. A collective quantity of a closed system comes
back around; the full state of the gas wasn't shown to return, and couldn't be. \\[3pt]
Heat of collapse & Vinante and others (2017): a cantilever at millikelvin shows a
non-thermal force noise of unknown origin, compatible with the collapse heating one model
predicts. Donadi, Piscicchia, Curceanu, Di\'osi, Laubenstein and Bassi (2021): the radiation a
gravity-driven collapse would emit was not seen at Gran Sasso, ruling out that model's
parameter-free form. These test the noise a collapse model adds, not a Landauer price on a kept
outcome. \\[3pt]
The ceiling on fresh states & Bekenstein's proposed bound uses energy and size; it is not a measured state count. \\
\end{longtable}

The cited erasure experiments test specified memories and protocols, including
symmetric-memory Landauer bounds and heat versus duration. They are not an
exhaustive survey of physical computation, nor an experimental proof of the
abstract state-space dichotomy. \href{https://doi.org/10.1103/PhysRevD.23.287}{Bekenstein's bound}
relates entropy to energy and size under its physical assumptions; it does
not supply a finite enumeration of arbitrary device states. The finite-state
theorems here take their finite set as a mathematical hypothesis.

%===========================================================================
\section{From ``A record and a toll''}
%===========================================================================
\subsection{Clay, seals, and seven witnesses}

Start a little over five thousand years ago, in the Near East, with people who counted sheep with
little clay tokens. To send a count across a distance without the courier changing it, they
put the tokens inside a hollow clay ball, pressed the shapes of the tokens onto the outside
while the clay was soft, and sealed it. The archaeologist Denise Schmandt-Besserat has
spent a career on these, and she puts the practice at around 3300 BC; the ball is called a bulla. The push is one you can do with
a cookie tin, with a paper seal glued across the lid and rim. Let it dry, then try
opening the lid without tearing or disturbing the seal. The bare lid would come
straight off; it is the seal that makes this a test. A substitute seal could fool
you, so this little trial shows a disturbance, not an unforgeable history. In the
clay ball, the outside shows the count inside and opening the ball breaks it.
That is the property the tin and seal are trying to reproduce.
The idea is a real one and it will show up again in a form you won't expect: a record is worth
something only if it can't be silently undone.

Now law, which is where people get serious about who's allowed to say what happened. The
book of Deuteronomy, chapter nineteen: nobody is convicted on the word of one witness; it
takes ``two or three.'' The Roman jurist Gaius, writing in the second century, says a will the
praetor will honour is one sealed by seven witnesses. Why seven? Gaius doesn't say. The
legal historian Henry Maine, in 1861, had an answer: the older Roman will was a mock sale, done in
front of five witnesses, a man holding the scales, and the buyer of the estate, seven people in
all, and when the ceremony fell away the seven seals kept the count. Maine's is the only account
of it I found, so treat it as his. Gaius's own text does show they kept those old formalities
around a will, and it isn't hard to see why anyone would: one person can lie, and a will moves a
fortune, and the dead man can't object. So the record isn't the paper; it's the number of independent people who have put their own
seal, their own name, their own risk, next to the claim. Try it: write ``paid five dollars,
signed A'' on a napkin and ask a friend to forge it. It takes them ten seconds. Now require two
other people to initial it at the moment of payment and each keep a copy. If the readers compare those independently kept copies and trust their
holders, changing your napkin alone will be exposed. Stolen copies, forged
initials, collusion or a reader who never compares them can defeat the test. The claim didn't change,
but what stands behind it did.

Which brings me to seals, and the medieval historian Michael Clanchy, whose book is called
\emph{From Memory to Written Record}. England between the Norman Conquest and the early 1300s
was learning to trust writing. Until then it
had trusted people who remembered. The lesson was painful: seals were forged, charters were forged, and historians today have a whole craft of
telling real charters from fake ones because the fakes are so many. A seal is only as good as
whoever holds the stamp. Keep that one too, because in 2011 a certificate authority called
DigiNotar found out, at the cost of its existence, that a mathematically perfect signature is
worth exactly what the custody of the key is worth, and intruders had taken control of the
machines that signed. It is the same failure nine
hundred years apart, and I find that oddly comforting.

And then the obvious next question: if seals could be forged, what caught the forgeries? It wasn't
the seals. The most famous forged record in European history is the Donation of Constantine,
a document in which the emperor Constantine supposedly handed rule over Rome and the western
empire to the Pope. It was most likely written in the eighth century, some four hundred years
after Constantine died, and it was leaned on for centuries. Around 1440 a scholar named
Lorenzo Valla took it apart by its Latin: the language belonged to a much later age, and it
used words, ``satraps'' among them, that no Roman of Constantine's day would have used for his
officials. (Valla wasn't even
first. Nicholas of Cusa, a German churchman and philosopher who later became a cardinal, had
already called it a forgery in 1433, in \emph{De concordantia catholica}, a long book on how
the Church and the empire ought to hang together. I know that from the reference works, and I
haven't read the book. Valla's is the version people
remember, which is a pattern that keeps turning up.) Look at what kind of check that was.
The question wasn't ``is the stamp real.'' It was ``could these words have been written when they say they
were.'' A record can pass the first and fail the second, which is a hint that what I keep
calling one thing, a record, is doing more than one job.

\subsection{Two columns and a stamp}

The one I'd never have guessed mattered is bookkeeping. Around 1300, the ledger of a
Florentine company called Farolfi, kept at its branch in Provence, shows the earliest complete example anyone
has found of a trick merchants were working out: write every transaction twice, once where the money came
from and once where it went, and the two columns have to balance. The treasurers of the city of Genoa kept their ledger for
1340 that way, and accounting historians still count it among the most famous government
account books of the Middle Ages kept in double entry. And in 1494 a friar named Luca Pacioli printed the method in a mathematics book,
and it spread. The method is double-entry bookkeeping, and it's easy to take it for just a tidy
habit.

Do it for one day. Keep two columns of every dollar that moves, and total them at night. Now
have a friend secretly change one entry in one column. The totals disagree; you catch it. Now
have them change the matching entry in both columns. The totals agree, and you've caught nothing.
Push on that and you've found what double entry actually is: a check for slips, which a
careful liar walks straight through. Did Pacioli see that? I can only get at it secondhand. The scholar Forrester, in a paper on the University of Strathclyde's accounting-history pages, reads him as reporting that in many places the
books were taken to the local consul to be authenticated before a single entry went in. The
stamp is a witness again, from outside. The two columns are the discipline; the stamp is the
claim that someone other than you will stand behind the fact that these are the books.
That works only as far as the authentication and custody distinguish this
set of books from a substitute. Keep both halves separate in your head, because one is a method and the other is a
commitment.

And then I found the place where this got turned into a motto. The Royal Society, founded
in London in 1660 and still running, picked as its motto \emph{Nullius in verba}: take nobody's
word for it. (Where did they get it? From somebody's words, naturally. It's half a line of the
Roman poet Horace, from a verse letter to his patron, saying he isn't bound to swear by the
words of any master. So the society that takes nobody's word for it took its motto from a poet
and cut the line to fit. That's not a contradiction, and it's a nice picture of the method:
borrow what's useful, say whose it was, and don't swear by it.) The historian Steven Shapin wrote a paper about what that meant in practice,
and the answer is more interesting than the slogan. Robert Boyle, pumping the air out of a
glass globe to see what happens to a candle or a bird, had three problems: the machine, the
report, and the witnesses. He solved the third by having gentlemen present who'd sign their
names, and the second by writing the report so richly that a reader who wasn't there could
picture the scene and count as a witness anyway. Shapin calls that virtual witnessing, and he
notes, dryly, that the air pump was very expensive, seventeenth-century big science. So there
\emph{was} a cost; it just wasn't a cost of declaring the fact. It was the cost of the instrument
and of getting enough people into the room. Hold on to that distinction, because the rest of the book hangs on it.


\subsection{What they all agree on, and who they rule out}

Lay them down, and this time what shows through is not one thing. It's five different jobs
that the word ``record'' has been doing, and I'd been letting one word carry all five. Most of
them have a corpse next to them, and the corpses are different too.

Durability: the record can't be silently undone. The bulla, the write-once disk, the hash
chain and Bennett's card that has to be wiped all do this job. The corpse is the editable log, and the company
that kept its debts off the books: Enron hid its debt in side entities set up to keep it off
the balance sheet, which is how the Justice Department described the scheme in 2002, and its
auditor, Arthur Andersen, shredded the paper. Andersen was convicted for the shredding that
year, and the Supreme Court threw the conviction out in 2005 over how the jury had been
instructed; either way, the books concealed debts they should have shown. But notice that this job doesn't need a
witness at all. A copy will do: a hash chain nobody signed is still hard to change quietly, so
long as somebody else holds the latest link, because a chain with no copy anywhere can be
rewritten from end to end and recomputed, and nobody would know.

Authenticity: the record came from who it says it came from. Seals, the signature only a
key-holder can make, and the consul's stamp all do this one. The corpse is the forged charter and DigiNotar.
This job is about custody of a stamp, and it's a different job from durability; a perfectly
authentic entry can still be erased, and a perfectly durable one can still be a forgery.

Evidence: enough independent people stand behind the claim that lying is expensive. Two
witnesses, seven seals and Boyle's gentlemen do this one. It rules out anonymous assertion, and it rules out
``the universe keeps the receipt'', because information spread
into the environment is not a record anyone can read or rely on. This one is about numbers
of people and their risk. Paper has nothing to do with it.

Accounting: the record checks itself for slips. The two columns do this one. And one thing has to be kept
straight: the second column is \emph{not} an outside witness. It's you,
checking you. It catches a dropped digit. A careful liar gets past it, as you saw. It
belongs on this list as its own job, separate from the seven seals.

Cheap checking: a record whose correctness can be confirmed for much less than it cost to
make. Cook's certificates, the signature and the hash do this one. This one isn't universal,
and I thought it was. Whether checking needs to be cheaper than making depends on who you're
worried about and how much a wrong record costs you; plenty of records are checked by
redoing the whole thing, and that's fine when the thing is small.

Now for the thing none of them says. Not one of these jobs is
pricing the moment a record goes up and stays. Bennett prices merging and erasing, which that
moment might or might not be. Proof-of-work prices a search the designer chose. Bookkeeping prices
a discipline. Witnesses and seals price showing up and guarding the stamp. I looked, with
everything I had, and that's where it ran out. If there is a cost charged for that moment itself,
nobody in this lineage wrote it down, and that's either because it isn't there or because the
one person who did write it down is on a shelf I never found. The literal one-bit reset already has a thermodynamic analysis: Landauer's restore
to one. The question here is which premises force an instruction that raises a
record to merge full states, and which further premises put a price on that loss. So I'm going to have to try it myself, on the
smallest thing I can build, and let you watch. And I'll have to be honest about which of the
five jobs the small thing even attempts, because the toy has a permanent flag under its three moves, but no witnesses, authenticated identity or check of a claim.


Take an ordinary computing machine and put a lamp and tally beside it.
Let computing run as before, leaving both alone. A separate mark leaves the
computing state fixed, adds one to the tally and turns the lamp on when its
named event holds. Forget the lamp and tally and the old computation is
still there, step for step. Below I build that box, compare different ways
to implement it, and ask exactly which differences its questions require.
The old machine may already have prices or a reading of its own; this
construction supplies the new lamp and its zero/one account. With effective
rules an ordinary computer can store all three parts.


\subsubsection*{What the next move could tell us}

Put two boxes beside each other and read everything the task lets you read.
Suppose the answers agree. That is not yet enough to put the boxes in the
same pile. Press the same button on each. Perhaps one lamp now comes on
and the other stays dark. The difference mattered, even though the first
reading missed it.

Keep two states together only when every permitted finite continuation
gives the same required answer. Include doing nothing: the current answer
counts too. This tells us which differences the task can afford to forget.

\begin{theorem}[Keep what a later question can separate]\label{thm:future-reading-equivalence}
For a fixed deterministic machine, permitted instructions and required
observation, equality of observations after every finite continuation is
the unique coarsest equivalence that preserves the current observation and
is stable under every instruction. A deterministic implementation preserving
those steps and answers cannot identify two states separated by such a
continuation.
\end{theorem}
\begin{proof}
Agreement with oneself, reversal of an agreement, and joining two agreements
all hold one continuation at a time, so this is an equivalence. Doing nothing
checks the current observation. To check stability after a move, put that
move in front of each continuation. If two states were already grouped by
any observation-preserving, stable relation, apply its stability once for
every move in a continuation, then read the answer. Those states agree on
every continuation and belong to our equivalence. Our grouping therefore
contains every permitted grouping. Finally, if an implementation put two
separated states in one full state, the same continuation from that state
would have to produce two different answers. Determinism gives only one.
\end{proof}

Change the questions and the piles can change. A box that toggles between
two positions needs two distinguishable states if I ask its position. If
my only required answer is always ``ready'', and every later answer is also
``ready'', those positions can share one description. The task tells me
what has to survive.

\subsubsection*{The picture, the lamp and the tally}

Take an ordinary computing box. Put a lamp beside it and a tally beneath it.
A computing button moves the box by its old rule, costs zero, and leaves
the lamp and tally alone. A mark button leaves the box where it is, adds
one to the tally, and turns the lamp on if its named event holds there.
Once on, the lamp stays on. Start with the lamp down and the tally at zero.

Those instructions tell us what to do without telling us how to build the
box. One builder may keep every button press on a tape. Another may keep an
extra turn counter. I require each builder to show me the same three things:
the computing picture, the lamp and the exact tally. Its computing moves
must follow the base rule through that picture. Its marks must leave the
picture fixed and latch exactly the named event. Its prices are zero for
computing and one for marking; its reading obeys the toll and permanence;
its tally increases by exactly the charge on every move.

\begin{theorem}[Every such box has the same three-part description]\label{thm:lift-terminal}
Fix those base rules, mark events and accounting requirements. Every
implementation has exactly one map to the product of base state, Boolean
reading and natural tally that preserves the three readings. That map
preserves every instruction and every finite run. From a start with lamp
down and tally zero, it maps the reachable implementation states onto all
reachable states of the product lift with the matching base start.
\end{theorem}
\begin{proof}
Read the picture, lamp and tally, and write them down in that order. There
is no other choice for a map required to preserve those three readings.
Press a computing button. The picture follows the base rule. The lamp stays
put by the toll and permanence, and the tally stays put because the charge
is zero. Press a mark instead. The picture stays put, the lamp takes its old
value or the event's yes, and the tally rises by one. These are exactly the
two rules of the product lift. Repeat for every move in a finite run.

The clean starts agree in all three readings. Any reachable product state
has some finite list of buttons reaching it. Press that same list on the
implementation. The step argument makes its three readings the desired
product state. Every reachable product state therefore has a reachable
implementation state above it.
\end{proof}

Here are the three stages. First write a list of buttons. Then run it on
the builder's box. Then read the picture, lamp and tally. You get the same
result as running the list directly on the three-part description. Written
as arrows, after we have done it:
\[
\text{button lists}\ \longrightarrow\ \text{implementation states}
\ \longrightarrow\ \text{picture, lamp, tally}.
\]

The first arrow has only one possible meaning once the start and buttons
are fixed. The empty list gives the start; appending a button takes one
more step. This is the free history construction. The second arrow can
forget history. Its direction matters.

Turn a box halfway around with a free button. Turn it again. It is back
where it started, with the same lamp and tally. Now give another builder's
box a counter that goes up on every turn, without charging for it. After
two turns its counter is two larger. It meets all the requirements above:
the extra counter is separate from the tally of charges.

\begin{proposition}[Forgetting the turn count goes one way]\label{prop:lift-history-direction}
Even on reachable states, the three-part lift need not have a
step-preserving map into every implementation satisfying its contract.
\end{proposition}
\begin{proof}
In this example two turns return every product state to itself. A map
preserving turns would have to send it to an implementation state that
also returns after two turns. No such state exists there: its extra
counter has increased by two. Forgetting that counter works in the other
direction and preserves both turning and marking.
\end{proof}

There is an exact way to tell when the reverse map would work. Whenever
two button lists reach the same source state, they must also reach the
same target state. Otherwise the map would have to send that one source
to two targets. Supply one access list reaching each source state, and
this condition is sufficient: run its access list on the target. Any other
list reaching that source gives the same target by the condition. The
empty list preserves the start, appending a button preserves a step, and
reachability leaves no other possible map. The turn counter fails exactly
this test: the empty list and two turns agree in the source and disagree
in the target.

\begin{center}
\begin{tikzpicture}[thiele]
\node[pipe,text width=2.8cm,minimum width=0pt] (a) at (0,0) {Front\\Lamp down, tally 0};
\node[pipe,text width=2.8cm,minimum width=0pt] (b) at (4.2,0) {Back\\Lamp down, tally 0};
\node[pipe,text width=2.8cm,minimum width=0pt] (c) at (8.4,0) {Front\\Lamp down, tally 0};
\node[pipe,text width=2.8cm,minimum width=0pt] (d) at (0,-3) {Front; turns 0\\Lamp down, tally 0};
\node[pipe,text width=2.8cm,minimum width=0pt] (e) at (4.2,-3) {Back; turns 1\\Lamp down, tally 0};
\node[pipe,text width=2.8cm,minimum width=0pt] (f) at (8.4,-3) {Front; turns 2\\Lamp down, tally 0};
\draw[pipearrow] (a) -- node[above=6pt]{turn} (b);
\draw[pipearrow] (b) -- node[above=6pt]{turn} (c);
\draw[pipearrow] (d) -- node[below=6pt]{turn} (e);
\draw[pipearrow] (e) -- node[below=6pt]{turn} (f);
\draw[proj] (d) -- (a);
\draw[proj] (e) -- node[right=8pt,align=left]{forget\\turn count} (b);
\draw[proj] (f) -- (c);
\end{tikzpicture}

\small The top row returns to the same full state. The bottom row does not.
Forgetting the turn count respects each turn; recovering it from the top row cannot.
\end{center}

\begin{theorem}[The smallest description that answers these questions]\label{thm:lift-reachable-minimal}
Under the preceding contract, two implementation states give the same three
readings after every common finite continuation exactly when their current
three readings agree. If every two distinct reachable implementation states
can be separated by such a continuation, projection is a bijection onto the
reachable product states. It preserves starts and steps and is the unique
start- and step-preserving map there.
\end{theorem}
\begin{proof}
If the current readings differ, read them now. If they agree, the two
product states are equal. Equal states obeying the same rule give equal
next states; repeat for every move. Projection preserves those moves, so
the implementation's readings agree throughout too.

Now require every distinct reachable pair to be separable. Two states
projecting to the same triple could never be separated, so they must be
one state. Projection is therefore one-to-one; the previous theorem makes
it onto the reachable product states. Any start- and step-preserving map
is fixed on a reachable state by a list reaching it: map the start, then
follow the list. It must agree with projection there.
\end{proof}

The same reasoning explains uniqueness more generally. Take two reachable
machines with the same answers to every button list from their starts, and
suppose each keeps no two states that all future answers fail to separate.
Supply an access list for each state of the first. Run that list on the
second. If two lists reach the same first state, append any continuation:
their answers agree on the first machine, hence on the second. Separation
in the second makes their endpoints equal. The map is well-defined and
preserves steps and outputs. Reachability makes it onto; separation in
the first makes it one-to-one. Following an access list forces uniqueness.

For our box, that smallest description keeps the tally because I asked
for the tally. Stop asking for it and repeated failed marks may become
indistinguishable. Ask for every past button and the extra tape starts to
matter. With a clean start and an event that never holds, no reachable
state has its lamp up. Counting unreachable lamp-up states would answer
a different question. ``Smallest'' here means no two retained states are
indistinguishable by the required future readings, including on an
infinite machine.

This way of grouping states is the familiar behavioral equivalence behind
minimal automata. The initial history construction and final observation
construction are treated in Rutten's \emph{The Method of Coalgebra},
Chapters 2 and 4 and Section 8.6.\footnote{Jan Rutten, \emph{The Method of
Coalgebra} (2019), \url{https://www.cs.cornell.edu/courses/cs4810/2025sp/Handouts/Rutten.pdf}.}
Here the box tells us exactly which observations and operations the lift
preserves. Choosing the event, the window and the accounting is where the
application enters. Once they are fixed, the argument above determines
which implementation differences can be forgotten.


\subsubsection*{What can choose the question?}

Take a dial with two faces and one button. Each press turns it to the other
face. For this question, supply no starting face, no lamp and no names with
an outside meaning. Draw its two arrows. Now exchange the names on the two
faces. The arrows describe exactly the same machine.

Suppose those arrows alone are meant to tell us which face deserves yes.
A rule that does not depend on our choice of names must give the exchanged
face the same reading. So both faces get the same answer. A nonconstant
reading cannot be selected that way. Give one face a role---the loaded
start, a completed job, a passed check---and we have supplied more to work
with. The role has to be stated and defended. Renaming the arrows did not
supply it.

The argument is two lines: the exchange takes the first face to the second;
respecting that exchange makes their readings equal. Those are all the
faces. It does not make every machine's reading constant. It rules out a
nonconstant reading selected from the transition structure alone on every
machine. \footnote{Checked in \texttt{coq/audit/IntrinsicBoundary.v}:
\path{toggle_symmetry_forces_constant_reading} and
\path{no_nonconstant_renaming_invariant_toggle_reading}.}

That fixes what the uniqueness claim can say. The box, lamp and tally have
a unique observable realization under the contract already proved. The
arrows by themselves do not choose that contract. Adding the word
``intrinsic'' cannot do the choosing for us.

\subsubsection*{Does the picture keep enough of the device?}

Put a lamp on a box and a hidden switch inside it. The button copies the
hidden switch to the lamp. Start with the lamp dark in both boxes, and
put the hidden switch down in one, up in the other.

\begin{center}
\begin{tabular}{@{}lll@{}}
Lamp before & Hidden switch & Lamp after the button \\ \hline
Dark & down & dark \\
Dark & up & lit \\ \hline
\end{tabular}
\end{center}

A description that keeps only the lamp has one starting state, dark,
and is asked for two different next answers to the same button. It cannot
do it. Keep the hidden switch too and the problem disappears. This is a
reason the additional distinction is necessary for this task: the next
press can expose it. It is not a reason every task needs that switch.
\footnote{\path{hidden_bit_blocks_visible_only_dynamics},
\texttt{coq/audit/ObservationContract.v}.}

There is a complete test for this kind of description. First write down
the underlying configurations, the permitted buttons, every permitted
successor, the picture being retained, and the answers the task requires.
This is the operational model of the device. Then check two things.
Configurations with the same picture must give the same required answer
now. And the same button, from any configurations with that picture, must
leave the same next picture, for every permitted successor.

Both conditions are necessary. An answer decoder has only the picture to
read. A deterministic next-picture rule has only the picture and button.
Neither can return two different results on one input.

They are sufficient too, with explicit choices. Supply a representative
configuration for each picture and a permitted successor for each button
at each configuration. Decode a picture by reading its representative.
Compute its next picture by pressing the button on that representative
and taking the chosen successor. The two conditions ensure that every
other configuration and every other permitted successor agrees. Repeat
that argument at each step: all finite button lists preserve the required
answers. On a finite listed model this can be checked against its complete
table. For arbitrary state types the choices and their effectiveness are
separate data; the proof does not manufacture a search procedure.

This test allows several successors of an underlying configuration.\footnote{
\path{operational_adequacy_iff}, \texttt{coq/audit/ObservationContract.v}.} If two of them give different
required answers, no exact deterministic encoding of that same
configuration and button can represent both. \footnote{\path{branching_answers_exclude_exact_deterministic_encoding},\newline
\texttt{coq/audit/ObservationContract.v}.}
A probability distribution over outcomes, or a finer state including a
chosen input or seed, is a different contract and must be described as one.

The test starts with a model. Establishing that a real apparatus follows
that model still needs evidence about the apparatus: preparation, controls,
readout and the range over which the description is claimed. Passing the
mathematical test does not supply those measurements.

\subsubsection*{How long did we watch the lamp?}

Watch two boxes for three ticks. Both start dark and light on the first tick.
Both stay lit on the second and third. One stays lit forever. The other
has a counter inside that turns it off on the fourth tick.

\begin{center}
\begin{tabular}{@{}lll@{}}
Tick & Enduring lamp & Delayed reset \\ \hline
0 & dark & dark \\
1 & lit & lit \\
2 & lit & lit \\
3 & lit & lit \\
4 & lit & dark \\ \hline
\end{tabular}
\end{center}

The first three ticks do not tell the boxes apart. Pick any longer finite
watch and I can move the reset one tick beyond it. Both boxes can pay one
unit per tick and satisfy the basic toll. Only the first has permanence.
The permanent interface excludes the second by its rule, not by pretending
that the finite observation already excluded it.

The construction works for every finite horizon.\footnote{
\path{finite_observations_do_not_establish_permanence},
\texttt{coq/audit/IntrinsicBoundary.v}.} It is not a prediction that every
real record fails. It shows why a finite run alone cannot establish an
unrestricted forever claim. A proof from an adequate mechanism, or an
explicitly bounded reliability claim supported by measurements, has work
to do beyond reporting that the lamp stayed lit while we watched.


%===========================================================================
\section{From ``Why the toll isn't arbitrary''}
%===========================================================================
There is also a limited statement about how long a noisy record lasts, using a different
quantity: total entropy production during relaxation. Take a device with no outside hand:
a random process hopping among finitely many states at rates that never change, with a
reverse for every hop and local detailed balance, able to get from any state to any other.
Require at least one up-state and one down-state. Left to run, it settles into a long-run
spread $\pi$. Measure each up or down stretch from its entrance, with entrances sampled
in the stationary process, and call the mean lengths $T_{\mathrm{up}}$ and
$T_{\mathrm{down}}$. Start the device in a known no-state $s$. The total entropy it
produces over relaxation is at least
$\ln(1+T_{\mathrm{up}}/T_{\mathrm{down}})$, in units of $k_B$.

First compare the starting spread with the spread the device settles into. A state
that is common at the start but rare in the settled spread is a large mismatch.
For each state, take the ratio of those two chances, take its natural logarithm,
and average using the starting chances. With $p$ the starting spread, that number is
\[
D(p\Vert\pi)=\sum_{x:p(x)>0}p(x)\ln\frac{p(x)}{\pi(x)}.
\]
It is called relative entropy. Here every $\pi(x)$ is positive, so the ratios are
available. The number is nonnegative: apply $\ln y\le y-1$ to $\pi(x)/p(x)$,
multiply by $p(x)$, and add. The right side is at most zero, since the included
$\pi$ chances total at most one and the $p$ chances total exactly one.
Negating that comparison gives $D(p\Vert\pi)\ge0$.

For a known start $s$, all the starting chance is at $s$. There is only one term:
$D=\ln(1/\pi(s))=-\ln\pi(s)$. A state with settled chance one tenth therefore
starts with mismatch $\ln10$. Since $s$ is a no-state, its settled chance cannot
exceed the total chance of all no-states. Its mismatch is at least
$-\ln\pi(\mathrm{no})$.

The physical step in the argument is the fixed-rate entropy-production split of
Massimiliano Esposito and Christian Van den Broeck (2010). In that setting the total
entropy produced is at least the decrease in this relative entropy. Once the spread
has settled to $\pi$, the relative entropy left is zero. So the initial mismatch
supplies the lower bound for the whole relaxation.

Now read the same number from the record of up and down stretches. In a long settled
run, each entry into YES is followed by an exit; apart from the two ends of the record,
there are equally many up-stretches and down-stretches. Total time up is the number of
up-stretches times their mean length, and likewise for time down. Dividing gives
\[
\frac{\pi(\mathrm{up})}{\pi(\mathrm{no})}
 =\frac{T_{\mathrm{up}}}{T_{\mathrm{down}}}.
\]
The two chances add to one. If the ratio on the right is $r$, the no-share is
$1/(1+r)$, and its negative logarithm is $\ln(1+r)$. That gives the stated bound.
The settling and mean-length arguments below justify passing from a long record
to these limiting quantities.

For a large up-to-down ratio, doubling that ratio raises this lower bound by about
$k_B\ln2$. Making both stretches longer by the same factor leaves it unchanged. Across
a family of devices whose ratio tends to infinity, the bound diverges; no single finite
irreducible device in this class keeps its flag up forever. This is a relaxation bound,
not a heat charge on each flip. It uses neither a thermodynamic reverse-path construction
nor time-symmetric control, so it does not identify the quantities in those other results.
It also leaves the known-input case of the deterministic entropy-loss bound where it was:
this noisy evolution and its total entropy production are a different question.
The entropy is added up over one whole relaxation, from the known no-state $s$ until the device
has settled, and local detailed balance is what ties it to heat exchanged with the reservoirs.
Total entropy production includes the change in the device's entropy too; it is not heat alone.
It also isn't the permanence of the finite-state theorems on a
noisy device: on this device the flag comes back down, and $T_{\mathrm{up}}$ and
$T_{\mathrm{down}}$ are average stretches.

The core of it for a chain that ticks, worked out, runs as follows in the case where the long-run spread
$\pi$ gives every state a positive chance and satisfies detailed balance: in the long run as much
weight flows from $x$ to $y$ in one step as from $y$ to $x$, $\pi(x)P(x,y) = \pi(y)P(y,x)$, where
$P(x,y)$ is the chance of stepping from $x$ to $y$. Count the entropy produced in one step as
the change in the device's own entropy plus what goes out to the surroundings, which for a hop
from $x$ to $y$ is $\ln(P(x,y)/P(y,x))$. Detailed balance turns that ratio into
$\pi(y)/\pi(x)$, and the sum collapses: the entropy produced in one step is exactly the drop in
relative entropy to $\pi$. It's never negative, because a step that leaves $\pi$ where it is
can't make any other spread harder to tell from $\pi$ (that's the data-processing inequality, a
standard fact of information theory). Add the steps up and the sum telescopes. A known start $s$
has relative entropy $-\ln\pi(s)$, so the entropy produced in the first $T$ steps is $-\ln\pi(s)$
minus the relative entropy left at step $T$: it never exceeds $-\ln\pi(s)$, and it reaches it in
the limit exactly when that remainder goes to zero. Next, $-\ln\pi(s) \ge -\ln\pi(\mathrm{no})$,
because $s$ is one of the no-states and so carries at most their total chance. The flow from the
no-states to the yes-states equals the flow back, because in the long run the weight on the
yes-states doesn't change, so what flows in equals what flows out. Last, take each mean stretch
as its mass over that flow; then $T_{\mathrm{up}}/T_{\mathrm{down}} =
\pi(\mathrm{up})/\pi(\mathrm{no})$, and since the two chances add to one,
$-\ln\pi(\mathrm{no}) = \ln(1 + T_{\mathrm{up}}/T_{\mathrm{down}})$.

Try the stretches with a bag of ten counters: nine marked YES and one NO. Draw,
record the answer, replace the counter and mix again. Each draw gives YES with chance
$9/10$, whatever the last answer was. A YES stretch ends on the next draw with chance
$1/10$, so its mean length is ten draws: the mean $L$ obeys
$L=1+(9/10)L$. A NO stretch similarly averages $10/9$ draws. Their ratio is nine,
just like the ratio of YES to NO counters. In a long record, count all the YES entries
and divide by the number of YES stretches; that is the measured mean stretch length.
Divide both counts by the total number of draws and you have mass divided by flow.

For this rule the long-run flow from NO to YES is $(1/10)(9/10)=9/100$ per draw,
and the reverse flow is the same. A start known to be NO has relative entropy
$-\ln(1/10)=\ln10$ against the long-run spread; after one draw the spread is already
$1/10,9/10$, so none of that relative entropy remains. This bag shows the probability
accounting. The entropy-production interpretation still needs the stated physical model;
shaking counters in a bag does not measure that heat.

Two things need establishing for a general device: that it actually settles, and that
``mean stretch length'' means what it sounds like. I'll say what each says, and why it holds.

Settling comes first. A chain doesn't have to settle. It can get stuck in a corner, or cycle round
forever. What rules that out here is a condition named for Wolfgang Doeblin, a German-born French
mathematician, son of the novelist Alfred Döblin, who took his doctorate at the Sorbonne in 1938
for work on chains like this one and died in June 1940, at twenty-five. In words: from wherever
the device is, its next step lands on each state with at least some fixed fraction $\delta$ of
that state's long-run chance. No state is ever out of reach in one step, and none is favoured
so heavily that the others starve. Under that condition each step shrinks the distance between
where the device is and where it's going by the factor $1-\delta$ (the distance here is the total
of the differences in chance, state by state). To see why, split each step in two: with chance
$\delta$, forget where the device was and draw a fresh state from $\pi$; otherwise, move by
whatever is left over, and the condition is exactly what makes the leftover a genuine set of
chances. The fresh draw puts the device's spread and $\pi$ in the same place, so only the other
share, $1-\delta$, can keep them apart, and no step pushes two spreads further apart than they
were. So the distance dies off geometrically and the spread converges to $\pi$ at every state.
The payoff for the bound is that the entropy produced from the known start stays under
$-\ln\pi(s)$, and it also gets there: after $T$ steps it falls short by at most
$(4/\pi_{\min})(1-\delta)^{2T}$, where $\pi_{\min}$ is the smallest long-run chance of any state.
The shortfall is the relative entropy left, and $\ln y \le y - 1$ once more puts that under the
sum of $(p(x)-\pi(x))^2/\pi(x)$, which is at most the square of the distance over $\pi_{\min}$;
the distance starts at no more than $2$ and shrinks by $1-\delta$ a step. So $-\ln\pi(s)$ is the
exact total entropy production of the whole relaxation, in units of $k_B$, and the
bound $-\ln\pi(s)$ is reached.

For rates satisfying detailed balance with $\pi$, the same equality holds when the device
hops at random moments, with no clock ticking. Write out how fast the relative entropy to $\pi$
changes and pair each hop with its reverse. Detailed balance makes each pair's share minus
(the flow one way minus the flow back) times the logarithm of their ratio, and that is exactly
the pair's share of the rate entropy is produced; a difference times the logarithm of the same
ratio is never negative. So at every instant the relative entropy falls at exactly the rate
entropy is being produced, and that rate is never negative; add the rate up between two times
and you get exactly the drop in relative entropy between them. For settling, Doeblin's condition
becomes a condition on rates: the rate of hopping from any state $x$ to any state $y$ is at least
$c$ times $\pi(y)$, for some fixed $c > 0$. Split the rates in two again: a part that, at rate
$c$, forgets where the device was and redraws it from $\pi$, and the rest, which still satisfies
detailed balance. The production rate adds over the two parts, because with detailed balance it's
the rates times quantities that don't depend on the rates. The rest produces nothing negative,
and the redraw part alone produces at least $c$ times the relative entropy still left. So
relative entropy falls at least as fast as $e^{-ct}$. One wrinkle needs its own line. A known
start puts all its weight on one state and none anywhere else, and the logarithms in the
relative entropy go bad at zero. The relative entropy is still continuous at the starting
instant, because each of its terms has the form $u \ln (u/\pi)$, and $u \ln u$ goes to zero as
$u$ does. So adding up the production from the start out to infinity gives exactly
$-\ln\pi(s)$.

I should say what Doeblin's condition costs. It asks for every hop to be possible directly, in
one step or at a positive rate. A device whose states reach each other only through other states
can still settle, and the standard theory of these chains says when: always, for a device that
hops at random moments and can get from any state to any other; and for a chain that ticks, as
long as it also doesn't march round a cycle. Those devices meet the condition for some fixed
stretch of time or number of steps instead of one, and the same argument runs on blocks of that
length, because the relative entropy never goes up in between.

Then come the stretches. The bound used the fact that the long-run chances of up and down stand in the
ratio of the mean stretch lengths, and I took the mean length of a stretch with the flag up to be
the long-run weight of the up-states divided by the flow into them. For a chain that ticks, that's right for the following reason. Start each stretch the way the flow starts them, count the steps until the chain
leaves the up-states, and the average is exactly $\pi(\mathrm{up})$ over the flow. Look at the
long-run weight that has been up at each of the last $k + 1$ ticks. It splits in two: the part
whose stretch began exactly $k$ ticks ago, which is the flow times the chance that a fresh stretch
lasts more than $k$ ticks, and the part that was up one tick earlier still. Add those splits from
$k = 0$ up to some $K$ and the sum telescopes: $\pi(\mathrm{up})$ is the flow times the sum of
those chances, plus the weight that has been up for more than $K$ ticks. The sum of the chances
that a stretch lasts more than $k$ is its mean length. Doeblin's condition makes the leftover die:
from any up-state the next step lands in the down-states with chance at least
$\delta\,\pi(\mathrm{down})$, so the weight that has stayed up for $K$ ticks is at most
$(1-\delta\,\pi(\mathrm{down}))^K$.

With the stretches measured in seconds, when the device hops at random moments, it comes out the
same, by a different route. Ask, for each up-state, how long on average the device takes to
leave the up-states from there. Call that its leaving time. The leaving times are finite: from
every up-state the device hops into the down-states at a total rate of at least
$c\,\pi(\mathrm{down})$, so the chance a stretch is still running at time $t$ falls at least as
fast as $e^{-c\,\pi(\mathrm{down})\,t}$, and no leaving time is more than one over
$c\,\pi(\mathrm{down})$. And they obey one plain rule, which you get by asking what happens in the
first instant: from any up-state, the rate out of it times its leaving time, minus the rates into
the other up-states times their leaving times, comes to exactly one (per second times seconds is a
pure number). Average the leaving times over the way the flow starts stretches, multiply by the
flow, and the rule plus the balance of the steady spread (as much weight flows into each state as
out of it) rearrange the sum into exactly $\pi(\mathrm{up})$, so in seconds too the mean stretch is its mass over the flow. Both of
the bound's premises are proved now, in steps and in seconds, under Doeblin's condition.


%===========================================================================
\section{From ``Learning borrows, committing pays''}
%===========================================================================
\subsection{Somebody got there first, goddamn it}

I'd been sitting on this idea that \emph{committing}, making a result stay so it can't be handed
back, is what pays, and learning, writing on a blank page you could hand back, only borrows, and I thought: this
one's mine, I finally had an original thought. Then I found Carlo Rovelli's 2024 paper, ``The
Thermodynamic Cost of Choosing.'' I hadn't. Somebody got there first, goddamn it: a real physicist, with
a real paper, saying that making a choice that sticks dissipates at least the energy Landauer charges
for erasing a bit. He even quotes a German saying, \emph{Wahl ist Leiden}, to choose is to suffer, and at
that point what are you going to do but laugh.

And then, taking him apart properly, it got good, because on its face he contradicts Bennett. Bennett
says measuring owes no floor; Rovelli says any measurement costs. So which is it? Look at what Rovelli
calls a choice. A coin in the air hasn't chosen anything; it has to hit the ground and stop. And
then comes the sentence I'd put on my wall: ``stabilization must be irreversible, otherwise it is not
stabilization.'' Bennett's floorless measurement is one that could still be undone; that's \emph{why} it
owes no floor. It only borrows the page. Rovelli's choice is one that's been made to stay; that's why it costs. The way I read them,
they're standing on opposite sides of the same line: learning, which borrows, on Bennett's side;
committing, which pays, on Rovelli's. (Bennett's demon still pays in the end, when it wipes the page.) Neither of them puts it
that way, so push on it. But if it holds, it's the best sign I've got that the line I'm drawing between finding
something out and making it stay isn't something I made up.

Push on it with his paper open and it gives in four places, which I go through at the end of
this chapter, once the demon has been built. So the two-sides-of-one-line picture is my reading,
not his, and the piece of it that's really his is the one sentence about stabilizing. That
sentence I'll still put on my wall.

So he's my closest neighbour, and the difference between us is exact. His is physics: a dissipation bound,
at least $kT \ln 2$ at the temperature of the recording parts, with no outside heat bath, for a
choice that sticks. Mine is a rule in a model of computation: a charge in
the machine's own count, on the step that flips a permanent reading on, with no temperature anywhere
in it. It's the same instinct in a different place, and I'll take it, because it's annoying to be second
but worse to be alone and wrong. And he's the closest only on choosing. On copying there's an older neighbour, in
the chemistry of copying molecules, and the end of this chapter gets to it.

\subsection{A small being at a small door}

You met Maxwell's demon briefly in Section~\ref{sec:toll}. It gets a whole section here,
because the step it seemed to pay for and the step it actually pays for are what this section
is about.

In December 1867 James Clerk Maxwell, who had published his equations for
electricity and magnetism two years before, in 1865, and to my mind had as good a claim as anyone to being the most important
physicist alive, wrote a
letter to his friend Peter Tait describing a problem he couldn't get rid of. Take a box of gas
at one temperature, with a wall down the middle and a tiny door in the wall. Now imagine, he
said, a ``very observant and neat-fingered being'' sitting at the door, who can see individual
molecules coming. When a fast one approaches from the left, it opens the door and lets it
through. When a slow one approaches from the right, same. Otherwise it keeps the door shut.
It never pushes anything; it only opens and closes a frictionless door at the right moments.
After a while the right side is hot and the left side is cold, and you've got a temperature
difference you could run an engine off, made out of nothing but \emph{looking}. That's supposed to
be impossible; it's the thing the second law of thermodynamics says can't happen, and the
second law is the law physicists are least willing to give up. Arthur Eddington, the English
astronomer, wrote in a 1928 book of popular lectures that a theory
found to contradict it has nothing for it but to collapse in deepest humiliation, and nobody
has had to test that yet.

Maxwell put it in his textbook in 1871, and it got named Maxwell's demon, though not by
Maxwell. William Thomson, later Lord Kelvin, called it a demon in \emph{Nature} in 1874, and
Maxwell, in an undated note to Tait headed ``Concerning Demons,'' tried to take the name back:
``Call him no more a demon but a valve \ldots'' Nobody listened, which is probably for the best,
because ``Maxwell's valve'' would not have lasted a century. But keep his correction, because
it's exactly right and it's the whole problem: the being does no work. It opens and shuts a
door. So physicists
spent the next century trying to find where the being was cheating. I want you to feel the
pull of the obvious answer, because I felt it and so did most of them: the being has to
\emph{see} the molecules, seeing costs, and the case is closed. In 1929 Leo Szilard built a one-molecule
version of the box and did the arithmetic, and said the second law survives if finding out
which side the molecule is on costs at least as much entropy as the being gains. In 1951 the French physicist Léon
Brillouin made it concrete: to see a molecule you have to bounce light off it, and the way the
philosopher of physics Jeffrey Bub sums up the result, finding out that way costs at least
k ln 2 of entropy for every bit, which is kT ln 2 of heat, Landauer's number before Landauer, and there was your charge on
looking. For thirty years that was the answer. Information costs because getting it costs.

That's a price on one way of looking, and it doesn't show that every possible way of
making a record has to pay it. A device can burn energy without that being a law.

Szilard is worth one more minute, for a reason that has nothing to do with the demon. Four
years after the one-molecule engine, in September 1933, by his own later telling, he was
waiting at a traffic light on Southampton Row in London, having just read a newspaper report
in which Ernest Rutherford, the physicist who had announced the atomic nucleus in 1911, called the hope of getting power out of atoms moonshine, and as the light
changed it occurred to him that an element which gave off two neutrons when it swallowed one
could keep a reaction going by itself. That's the nuclear chain reaction, thought up crossing
a street. It's his memory of himself, and Richard Rhodes's history of the bomb opens with it, in a
first chapter called ``Moonshine,'' which is where most
people meet it, so hold it as the man's own story, not a record made at the time. Does it change
anything about the demon? No. What it adds is a picture of the man: the same mind looked at one
molecule in 1929 and one neutron in 1933, and both times the question was what a single small
step forces downstream.


This is also where I settle up with Rovelli, my closest neighbour from the start of this chapter,
because I put him and Bennett on opposite sides of one line: a measurement that can still be
undone borrows and owes no floor, and a choice that sticks pays. I owe you how much of that is
mine. All of the reconciliation is mine. Rovelli's paper doesn't make it, and in four places
his own framing points somewhere else. He reserves ``choice'' for an outcome that is stable and
\emph{not predictable}, and he rules out a thermostat because it runs on a predictable
mechanism; the commit in this book is a step a rule fixes in advance, so by his own definition
his result doesn't cover it. His choice is a branching, one possibility splitting into two,
and he calls choosing in a sense the inverse of erasing, because branching is the inverse of
merging; the toll comes from a merge, the opposite shape. His stability is for a while: he says
stabilization means something only on times that aren't too long next to the time it takes
things to thermalize, and past that nothing stays chosen; my permanence is forever (and the
last chapter owned up that forever is an idealization). And he doesn't stand beside Bennett at
all: he holds that any measurement is dissipative, credits the point to the philosopher of
science Hans Reichenbach, and
doesn't cite Bennett. So what Rovelli and I share is a sentence, that a choice that sticks
costs something like Landauer's heat. There's no shared mechanism. Push on it and his coin hitting the
ground, a fact gained from something nobody could predict, sits closer to what this section
calls learning than to committing. I still think it's the same instinct. This is my reading,
and he doesn't make it; you now have everything you need to disagree with it.

And on one point Rovelli isn't the nearest physics. The nearest I've found is in biology. Thomas
Ouldridge and Pieter Rein ten Wolde (2017) derived what they call the fundamental costs of making a
\emph{persistent} copy of a polymer, one whose sequence stays correlated with its template after
the two have physically come apart, and of destroying one. On their account a copy like that is a state
out of equilibrium, and making it has a thermodynamic cost they measure efficiency against. A companion
paper the same year, with Christopher Govern, does the same kind of accounting for the molecules a
cell uses to read out the state of its receptors, which is a measurement in all but name. The way
I read them, the copy that has to stay is the one that pays, and that's the closest thing in
physics I've found to the line this section draws between finding out and committing. The
differences, so it doesn't get borrowed for more than it is: theirs is work and stored free energy, in chemistry, for a
correlation between two molecules, and their title has the destruction of copies in it too, so
their copies don't last forever; mine is a count in a model of computation, for a reading that
never turns off, with no temperature anywhere. It's the same instinct again, from a different direction.

The chemistry has the sharp form of it, too, and it came two years later. Jenny Poulton, Pieter
Rein ten Wolde and Thomas Ouldridge (2019) studied minimal models of polymer copying, and in
those models, when the copy grows at equilibrium, with the weakest chemical driving, a persistent copy's
accuracy advantage over random copying and the efficiency of storing that
information both vanish. This does not mean every copied site is wrong.
Driving growth out of equilibrium can produce correlations that remain
after the two polymers let go, with chemical work supplied to the process. This is finding out against committing, in chemistry, with the dissipation attached, and I'd
rather say plainly that biology had the distinction before I did. What stays different is the
kind of thing the record is: theirs is a correlation in a chemical ensemble, mine is a reading
enforced by a step in a model of computation.

Two more neighbours have ``learning'' and ``commit'' in their titles, so the slogan shouldn't be
read against them by accident. Sebastian Goldt and Udo Seifert (2017) did the stochastic
thermodynamics of a neural network learning a rule, and argued from it that learning has a thermodynamic cost: on
their account, the entropy the network produces bounds how much it can learn. That doesn't contradict
this section, and the reason is that their learner isn't copying onto a blank page. Its weights are a
driven thermal system, jostled by noise and dragged toward what the data push them to, and what
they price is that physical process. This section's claim is narrower: one way of finding out,
a copy onto a blank, is \emph{allowed} to owe no floor and no cost. It can dissipate under the stated thermal-model assumptions, and it still spends the blank. Nothing here says every way of learning is
cheap, and theirs isn't. From statistics, Aur\'elien Garivier, Tor Lattimore and Emilie
Kaufmann (2016) compare two arms with unit-variance Gaussian rewards and
different unknown means. Their explore-then-commit strategies sample the
arms equally during exploration, then choose one for the remaining horizon.
Within the uniformly efficient classes they compare, the best leading
logarithmic regret is twice that of fully sequential strategies. Whether
the gap between the means is known changes both constants; fixing the
exploration length in advance changes the comparison again. The factor
two belongs to that model and strategy class, not every act called commitment.
That's the closest
published sentence I've found to ``committing costs something'' with a number on it. Their price is
reward lost, which is a different unit from a unit of cost on a ledger, so it's a cousin. A twin would
have to count cost.


%===========================================================================
\section{From ``What structure is''}
%===========================================================================
\subsection{A string, a square, and a drowning nobody can place}

Start with the oldest one, and start by doing it. Take a string. A guitar string, a rubber
band over an open box, anything that twangs. Pluck it. Keep its tension as nearly
unchanged as you can when you stop it at a new point; stretching it tighter changes
the pitch for a second reason. Now stop it exactly at the middle and pluck half.
With the tension held steady, that's the same note an octave up, the one that sounds like the same
note sung higher. Press at a third of the way along and pluck the long side: that's the fifth,
one of the handful of intervals a bugle can play at all, which is why you hear it all through
bugle calls. Half, two thirds: small whole numbers,
and your ear says yes before your brain does. About twenty-five hundred years ago, in the
south of Italy, the followers of a man named Pythagoras are credited with finding exactly that, and it's as true
in your kitchen tonight as it was on their hillside.

I'll tell you what they did with it, as best anyone can tell. Nothing survives that's surely from
Pythagoras's own hand. The earliest fragments are from a later member, Philolaus, and scholars
argue over those. Mostly what's left is Aristotle, one to two centuries later, writing in his
\emph{Metaphysics} about them the way you'd write about the neighbours' religion, and he says they
took the elements of numbers to be the elements of all things: everything is number. That doesn't mean everything can be counted; it means the parts of numbers are the parts of all
things.

Stop and ask what I asked. What does that even mean? Is a chair a number? Does a chair have
a number tucked inside it somewhere? Aristotle doesn't say, and he's reporting a belief he
doesn't hold. So don't take it on his word, or theirs, or mine. But before you push on it,
know what it was holding up, because you can't test a load-bearing wall without knowing what
sits on it. These weren't people doing math for fun on weekends. They were a brotherhood,
with rules about what you could eat and when you had to keep quiet, and that plucked string
was their scripture. Halve the string, hear the octave: to them that was proof that under the
mess of the world there were whole numbers in clean ratios, that the cosmos was a chord and
they'd heard it. ``Everything is number'' wasn't a slogan; it was what they'd organized their
lives around. And that's exactly the kind of sentence you should be suspicious of, mine
included.

Now push, and push where they loved most: the square. Draw one, with a side one unit long,
any unit you like, and draw the line from one corner to the opposite corner. Call that the
diagonal. If everything is number, then the diagonal is some whole number of parts long when
the side is some whole number of parts, so the diagonal divided by the side is a fraction, a
whole number over a whole number. Cancel anything the top and bottom share, the way
four-sixths becomes two-thirds, until they have no common factor left. I'll call the top $p$
and the bottom $q$, and the letters are purely so I don't have to keep writing ``the top'' and
``the bottom.''

You need one fact about squares, and it's the one Pythagoras is famous for. The square built
on the diagonal has exactly twice the area of the square built on the side. If you want to
see it for yourself, cut out two identical squares, cut the second one along both diagonals into
four triangles, and lay one triangle against each side of the first, its long edge on that side. The outline is
the square on the first one's diagonal, and it's made of exactly two of the small squares; it's a
five-minute job with scissors. So $p$ times $p$ is twice $q$
times $q$. People write ``times itself'' as a small raised two, so that's $p^2 = 2q^2$, and that
is the only piece of notation in this argument.

Twice anything is even. So $p$ times $p$ is even. An odd number times itself is always odd: draw that many rows of that many dots,
pair off the columns, then pair the dots in the one column left over, and one dot
remains. The same pairing works for any odd side, so if $p$ were odd, $p$ times $p$ would be
odd too. It isn't, so $p$ itself must be even: it's two
times something, and I'll call that something $r$. Put it back in. Two $r$ times two $r$ is
four times $r$ times $r$, and that equals twice $q$ times $q$. Halve both sides: two times $r$
times $r$ equals $q$ times $q$. So $q$ times $q$ is even, so $q$ is even too, by the same step
as before. But you already cancelled everything the top and bottom had in common, so they
can't both be even. The argument has eaten itself, and when an argument eats itself the
thing you supposed at the start is false. There is no such fraction. The diagonal of a square
can't be named by any ratio of whole numbers, ever. Aristotle reaches for this very argument in
his \emph{Prior Analytics} as his example of proof by contradiction (suppose the diagonal can be
measured, and odd numbers come out equal to even ones), and you just did it with a pencil.

There's a story that a member of the brotherhood named Hippasus let this secret out, and his
brothers took him to sea and drowned him for it. I went and checked, because it's a good story and
good stories are the ones to check. The oldest drownings I could find are in Iamblichus, writing
around AD 300, seven or eight hundred years after Hippasus. In Iamblichus the man drowned by name is
Hippasus, and what he gave away was how to fit twelve pentagons together into a ball, the
dodecahedron. A few pages on, Iamblichus tells it again about whoever first let out the
incommensurable, and that time he names nobody. Carl Huffman, who wrote the Stanford Encyclopedia of
Philosophy's entry on Pythagoreanism, follows the historian Walter Burkert there and finds no ancient
writer who connects Hippasus to this discovery at all. So the version with Hippasus drowned for the
diagonal looks like two old stories stitched together later. It's a legend, but notice that the legend exists.
Somebody, long after, felt that a man \emph{should} have been drowned for this. That's how much
was riding on it. To whoever stitched the story together, the diagonal wasn't a math problem; it was the end of the world.

So the claim is dead, and you killed it. I only handed you the pencil. Now the question that matters: what's left?
Look at your square again. The diagonal is still there. It still stands in a perfectly
definite relation to the side; you drew it, and nobody drowned. That relation was real before
anyone tried to express it as a fraction, and it stayed real after that attempt failed. The relation
came first. The number, when there is one, is a name for the relation, and sometimes there
isn't a name. I'm going to keep that, and I want you to keep it too, because you're going to
watch the same move get made again and again as this book goes on: somebody says ``everything
is X,'' and it's true enough to do real work with, and it sounds right when you pluck it, and
then somebody draws the diagonal. Whether the thing I built survives its own diagonal is a real
question, and Section~\ref{sec:diagonal} answers it.

\subsection{Shapes poured into nothing}

About a century and a half after the brotherhood, Plato, in a dialogue called the \emph{Timaeus}, tried to
build the whole physical world out of geometry. Fire, air, water and earth were four
regular solids (the pyramid, the eight-faced one, the twenty-faced one, and the cube), and the
faces of those solids were built from exactly two kinds of right triangle. You can make one
of each in a minute: cut a square from corner to corner for the first, and cut an equilateral
triangle from a corner to the middle of the opposite side for the second. Everything else, in
Plato's story, was a receptacle, a featureless medium that takes on shape the way wax takes a
seal. Structure, for Plato, is shape laid into something that has no shape of its own.

Why did he care? He wanted a world made on purpose, by a craftsman working from a plan, and
triangles were the plan. Is it true? No, but not for the cheap reason. You can't refute a
claim about what things are made of at the smallest scale by picking one up; the smallest
scale is exactly where your hands don't reach, and Plato knew that. What sank it is that it
made no prediction anyone could check, and the people who eventually did reach that scale,
two thousand years on, found atoms, and atoms aren't triangles. The idea was a guess with no
way to lose, until it lost. But push on it and two things survive. One: structure is something \emph{laid in}, and the stuff it's laid into doesn't carry it
by itself. That's a real idea, and I'd ask you to keep it loosely for now; when the machine
arrives in Part~III, you can ask whether the state it holds is the wax or the seal. Two: Plato put a
crack in his own forms himself, in another dialogue, the \emph{Parmenides}, with an objection that Aristotle later
nicknamed the third man, and it's the most useful mistake in this section. Take all the large
things. Say they share one thing, Largeness. But Largeness is itself large, so now
there's a further thing that all the large things \emph{and} Largeness share, and that one is
large too, and it never stops. The extra move is treating Largeness itself as one more large thing. Grant that
again at each stage and it never stops; refuse it and this particular regress
never starts. The lesson, and I think it's right though Plato never spells it out: the
structure of a collection is not one more member of the collection. The plan isn't a brick.
Hold on to that one.

\subsection{Prove it from what you and I both accept}

Around 300 BC, in Alexandria, Euclid wrote the \emph{Elements}, and the thing to take from it is
not the geometry. It's the method, and the method is the whole reason I trust anything in
this book, so give me a paragraph on it. He starts with definitions (``a point is that which
has no part''), five postulates (things he asks you to grant, like ``to draw a straight line from any
point to any point''), and five common notions (things nobody argues with, like ``things which
are equal to the same thing are also equal to one another''). Then he derives everything else,
four hundred and sixty-five propositions, from only those, in order, each meant to use nothing but
what came before (he sometimes leans on things he never stated, and later geometers filled
those in). He never says ``trust me'' or ``it's well known''; he says, in effect, here are the rules you and I both accept,
now watch. If you've ever been in an argument that ended with ``well, that's just your
opinion,'' Euclid is the man who found a way out of that room, and I'm not embarrassed to say
I learned the way out from someone dead for two thousand three hundred years. Four hundred and
sixty-five times he does the same trick, and the trick is everything.

And then there's the one rule Euclid seems almost embarrassed by, and it's a hole worth falling into. His
fifth postulate, the one about parallel lines, is long and awkward next to the others, and for two
thousand years people tried to prove it from the rest so it wouldn't have to be assumed. Proclus, the fifth-century philosopher who wrote
a commentary on Euclid's first book, tried. Mathematicians in Persia and the Arab world tried. In 1733 an Italian priest named Girolamo Saccheri spent a
whole book assuming it was false so he could watch the contradiction arrive, and he got pages deep
into what is now a perfectly good geometry before declaring that it was repugnant and stopping. A
century later three people, separately, took the step he wouldn't: drop the postulate, and what
you get is other geometries that hold up just as well, and nonsense never arrives. Farkas Bolyai begged his
son János to stay away from the problem, warning that it had swallowed years of his own life; the son went in anyway
and came out with one of those geometries. The part I keep is that they found something real by
refusing to believe one of the rules everyone accepted, and checking what happened without it.
Refusing a rule and checking what happens without it is the push, and this one went on for two
thousand years.

What is structure, to Euclid? Whatever the starting rules force. And when are two shapes the
same? You can do this one yourself: cut out two paper triangles and lay one on the other. If
they coincide, they're equal. That's his fourth common notion, and it's how he proves his
fourth proposition, by picking up one triangle and setting it down on another. Sameness is
coinciding under a move. Keep that; it's about to get very general, and I didn't see that
coming either.

\subsection{The statue and the bronze}

Aristotle, Plato's student, pulled the forms down out of heaven and put them inside things,
which was a bit rude to his teacher and also, I think, correct. His example, from the
\emph{Metaphysics}: a bronze statue is bronze plus a shape. And a circle can be made in bronze or
stone or wood, so, he says, the bronze is no part of what it is to be a circle. The same shape in
marble makes a different statue with the same form.
The bronze is the matter and the shape is the form, and the form is what the thing \emph{is}.

Push on it with two things from the kitchen. A cup can be ceramic or metal and
still hold your coffee. A sandwich carved in wood can keep the outline and fail
to be lunch. The two swaps preserve different things. If the question is shape,
the carving may pass; if the question is whether you can eat it, it does not.

Aristotle already makes a distinction of this kind. In Book VII, chapter 11,
he warns against treating a living animal just like the circle whose bronze can
be changed: its parts must be in a condition to do their work. The sandwich is
a way to make that difficulty visible, not a correction he never considered.

That is the bet I need for the record. Change the material, then check the
operations and the answers again. If they still agree, the record has survived
the change in the respects I promised. Section~\ref{sec:latch} states those
promises over a chosen base; Part~III builds an instance. Neither theorem says
that every material can do every job, or that the price in one material is the
price in another.

\subsection{Where things are, without numbers}

In 1679 Gottfried Leibniz, who co-invented calculus and had opinions about everything, wrote
to the physicist Christiaan Huygens asking for something that didn't exist: a way to reason
about \emph{where} things are directly, the way algebra reasons about \emph{how much}, without ever
laying down coordinates. He called it the analysis of situation. Try it yourself. Describe
where the furniture is in the room you're sitting in without using a single number. ``The couch
is against the wall. The table is between the couch and the window. The lamp is next to the
chair.'' You just did it, and what you used was relations between places and nothing else.
Huygens wrote back that November and was not impressed. It took roughly two hundred years for
Leibniz's idea to become a field, and the field is called topology, and Huygens is remembered
for other things.

What survives the push is the same thing that survived the square, coming at you from the
other side. Position is a relation among places. The numbers you lay on top are a
convenience, and they're a convenience you can take away and still have the room. I'll be
honest about why this one landed for me: the thing I built keeps track of which things
belong together, and I spent a long time thinking that was a list of numbers before I
understood it was a relation that the numbers were merely naming. Leibniz had that in 1679.
Huygens told him no. I'd like to think I'd have sided with Leibniz. I probably would have
sided with the famous guy.

\subsection{Naming a thing by saying when two are the same}

In 1884, in a short book called \emph{The Foundations of Arithmetic}, Gottlob Frege noticed
something you can check in your head right now. ``Line $a$ is parallel to line $b$'' and ``the
direction of $a$ is the direction of $b$'' say exactly the same thing. There is no fact in one
that isn't in the other. But the second sentence has a new object in it, a \emph{direction}, that
the first one doesn't. Where did it come from? From the sameness. Once you can say when two
things count as the same, you've earned the right to name the thing they share. Frege used
exactly this move to say what a number is: the number of things in a pile is whatever every
pile you can match up with it, one to one, has in common. The residue has teeth:
a sameness relation lets you name a new object.

Now push, because Frege pushed too hard, and it cost him. In his formal system, Basic Law V treated value-ranges as objects and
identified them exactly when the corresponding functions agreed everywhere.
Together with his unrestricted comprehension principles, that let the
self-membership construction in. The familiar set version is easier to tell: In June of 1902, a young British philosopher and logician, Bertrand Russell, wrote him a short letter. Consider
the collection of all collections that are not members of themselves. Is it a member of
itself? If it is, then by its own rule it isn't. If it isn't, then it qualifies, so it is. Do
it in your head; it flips every time you look at it, like one of those drawings that's a vase
and then two faces and then a vase. Frege's floor had a hole in it, and he knew it the day he
opened the letter, and his second volume was already at the printer. One repair was Zermelo's 1908 axiom of separation: a property selects
a subset of a set already given. His system also had other existence
axioms, including power set and infinity; separation was not its only way
to obtain a set. What it removed was unrestricted sorting of ``everything.''
You sort a pile. I'll tell you now that this rule shows up again, later, in a place where I
had to obey it or fall through the same hole, and I'll point at it when the book gets there.

Frege's reply is worth hearing, because it's the best example I know of how to lose. He wrote
back within days. He didn't hide it or argue around it. He said, more or less, that arithmetic was
tottering, and in an appendix to the volume that was about to come out he wrote that hardly
anything worse can happen to a scientific writer than having a foundation give way right after the
work is finished. Then he tried to fix it. If I'm wrong about something in this book, and I'd bet
I am somewhere, that's how I'd like to take it.

\subsection{Choose the changes, keep what survives}

By 1872 there were too many geometries. Euclid's, and ones where parallel lines meet, and
ones with no such thing as distance, and nobody could say what made them all geometry and kept
them from being a pile of unrelated hobbies. Felix Klein, twenty-three years old and taking up a
professorship at Erlangen, published an answer in a program paper that people still call the
Erlangen program. Pick a set of allowed changes: rotations, stretches, whatever you decide.
Then a geometry is the study of the properties that \emph{stay the same} under those changes. Pick
a different set of changes and you get a different geometry on the very same page. Euclid's, in Klein's reading, is the one where you may slide,
turn and flip a shape and scale it evenly, and stretching it more one way than another is off limits. And I'd
underline, if this were my book, which it is, that in Klein's scheme any group of changes can define
a geometry. The choice is yours. Nobody tells you which changes to allow.

You can do this one at your desk. Take any drawing and ask what's still true after you turn
it upside down (lengths, angles), after you stretch it (lengths are gone by then; which lines cross which is still there), after you bend and crumple the sheet without tearing it (which drawn paths remain
connected survives). Two distant ink marks may touch across a fold; that new
physical contact is not a new connection in the drawing. You have to say which
relation you are following. The three changes give three answers and three structures, all from one drawing. The residue is that structure is what survives a chosen set of
changes, and ``chosen'' is doing real work in that sentence. There's no structure of a thing
from nowhere. You have to say what you're allowed to do to it first. I didn't know that when
I started, and I'd been saying ``the structure of'' things for years.

Klein's move has a descendant in logic, and it goes all the way. In a lecture in London in
1966, published only in 1986, three years after his death, the Polish-American logician Alfred Tarski
asked what makes a notion logical, and answered with Klein's recipe at its far end: allow every
one-to-one transformation of the world onto itself, and call logical the notions that survive
all of them. The more changes you allow, the fewer notions are left standing, and with every
change allowed what's left is logic. Tarski's answer is the sentence above, ``structure is what survives a
chosen set of changes,'' with the choice made as large as it can be.

\subsection{An invariance comes with a conserved quantity}

In 1918 Emmy Noether, who at the time wasn't allowed to qualify for a university teaching post
because she was a woman, and for a while taught courses advertised under David Hilbert's name, proved something that physicists still
treat as close to holy. Take a physical system whose laws can be written the way most of
physics is written, as ``of all the paths the system could take, it takes one that makes a
certain total stationary'' (the total is called the action, and stationary means a small change
to the path doesn't change the total to first order; the path needn't make it smallest, and
that's all I'll say about it here). For such a system, every continuous symmetry that leaves the
action unchanged comes with a quantity that is conserved. If the action doesn't change when you
shift everything in space, momentum is conserved. If it doesn't change when you shift everything
in time, energy is. A symmetry isn't decoration. It has a consequence downstream that
you can measure with instruments. I can't hand you this one to do at the table; the proof
needs calculus I'd have to build first and I'd lose you and probably me, and the ``written
as an action'' condition is doing real work that I'm not going to unpack here. What I can hand you is the
shape. Structure, once it's really there and the setting is the right kind, has consequences
you can measure. That's the sentence that made me think a machine could carry a
cost for structure in the first place. Noether makes it reasonable to ask. She doesn't make it true, and
nothing I build in this book is a system with an action.

\subsection{All of mathematics as structures}

Starting in the 1930s, a group of young French mathematicians began publishing a vast
treatise under one made-up name, Nicolas Bourbaki, as if he were a single person. He wasn't.
In 1948 he also applied for membership in the American Mathematical Society, and the application
was quietly turned down, which I find comforting. (The historian Michael Barany has gone through
that application, a second one a year
later, and the letters they set off.) Their goal was to rebuild all of mathematics on one idea: a
structure is a set of elements ``whose nature has not been specified,'' together with one or
more relations on them that satisfy stated axioms. They believed three kinds of structure
were the mothers of all the rest: algebraic, order, and topological.

Push on it. The historian Leo Corry went through the entire
treatise looking for the place where this formal definition of structure actually does any
work, and what he found is that it ``plays no actual role in the presentation of theories within the treatise.'' Outside
the chapter that defines it, it does work almost nowhere. And I found no argument in their writing for why
the mother structures number three. And the chapter on category theory,
which they promised their readers, never appeared; Corry is where I found that too, and Ralf Kr\"omer's 2006 history of Bourbaki and categories in the 1950s follows the attempt. That does not mean Bourbaki left maps out. Their \emph{Theory of Sets}
has sections on morphisms, derived structures and universal mappings
(Chapter IV, sections 2 and 3). Corry's criticism concerns how much their
general formal apparatus was actually used. The lesson I can take for
my own work is smaller: naming a structure is not enough; show which
constructions and comparisons it supports. I learned that one the
slow way, and the origin story at the back of the book (Section~\ref{sec:origin}) tells it.

Philosophers of mathematics took the same idea further and gave it a name, structuralism. In 1965
the philosopher Paul Benacerraf pointed out that the whole numbers can be built out of sets in two
equally good ways that disagree about which set the number three is, and concluded that numbers
aren't objects at all: what arithmetic is about is the pattern, and anything standing in the right
relations to everything else can play the part of three. The philosophers Michael Resnik
(\emph{Mathematics as a Science of Patterns}) and Stewart Shapiro (\emph{Philosophy of Mathematics:
Structure and Ontology}), both in 1997, made the patterns themselves what mathematics is about. The
philosopher Geoffrey Hellman (\emph{Mathematics without Numbers}, 1989) kept the talk of patterns
and did without the patterns, reading each statement as a claim about what would hold in any
system with the right structure, if there were one. The logician Steve Awodey (1996) argued that
category theory says it best, a structure being what its maps keep, which is where this section
goes next. And the philosopher of science James Ladyman (1998) carried the idea into physics: what a
successful theory gets right about the world is its structure. This section's thesis, that
structure is the relations and the things in them don't matter, is theirs. None of the formal
results takes their arguments as a premise, and the bibliography lists them under the same heading
as Bourbaki.


\subsection{The structure is in the arrows}

I met this one first, and backwards. I'd built it by feel, in a rendering engine, as
``digital DNA,'' and only went looking for its name after the thing already ran. In 1945
two mathematicians, Samuel Eilenberg and Saunders Mac Lane, published a paper that founded a field now called
category theory, and they said in plain words that it was a continuation of Klein, which I'd
have loved to know at the time. The move they made was to stop asking what a thing is, and ask how it's
carried. A category is some objects and
some arrows between them, where arrows can be followed one after another and every object has
a do-nothing arrow to itself, plus two rules that make the following dependable: following
three arrows in a row gives the same result whether you pair up the first two or the last
two, and the do-nothing arrow really does nothing when you tack it onto either end, and that's the whole definition. On the page the two rules look like nothing, and they're the part people
skip, so slow down on them for a second, because each one is there for a reason you can see.

Picture three dots in a row, A, B and C, an arrow $f$ hopping from the first to the second and
an arrow $g$ from the second to the third. Composition says there's an arrow taking the whole
stride from A to C in one, $g$ after $f$. Without it, every multi-hop path has to be named on its
own, and the structure falls apart into a pile of unconnected pairs, a phone book without a
search function. Associativity says that with a third arrow $h$ on the end, pairing up the first
two and then adding $h$ gives the same arrow as adding $h$ to the last two first. Both spell the
same three-hop path. With your paper triangle, bracket ``slide, then turn'' before
``flip,'' or bracket ``turn, then flip'' after ``slide'': you still carry out the
same three moves in the same order. Swapping the turn and the flip would be a
different question. If the two bracketings came out different, two people who built the same chain in
different orders would hold two different arrows, which is a bug, because the chain is the thing and
how you assembled it isn't. The identity laws say that tacking the do-nothing loop onto either
end of an arrow gives back the arrow. If the loop changed things, it wouldn't be the loop. Those six
terms, objects, arrows, composition, identities, associativity and identity laws, are all of
Eilenberg and Mac Lane's definition.

Composing relations is older than the definition, too. Who put a name on that move first? I went looking. Augustus De Morgan, writing on the logic of relations around
1860, took the ``is'' out of the old syllogism, put a general relation in its place, and composed
two relations to get the conclusion. The example that gets quoted from that tradition, Peirce's
a decade later, is ``$x$ is the lover of a servant of $y$,'' which is relational composition with better gossip. When exactly
De Morgan read that paper to the Cambridge Philosophical Society, and on which pages it was printed, the
sources disagree, which is a small reminder that history is where you start digging. Then in 1870
Charles Sanders Peirce, an American scientist and philosopher who turns up again in Section~\ref{sec:established}, carried George Boole's algebra of logic (the arithmetic of true and false that Part~III comes
back to in Section~\ref{sec:smallest},
along with the story of Boole himself, a largely self-taught schoolmaster) over to relations, with
composition as one of its basic operations, and he's usually counted as
the real founder of the theory. So following one arrow after another is older than computers,
older than categories, older than Turing's machine: it is a nineteenth-century logician's product of
relations.

And where did they get the word? They didn't get it from mathematics. They borrowed \emph{category} from the
philosophers, from Aristotle and Kant, and \emph{functor} (a map from one whole category to
another) from the philosopher Rudolf Carnap, and then gave both words exact meanings none of
those philosophers had in mind. So Aristotle walks back into this chapter without being asked:
the man who put the form inside the bronze lent his word to the field that keeps track of forms
by their arrows. That doesn't change the laws. It's a reminder that a word can start in
philosophy and arrive in mathematics with a precise meaning bolted on, which is the thing I've
been trying to do with \emph{established} since the start of the book.

What a thing \emph{is} is how it maps. For now the residue is the emphasis I needed: structure
lives in the maps. Nothing later in this book needs more category theory than that sentence, and
the latch of Section~\ref{sec:latch}, where a record is pinned down by the maps between a machine
and the base it extends, is where the sentence does its work.

\subsection{Structure is the part that compresses}

Play this one with me. I write ``th'' and you guess the next letter. You said ``e,'' and you were
right. Now ``the q'': you said ``u.'' Congratulations, you just measured the structure of English.
In 1948 Claude Shannon measured it properly, in a paper that founded information theory, and
found that about half of written English is ``determined by the structure of the language.''
It's the reason a zip file of a long letter is so much smaller than the letter.
Be careful what that says. It does not say that every other letter is fixed by the ones
before it; you just guessed ``e'' after ``th'' and you could have been wrong (think ``than''). It
says that, on average and over long stretches, about half of what's written could be
rebuilt from the rest, so that the text is roughly twice as long as the choices in it. And he said
why that matters in one line that I'd frame if I framed things: the structure ``allows one to
make a saving in time'' (he added ``or channel capacity''). The more you can predict, the less you have to read.

Thirty years later Jorma Rissanen, a Finnish information theorist at IBM's research lab in San
Jose, turned that saving into a price. His idea is called the
minimum description length principle, and it says that the best model of some data is the one
that makes the total description shortest, where the total is the length of the model \emph{plus}
the length of the data once the model has explained what it can. A model is paid for in bits,
and it's worth stating only if it saves more than it costs. Peter Grünwald, a statistician who leads a machine learning group at CWI, the Dutch national
institute for mathematics and computer science in Amsterdam, in a tutorial you can
find online, adds the footnote that keeps this honest: the price of a model depends on the
code you chose to write it in, so the price itself needs a principle behind it, and that principle
is something you have to argue for, not something you're handed.

That's a lot of ``length of'' to swallow in one paragraph, so do one with me, with
numbers small enough to check on the back of an envelope. Every number in it comes from a
code I'm about to make up on the spot, and I'll say every rule of the code out loud, because
Grünwald's footnote is the whole point: change the rules and every number moves.

The data is a row of thirty-two digits, each a 0 or a 1, and you and I agree in advance
that the row is thirty-two long, so nobody has to spend anything saying so:

\[\texttt{0101010101010101 0101010101010101}\]

(The space is only there to rest your eyes.) Every description in the code starts with one
digit that says which kind of description follows: 0 means ``here's the row, written out,'' 1
means ``here's a pattern, then the exceptions to it.'' Each digit is one yes-or-no answer, one
bit, which is what Shannon's unit is. Written out, the row costs that first digit plus the
thirty-two digits themselves, thirty-three bits, and there's no model in it at all. A pattern
description goes like this. First come five digits giving the block length minus one: 00000 means one digit, and 11111 means thirty-two. That covers every allowed length. Then comes the
block itself. The rule is ``repeat the block, then stop at exactly thirty-two digits, cutting off the last repetition if necessary.'' Then come five digits
for how many places the row disagrees with the pattern (up to thirty-one, which is plenty here), and
five more digits for each of those places, numbered zero through thirty-one and listed in increasing order without repeats. At every listed place, flip the predicted digit.

Do it for the row above. The block is $\texttt{01}$, two digits long. So the model costs 1 (which
kind) plus 5 (the block's length) plus 2 (the block), which makes eight bits. Then comes what's left to say
about the data once the model has said its piece: the row disagrees with ``repeat 01''
nowhere, so that's five digits saying ``zero exceptions'' and nothing else, five bits. The total is
thirteen, against thirty-three for writing it out. The pattern pays for itself, with twenty
bits to spare, and Rissanen's sentence has just happened on paper: a model is worth stating when
it saves more than it costs.

Now spoil the row. Flip two of its digits anywhere you like. The model doesn't change, eight
bits, but now the leftover says ``two exceptions, here and here'': five plus five plus five,
fifteen, for a total of twenty-three. It's still cheaper than thirty-three; the pattern still pays, it just
pays less, because the data has more to say that the model can't. Flip four, and the total is $8 + 5 + 4 \times 5 = 33$, a dead tie, and a pattern that saves nothing isn't worth stating. Flip five and it's thirty-eight, and now writing the
row out is cheaper than explaining it, and this particular pattern has lost to writing the row out. Another pattern might
still win. Flip places 6, 12, 18, 24 and 30 of the original row: now the block
\texttt{010100}, repeated and cut off at thirty-two digits, fits exactly.
Its description costs $1+5+6+5=17$ bits. Five changes defeated ``repeat 01''
and produced a better description of their own. The comparison is between
descriptions in the stated code, not a count of spoiled digits alone.

\begin{center}
\begin{tabular}{lrrr}
\toprule
Message & Original row & Two flips & Five specified flips\\
\midrule
Write all 32 digits & 33 & 33 & 33\\
Repeat \texttt{01}, list exceptions & 13 & 23 & 38\\
Repeat \texttt{010100}, no exceptions & does not fit & depends & 17\\
\bottomrule
\end{tabular}

\small For the last column, flip positions 6, 12, 18, 24 and 30.
That defeats one short description and supplies another. Pattern descriptions
include the type bit, length field and exception-count field. The raw row
uses one type bit followed by its thirty-two digits.
\end{center}
 So ``is there structure here?'' has
turned into a question with a number on each side, and a break-even point you can find yourself: four spoiled digits, in this code.

And now change the code, because I promised you'd see the numbers move. Say I'd used seven digits for each exception's place (wasteful, but nobody stopped me). Then three spoiled digits cost $8 + 5 + 3 \times 7 = 34$, one more than writing the row out, where the five-digit code charged twenty-eight and the pattern paid. The pattern stopped paying one spoiled digit sooner. The row is the same, the pattern is the same, the verdict is different, and the only thing that
changed was a convention I chose. That's Grünwald's footnote, and it's why I won't let you
take any of those numbers anywhere else. They're the price of a description in one made-up
code, a length of a written-down message, and nothing more. They are not the charge I opened
this book with. That charge, as Section~\ref{sec:toll} states it, is on one step a machine
takes. How long anybody's description is doesn't enter it, and if the two ever turn out to be related,
that has to be shown. Both of them being counted in whole numbers doesn't get to stand in for the
showing.

Now push, because this is where ``structure means information'' gets born, and it's wrong, and
it's an easy thing to believe. Flip a fair coin two hundred times and write down the string.
Usually there is no substantially shorter way to say that string than to say the whole thing. You can
watch it happen in the little code above, cut to thirty-two digits: a coin-flip row disagrees
with ``repeat 01'' in about sixteen places, which would cost eighty-odd bits to list,
and even the one block guaranteed to fit with no exceptions, the whole row used as its own
block, costs $1 + 5 + 32 + 5 = 43$. The written-out row at thirty-three beats those two descriptions. Another block
could do better; a lucky run of coins can make a pattern, so check it. The general
reason long random rows usually resist substantial compression is a count: put
each short description beside the one row it prints, and there are too few short
descriptions to stand beside more than a small fraction of all the long rows. Nikolai
Vereshchagin, a mathematician at Moscow State University, and Paul Vitányi, a computer scientist at
CWI in Amsterdam, in a 2004 paper on exactly this question, put it without
mercy: every bit of a string like that is ``meaningless information, random noise.'' Such a string carries maximum information and zero
structure. So in this sense structure is the part that compresses, the part where the ``th''
tells you about the ``e,'' and most information and most structure are different things. Hold
that one tightly, because if you built something that charged for information, you'd be
charging for the wrong thing, and for a while I wasn't sure which one I'd built.

And now notice that the word just changed meaning on you,
and on me. I'm pointing at
it so it doesn't slide past. Shannon and Rissanen and Vereshchagin mean \emph{regularity}: a pattern that
lets you say something shorter. Tarski, a few pages back, and Paul Halmos, a Hungarian-American mathematician whose short
textbook on sets is coming up, mean a set with named relations on it, and
those two are not the same thing. You can run the test yourself: number a thousand
objects, flip a coin for each, heads goes in group one, tails in group two. The result is a perfectly
good sorting in the set theorist's sense, every object in exactly one group. And the list of
who went where is a thousand coin flips, which is to say it compresses about as well as a
thousand coin flips, which is to say not at all. Feed it to a code built like the little one,
with the fields widened to fit a thousand, and you get the coin-flip row again, only longer: typically most of the assignment still has to be written down. A lucky pattern
can save some of it; the coin tosses alone do not guarantee there is none. By the set theorist's rule, the sorting
\emph{is} the structure, and it's every bit as much a sorting as ``evens in one group, odds in the
other,'' which compresses to almost nothing. The two are the same kind of object with wildly different
description lengths. So the two meanings don't even rank things the same way: one asks ``is
there a sorting?'' and the answer is yes for both; the other asks ``does a short rule say who
went where?'' and the answer is yes for evens and odds and no for the coins. The coin sorting has structure in the
set theorist's sense and none in Rissanen's. I'll keep both meanings, and I'll try to say which
one I'm using, and when I slip you'll know what to look for.

\subsection{The four parts of every computer}

In 1936 Alan Turing described the machine that every computer you've ever touched is a
version of, and he did it before any of them existed. It's a clerk with a pencil, following a rulebook,
written down as a machine (that clerk comes back in Section~\ref{sec:established}). Write down its parts; there are four.
There's a tape, divided into squares; a head that reads one square at a time and can write on it and
move left or right; a finite list of states the head can be in; and a table of rules: in this
state, reading this symbol, do this. Nothing else is in it, and two different things are true about it
that get run together. The standard effective models compared with it, including register machines and lambda calculus, compute the same partial functions; those are theorems, proved one pairing at a
time, and Turing himself sketched one of them, the pairing with the logician Alonzo Church's lambda
calculus, in an appendix to the 1936 paper. That it also captures everything a person
could ever carry out by rote, everything ``computable'' in the plain-English sense, is a
claim, called the Church-Turing thesis, and nobody has proved it, because the plain-English
side isn't a mathematical object. Turing himself kept those two apart in the 1936 paper, and
so will I. There's a third claim that gets stacked on top, and it's a different kind again:
the physical Church-Turing thesis, that every device you could build out of real stuff
computes no more than this machine does. That one is a claim about physics. No mathematical proof can settle it.
Nobody has built a device that breaks it. Your laptop, your phone, every server in every data
center compute nothing a Turing machine can't; what differs between them is speed, space,
energy, and cost. Hold on to that last word. Cost is where this book lives.

Now do the thing I did, and do it honestly. Look through those four parts for a place that says ``these squares belong
together.'' There isn't one, and that's not me failing to look. Everything the machine is at any
moment is which state it's in, where the head is, and what's on the tape, and the table says what
happens next, and that's the full list. So any grouping has to live in the states and the table, as a
rule, or on the tape, as symbols, which is to say as bytes like any other bytes. The machine can
hold a grouping. It just has no place where holding one is different from holding anything else.

In 1972 David Parnas, a software engineer, wrote a short paper that every programmer should
be made to read and most haven't, on how to cut a program into modules. He took one small
program and cut it two ways. Both worked. But one was easy to change later and the other
wasn't, and the easy one could be ``much less efficient'' if talking between modules cost
anything. His point was that a module is ``a responsibility assignment,'' a decision about who answers
for what, and how you cut decides what every later change costs. If you've ever maintained
someone else's code, you've paid for their cut the way Parnas said you would, probably at two in the morning. What survives
the push is that a division is a commitment: it changes what the work costs, and it is no caption, no
label stuck on afterwards. I'd like to say the thing I built understood that from the
beginning. It didn't. I did, eventually.


%===========================================================================
\section{From ``Windows and shadows''}
%===========================================================================
\subsection{People who found this on their own}

The fact is so plain that it got discovered separately in every field that has to deal with
seeing part of something, and I want to lay them beside each other, because that's the kind
of evidence I trust: people who'd never met kept walking into the same wall.

In 1956 Edward Moore, working on the small switching machines that became computer
hardware, asked when an experimenter who can only feed a machine inputs and watch its
outputs can tell two of its internal states apart. His answer gave a precise meaning to
``indistinguishable from outside,'' and a theorem with a sting in it: even when every pair of
states can be told apart by \emph{some} experiment, there are machines where no single experiment
tells you which state you started in. What you can always do is run an experiment after
which you know which state you're in \emph{now}. In 1960 Rudolf Kalman,
whose filter went on to aerospace and industry, asked the same question for continuous
systems, called it \emph{observability}, and gave a test (a matrix rank, which you can compute)
for whether watching the output over time pins down the internal state. His unobservable
states are exactly a fibre: states whose difference never reaches the output, and never will,
however long you watch. That second half took me longest to take in. More samples help only when the states eventually show different outputs, and states in
the same fibre of the \emph{whole observation history} remain indistinguishable. Try it with two toy cars on identical roads, a mile apart. Give them the same
speed at every tick and let neither one's motion depend on its mile marker.
Their speedometers then tell the same story throughout. No amount of rereading
that story supplies the missing starting position. Those assumptions do work:
a landmark sensor, or motion affected by a known hill, could expose the position.

Read the mile markers at two known times instead. Their difference divided by
the elapsed time gives the average speed between the readings. To recover an
instantaneous speed, you need a model of the intervening motion or a limiting
sequence of exact readings. Real noisy instruments give estimates with uncertainty.

Where did Kalman's filter go first, before aerospace and industry at large? It went toward the Moon. In the autumn
of 1960 Kalman called an acquaintance, Stanley Schmidt, at NASA's Ames Research Center, and came
to visit, not knowing what Schmidt's group was working on. The group, a branch with eight people
doing the analysis, was stuck on how a spacecraft going around the Moon and back could work out
where it was from what its crew could measure through an optical sight: a few angles to the Earth
and the Moon. That question is a window problem with a capsule in it. The state is position and velocity,
the view is a handful of angles, and the job is to keep a running best guess at the state from a
sequence of views. Kalman presented his new filtering paper to them, and with real work (their
problem wasn't linear and his theory was) it fit, and it went into the feasibility studies for
navigating Apollo. Leonard McGee, of Schmidt's group at Ames, and Schmidt, telling the story in 1985, call it a fortunate
sequence of events and say the solution might have gone undiscovered but for that meeting. The
story shows the limit from the useful side. The filter didn't make the angles show more.
It used the dynamics and a sequence of noisy views to estimate the state.
Observability asks whether ideal observations can distinguish the states at all;
filtering also has to handle noise and uncertainty. More computation cannot
recover a distinction that the entire specified observation history loses.

In 1977 Patrick and Radhia Cousot built a whole discipline of program analysis, called
abstract interpretation, on throwing information away \emph{on purpose}. You run the program over
coarse descriptions of its values, and the exact values never enter. And you accept that the answer you get
has to be true of \emph{every} real value the coarse description could stand for. For example, an input known only to lie between $-1$ and $1$ has a square between
$0$ and $1$; the interval does not tell you its exact square. That is a sound bound
over the hidden possibilities. Exact recovery would require the stronger condition
of fibre constancy. And in 1970 Edgar Codd, inventing the relational
database, made programs independent of how the data is stored, and the relational systems
that followed gave users \emph{views}, deliberately partial windows onto the stored data. A
decade after Codd, in 1981, two database researchers, Fran\c{c}ois Bancilhon and Nicolas Spyratos, gave the view-update problem an exact semantics. The problem asks, if a user changes what they
see through a view, which of the many underlying states that cast that view should the
database move to? Several are legitimate, and that's the fibre again, seen from the writing side.

Two more fields have the fact as a founding definition. In the theory of knowledge, what an agent
knows is modelled by a partition of the possible states into cells, the states it can't tell
apart, and it knows a fact at a state exactly when the fact holds at every state in that state's
cell. The economist and game theorist Robert Aumann used that model in 1976 to prove that two
people who start from the same prior, and whose estimates of an event's chance are common
knowledge, can't agree to disagree; the computer scientists Ronald Fagin, Joseph Halpern, Yoram
Moses and Moshe Vardi built a whole logic of knowledge on the same partitions in \emph{Reasoning
About Knowledge} (1995). Theorem~\ref{thm:fibre} is that model with the cells written as fibres: a
window knows the answer to a question exactly when the question is constant on its cell. And in
computer security, the computer scientists Joseph Goguen and Jos\'e Meseguer defined
\emph{noninterference} in 1982: one group of users doesn't interfere with another when nothing the
first group does changes what the second can see. Hiding the counter in today's view is only the first step toward that property.
If tomorrow's move reads the counter and writes it onto the display, the secret
comes out. To keep it hidden through runs, the permitted moves must preserve
the agreement of the views too. The adaptive-observer argument earlier in the
book supplies that extra condition for its horizontal moves.

And Parnas, from Section~\ref{sec:structure}, hid things on purpose: a module is characterised by the design decision
it \emph{hides} from its callers. Hiding is a window. The caller can't later read back what was
hidden, including the things it later wishes it could check, and that's the price side of
the coin whose benefit side Parnas was selling. None of these people were doing philosophy.
They were trying to build filters, databases, compilers, and modules, and the wall is the
same wall in each.

There's one more from information theory. Oliver Kosut, at Arizona State, and J\"org Kliewer, at the
New Jersey Institute of Technology, two information theorists, asked in 2018 when a receiver
on a channel with a possible attacker, and no shared key, can always either get the message
right or notice the attacker. The answer is all or nothing: the rate is zero, or it's everything
the channel would carry with no attacker at all, and which one depends on a condition they call
\emph{overwritability}, that the attacker can make what the receiver sees look exactly like an
honest transmission of some other message. Overwritability is a printer, a way to put the same thing in the
window without the event behind it. The fibre fact is their condition in its plainest form, with
the probabilities taken out.

And there's one from machine learning, where the two computers of \hyperref[sec:point]{The point}, same bits and different histories, turn up with a lawyer attached. Four
machine-learning security researchers at the University of Toronto and the Vector Institute, Anvith
Thudi, Hengrui Jia, Ilia Shumailov and Nicolas Papernot, asked in 2022 how a company could prove a
trained model had forgotten someone's data, and found that the same model can come out of
different training sets, so the finished model can't show whether a point was used or removed.
The only claim they found you can audit is about the procedure that was run. It's the same wall: the
finished state is a window, and what happened is in the run.

Philosophers who study testimony, which is the trade of being told things, have the same wall from
the other side. Read one after another they blurred into a list of opinions for me, and the list
didn't tell me anything until I built a picture that showed where each one sits, so I'll give the
picture first.

Somebody passes you a note across a room. It says the answer to question seven is 42. What crossed
the room? Ink. Whatever made the writer believe 42 stayed on their side: the working-out, how careful
they were, whether they knew or guessed. Pull that apart and there are five things in play: the
claim on the paper, the writer's head, whatever process got the writer there, the trip across the
room, and your head at the other end. Every position I found puts its weight on one of those five.
Once I saw that, they stopped being a list and became a map, and the map has connections the list
hid.

The two ends of it are Locke and Burge, and they're the same dial turned opposite ways. The
philosopher John Locke, in the first book of his \emph{Essay} in 1690, says believing something
because someone with authority said it isn't knowing it. If the reason didn't cross the room, you
don't have it. Tyler Burge, a philosopher of mind and knowledge at UCLA, argued in 1993 that you're
entitled to accept what you're told, if it's presented as true and it makes sense to you, unless you
have a reason not to. Locke says only ink crossed, so you've got nothing. Burge says the ink comes
with a default, and the default holds until something knocks it over. I get the pull of Burge.
Nobody could get through a day checking every note. But a default anybody can trigger by writing
something sensible is a record nobody paid for, the thing this book keeps refusing, so on that dial the book
sits at Locke's end. And sitting there has a price: staying there means carrying the check across
the room every time.

Jennifer Lackey, a philosopher at Northwestern who has spent much of her work on testimony, cuts the wire between the writer's head and what arrives. Her case, from 2008, is a teacher who doesn't believe the theory
she's required to teach, and her students come away knowing it anyway. I went back and forth on
whether those students really know it, and I haven't settled it. But look at where the case cuts. If
it's right, what arrives at your end never had to be believed at the other end, and the writer's head
drops out of the map. Sanford Goldberg, also at Northwestern, moves the weight one more place over,
onto the process. In 2010 he argued that whether your word justifies me can depend on how reliable
the processes behind your word were, which sit outside my head. That one took me a few reads. What I
think he's saying is that part of my reason for believing you lives in you, in machinery I can't
see, and he's fine with that. I'm not, quite, and why I'm not is why this book exists: machinery I
can't see is machinery I can't check. Christoph J\"ager and Nicholas Shackel, philosophers at
Innsbruck and Cardiff, built a model of degrees of belief in 2025 in which the writer knowing is neither
necessary nor sufficient for the reader to know. That lands where Lackey's teacher landed, from a
completely different direction, a story on one side and a calculation on the other. Agreement like
that I trust more than either piece alone, as long as I remember a model only shows what can happen
inside its own rules.

Then come two sets of people who drew arrows. Wesley Holliday, a logician then at Stanford, put testimony
into a logic of belief change in 2009, with an authority graph recording whose word each agent takes
and rules for what hearing a speaker does to a belief. I've only read a summary of Holliday's paper,
so I'm taking that on somebody's word, which is a funny thing to admit in this paragraph. Security
people had drawn the same kind of arrow two decades earlier. In the logic of authentication of
Michael Burrows, Martin Abadi and Roger Needham (1990), two computer scientists at Digital's research
lab in Palo Alto and one at Cambridge, if I believe an authority has jurisdiction over a statement
and I believe the authority holds it, I believe it. Set that beside Burge and it's his default with
one change: the authority is named before the note arrives. That one change is the whole difference.
Burge's default trusts whoever is writing. Theirs trusts only a writer you picked in advance.

Kenny Easwaran is the one who goes after what crosses the room itself. He's a philosopher who works on
probability and on how mathematicians know things, and in 2009 he took up a question I'd never have
thought to ask. The philosopher Don Fallis had pressed it in a series of papers: why will
mathematicians accept a proof that's long and complicated, or partly carried out by a computer,
and turn down a probabilistic one that can be far more certain? That looked backwards to me at first,
because more certain should win. Easwaran's answer is a property he calls \emph{transferability}. A deductive proof puts the
reason on the paper. Every step crosses the room, so you can learn from the writer without trusting
the writer. A probabilistic proof leaves the reason behind and sends only the writer's word about
what the dice did. Once I saw it that way it stopped looking backwards, and it's the cleanest
statement I found of what the machine of Part~III does with a check: make the reason cross the room,
carry the thing, and run the check on your own side. That reading is mine. His subject is proofs passed between mathematicians, so if it's wrong, it's
wrong on me.

The machine carries the check across the room and runs it on a specified checker. Its
correctness proof says that YES means the claim holds. That proof is supplied once, outside
the individual runs. In Burrows, Abadi and Needham's terms, the checker is the named
authority. Here you can examine its rules and the proof of what they mean. That is the
connection between the event the machine records and the claim the event supports.

Only Burrows, Abadi and Needham write about machines passing messages, and even
they write about the computers as if the computers held beliefs. The rest are about what survives the trip from one head
to another. Either way it's the copier's problem (the computer in \hyperref[sec:point]{The point} that was handed the sorted list and copied it in), with philosophers' names on it and, in one
case, computer scientists'.

Distributed-systems people hit the wall too, and one pair of them measured how far it goes. Their
papers were a pile of acronyms to me until I put them all in one room. People at desks work from one
rulebook and pass slips to each other. Nobody sees into anybody's head; all anyone gets is the slips
handed to them. Somebody may be cheating. What can the honest ones do about it with nothing but slips?

Andreas Haeberlen, then at the Max Planck Institute for Software Systems in Saarbr\"ucken, and Petr
Kuznetsov, then in Berlin, asked in 2009 which faults an outside observer can
detect at all. As I read their answer, it sorts cheating by where it sits. A cheat can break the
rulebook without changing any slip. A cheat can write a slip an honest player could have written too,
for all anyone holding it can tell. A cheat can fail to send a slip that's due, an omission fault, or
send one the rulebook would never produce, a commission fault. The first two, they argue, no algorithm
can detect. The last two can be caught, and they put a floor under how many messages the catching
takes, which I liked for obvious reasons. My first guess was that the classes ranked how bad a fault
is. They don't. They rank how far the fault gets toward a slip, and the first two are a fibre: an
honest room and a crooked one that hand out the same slips. That pair of classes is the fibre fact with an auditor
attached.

Their group's PeerReview system (Haeberlen, Kouznetsov, the same Kuznetsov spelled another way, and
Peter Druschel, the institute's founding director, 2007) is the ladder up to that wall:
every node keeps a secure record of the messages it sent and received, and a checker replays
the record against a reference implementation, a copy of the rulebook it can run, which exposes the
faults a correct node got to see. It comes close to the whole-state escape from the three escapes
below, built: ship the log, and the checker re-runs it. I say close because a log of slips stands in for a
whole state only if the rulebook can rebuild what a desk did from them, and only if the checker's
copy of the rulebook is the right one.

The next three ask what you can say once you're at the wall. Ralf K\"usters, Tomasz Truderung and
Andreas Vogt, three security researchers at Trier, asked what a verdict should hold. They gave in
2010 a general definition of accountability, a verdict naming who misbehaved, and argued, with a
proof inside their own definitions, that verifiability, the property that an outsider can check a
result, is a restricted form of it. The way I read the restriction: verifiability says something went
wrong, and accountability points at a desk. Radha Jagadeesan, Alan Jeffrey, Corin Pitcher and James
Riely, programming-language people at DePaul in Chicago and, for Jeffrey, Bell Labs, asked when the
verdict comes. They built in 2009 an operational model of systems that enforce a policy after the
fact by audit, and studied the trade between the power of the auditor, the cost of the audit and
what's asked of the agents. Robert K\"unnemann, Ilkan Esiyok and Michael Backes, at the CISPA
security center in Saarbr\"ucken, in 2019 implemented a check for it, with Certificate
Transparency (CT), the public logging of the certificates web browsers trust, among their examples. As I understand it, they did it by turning accountability into a
property of the trace, the whole run. It's Thudi and his coauthors' wall again: what happened is in
the run.

In September 2026 Timo Treitz, a student at Saarland University, and K\"unnemann gave CT the full
treatment (Section~\ref{sec:models} has my fragment of it) and found, in their model, that plain CT
gives accountability only if the log is assumed honest, an assumption the SCT auditing extension
(SCT is the signed certificate timestamp a log hands out) can drop and the gossiping extension can't. In my room, CT hands PeerReview's log of slips to a third party
to keep. An honest log is the commitment bit from the three escapes below, with a
name on the contract: somebody's promise that the bit is set when, and only when, the claim
holds, trusted at the step.

Turn the room on me and it gives twice. A raise of the flag that costs nothing and never reaches the window is
Haeberlen and Kuznetsov's first class, and no auditor outside catches it; that's the clerk's trouble,
next. My answer to that is PeerReview's, and it trusts a copy of the rulebook too: the checker, proved
honest once, outside every run. And every verdict in this book is accept or reject on one claim. It
never points at a desk. In K\"usters, Truderung and Vogt's terms that's the restricted form, and the
rest of accountability isn't here.


%===========================================================================
\section{From ``Structural entitlement''}
%===========================================================================
\subsection{The count, exactly}

I'll set out the words and then the count, for any system with Section~\ref{sec:counter}'s counter. A \emph{list} is a finite list
of candidates, repeats allowed, and $|\Omega|$ is the length of the list $\Omega$, repeats
counted. $\lceil \log_2 x \rceil$ is the base-two logarithm of $x$ rounded up: the smallest
$k\in\mathbb{N}$ with $2^k \ge x$. At $x=0$ it is $0$, since $2^0 \ge 0$. And $a \mathbin{\dot{-}} b$ is subtraction on whole numbers that stops at
zero, so $3 \mathbin{\dot{-}} 5 = 0$.

\begin{definition}[Decision tree]\label{def:u08-dtree}
A \emph{decision tree} is a leaf, or a branch with two subtrees. Its depth is $0$ for a leaf and
one more than the deeper subtree for a branch; its leaf count $L(T)$ is $1$ for a leaf and the
sum of the two counts for a branch. Nothing is written on a tree: no question, no answer, no
candidate.
\end{definition}

\begin{lemma}[A tree is no wider than two to its depth]\label{lem:u08-tree-width}
For every decision tree $T$, $L(T) \le 2^{\mathrm{depth}(T)}$, and so
$\lceil \log_2 L(T) \rceil \le \mathrm{depth}(T)$.
\end{lemma}

\begin{proof}
Start at the root. At the first level there can be at most two places, at the
second at most four, and each further level at most doubles the places available.
To count leaves that stop early as well, extend each such leaf with branches
until all tips reach the tree's greatest depth. This only adds leaves.
At depth $d$ the filled-out tree has $2^d$ tips, so the original tree has at most
that many leaves. For depth zero there is just the root, which is one leaf.

Thus $d$ bits are enough to number all the leaves. The least number of bits
that will do, $\lceil\log_2 L(T)\rceil$, cannot be greater than $d$.
\end{proof}

Nothing about any machine is in the definition or the lemma. They're about trees, a leaf or a
branch, and nothing else.

\begin{theorem}[Paid narrowing, counted]\label{thm:u08-count}
Take a system with a counter $\mu$ and a price list $c$ that obey conservation
(Theorem~\ref{thm:conservation}): each step raises $\mu$ by exactly its charge. Let $\Omega$ and
$\Omega'$ be lists with $|\Omega'| > 0$, let $T$ be a decision tree with
$|\Omega| \le L(T) \cdot |\Omega'|$ (the tree \emph{covers} the narrowing), and let $t$ be a run from
$s$ ending at $s'$ in which at least $\mathrm{depth}(T)$ steps are charged at least one each.
Then
\[
|\Omega| \;\le\; 2^{\mu(s') - \mu(s)} \cdot |\Omega'|, \qquad\text{so}\qquad \lceil \log_2 |\Omega| \rceil \mathbin{\dot{-}} \lceil \log_2 |\Omega'| \rceil \;\le\; \mu(s') - \mu(s). \] \end{theorem}

\begin{proof}
Count the bill first. There are at least as many paid steps as levels in the
tree. Each paid step adds at least one and every other charge is nonnegative.
The total bill is therefore at least the tree's depth. Conservation says
that bill is exactly the counter's rise, $\mu(s')-\mu(s)$.

Now count the room the tree supplies. It has at most two to its depth leaves.
The covering premise allows at most $|\Omega'|$ prior entries per leaf, so
\[
|\Omega|\le L(T)|\Omega'|
 \le 2^{\mathrm{depth}(T)}|\Omega'|
 \le 2^{\mu(s')-\mu(s)}|\Omega'|.
\]
That is the first bound.

For the bit version, let $k=\lceil\log_2|\Omega'|\rceil$.
The posterior has at most $2^k$ entries. Insert that count in the first bound:
the prior has at most $2^{\mu(s')-\mu(s)+k}$ entries. So numbering its entries
needs no more than the counter's rise plus $k$ bits. Subtracting the posterior's
$k$ bits gives the claimed bound; if the subtraction would go negative, the
stipulated subtraction stops at zero. An empty prior already has bit count
zero and satisfies the same conclusion.
\end{proof}

Two arguments are stuck together in that proof, and they come apart. The tree half (the
covering and the lemma) bounds the index bits by the sum of the charges of the run, on any price
list at all. Conservation is only what turns that sum into the counter's rise. So the form in terms of cost,
with the sum of the charges in place of the counter, holds in any Thiele machine, and it's
part (3) of Theorem~\ref{thm:ent-cs} below. The shape of the bound shows in what's missing from
the hypotheses: anything tying the lists to the run.

There are three honest pieces of small print, all of them visible in the statement. Both logarithms round up
separately, and the subtraction stops at zero, so when the prior and the posterior round up to
the same number of index bits (four candidates narrowed to three, say) the rounded bound says nothing, though the first form of the theorem charges the run one unit: one leaf can't cover four candidates with three.
The lists are plain lists, repeats allowed, so ``index bits'' is a difference of rounded list lengths
with repeats counted; a list can get shorter by dropping a duplicate without losing a single
candidate, and ``a candidate was ruled out'' and ``the list got shorter'' are two different
claims. And the difference isn't the number of questions anyone asked.

The covering condition, $|\Omega| \le L(T) \cdot |\Omega'|$, is the nearest thing the count
has to an anti-stopped-clock clause, and I have to be exact about how near. Picture each survivor standing for a little crowd
of the prior's candidates, a crowd no bigger than the tree has leaves. The condition says the
crowds can hold the whole prior between them: share the prior's candidates out among them and none
overflows. That limits how much can be deleted per leaf you paid for, but it doesn't say
why anything was deleted. So the condition rules out deleting too much for the tree you've
paid for, and it doesn't rule out deleting for no reason. A run can satisfy every hypothesis of
the count and have ruled nothing out on purpose; the count prices the size of the drop and
says nothing about the honesty of it. And the covering can't just be dropped. On the small machine of Part III (Section~\ref{sec:smallest}), whose record moves CHECK, COMMIT and CERTIFY cost one each, take the sixteen starts
with $A$ at $0$ through $15$, and the run CHECK ``$A \ge 15$,'' COMMIT, CERTIFY from the start
with $A$ at $15$. One yes-or-no question with a big constant in it narrows sixteen candidates to
one, four index bits, and the run has three paid steps, a ledger of three. A tree no deeper than
three has at most eight leaves, and eight times one survivor is less than sixteen, so the
covering fails, and without it the bound would say $4 \le 3$. More candidates only widen the gap:
$2^n$ starts and the constant $2^n - 1$ give $n$ index bits for the same cost of three
(Section~\ref{sec:part3-on-machine}, Proposition~\ref{prop:ent-share-needed}).

A side trip is worth taking here, because it looks like it breaks the count and doesn't. Two
information theorists, Rudolf
Ahlswede and Gunter Dueck (1989), noticed that asking ``is it this one?'' is a different job from
asking ``which one is it?'' If each receiver only wants to know whether the message is theirs, a
noisy channel used $n$ times can serve a number of receivers that grows like two raised to two to
the $nC$, with $C$ the channel's capacity, where naming the message only gets you two to the
$nC$. The task is called identification, and on its face it says checking can be exponentially cheaper than
naming. But their tests are allowed to say yes wrongly once in a while. Make the yes sound, a
test that never says yes to the wrong message, which is the only kind of check this book lets
raise a record, and the double exponential collapses. The argument is two lines of pigeonhole.
Each receiver's test, to be any use, says yes on some
output its own message can produce, and soundness means no other message can produce that output;
so those outputs are all different, one per receiver, and there can't be more receivers than
outputs, $|Y|^n$ for an output alphabet $Y$. That bound is single exponential, the same order as naming. So for a
check that never lies, the card count is the right scale after all, and the cheap trick needs a
check that can be fooled.


\subsection{When the run does the narrowing}

Theorem~\ref{thm:ent-cs} takes the lists in. In the version below, the run makes the posterior, and the count stops being an item on the list and becomes a consequence. That version needs one thing a bare Thiele machine doesn't have: which steps are tests, and what each asks. So a reading names them. For each instruction it gives either nothing or a yes-or-no question about the
state the run is in when it gets there, and a test step \emph{passes} when its question says yes
there. The candidates are states the run could have started from. A passed test is asked of a
candidate by running the same moves from the candidate up to that step and asking the same
question there. The \emph{posterior} is the candidates that pass every test the run passed, asked
that way, in the order the run asked them.

\begin{theorem}[Structural entitlement, with the narrowing made by the run]\label{thm:ent-run}
Take a Thiele machine and a reading that names its test steps, every test step charged at
least one. Take a state $s$ that reads no, a run $t$ from $s$ that ends in a state that reads yes,
and a list $\Omega$ of candidate starts that contains $s$. Let $\Omega'$ be the posterior $t$ makes
from $\Omega$, and let $b_i$ and $a_i$ be the numbers of candidates standing just before and just
after the $i$th test $t$ passed. Then (1) $\Omega'$ sits inside $\Omega$ and contains $s$; (2)
every candidate of $\Omega$ that isn't in $\Omega'$ fails, asked that way, a test $t$ passed; (3)
$|\Omega| \cdot \prod_i a_i = |\Omega'| \cdot \prod_i b_i$; (4) if every test $t$ passed keeps at
least a $2^{-k}$ share of the candidates standing before it, $b_i \le 2^k a_i$, then
$\lceil \log_2 |\Omega| \rceil \mathbin{\dot{-}} \lceil \log_2 |\Omega'| \rceil$ is at most $k$ times
the number of tests $t$ passed, that number is at most the number of steps charged at least one,
and the index bits are at most $k$ times the sum of the charges; and (5) some step of $t$ took the
reading from no to yes and was charged at least one.
\end{theorem}

\begin{proof}
(1) Start with the prior list. Each passed test only removes candidates; it
never inserts one. The actual start $s$ passes every test used to filter the
list, because these are precisely the tests the actual run passed. So $s$
remains throughout, and the final list is a nonempty part of the original.

(2) Follow any candidate that is absent at the end. It was present at first,
so there is a first test at which it disappears. The filter removes it only
when its version of that test says no while the actual run's says yes.
That test is the required witness.

(3) Write the list lengths on the bench: $q_0$ before any passed test,
$q_1$ after the first, through $q_n$ after the last. Here $q_0=|\Omega|$
and $q_n=|\Omega'|$. Every before-count $b_i$ is $q_{i-1}$ and every
after-count $a_i$ is $q_i$. The left product in the statement is therefore
$q_0q_1\cdots q_n$; the right product has the same factors with $q_n$ placed
first. They are equal, including when there are no tests and both products
are empty products, equal to one.

(4) At the first test the share premise gives $q_0\le2^kq_1$.
At the second it gives $q_1\le2^kq_2$. Substitute one into the other,
and continue. After $n$ tests, $q_0\le2^{kn}q_n$.
If $j$ bits number the final candidates, then $q_n\le2^j$, so $kn+j$
bits suffice to number the original ones. The rounded bit drop is at most
$kn$. Each of the $n$ passed tests is a separate paid step, and each adds at
least one to the nonnegative bill. Thus $n$ is at most the number of paid
steps and at most the total charge, giving the two cost comparisons.

(5) Read the actual run from its initial no until its first yes.
The preceding reading was no, so that step raised it. The toll charges
that step at least one.
\end{proof}

Part (3) is the exact count: take
base-two logarithms and $\log_2 |\Omega| - \log_2 |\Omega'|$ is the sum, over the tests the run
passed, of $\log_2 (b_i / a_i)$, the bits each one took. Part (2) answers the complaint I made
about the covering, that it doesn't say why anything was deleted: here everything deleted was
deleted by a test the run passed and paid for. Two premises stay, the prior list and the share, a bound $k$ fixed before the run on what one test can take. Every run meets some bound, its own biggest test's, so what the share adds is fixing it in advance, and that is needed. On the small machine one CHECK of ``$A \ge 15$,'' asked of sixteen
candidates, keeps one, a $2^{-4}$ share, and narrows by 4 index bits on a run of 3 record moves
(Proposition~\ref{prop:ent-share-needed}); with $k = 1$ part (4) would say $4 \le 3$. The same
statement on every certified Thiele-complete machine, where the tests are the CHECK moves and the charges are
counted by the ledger, is Theorem~\ref{thm:ent-run-machine}.


%===========================================================================
\section{From ``Records over any base: the latch''}
%===========================================================================
\subsection{Bases from outside, and a relation I have to be honest about}

The theorem is stated for any base, so let me put it on three models of computation that other
people built: a small Turing machine (any program); a register machine
of the kind the computer scientists Stephen Cook (of the certificates) and Robert Reckhow described in 1973, with unit-cost indirect addressing (any
program); and a lambda calculus (call-by-value, meaning an argument is worked out before it's
plugged in, and made deterministic by letting a term with no step left to take step to
itself). On each, the latch form holds, by the one generic theorem applied to that base
(the bases are defined below). So a register machine
with a base-driven record is a register machine plus a cup, and so is a lambda calculus. (There's one
more instance, the machine I build later, where a particular flag is shown to
be a base-driven record and the same theorem makes it a latch, in
Section~\ref{sec:smallest}.)

I'll give the three bases explicitly. A base needs a state, a next-state function,
a set of starting states and a halting predicate; the latch theorem asks for nothing
else about its instruction vocabulary.

\begin{definition}[Three concrete bases]\label{def:u11-concrete-bases}
The \textbf{tape base} has a finite Boolean list $t$, a head index $h$, a program counter
$q$ and a natural-number cost field. Reading beyond the list returns false. Writing
at $h$ extends the list with false cells if needed, then changes that cell.
Moving right increments $h$ and extends the list through the new head position.
Moving left decrements a positive $h$; at $h=0$, it prepends a false cell and keeps
the head at zero. A fixed finite program contains these instructions:

\begingroup\small
\begin{longtable}{@{}>{\raggedright\arraybackslash}p{0.482\linewidth}>{\raggedright\arraybackslash}p{0.482\linewidth}@{}}
\toprule
\textbf{Instruction} & \textbf{Next state} \\
\midrule
\endhead
\bottomrule
\endlastfoot
Write a bit & Perform the write and advance $q$ by one. \\
Move left or right & Perform the move and advance $q$ by one. \\
Branch to $j$ & Set $q:=j$ if the current cell is true; otherwise advance $q$ by one. \\
Halt & Leave the entire state unchanged. \\
ClaimTapeIsZero with a label $\ell$ & Advance $q$ by one and do nothing else. \\
\end{longtable}
\endgroup

An out-of-range program fetch is Halt. Every state is allowed as a starting state,
and halting means the fetched instruction is Halt. The cost field is preserved by
every instruction in this base. In particular, its ClaimTapeIsZero instruction
neither erases the tape nor charges the cost field; neither the name nor the label supplies any
semantics.

The \textbf{random-access base} has a register function $X:\mathbb N\to\mathbb N$ and a
program counter $q$. For a fixed finite program its instructions are:

\begingroup\small
\begin{longtable}{@{}>{\raggedright\arraybackslash}p{0.482\linewidth}>{\raggedright\arraybackslash}p{0.482\linewidth}@{}}
\toprule
\textbf{Instruction} & \textbf{Next state} \\
\midrule
\endhead
\bottomrule
\endlastfoot
Constant $(i,c)$ & Set $X(i):=c$. \\
Add $(i,j,k)$ & Set $X(i):=X(j)+X(k)$. \\
Subtract $(i,j,k)$ & Set $X(i):=\max(0,X(j)-X(k))$. \\
Indirect load $(i,j)$ & Set $X(i):=X(X(j))$. \\
Indirect store $(i,j)$ & Set $X(X(i)):=X(j)$. \\
Positive branch $(j,a)$ & Set $q:=a$ if $X(j)>0$; otherwise advance $q$ by one. \\
Jump to $a$ & Set $q:=a$. \\
Halt & Leave the entire state unchanged. \\
\end{longtable}
\endgroup

Each register-writing instruction advances $q$ by one; its operands are evaluated
in the old state, and all other registers are preserved. A fetch outside the program
also leaves the state unchanged. Starting states have $q=0$ and arbitrary register
contents. A state is halted when it fetches Halt or is outside the program.
Unit cost here counts instructions. It has nothing to do with the record's ledger.

For the indirect instructions, draw boxes 1, 2 and 3 and put 2, 7 and 0 in them.
Indirect load $(3,1)$ first reads the 2 in box 1, uses it as an address, and copies
box 2's 7 into box 3. Now put 9 in box 3 and execute indirect store $(1,3)$:
the 2 in box 1 tells it where to put the 9, so box 2 changes and box 1 stays 2.
An address is being read from a box; it is not necessarily the box's own number.

The \textbf{lambda base} uses the terms and weak call-by-value rule defined in
Definition~\ref{def:u12-L}. To compute its next state, first try to reduce the function
part of an application. If that part cannot step and is an abstraction, either
substitute an argument that is already an abstraction or take a step in the
argument. If neither applies, the term cannot step. Variables and abstractions
cannot step by themselves, and there is no reduction under an abstraction.
The base leaves a term unchanged when it cannot step and calls precisely those
terms halted. Every term is allowed as a starting state, including open terms.
\end{definition}

So all three are bases in the sense of Definition~\ref{def:u11-base}, and the latch theorem
applies to each. It says nothing else about them, nothing about what they can compute or how fast.

There's also a relation defined between bases, \emph{weak base equivalence}: two bases are
equivalent when some correspondence covers their start states both ways, related states agree
on an observation and on halting, and one step on one side is matched by some number of
steps, possibly zero, on the other. That ``some number of steps'' is the \emph{stuttering} idea
Leslie Lamport (the same Lamport as Section~\ref{sec:counter}'s message counters) put forward in 1983: a fine-grained machine and a coarse one shouldn't count as
different because one takes three ticks to do what the other does in one. The relation
carries no cost; a machine that multiplies huge numbers in one step and a Turing machine
differ in price. The relation never looks at price, so on what it does look at they can match. It's proved to be an equivalence relation, and the latch
form is preserved across it.

Draw a coarse machine $0\to1$, with 1 halted, and a finer one
$a\to b\to c$, with $c$ halted. Let $a$ and $b$ both show 0 and $c$ show 1;
halted states stay where they are. Relate 0 to both $a$ and $b$, and 1 to $c$.
The fine move $a\to b$ changes nothing visible, so the coarse machine matches it
by taking no step. Its move $b\to c$ matches $0\to1$. Starting from $a$, the
coarse move is matched by two fine steps. ``Possibly zero'' allows this;
the relation is about matching observations and halting, not matching the number
of ticks or their cost.

There is a limit to that statement. The ``preserved across it'' proof throws away its equivalence hypothesis; the latch holds
on every base regardless, so the preservation is true for a boring reason. And nowhere in this
book are \emph{two} of the bases actually related by this relation. The relation is
defined, it's an equivalence, and it is not used to connect the Turing machine to the
register machine to the lambda calculus, or any of them to the machine of Section~\ref{sec:smallest}. It would be
easy to read that list as if it did. The honest sentence is the one that's proved: each base, separately,
carries a latch.

\begin{definition}[Weak base equivalence]\label{def:u11-weak-equiv}
Bases $B_1$ and $B_2$, with observations $o_1$ and $o_2$ into one common set, are \emph{weakly equivalent} when some relation $R$ relates every starting state of each to some starting state
of the other, and any related $x$, $y$ have $o_1(x) = o_2(y)$, the same halting, some $n \ge 0$
with $R(\mathrm{next}_{B_1}(x), \mathrm{next}_{B_2}^n(y))$, and some $n \ge 0$ with
$R(\mathrm{next}_{B_1}^n(x), \mathrm{next}_{B_2}(y))$.
\end{definition}

\begin{theorem}[Outside bases, and the relation]\label{thm:u11-four-bases}
(1) Weak base equivalence is reflexive, symmetric and transitive. (2) If two bases are weakly
equivalent, Observation~\ref{thm:u11-latch} holds over one exactly when it holds over the other.
(3) Observation~\ref{thm:u11-latch} holds over each of the three bases in
Definition~\ref{def:u11-concrete-bases}, for every fixed program where the base has one.
\end{theorem}

\begin{proof}
(1) A base matches itself by pairing each state with itself and matching one
step with one step. Reverse any matching between two bases and it gives a
matching in the other direction.

For three bases, pair a state of the first with a state of the third whenever
some middle state matches both. A start on either outside has a middle start
beside it, and that middle start has a start on the far side. Observations
and halting agree across the first pair and then across the second, so they
agree across the outside pair too.

Follow one step on the first base. Its match supplies a finite sequence of
middle steps. Match the first of those steps on the third base, then the
second from where that match ended, and continue. Joining the finite
sequences gives a match for the original step. If the middle sequence is
empty, take no step on the third base. Do the same in the other direction
for a step on the third base. This checks both stepping requirements and
proves transitivity.

(2) The latch argument uses only a base step and the stated record rules.
It therefore holds over both bases, whether or not they are matched.
(3) Each of the three displayed constructions supplies that step, its starts
and its halting reading. The same latch argument applies separately to each.
\end{proof}


\subsection{What it takes to be complete}
\label{sec:complete}

The latch left two choices open, the price list and the event. A machine I can hold up against
everything else has to make both, and a name for ``this machine computes, and its record keeps
something'' has to ask for more than the toll does. The toll alone is a low bar. A machine that
charges one for every step it takes is a clock bolted to a computer, and it obeys the toll for any
record at all: every raise costs at least one because every move does. And a machine whose record
never goes up obeys it too, because a toll on a raise that never happens is never broken. The clock-priced machine may keep useful records; its prices simply do
not distinguish computing from recording. The constant reading can keep no
new yes/no distinction. The toll permits both.

I'll give the names now, all of them, once, with the reason for each, and they work the way
Turing's do.
\emph{The Thiele Machine}, with capitals, is the model: a computer with a reading that can go from
no to yes, natural instruction prices and the toll. The reading may also
be constant or revocable. An internal ledger is extra structure. \emph{A Thiele machine}, lower case, is any particular system
that fits the model, the way one rule table is a Turing machine. That bar is low on purpose, and a
clock clears it. \emph{Thiele-complete}, defined next, is the bar worth meeting, the way
Turing-complete is: a machine that computes anything, with its record kept apart from its computing
and able to tell starts apart. Part~III adds one stricter bar, \emph{certified Thiele-complete},
for a record that, from a clean start, goes up only on a claim that was checked, committed and certified. And
Appendix~\ref{sec:universal-machines} names a \emph{universal Thiele machine}: one machine that
runs the others, the way Turing's universal machine runs every Turing machine. Those are all of
them: the model, the systems that obey it, the ones that compute and keep, the ones whose keeping
is earned, and the one that runs them.

The bar comes in four clauses, and nothing in it mentions what the record is about.

\begin{definition}[Thiele-complete]\label{def:thiele-complete}
A machine with the five parts is \emph{Thiele-complete} when it can be read through an \emph{interface} that
hands over three things: a universal base, a split of its moves into \emph{computing moves} and
\emph{marks}, and a \emph{ledger}, a whole number for each state. A \emph{universal base} is a way
of reading each state as a configuration of a machine that runs every program (I use Minsky's two
counters, Section~\ref{sec:counters-origin}), a class of \emph{live} states, a computing move for each of that machine's instructions that does to the configuration exactly what the instruction does from every live state and keeps it live, and, for each start, a live loaded state showing that start. The interface has to meet four clauses.
(A) \emph{Computing can't reach the record, and the record stays.} No computing move changes the
record, and a record that's up stays up under every move.
(B) \emph{Starts are blank.} Every loaded state has the record down.
(C) \emph{Computing costs nothing, marks cost one.} Computing moves cost 0, every other
move costs exactly 1, and each move adds its cost to the ledger.
(D) \emph{The record tells starts apart.} Some run, the same list of moves, ends with the record
up from one loaded state and down from another.
\end{definition}

Read the clauses as four plain sentences, with the universal base underneath them: the machine
computes anything a Turing machine can. (A) The computing never touches what's kept, and what's
kept stays. (B) Nothing is kept before the run starts. (C) The
ledger counts the marks, and only the marks. (D) What's kept depends on where you started: there's
a run whose record says something about its start.

(D) is the clause where distinction comes in, and it's a low bar: a mark that reads one bit of the
input meets it. What it excludes is a record that ignores the start. A record that goes up on every
run from every start keeps a before and an after, and nothing about where the run began. (D) doesn't
ask that every raise tell starts apart; it asks that the machine can. A machine may have a mark that
raises the record from any state at all; it pays one for it, like every raise, and (D) is still met
by some other run. Which runs a machine wants its record to tell apart is its own choice, and
Part~III makes one.

The name echoes Turing-complete, but this is a specification a machine meets, not a class closed
under simulation: a machine is held up against it through an interface, and changing the unit of
price drops a machine out.

\begin{theorem}[What every Thiele-complete machine has]\label{thm:tc-has}
Let a machine be Thiele-complete through an interface. (1) It obeys the toll, and the move that
raises its record is a mark, never a computing move, and costs exactly one. (2) Along any run, the
ledger rises by exactly the number of marks in it. (3) From any state with the record down, every
run that ends with the record up raises the ledger by at least 1. (4) Some run from some loaded
state raises the record. (5) Every move costs 0 or 1, and some move costs 0. (6) Its base runs
every two-counter program step for step from every live state, and halting from a loaded state
matches halting from the configuration it shows.
\end{theorem}

\begin{proof}
(1) If a computing move raised the record it would change it, against (A); so the raising move is
a mark, and a mark costs one by (C). (2) is (C) added up along the run. (3) Somewhere along the run
the record goes from down to up, and that move is a mark, by (1), so it adds one. (4) is (D), and by (B) the record started down. (5)
Every move is a computing move or a mark, and the universal base supplies computing moves. (6) is the universal base, one instruction at a time, from a loaded state that is live and shows its start.
\end{proof}

The floor is one, and one is reached: the lift below gives a machine whose record goes up after a
single mark.

\begin{observation}[A clock is a Thiele machine and isn't Thiele-complete]\label{obs:tc-clock}
Take the two-counter machine, let every move cost 1, and take as the record any yes/no question
about its configuration; or a flag that switches on, and stays on, the first time a chosen test of
the configuration says yes. It obeys the toll and has a two-counter machine inside it, and it isn't
Thiele-complete. No machine whose every move costs at least 1 is.
\end{observation}

\begin{proof}
A Thiele-complete machine has a move priced 0, by Theorem~\ref{thm:tc-has}(5), and this one has
none.
\end{proof}

A two-counter base priced 0 with a record that's down at every state obeys the toll and meets (A),
(B) and (C), and isn't Thiele-complete either: (D) asks for a run that raises the record.

\begin{theorem}[Each of the four clauses is needed]\label{thm:tc-clauses}
For each clause there's a machine that meets the other three through some interface and is
Thiele-complete through none. Two of them are built on the two counters with two marks, one that raises
the record when $A$ is 0 and one that always raises it, and a ledger beside them that each move adds its cost to. (A) fails: that machine, and one more move,
priced 0, that puts the record back down. (B) fails: the two counters, every move priced 0, with
a record that no move ever changes, up at the starts where $A$ is 0 and down elsewhere. (C) fails:
the two counters with two marks, every move priced 1. (D) fails: the two counters, every move priced 0,
with a record that's always down.
\end{theorem}

\begin{proof}
Each machine meets the three by construction. The first fails (A) through every interface, because
a move lowers the record, so the record doesn't stay. The second: every one of its moves costs 0,
so (C) makes every move a computing move under any interface; then by (A) no run changes the
record, and with every start blank by (B) no run raises it, so (D) fails. The third has no move
priced 0, so it fails by Observation~\ref{obs:tc-clock}. The fourth never raises its record.
\end{proof}

\begin{theorem}[The record can't be read off the window]\label{thm:tc-window}
Let a machine be Thiele-complete through an interface. (1) Every configuration is shown, at no
cost, with the record down, by one computing move from some loaded state. (2) Some loaded state
has two runs that end on the same window, one with the record up and one with it down, and their
ledgers differ by at least 1. (3) So no rule that reads only the window says, at the end of every
run from every loaded state, whether the record is up, and none says what the ledger reads.
\end{theorem}

\begin{proof}
(1) From the loaded state $(1, x+1, y)$, the computing move ``DEC $A$ to line $p$'' shows
$(p, x, y)$, costs nothing by (C), and leaves the record down by (A) and (B). (2) Take the run (D)
gives, from the loaded state where it raises the record. Its end shows some configuration. From
the same loaded state, computing moves alone reach any configuration at all: count each counter
down to zero, count it up to the value wanted, and jump to the line wanted. That run ends on the
same window, and it's all computing, so by (A) and (B) its record is still down and by (C) its
ledger hasn't moved, while the first run's ledger went up by at least one, at its raise. (3) A rule
that reads the window gives one answer for one window, and (2) has one window with two answers.
\end{proof}

Part~III proves the same blindness for its certified machines (Observation~\ref{thm:u34-tc-window}),
with a gap of three, because a certified raise costs three. Here the gap is one, and one is
reached.

\begin{theorem}[Every universal base becomes Thiele-complete]\label{thm:tc-lift}
Take any machine with a universal base and any events on its states, and add a record, a ledger and, for each event, a
mark that costs one and raises the record exactly when the event holds; computing moves cost nothing
and leave the record alone, and every move adds its cost to the ledger. If some run of computing moves makes some event hold from one loaded
state and not from another, the result is Thiele-complete, its computing is the original machine's
step for step, and its record goes up after a single mark from any state where the event already
holds.
\end{theorem}

\begin{proof}
(A): computing moves are the original's and don't touch the record, and a mark only adds to it.
(B): the record starts down. (C): by construction. (D): run the moves that make the event hold
from one loaded state and not the other, then the event's mark. Step for step: a computing move
acts on the original part and nothing else. And a mark from a state where its event holds puts the
record up.
\end{proof}

One mark on ``$A$ is at least 1'' always works: the starts with $A$ at 0 and at 1 already differ
on it, so that mark alone, run from each, tells them apart. So Thiele-complete isn't a property only a special
machine has. Any model with a universal base becomes Thiele-complete by adding marks, and needs no
claims and no checker to do it; a model that simulates the two counters some other way does so once
its moves are grouped into ones that carry out a counter instruction each (Theorem~\ref{thm:lift-models}). Two kinds of base don't lift, whatever the
events: one whose moves don't depend on the state it's in, and one whose moves lead from each state
to only finitely many next states. A universal base needs, from one state, a move at no cost to
every line of a program, and that's infinitely many places to go.

Some Thiele-complete machine runs every Thiele machine a program can describe.
Appendix~\ref{sec:universal-machines} defines a \emph{universal Thiele machine} and builds hosts
that run any machine a program can describe, in order, carrying its reading as a latch and matching its ledger up to a surcharge of at most three, with every
computing step done by counter moves on the host and every rise of the host's record placed by its
program. None of that needs certification: the presented host's own program, read on a host with no claims, no checker and no table of facts, and one paid mark that raises the record when a register is at least 1, meets the definition too.

The rest of the book holds machines to this bar. Part~III gives one example that clears it,
and then asks for the stricter certified bar: that from a clean start its record go up only on a claim that was checked, by a
checker proved to mean it, and then committed, and that some claim's check can come out either
way. Asking that is a choice of what the record answers to, the second of the latch's two open
choices, made one way, with clean starts and three priced acts added (Theorem~\ref{thm:tc-decomposition}).


%===========================================================================
\section{From ``The whole direction''}
%===========================================================================
People got here before me, and from a direction I didn't expect, which was programming for many
processors at once. Lindsey Kuper and Ryan Newton's LVars (2013), ``lattice-based data structures
for deterministic parallelism'' in their own title, are memory cells whose value only climbs in an
order, where a read can't ask for the value and can only ask whether it has climbed past a
threshold. That's a cup. Writers running side by side can push the value up in different orders,
and a threshold once passed stays passed. Their deterministic behaviour also depends on
which reads and updates the language permits. Permanence alone would not do it: a program
allowed to ask whether the cup is up \emph{yet} could still see another writer's timing. Amin Timany and Lars Birkedal (2021), coming at it from proofs about programs,
show, as I read their paper, that any preorder at all can be carried as monotone ghost state
in a program logic. That's the closest thing I've found to somebody else arriving at the decomposition into
lines.

The two rest on the same thing, a value that only climbs and a question that only asks whether it's
past a point, and they use it for opposite jobs. In an LVar the threshold read steers the program.
Ghost state, as far as I understand the term, lives in the proof, and the running program never
reads it. Here the record is never allowed to steer the base computation, so that particular restriction sits closer to the ghost state. Push on
that and it gives a little: Section~\ref{sec:structure}'s glance, the whole reason to keep a record,
is a read. What keeps me on that side is that a glance worth trusting gives the answer the long way
would have given. Their purposes aren't mine, but the shape is theirs first. What I can't find in
either is a price on crossing a line.


%===========================================================================
\section{From ``The shadow, exactly''}
%===========================================================================
\subsection{How much gets destroyed, on average}

The halving price is a worst case, because it prices the largest landing pile a move can have. The next question is
the average one: if I don't know which state the machine was in, cups and all, but I have some
spread of chances over the states, how much do I lose, on average, about which state it was, once
I know where it landed? Shannon's conditional entropy measures exactly that (Section~\ref{sec:a2-not-arbitrary}
built the entropy from scratch, so I won't do it again here), and I'll write it in bits.

\begin{theorem}[The halving price in bits bounds the average loss]\label{thm:dir-bits-loss}
Take any spread of non-negative weights over the states, uniform or not, with finite total weight, on a machine
of any size. (1) If every landing pile has at most $N$ states, with $N\ge1$, the weighted conditional entropy of
the state given its landing spot is at most the total weight times $\log_2 N$ bits. (2) So on a machine with the
halving price, that weighted information loss is at most its cost in bits times the total weight, even when the
spread's own entropy is infinite. For a probability distribution, the total weight is one and the average loss is
at most the cost. (3) If a landing pile has $2^c$ states, the uniform probability distribution on it loses exactly
$c$ bits and attains that bound.
\end{theorem}

\begin{proof}
Look at one landing spot. Suppose the sources that land there carry total
weight $w>0$. Divide each source's weight by $w$. Those shares now add to one:
they describe the uncertainty about the source once you know this landing spot.
There are at most $N$ possible sources. The logarithm argument in
Section~\ref{sec:a2-not-arbitrary} says the uncertainty is greatest when those
possibilities have equal shares, and is then at most $\log_2 N$ bits.
Multiply by $w$ to put this pile back into the whole spread. Its contribution
is at most $w\log_2 N$. A pile with zero weight contributes nothing.

Add the contributions. The pile weights add to the original total weight,
so the same $\log_2 N$ multiplies that total. Under the halving price,
$N=2^{\mathrm{cost}}$, whose base-two logarithm is exactly the cost.

Infinitely many piles cause no subtraction of one infinite entropy from
another. I am adding the nonnegative uncertainties inside the piles.
Every finite collection contributes at most its weight times the bound,
and its weight is at most the total. Taking the largest value approached
by those finite sums keeps the same bound.

Finally put all probability on one pile of $2^c$ sources, equally shared.
Before landing there are $2^c$ equally possible origins; afterwards the
landing spot distinguishes none of them. The uncertainty left about the
origin is $\log_2(2^c)=c$ bits. That attains the bound.
\end{proof}

Where does it stop? My first guess was one place, a landing pile with infinitely many states, and
that the average would blow up there right along with the worst case. That's half right, and the
half that's wrong surprised me. No finite charge satisfying the halving price can price a move that sends infinitely many states
to one spot. The weaker merge price could still charge one, since it does not measure the pile's size. But the average depends on the spread. Put weight $\tfrac12$ on
the first state, $\tfrac14$ on the next, $\tfrac18$ on the next, and so on, and the information lost
by crushing the pile never reaches 2 bits, however many of its states you count. Under that spread,
a two-bit information budget covers the average, though no finite budget covers the worst case. So I went looking
for a spread that does blow up, and it took a nastier one: blocks of states, block $k$ holding
$2^{2^k}$ states and a share $\frac{1}{(k+1)(k+2)}$ of the weight, which adds up to one. Under that
spread the loss grows past every bound, and no finite information budget covers that average loss.

You can generate the first spread by tossing a fair coin until the first head.
Zero tails names state 0, one tail names state 1, and so on. State $j$ has chance
$2^{-(j+1)}$, and identifying that outcome takes $j+1$ toss-results, so its entropy
term is $(j+1)/2^{j+1}$. Add these by counting each possible toss position:
the first is reached with chance 1, the second with chance $1/2$, the third with
chance $1/4$, and so on. The total is $1+1/2+1/4+\cdots=2$; every finite partial
sum of the state terms is smaller. Resetting all those outcomes to one box can
lose only those two bits on average, despite having infinitely many possible starts.

For the block spread, split each block's weight equally among its states. The block
weights sum to one because $1/((k+1)(k+2))=1/(k+1)-1/(k+2)$. Within block $k$,
there are $2^{2^k}$ equally likely choices, so just identifying a state within that
block contributes $2^k/((k+1)(k+2))$ to the average bits. Those contributions grow
without bound as $k$ grows, and forgetting which block it was only adds to the loss.
And no finite worst-case price can handle an infinite pile: for any proposed cost
$c$, select $2^c+1$ different states in it. They all land together, violating the
halving rule's maximum of $2^c$.

\begin{theorem}[An infinite pile]\label{thm:dir-infinite}
(1) A move that sends infinitely many states to one spot isn't priced by the halving price at any
finite cost. (2) Under the halving spread $2^{-(j+1)}$, crushing the pile loses exactly $2$ bits, and counted over its first $n$ states, less than $2$ for every $n$. (3) Under the block spread, the information lost
grows past every bound, so no finite bit budget covers it.
\end{theorem}

\begin{proof}
The three computations above.
\end{proof}

And heat? Part~I already said where the line is, and it's the same line here. Landauer's
principle, that crushing $c$ bits' worth of possibilities gives off at least $c\,kT\ln 2$ of heat on
average, is a premise from physics, and a contested one (Section~\ref{sec:a2-not-arbitrary}). Nothing in
this section proves anything about heat. What the counting does give is that, under that premise,
on a machine with the halving price, a move's cost priced at $kT \ln 2$ per unit,
times the weight of the piles, is always enough
to cover the minimum the premise would
demand, for every spread of finite total weight.  If the premise falls, that sentence falls
with it, and every theorem above it stands, because none of them used it.


%===========================================================================
\section{From ``The diagonal''}
%===========================================================================
\subsection{Five names and a table}

Do this first, with a pencil, because the whole family of arguments is this one trick and
you should own it. Write five friends' names down the left of a page. Next to each name,
write a five-digit string of zeros and ones, anything you like. Now build a new string: its
first digit is the \emph{opposite} of the first digit of the first row; its second digit is the
opposite of the second digit of the second row; and so on down the diagonal. Look at what
you've made. It differs from row one in position one, from row two in position two, from
every row at that row's own position. So it's nobody's string, and your table of five missed
something.

With five rows and five digits that's a cute fact. Georg Cantor, a mathematician working out how to compare infinite
sets, made the point in 1891 that it works for \emph{any} list, however long, even infinite: hand me a complete list of all infinite
strings of two symbols, I'll walk the diagonal flipping as I go, and I'll hand you back a
string that's on no row. So there is no complete list. There are more such strings than there
are places on any list, and that's where ``some infinities are bigger than others'' comes
from. (He wrote the two symbols as m and w, where I'd write 0 and 1, and that paper wasn't even his
first proof of the fact. The first was in 1874, by a harder road through the real numbers, and in
1891 he called the new one much simpler. The diagonal was him finding the easy way.)

Now look at what the argument \emph{did}, because it's going to do the same thing again and again,
in new clothes, and I want you to recognise it each time. (Not every name below is a
flipped diagonal in the table sense. Kleene's theorem is the tool that makes the diagonal
buildable, Rice's is the general form, the one from Ken Thompson, who built the
original Unix operating system at Bell Labs, is an attack that uses the same
trick, and the last, Lawvere's, is the theorem that makes them all one.) There was a table: rows are the things being
listed, columns are questions, each cell says ``does row i do column j.'' There was the
diagonal: the cells where a thing is asked about \emph{itself}. There was a flip: something with no
fixed point, like turning yes into no, applied along the diagonal. And the thing that got
refuted was \emph{the assumed object}, the complete list. Nothing outside the assumption was
touched. Hold on to those four: a table, a diagonal, a flip, and the assumption eating itself.

\subsection{The same move, again and again}

Russell, 1902, from Section~\ref{sec:structure}: the collection of all collections that don't contain themselves.
It's Cantor's flipped diagonal with ``contains'' as the table, and Russell said as much, in
\emph{The Principles of Mathematics} in 1903. He got there, he wrote, by running Cantor's argument
on a supposed collection of everything, which ought to be the biggest collection there is. What got
refuted was the assumed object, the collection of everything, and Zermelo's fix was to stop
assuming it. I promised in Section~\ref{sec:structure} that I'd come back to this here. Every reading in this book is a yes or
no question on the states of one machine, a collection cut out of a collection I was already
given, and never a sort of everything. When Part~I asked for one state that codes every question
about the states (Section~\ref{sec:what-a-chair-is}), the same flip refuted it.

Jules Richard, a French mathematician who taught at the lycées of Tours, Dijon and Châteauroux,
took his doctorate in Paris in 1901 and wrote mostly on geometry. In 1905, in a short note, a
letter to the editor of a French science review, he observed: some phrases (his were in French) define real numbers (``the ratio of a circle's
circumference to its diameter''); there are only countably many phrases, so list the numbers
they define in order; run Cantor's diagonal down their decimal expansions and you get a real
number no phrase defines, except that you just defined it, in the phrase that describes the
diagonal. The lesson I take from that one is the lesson this whole
section turns on: keep apart the language a system \emph{talks in} and the language you use to
talk \emph{about} it. Mix them and the diagonal bites you.

Richard's paradox has a cousin small enough to do at the table, and it comes with a footnote I
can't resist. Russell, the first to put it in print, in 1906, credited it to G. G. Berry, a
junior librarian at the Bodleian Library in Oxford. The paradox has outlived empires, and it
came from the library desk. In the version in \emph{Principia Mathematica} it's about syllables. Take
the phrase ``the least integer not nameable in fewer than nineteen syllables.'' There are only
so many phrases shorter than nineteen syllables, so they name only so many numbers, so some
least number is left over, and the phrase picks it out. Now count the phrase: the, least,
in-te-ger, not, name-a-ble, in, few-er, than, nine-teen, syl-la-bles, and it comes to eighteen. In eighteen
syllables it has named the number that can't be named in fewer than nineteen. It's the same lesson as Richard's, in a smaller coat: ``nameable'' is a word about the
language, used inside the language, and that's where the bite is.

Kurt Gödel, 1931, with an improvement in 1936 by Barkley Rosser, a logician who had studied under
Alonzo Church at Princeton and whose trick lets the theorem assume only that the system is
consistent: the big one, and I'll say exactly
what it says, because every word of it is carrying weight. Take any consistent formal system with a mechanically checkable list
of axioms that can do a bit of arithmetic. Then there are sentences it can neither prove nor
refute. The engine is a lemma that lets you build, for any property you can write down, a
sentence that the system proves \emph{equivalent} to that property holding of the sentence's own
code number. Apply it to ``not provable'' and you get a sentence equivalent, inside the
system, to its own unprovability. The table is the sentences listed by code, the diagonal is a sentence about its
own code, and the flip turns provable into not-provable. The Stanford Encyclopedia stresses two cautions, and
I'll repeat them. The sentence doesn't literally ``say'' anything about itself. It's provably
\emph{equivalent} to a claim about its code, which is less mystical and more exact. And the
theorem doesn't say the sentence is unknowable. It says \emph{this one system} can't settle it.

And how did the world take it? Mostly, at first, it didn't. Gödel, twenty-four years old,
announced the first theorem during a round-table discussion on the last day of a conference in
Königsberg in September 1930, and by the account of Solomon Feferman, one of the editors of
Gödel's collected works, the room met it with general incomprehension, except for one man. John
von Neumann pulled Gödel aside to talk about it, kept thinking about it, and a couple of months
later saw that the argument gave something more: a consistent system of this kind can't prove
its own consistency. He wrote to Gödel with it in a letter dated 20 November 1930. Gödel had
already got there. His paper had reached the journal three days earlier, with that
second theorem announced in it, and von Neumann left the credit with Gödel. It's the
``somebody got there first'' shape from Section~\ref{sec:narrowing}, with one difference: here the man
who got there second was von Neumann.

Alan Turing, 1936. I'd always taken it that Turing proved the halting problem undecidable, and
here I had to go and look. The set theorist Joel David Hamkins and his coauthor Theodor Nenu went
back to the 1936 paper and report that the word ``halting,'' the modern statement, and the
self-referential proof everyone teaches are all, strictly speaking, absent from it. What Turing
proved undecidable were two other questions about his machines (whether a machine prints
infinitely many digits; whether it ever prints a given symbol). The second is \emph{equivalent} to
halting and the first is strictly harder. So their answer cuts both ways, and they say so: Turing
never stated the halting problem, and he essentially proved it, because his second question is a
short translation away from it. The name came later. Martin Davis seems to have been the first to
write ``halting problem,'' in a book of 1958, and Stephen Kleene had already stated the problem,
without that name, in his 1952 textbook \emph{Introduction to Metamathematics}. The textbook version is still right as mathematics, and it fits on two index
cards. Card A is the Checker: hand it any program card and it announces ``halts'' or ``runs
forever.'' Card B is the Contrarian: ``Hand me to the Checker. If it says I halt, run forever.
If it says I run forever, stop.'' Hand B to A, and either answer is wrong, so there's no Checker.
The table is programs against programs, the diagonal is a program run on itself, and the flip turns halt into loop.

Stephen Kleene, 1938, the logician whose textbook I just mentioned, and this is the one that turns the paradox into a tool, so slow down.
Everything above needs the self-referring object to \emph{exist}: Gödel's sentence has to be
buildable, the Contrarian has to be able to ``hand me to the Checker.'' Kleene's second
recursion theorem says: for any transformation of programs you can compute, there's a
program that behaves exactly as if it had been handed its own source code. Any program may
be written as though it already had a copy of itself in hand. He stated it almost in passing, in a paper
on something else, and proved it in the last two lines of one of its sections. I have that count
from Yiannis Moschovakis, a logician at UCLA and Athens who wrote a 2009 paper on it, expanded in the Bulletin of Symbolic Logic in 2010, called
``Kleene's Amazing Second Recursion Theorem,'' and he calls it a little
gem. It's the reason ``programs that print their own source'' exist in every ordinary programming language: one that can compute whatever a Turing machine can and can build, from a program and an input, a program that runs the first on that input (they're called quines, after a philosopher I'll come to in a moment; try ``Print the
following twice, the second time in quotes: `Print the following twice, the second time in
quotes:'\,'' and obey it). Self-reference in computation isn't a paradox but a construction.
The paradoxes happened only when someone assumed an \emph{unrestricted} collection or language;
the theorem happens when you build the self-referrer honestly.

And where does the word quine come from? The philosopher didn't coin it himself. Douglas Hofstadter
coined it in 1979, in \emph{Gödel, Escher, Bach}, in honour of Willard Van Orman Quine, who had
made a long study of self-reference that works indirectly. Quine left a sentence you can take
apart at the table, from a 1962 essay on paradox, and in full it reads: ``yields falsehood when
preceded by its quotation'' yields falsehood when preceded by its quotation. Do what it says.
Take the part in quotation marks, put a quoted copy of it in front of itself, and look at what
you've built: it's the whole sentence, word for word. So the sentence says that it itself is
false. It's the liar paradox, but look at how it got there: no ``this sentence'' anywhere, no
pointing at itself, just an instruction to quote and copy, the same move as the print-twice
instruction above. Quoting and copying is Kleene's trick in plain English, and it's why the paradox and the tool
look so alike. The paradox is what you get when the copy asserts something about its own truth.
The tool is what you get when the copy just does something.

Henry Rice, 1953, whom you met briefly in Section~\ref{sec:windows}: the general form. Any property that depends only on what a program \emph{does},
whatever its text, and that's true of some programs and false of others, has no
algorithm that decides it for all programs. ``Does this program halt'' is one such property.
``Does this program ever output 7'' is another. ``Does this program contain an if-statement''
isn't, because that's about the text. One standard proof goes straight through Kleene.

And the one that's actually an attack: Ken Thompson, in his Turing Award lecture, published in 1984.
Build a compiler that, when it compiles the login program, plants a back door, and when it
compiles \emph{a compiler}, plants both habits into that one too. Then delete the trick from the
source. The compiled compiler still carries it, and recompiling a clean source with the
poisoned binary re-plants it, forever. Reading the source never reveals it. A locksmith
teaches an apprentice to cut a spare for every town-hall key and to teach \emph{that} to every
apprentice they train, and never to write it in the manual; years later the manual is
spotless and every locksmith in town cuts the spare. Thompson's compiler is the recursion theorem with a
knife in it, and his moral was that you can't trust code you didn't write entirely yourself.

And Thompson wasn't first either, and he said so. Pull on the lecture's acknowledgment the way
Section~\ref{sec:established} pulls on Bitcoin's footnotes: he writes that he first read of the
possibility in an Air Force critique of the security of an early implementation of Multics, a
time-sharing system from the 1960s, and that he couldn't find a more specific reference to it,
and he asks anyone who can to send him one. The critique is Paul Karger and Roger Schell's 1974
security evaluation of Multics, which already described a trap door planted in a compiler that,
because the compiler is written in its own language, keeps itself alive. So the knife was in the drawer
ten years before the lecture was printed. Then comes the obvious next question, whether there's any way out or whether
``trust no binary you didn't build'' is the bottom. There's a partial answer, and it took a while.
David A. Wheeler, in a 2005 paper and a 2009 dissertation at George Mason University, worked out
\emph{diverse double-compiling}: compile the suspect compiler's source with a second,
independently made compiler, use the result to compile that source again, and compare the
output bit for bit with the suspect binary. If they match under Wheeler's assumptions, that gives a source-to-binary
correspondence. The second compiler must correctly translate the relevant
source, the execution and comparison environments must be trusted for these
runs, and the build's language and environmental effects must match the
conditions of the argument. Two compilers having different authors is useful
evidence for trust, not a substitute for those premises. The question has
moved to a boundary where we can name what we trust.

Bill Lawvere, in 1969, proved a fixed-point theorem that, as I read it, makes all of these one theorem in disguise: if a set can be mapped
onto all the functions from itself to some target, then every function on the target has a
fixed point; so if the target has a function with no fixed point (a flip), no such covering
exists. Cantor, Russell, Gödel and halting all come out as instances of one shape. Noson Yanofsky retold it
in 2003 without the category theory. I mention it because ``it's the same wall'' is a claim I
make below, and I want you to know it's fair at the level of \emph{proof shape}, and that every
instance has its own boundary of \emph{who counts as a decider}, and that boundary is exactly
where most overclaims live.

\subsection{Applying the diagonal}

Before applying the diagonal to a system with a record, identify who the proposed decider
is and what it must decide. Four familiar applications show why those details matter.

``Gödel shows minds aren't machines.'' John Lucas, a philosopher who spent thirty-six years as a
fellow of Merton College, Oxford, argued it in 1961, in a paper called ``Minds, Machines and
Gödel.'' Roger Penrose, the Oxford mathematical physicist who later shared a Nobel prize for his
work on how black holes form, revived it twice, in two books, \emph{The Emperor's New Mind} in 1989
and \emph{Shadows of the Mind} in 1994. I'm sure of who argued it and when. Whether they're right is
the question, and the standard replies, as far as I can tell, have held up: the Gödel sentence is true \emph{only if the system
is consistent}, so a human who ``sees'' its truth has to know the system is consistent, which
is the hard part; and nobody can build the Gödel sentence of a system they can't specify, so
humans might be machines too complicated to diagonalise. The theorem has a conditional in it
and the argument drops the conditional.

And the part that's easy to miss is that Gödel himself didn't go as far as Lucas. In a lecture
on the day after Christmas in 1951, at a meeting of the American Mathematical Society at Brown
University, he drew from his own theorems only an either-or. Either the human mind, even inside
pure mathematics, infinitely surpasses the powers of any finite machine, or there are problems
about whole numbers, of a certain simple kind, that no proof the human mind can conceive will
ever settle. Feferman, writing about that lecture, says Gödel privately believed the first half
and still stated only the disjunction, because he didn't have a knock-down argument against the
machine side. Gödel wrote soon afterward that he meant to publish the lecture, and never did; the
text sat among his papers until after his death. So the man whose theorem gets quoted for
``minds aren't machines'' stated an either-or and left the deciding to whoever came next. The
conditional is there again, kept by the person with the most right to drop it.

``Undecidable means unknowable.'' Two senses share the word: a problem with no single
algorithm correct on every instance, and a sentence independent of one particular system.
Neither means nobody can know a given answer. Any single instance is decidable, trivially:
for any one program, one of ``say yes'' and ``say no'' is right. What's ruled out is one
\emph{inside} procedure that's right on \emph{every} case.

``Everything interesting is undecidable.'' Rice is about behavioural properties over \emph{all}
programs. It says nothing against checking the text, nothing against deciding for specific
programs, nothing against analyses allowed to say ``I don't know,'' nothing against finite-state
model checking or type checking. Real compilers and analysers live with Rice by being sound but
incomplete, and what limits them in practice is mostly cost. The diagonal shapes them; it doesn't stop them.

``Self-reference is a paradox.'' It's a theorem with a construction, as Kleene, the quines and Thompson show.


Two forms of the wall live in this section, and they're honest in different ways.

The first is the \emph{nat family}, a tiny substrate built so that program two is the
Contrarian for a given candidate decider d, and for every d, d is wrong about program two. It's
a \emph{family}, one instance per d, and its recursion theorem holds by construction. It's the
index cards made formal and nothing more.

\begin{theorem}[The nat family]\label{thm:u12-nat-family}
Fix any Boolean function $d$ on $\mathbb{N}$. Let the substrate $N_d$ have programs and
states $\mathbb{N}$; program 0 returns 1, program 1 returns 0, program 2 returns 0 if $d(2)$
is true and 1 if it is false, whatever the input state, and every other program gives nothing;
each program is encoded as itself; $\mu$ is constantly 0, so any step relation will do;
$\mathrm{Rep}(f)$ holds when $f(p) = 1$ for every $p$ with $d(p)$ true and $f(p) = 0$ for
every other $p$; and $A(p)$ holds when program $p$ returns 1 on input 0. Then it is not the
case that $d(p) = \mathrm{true} \iff A(p)$ for every $p$.
\end{theorem}

\begin{proof}
The recursion theorem of $N_d$ holds with the fixed witness 2: for a representable $f$,
$f(2)$ is 1 or 0 according to $d(2)$, and program 2 returns exactly what that program
returns. The witnesses for $A$ are 0 and 1, and $A$ is extensional because it reads only what a
program returns. The flip of $d$ is representable by the
definition of $\mathrm{Rep}$, so Theorem~\ref{thm:u12-no-inside-decider} applies to $d$ itself. Directly:
$A(2)$ holds exactly when $d(2)$ is false.
\end{proof}

Now push on it once more, because the nat family shows the outside decider in three lines,
and you can do them yourself. Stand outside $N_d$, knowing $d$. Write down a new decider
$e$: say yes on program 0, say ``not $d(2)$'' on program 2, and say no on everything else.
Check it against $A$. Program 0 returns 1, so $A(0)$ holds and $e$ says yes. Program 1
returns 0, so $A(1)$ fails and $e$ says no. Program 2 returns 1 exactly when $d(2)$ is false,
and that's what $e$ says. Every other program gives nothing, so $A$ fails and $e$ says no.
So $e$ is right about every program. Does that break the wall? It doesn't, and the reason is the part that
matters: in $N_d$ a flip is representable only if it says 1 exactly where $d$ says yes, so
the flip of $e$ is representable only if $e$ and $d$ agree everywhere, and they disagree at
program 2. The correct decider exists; its flip just isn't one of the substrate's own moves.
(It's three lines and a pencil, and you just did it.) One and the same property is undecidable from inside and decided cleanly from outside, at the
same time. That's not a paradox, because it's two different people asking.

The second form is the one where the recursion theorem is \emph{proved} from the reduction rules rather than assumed.
It's on \emph{L}, the weak call-by-value lambda calculus that Yannick Forster and Gert Smolka presented in
2017 as a model of computation, and which a group at Saarland has since used as the
seed of a whole library of undecidable problems. Here L is built from its reduction rules. The recursion theorem is proved, Rice's theorem follows, halting
follows, and the diagonal is reproved in the shape of the substrate
theorem, for any property of programs that depends only
on what they compute and that holds of one program and fails of another. The statement
mentions no shortcut. So on L the result is the diagonal for one fixed model, and I don't plug L into the
substrate theorem, for this reason: a substrate's run map has to answer, with a state or with nothing, and for L that map is the halting function itself, which no term of L computes. Using it would take something from outside L, and L is in this book as the model where nothing is taken from outside. So on L I run the argument again from
scratch instead of applying the general one. Both are true; they're different sentences.

I'll say what that string of words means, because L is the strangest thing in this book that
computes, and the smallest. The lambda calculus is Alonzo Church's, from the 1930s, and it's
a programming language with one idea in it: functions. A term is one of three things: a
variable; a function, written $\lambda$ followed by its body, the body saying what to do with
whatever you hand it; or one term applied to another, written side by side. Running a program
is one move, over and over: when a function meets an argument, put the argument into the body
where the variable was. The lambda calculus has no numbers, no memory, no instructions, and in 1937 Turing published a proof that
it computes exactly what his machines compute. ``Weak call-by-value'' is two house rules about
the order of moves: an argument is run down to a finished function (a \emph{value}) before it's
put in, and nothing runs inside a function that hasn't been handed anything yet. Variables
are written as numbers counting how many $\lambda$s back they were bound (a trick of Nicolaas
de Bruijn; it means two variables can never get mixed up by sharing a name), and a term is \emph{closed}
when every variable in it is bound inside it, so it leans on nothing outside. Last, to
hand a program to a program, L needs a way to write any term as data, the way a program's
source is a string you can pass around. The one used here is the \emph{Scott encoding}, named for
Dana Scott, a student of Church's who went on to build mathematical models of the lambda calculus
itself. As far as I can trace it, the encoding first appeared in lecture notes of his from Berkeley
in 1962 that he never published, as a way to write data; Torben Mogensen carried it over to
lambda terms themselves in 1992, and that's the form used here. All you need about it is that every term $t$ gets one \emph{code}, written
$\ulcorner t \urcorner$, and the code is itself a closed value. Two tools get built on top of
that in the proof: a \emph{quoting} term $Q$ that takes the code of a term to the code of that
code (a program that prints its own source needs exactly this move), and a \emph{fixed-point
combinator}, a term that lets a function call itself, which is how $Q$ works its way through
a code of any size.

Do two of these on paper before I give you the general rule. Write
$I=\lambda x.x$: its instruction is ``return whatever you were handed.'' Hand
$I$ its own card and you get $I\,I\to I$, and it stops. Now write
$D=\lambda x.x\,x$: ``make two copies of what you were handed and apply the first
to the second.'' Hand it its own card. Putting $D$ in for both copies of
$x$ gives $D\,D\to D\,D$. You're back at exactly the same instruction and
argument, so the rule never finishes. Both arguments were already
functions, so neither example sneaks any work under a function that hasn't been handed anything.

The variable numbers are what keep the copying from getting confused. In $\lambda x.\lambda y.x$,
the last $x$ points past the nearest $\lambda$ (the one for $y$) to the one before it, so its number
is 1; in $\lambda x.\lambda y.y$ the last variable has number 0. Rename $x$
and $y$ however you like and neither recipe changes. The codes below are a way of putting a whole
recipe on a card that another recipe can read.

\begin{definition}[L, as built here]\label{def:u12-L}
A term of L is a variable (a de Bruijn number), an application $s\,t$, or an abstraction
$\lambda s$; the values are the abstractions. One step is weak call-by-value, left to right:
$(\lambda b)\,v$ steps to $b$ with $v$ substituted when $v$ is a value, a step inside the
left part of an application is a step of the application, and so is a step inside the right
part once the left part is a value. A term $s$ \emph{evaluates to} $v$ when $s$ reaches $v$ in
some number of steps and $v$ is a value; $p \equiv q$ when $p$ and $q$ evaluate to the same
values. Every term $t$ has a code $\ulcorner t \urcorner$, its Scott encoding, which is a
closed value. A map $f$ on terms is \emph{L-computable} when some closed value $s$ has
$s\,\ulcorner p \urcorner \equiv f(p)$ for every $p$. A term $D$ \emph{decides} a property $P$ of
terms when, for every $p$, $D\,\ulcorner p \urcorner$ reaches $\lambda\lambda.1$ (true) if
$P(p)$ holds and $\lambda\lambda.0$ (false) if it fails. $P$ is \emph{extensional} when
$p \equiv q$ and $P(p)$ give $P(q)$.
\end{definition}

\begin{theorem}[Recursion, Rice and halting for L]\label{thm:u12-L}
(1) For every closed value $s$ there is a closed term $t$ that reaches
$s\,\ulcorner t \urcorner$; so every L-computable $f$ has a term $p$ with $p \equiv f(p)$.
(2) Let $P$ be extensional, with $P(y)$ and not $P(n)$ for closed terms $y$ and $n$. Then no
closed term decides $P$. (3) No closed term decides halting, the property of evaluating to
some value.
\end{theorem}

\begin{proof}
(1) First build a program that can assemble the code of the very application
I am about to run. The quoting program $Q$ takes a term's code and returns
the code of that code. Appendix~\ref{sec:quoting-construction} gives the idea of its
construction. The program
$\mathrm{mkapp}$ takes two codes and joins them into the code of an application.

Make $A$ do this when handed a code $x$: quote $x$, put that quoted code
beside $x$ using $\mathrm{mkapp}$, and hand the result to $s$. Written out,
\[
 A=\lambda x.\,s\,(\mathrm{mkapp}\ x\ (Q\,x)).
\]
Now run $A$ on its own code, and call that whole application $t$:
$t=A\,\ulcorner A\urcorner$. The first step puts $\ulcorner A\urcorner$
where $x$ stood. Quoting it supplies the code of $\ulcorner A\urcorner$;
joining the two codes supplies the code of
$A\,\ulcorner A\urcorner$. That is the code of $t$ itself. The computation
therefore reaches $s\,\ulcorner t\urcorner$, as promised. All the pieces
were closed, so $t$ is closed too. If $s$ computes $f$, the resulting
behaviour is exactly $f(t)$.

(2) Suppose a closed program $D$ decides $P$. Build another program that asks
$D$ about the code it receives, then runs the no-example $n$ after yes and
the yes-example $y$ after no. In this language you must keep the two branches
wrapped as functions until the answer chooses one; otherwise call-by-value
would try to run them before choosing. The following term does that:
\[
 \lambda x.\,D\,x\,(\lambda\_.\,n)\,(\lambda\_.\,y)\,(\lambda z.z).
\]
The final identity function is the harmless argument that opens the selected
wrapper. Thus a true answer runs $n$, and a false answer runs $y$.

Apply part (1) to this program, obtaining its self-matching $t$. If $P(t)$
holds, the judge says true and $t$ behaves like $n$, where $P$ fails. If
$P(t)$ fails, the judge says false and $t$ behaves like $y$, where $P$ holds.
Either answer contradicts the requirement that alike behaviour has the same
verdict. No such $D$ exists.

(3) Halting depends only on which values a term reaches. The identity
$\lambda z.z$ is already a value. The term
$(\lambda x.\,x\,x)(\lambda x.\,x\,x)$ takes a step back to itself each
time, and is never a value. These supply the yes-example and no-example
for part (2), so halting has no closed decider either.
\end{proof}

The construction in that proof isn't mine, and the credit belongs right here. Forster and Smolka's
L development of 2017 has the Scott encoding of terms, a closed quoting term with exactly this
specification, the code of a term to the code of its code, and a fixed point built the same way:
their $t$ is $C\,\ulcorner C \urcorner$ with $C = \lambda x.\, s\,(A\,x\,(Q\,x))$, so their $C$ is
my $A$, and their $A$, which builds the code of an application, is my $\mathrm{mkapp}$. They say they
follow the proof in the Dutch logician Henk Barendregt's textbook on the lambda calculus (1984),
where statement (1) is the full lambda calculus's second fixed-point theorem, and part (2) is what
they call Scott's theorem. The idea of the term construction is in Appendix~\ref{sec:quoting-construction};
the route is theirs, and Barendregt's before them.

So I've given the wall on paper, on a family built to fit it, and on one outside model.
The machine built in Part~III (Section~\ref{sec:smallest}) fills in this same shape: on
its host machine, the one the universal program of Appendix~\ref{sec:universal-host} runs on, which has a counter for every whole number,
programs are numbers, Kleene's recursion theorem is a theorem, and so is Rice's, and the fixed point
can be taken to agree on what the program computes and also on the ledger and the flag at the end of every run where they stop (Appendix~\ref{sec:small-diagonal}).

The table below lines up the picture all of that fills in against Turing's, one column per axis.

\begingroup\small
\begin{longtable}{@{}>{\raggedright\arraybackslash}p{0.482\linewidth}>{\raggedright\arraybackslash}p{0.482\linewidth}@{}}
\toprule
\textbf{Time axis} & \textbf{Structural axis} \\
\midrule
\endhead
\bottomrule
\endlastfoot
counts: steps & counts: $\mu$ \\
asks: does it halt? & asks: does it admit the shortcut? \\
the wall: no Turing machine decides halting & the wall: no decider whose flip is representable, under a recursion theorem \\
\end{longtable}
\endgroup

Read the table as two uses of one diagonal, each under its own premises, and as saying nothing more than that.
The structural wall never uses $\mu$; the $\mu$ entry only says what that axis counts. The
nat result is a family indexed by the candidate decider, and L's results are about one fixed
model and properties that mention no shortcut. Neither hands you a second fixed undecidable
problem, and neither separates Turing degrees: a property can sit behind the wall from inside
and be decided cleanly from outside, as the pencil argument on $N_d$ showed. And ``admits the
shortcut'' in the right-hand column is the theorem's sense, any behavioural property with a yes and
a no. ``Has a cheaper way to the same answer'' is a second question, and the column has no wall for it.

\subsection{Where the diagonal applies}

The three versions put the candidate procedure in different places. In the general
substrate argument, the substrate must represent the candidate's flip. In the natural-number
family, each candidate $d$ supplies the substrate whose predicate defeats it. In L, the
candidate is a closed term and the question is a nontrivial extensional property of closed
terms. The reduction rules and recursion argument give Rice's theorem there.

The pencil calculation on $N_d$ is useful precisely because it is outside that substrate.
To run the diagonal inside, I need the representability and recursion conditions stated
in the theorem.


%===========================================================================
\section{From ``What established means''}
%===========================================================================
\subsubsection*{Open the gauge wider}
\label{sec:history-capacity}

Go back to the Sorter. It checked the order of the list. Now ask it another question:
did it also check that none of the original numbers went missing? The list on the screen
can be the same whether that second check happened or not. One lamp saying ``checked''
will not tell me which work was done.

Put a sheet beside the scratchpad, with a named box for each check. Order gets a box.
Membership gets a box. Each other check you want to remember gets its own. When I come
back, I want to know which boxes were ticked. If two runs tick different boxes and leave
the same working value, the full machine has to leave me something that tells them apart.
That is more of the fibre above the same tabletop picture, brought within the question.

Here are the moves. WRITE puts a value $b$ on the scratchpad and starts a fresh sheet,
even if I write the same value again. CHECK $i$ tests its named predicate $P_i$ and records
its name when it passes. A failed check changes nothing in this monitor. IDLE leaves both
scratchpad and sheet alone. The reader can see the working value and ask about every
check name: did it pass since the last WRITE? The sheet must answer yes when it did and
no when it didn't.

Take a finite list of named checks, and let $m(b)$ count the ones whose predicates are
true at $b$. Each of those can be run or left unrun. The truths themselves may overlap;
it is my choice to perform each check that is independent. Four true predicates give
sixteen possible ticked sheets above the same value. The general count follows in the
same way.

\begin{theorem}[Exactly how many current sheets]\label{thm:current-sheet-capacity}
For a finite nonempty working set $B$ and this WRITE/CHECK/IDLE interface, an exact host
needs at least $2^{m(b)}$ reachable full states above each $b$. The minimum total reachable
state count is
\[
\sum_{b\in B}2^{m(b)}.
\]
\end{theorem}
\begin{proof}
Hold the working value fixed and start with two named checks that would pass
there. Run neither, the first alone, the second alone, or both. The working
value is the same in all four cases. The required answers are not. In particular,
the first-only case and the second-only case cannot finish at the same full
state: ask whether the first check ran, and that one state would have to answer
both yes and no.

Each further check that would pass doubles the choices. For every sheet you
already had, you can leave its new box blank or run that check and tick it.
Starting from one empty sheet, $m(b)$ such choices give $2^{m(b)}$ sheets. Every
sheet is reachable: WRITE the value $b$, then perform exactly the checks whose
boxes it ticks. Any two different sheets disagree at some named box, so they
need different full states. States showing different working values are also
different. Add the required states for each value to get the displayed sum.

Now build a machine with exactly that many states. A state consists of the
working value and its ticked sheet. Include every sheet that ticks only checks
true at that value. WRITE replaces the value and clears the sheet. CHECK ticks
its own box if its predicate is true, and otherwise leaves the sheet alone.
IDLE changes nothing. These rules never produce a forbidden sheet.

Start with the empty sheet. After a WRITE it is empty again, exactly as the
task requires. After a successful CHECK it has precisely that success added;
a repeated success changes nothing further. A failed CHECK adds nothing, and
IDLE adds nothing. If the sheet described the history correctly before a
command, it therefore describes it correctly afterward. Following any finite
list of commands proves the promised answers. The WRITE-and-check runs above
reach every state I included. The lower count is attained.
\end{proof}

Try it with just one box. The workspace holds false or true, and CHECK asks whether it
is true. The whole monitor fits in three states: false, true-unchecked, true-checked.
WRITE false goes to the first. WRITE true goes to the second. CHECK moves the second
to the third and leaves the other two alone. IDLE leaves all three alone.
Start at false. WRITE true reaches true-unchecked, and then CHECK reaches true-checked;
all three states are reachable. Now take any word of instructions. A WRITE makes the
verdict false and starts a new sheet, even if it writes the same value. A CHECK on false
leaves the verdict false; a CHECK on either true state makes or keeps the verdict true.
IDLE preserves it. These are all the cases, and each matches exactly what happened since
the last WRITE. Applying that argument after each instruction proves the promise for
every finite word. In particular, the first arrival at true-checked must use CHECK:
neither WRITE nor IDLE can enter it from another state.

The two true states also show you the merger. CHECK sends true-unchecked and true-checked
to the same next state, right above the same working value. A larger implementation can
move that merger later: let CHECK enter a finite region of accepting states in the fibre
and keep it there, and some two distinct states must share a CHECK successor. WRITE can
still clear the current sheet. Permanence under that one check and permanence under every
instruction are different promises.
 Call that finite accepting region $H$ and the entering state $u$.
The move sends all of $H\cup\{u\}$ into $H$. There are $|H|+1$ distinct sources
and only $|H|$ possible targets, so two sources share a successor. The move preserves
the working picture, so both sources and their common successor lie in that same fibre.

\subsubsection*{Keep the old receipts when the page is wiped}

Now wipe the scratchpad and use it again. I still want the receipt for the first page.
Then use it a third time. I want that receipt too. Think of carbon copies filed before
each sheet goes in the bin. On the tabletop, you can be back where you started. In the
file, you have a past you didn't have before.

On each of $n$ editions, give me the choice to certify it or leave it uncertified.
Bring every choice back to the same final working picture $b$. Then ask the machine
about any edition I name. It has to answer whether that edition was certified. Two
histories that disagree on that answer need two different full states, however alike
their tabletops look. There are $2^n$ possible yes-or-no lists to distinguish.

\begin{theorem}[Room for independently retained editions]\label{thm:edition-fibre-capacity}
Under those realization and exact-query requirements, the fibre above $b$ contains at
least $2^n$ reachable states. If one fixed host meets the requirements for every finite
$n$ at that same $b$, no finite set covers its reachable fibre.
\end{theorem}
\begin{proof}
Two different lists of length $n$ differ at an edition. Equal endpoints could not answer
that edition's query differently, so the endpoints are distinct. A putative finite cover
of size $K$ fails once $2^n>K$. Every endpoint being counted has the same working view $b$.
\end{proof}

You can build the witness directly. Keep a Boolean workspace, the answers for the
completed editions, and the current stage of CHECK, COMMIT and CERTIFY. On a true
workspace, run the chain for an edition you choose to certify, or skip it for one you
don't. Then WRITE true to start the next edition. Every choice brings you back to the
same true workspace and a fresh current stage. Open the file and the choices are still
there, in exactly the places you made them.

Put the whole machine on the table. There is the current page, holding false
or true. There is a marker beside it, with four positions: raw, checked,
committed, published. And there is a numbered file of yes-or-no answers for
pages already put away. Start with a true page, the marker at raw, and an
empty file.

CHECK can move the marker from raw to checked, but only if the page says true.
COMMIT can move it from checked to committed. CERTIFY can move it from
committed to published. Press any of those buttons at another stage and it
leaves the arrangement alone. IDLE always leaves it alone. Those are all the
ways to move the marker without replacing the page, so there is no shortcut
from raw to published.

WRITE does three things, in this order. It files yes if the marker is at
published and no otherwise. It puts the requested value on the page. Then it
returns the marker to raw. The answer goes into the file before the marker is
reset, and no command alters an answer already filed.

Check what that means after any command. CHECK, COMMIT and CERTIFY move the
current page through exactly the named stages; they do not touch earlier
pages' answers. IDLE leaves both alone. WRITE closes the current page with
the answer its stage warrants, then opens a raw page. Start with no earlier
pages and follow those rules one command at a time: the file always answers
the completed-page questions correctly, and the marker gives the current
page's stage.

To file a yes, run CHECK, COMMIT, CERTIFY, then WRITE true. To file a no,
use WRITE true alone. Either choice leaves a true page with the marker at
raw, ready for the next choice. After $n$ choices the page and marker are
identical in every run, while the file holds exactly the $n$ answers chosen.
Every yes-or-no list is possible, and different lists are different full states.

Keep asking for another edition, with the same freedom of choice and the same right to
read back every answer. Whatever finite number of full states you offer above $b$, a
long enough history outruns it. The tabletop can keep coming back to one picture.
The machine that keeps all those choices cannot fit its reachable fibre into a finite set.
That is where the infinity comes from.

The old yes-or-no gauge still has a job. Ask whether the file contains any certification
at all and it becomes a permanent lamp: once yes, always yes. Ask whether today's sheet
is certified and wiping that sheet starts the answer over. The small native machine in
the next chapter has the permanent lamp. I will look behind it at the whole native state,
then count which publication histories even that larger view leaves out. That camera sees
more than the working value here; the next count will name its fibre just as explicitly.

\subsection{The guy who got closest}

Somebody had that clue before me, and on that clue he's the closest neighbour I've found, so he gets the most space
and the hardest push.

Bennett again, the one from Part~I's counting of states and the demon, wrote a paper in 1988 that starts where I just ended up:
information isn't the same as value. His example is an almanac of where the moon and planets will
be every day for a hundred years. It has no more information than the equations it was calculated
from, but it saves you doing the calculating. So the value is the work. He called the amount of
work a thing carries its \emph{depth}.

Was this my idea, already done, by a guy at IBM? The only way to find out was to find where it
parts from mine, because if I couldn't find that place I'd just be renaming his idea. And it's one
word: \emph{plausibly}. Depth is the amount of work a thing was \emph{plausibly} made with. Bennett says
straight out that you can't know how a message actually came to be; for all you know it came from
flipping coins. Depth is read off the finished thing, the way a fossil's age is read off the
fossil.

Bring the two computers back. The sorted list in the first one has some depth. The identical list in
the second has exactly the same depth, because depth looks at the list and nothing else. A
photocopy of a deep thing is just as deep. Depth tells you what work a plausible way of making the thing would take.
It's not a record that \emph{this} machine did it, and it's not a receipt for what doing it
cost.

So Bennett's standing right next to the line, on the other side of it, and as far as I can tell he
got his own question right. He asked how much work a thing embodies. I'm asking who has the right to
lean on it and what they paid to get there. His answer lives in the object. Mine, if I've got one,
has to live in the machine.


\subsection{What putting it in the step buys}

The hardest objection is one I'd rather make myself. You defined a flag and a charge,
then proved the charge gets paid when the flag goes up. That's your definitions agreeing with each
other.

On the bare model, that's fair, and the copier's receipt for nothing is the proof. What makes the
record mean \emph{established} is the small machine's step, which won't let a commit through unless a
check of that very claim passed, meeting the first three of the four conditions from Gettier's
pocket, with the fourth, an honest checker, proved about the checker separately. And what it proves
isn't about one program: on \emph{every} list of instructions anybody can write, starting
clean, if the record is up, then somewhere earlier a check of some claim, about some edition,
passed; then a commit of that same claim; then the certificate; with nothing touching the
thing between the check and the commit.

So a raised flag, from a clean start, tells you some claim went through a passing check, a
commitment to it about the same edition, and a certificate. That piece of history is what the
flag keeps, and only that. The channel names the latest commitment that got through, which
needn't be the one the flag first went up on, and in general you can't read off the end state
which claim was certified at that first rise, or when. And the flag doesn't say the claim still
holds, at the certificate or at the end, because the counter can change after the commitment.

That ``every'' is what the operation-by-operation proof buys. A log written
by one program inherits that program's correctness conditions. A database
or reference monitor can instead enforce rules over every operation its
interface permits. The small machine does that too: its invariant covers
every string of its instructions from a clean start, including hostile
strings. The coverage comes from those restrictions and their proof; it
does not follow just from calling the record part of the step. And the tie
between a check and the thing it checked gets a job: the step reads it, and refuses a commitment
when it no longer fits, so the structure is doing something in the machine instead of sitting in a note
beside it.

There are two limits on that ``every,'' and I'll say them now so you don't have to find them later. First, it's every program,
not every everything. The guarantee covers any list of instructions, run from a clean start, one step at
a time, on a step that does what its definition says. Hostile programs are inside it. A hostile start
(a table of facts nobody checked, a reset), hostile hardware (a soldering iron on the flag's bit,
a second writer the step can't see) or a hostile reader isn't; Section~\ref{sec:scope} lists the
whole fence. The guarantee has the same shape as a database or reference monitor's:
it covers operations mediated by a faithful enforcing layer. How large that
class is depends on each interface; being a database does not make it a
one-program guarantee.

Second, the flag and the ledger on their own are what people who verify programs call ghost state:
erase them and every program computes exactly what it did before (Section~\ref{sec:counter} has
where that idea comes from). What isn't ghost on the small machine is
the table of checked facts, which the step reads and refuses on, at run time, for every program. And
the move of putting a rule under every program isn't mine either. It's the reference monitor's,
and its history is a few sections back.

\subsection{The record, the checker and the host}

Take the flag and ledger off the small machine's drawing. The core follows the same steps:
neither field tells it what to do next. Put them back and you can read whether certification
occurred and what the instruction schedule charged. That is the passive record's job. The
latch argument in Section~\ref{sec:latch} gives its form over any fixed base when the record
is driven by that base and never comes down.

The checker has another job. CHECK of ``$A\ge m$'' tests the counter in one instruction;
a program using only the two counter moves would have to work through their operations.
A step comparison therefore starts with the instructions each machine is given. Every
checker built in this book is a program. The abstract certified interface also admits a
supplied checker, so its computational power includes whatever that checker can do: supply
a halting oracle and CHECK can answer halting questions. The effective implementations here
use effective checks.

Now put a guest program inside the host. A guest that can write any value into its own
``checked'' variable can write YES without checking. The host I build gives the guest three
entry points instead: CHECK, COMMIT and CERTIFY. Its step rule controls those entry points
for every allowed guest program. The fact table and channel determine whether the next act
is permitted; the flag and ledger record the result. An ordinary universal machine can
implement these rules. The work is in preserving them through every guest instruction.

The universal program in Appendix~\ref{sec:universal-host} does that work with the small
machine's six instruction kinds on a host with more counters. It fetches a guest instruction,
decodes it, and carries it out. At matching points the guest and host ledgers agree. The
host halts exactly when the guest halts, keeps the guest's answer in its counters, and raises
its own flag exactly when the guest raises its flag, through its own earned chain. Its guests
are the small machine's programs from clean starts, with Minsky's encoding supplying their
universal computation. Appendix~\ref{sec:simpler-host} gives a simpler picture with a guest-step
operation supplied directly; the instruction-by-instruction construction shows how to build it.

Add a one-unit PAY instruction and the host can also follow a computably presented Thiele
machine: its selected moves, transitions, prices and total yes-or-no reading are all given
by programs. Appendix~\ref{sec:presented-universal} constructs that interpreter. Starting
with the guest's reading at no, the host follows its ledger with at most two extra units at
the raise. Starting with the guest already at yes, the clean host must establish its own
flag immediately and spends three. The same clean-start three-act count explains the
surcharge when the guest certifies with a whole ledger below three.


%===========================================================================
\section{From ``The earned-chain counter example''}
%===========================================================================
\subsection{Where the counters come from}
\label{sec:counters-origin}

The base first, and the base isn't mine. It's Marvin Minsky's. In a 1961 paper in the
\emph{Annals of Mathematics}, and then in his 1967 book \emph{Computation: Finite and Infinite
Machines}, Minsky studied a machine so bare it's almost a joke. Picture a few bins of
pebbles and a numbered list of instructions. There are only two kinds of instruction: put
a pebble in bin $c$ and go to the next line; or take a pebble out of bin $c$ if there's one to
take, and go one way if there was and another way if the bin was empty. There's nothing else: no memory but
the bins, no tape, no arithmetic beyond plus one and minus one. When the line you're on has no
instruction, the machine stops.

What Minsky proved is that two bins are enough to compute anything a Turing machine can
compute, as long as you're willing to pack the input in a strange way. The trick is to keep
many numbers inside one number, as the powers of different primes ($2^x 3^y 5^z$ holds
three numbers in one), and to do arithmetic on the exponents by multiplying and dividing
with the second bin as scratch space. Packing numbers that way can turn a modest input
into an enormous pile of pebbles, handled one at a time, so it's no way to build a fast
machine. It's a way to settle a question: with that packing, the two bins can do every
Turing computation. What the bins give me here is a base whose ordinary moves you can see
all of. Each of the earlier results says which properties of a base it needs, and
Appendix~\ref{sec:every-model} carries the earned layer over to other bases once you've named
what they have to supply.

Who cares? Minsky did, because the question then was where the edge of computation sits and
how little machinery reaches it. You should, because it means the base of this chapter's
machine hides nothing: if it can compute everything with plus one, minus one and a jump,
then nothing about the record can be smuggling computing power in through some instruction
I made fancy. For this construction I use Minsky's two-counter halting result and the counter constructions of Forster, Larchey-Wendling and their coauthors. Appendix~\ref{sec:counter-constructions} sets out the reduction; the step-for-step connection to this machine is the part I have to establish.

And the pebbles aren't only my picture. They were somebody's actual machine, the same year.
In one issue of the Canadian Mathematical Bulletin, September 1961, Z.~A. Melzak, a
mathematician at the University of British Columbia, described a
computer made of distinguishable holes and an unlimited supply of pebbles, with one
three-place operation, and a few pages later Joachim Lambek, a mathematician at McGill
University, trimmed it to two moves: put a
pebble in hole $X$; take one out of hole $X$ if there's one to take. He called it an
\emph{infinite abacus}, and early on he writes ``Our method is probably not new,'' which is
a hard sentence to beat. He also tells anybody who distrusts infinity not to worry
about all those holes, since any one computation only ever uses finitely many of them.
Is that two inventions or one? Lambek says his was a modification of Melzak's, so call it one
and a half, with Minsky's as a separate find the same year. And then Lambek went somewhere
you wouldn't guess: in 1986, with Philip Scott, a logician who had come to McGill to work with
him and later taught at the University of Ottawa, he wrote a book whose first part, as the
publisher describes it, shows that typed lambda
calculi and a certain kind of category are essentially the same thing. He went from pebbles to category
theory in one career, twenty-five years end to end. Does any of it change the argument? No. The
machine uses Minsky's version and Minsky's name. But a bare machine that several people
reach for in the same year looks like a shape that was sitting there. It doesn't look much like
anybody's invention.

The title of Minsky's 1961 paper is the clue I can't leave alone. It's about ``Post's problem of
tag.'' Who's Post, and what's tag? Emil Post was a logician, a fresh PhD from Columbia, who in
1921, on a fellowship at Princeton, was playing
with about the simplest rewriting game there is. Take a word of 0s and 1s. Look at its first
letter, chop the first three letters off the front, and stick something on the end:
\texttt{00} if the letter you looked at was 0, \texttt{1101} if it was 1. Repeat, and stop when
fewer than three letters are left. Try it at the table. Start from \texttt{1000}: you get
\texttt{01101}, then \texttt{0100}, then \texttt{000}, then \texttt{00}, and it stops. Start from
\texttt{100}: you get \texttt{1101}, \texttt{11101}, \texttt{011101}, \texttt{10100},
\texttt{001101}, and then \texttt{10100} again, and it goes round forever. Post wanted a rule
that tells you, for every starting word, which of those happens. He worked at it for months and
later called it hopeless, and as far as I can tell it still is: a hundred years on, nobody knows
whether that one little game, with that one rule, can be decided. Stephen Wolfram, the physicist
who built Mathematica, wrote a 2021
article on it that tells the Princeton story, and it's where I'd start. In his 1961 paper Minsky proved
that tag games in general can compute anything, which is why Post's name is on it. It didn't settle Post's
own game. And the chase goes one step further. In 2008 Liesbeth De Mol, a historian and philosopher
of computing who wrote her doctoral work on tag systems, turned
the famous $3n+1$ problem (halve an even number, triple an odd one and add one, and ask whether
you always get down to 1) into a small tag game of the same kind. So is Post's game secretly as
hard as everything, or just stubborn? I can't tell you, and as far as the published work goes,
nobody can yet.

None of that changes the machine below. It changes the picture. The two-counter machine came
out of a hunt for the smallest thing that can't be predicted, and the smallest things are exactly
where the open questions are still sitting.

There's one word clash to settle before it bites. Minsky's machine is usually called a \emph{counter machine},
and its bins are \emph{counters}. Section~\ref{sec:counter} also had a counter: the ledger $\mu$, the number that only
climbs. It's the same word for two completely different things, and I'd love to say I set that up on
purpose. In this chapter, \emph{counter} means one of Minsky's two bins, called $A$ and $B$, and
the climbing total is always \emph{the ledger} $\mu$.


\subsection{Nothing gets forged}

The theorem and its proof are in the book; this is what stood beside them.

The closest earlier result is in capability hardware, and it's much bigger. The computer
scientists Jack Dennis and Earl Van Horn, at MIT, introduced \emph{capabilities} in 1966: tokens a
process has to hold to touch an object, which it can pass on and can't forge. CHERI, a
capability extension of real instruction sets built at the University of Cambridge and SRI
International, rests on that idea, and in 2020 Kyndylan Nienhuis and twelve coauthors proved, for the whole CHERI-MIPS instruction set, that arbitrary code given
some capabilities can't increase them as it runs (up to any change of protection domain). Their
result says nothing gets forged across a full instruction set. Theorem~\ref{thm:u34-no-forging} is the same
shape, one instruction at a time, for six instruction kinds and a table of checked facts.

\subsubsection*{The actual counter covers every added fact}

Put the table beside the bill. Every passing CHECK writes one more entry and adds one
to the ledger. A failed CHECK writes no entry. COMMIT and CERTIFY can add to the bill
without adding a table entry. Walk down the run one instruction at a time: the table
can never grow faster than the bill.

\begin{theorem}[The native bill bounds the table]\label{thm:native-fact-bill}
For every clean run, whether driven by a finite instruction word or by the stored-program
step for a finite number of ticks,
\[
\operatorname{len}(\mathrm{facts})\le\mu.
\]
More generally, table growth is exactly the number of passing CHECKs, and is at most the
increase in the ledger from any starting state.
\end{theorem}
\begin{proof}
A passing CHECK adds one entry and one unit to the bill. A false CHECK, a
full-table CHECK or a trapped CHECK adds no entry but still adds one unit.
COMMIT and CERTIFY add one unit and no entry. INC, DEC and HALT add neither.
No instruction removes an entry.

Place a mark beside each passing CHECK in the run. The final table has
exactly one more entry for every mark. Each mark also accounts for one
unit on the bill; unmarked instructions can add further units, never take
them away. So table growth is exactly the number of marks and at most the
ledger's increase. This works from any starting state.

At a clean start the table is empty. The initial ledger is a nonnegative
whole number, so the table length is at most the final ledger. At the
specified loaded starts the ledger begins at zero as well. For a stored
program, every running tick executes one of the six instructions just
counted, and every stopped tick leaves both table and bill alone. The same
count therefore applies to any finite number of ticks.
\end{proof}

One passing CHECK from a clean start gives one entry and a bill of one. Repeat a passing
check and its duplicate entry counts as another line in the table. The tally follows
what was written there. COMMIT and CERTIFY can then make the bill larger while the
number of lines stays put.

\subsection{A full table traps}

The table has room for sixteen facts and evicts none of them, so there are two halves to prove:
nothing ever leaves, and nothing ever overflows.

\begin{theorem}[Bounded table, no eviction]\label{thm:u34-bounded}
Every instruction keeps every fact already in the table. If the table holds 16 or more
facts, CHECK $p$ $c$ puts the trap up and leaves the table as it was. If the table holds at
most 16 facts, it still holds at most 16 after any instruction.
\end{theorem}

\begin{proof}
With a raised trap no instruction changes the table. With the trap down, only CHECK
can change it, and a successful CHECK adds one fact without removing an entry.
Success requires a previous length strictly below $16$, so the new length is at most
$16$. At length $16$ or greater that branch is unavailable: CHECK raises the trap
and preserves the table. Every other instruction preserves it as well.
\end{proof}

Why trap? Why not make room? Because ``making room'' means deciding which fact dies, and a
fact that can die is a fact whose absence you have to explain later. A table that only grows keeps every checked fact available to read back, and the shared-memory result above leans on exactly that, since its oldest fact can never move. Provenance doesn't: a table that dropped its oldest fact to make room would still hold only facts a CHECK wrote. What
the choice gives up is real, though, and it's in the subsection near the end of this chapter on
what the machine doesn't do.


\subsection{What it takes to be certified}
\label{sec:certified}

Section~\ref{sec:complete} set the bar worth meeting, Thiele-complete,
and it says nothing about what the record is about. A clock fails it, and so does a machine whose
record never goes up. The small machine clears it (Theorems~\ref{thm:u34-tc-instances}
and~\ref{thm:tc-decomposition}(2) below), and
it does more: from a clean start, its flag goes up only on top of a check of a claim, by a checker proved to mean what
the claim says, then a commitment to the same claim with nothing changed underneath, then a
certificate. A Thiele-complete machine can raise its record with one mark that goes up from any
state at all, a receipt for nothing (Section~\ref{sec:established}). A certified one can't. The
second definition is the first one with three things added, and it's written so that any machine
can be held up against it. You name a language of claims first, then hand over a reading of the
machine in that language, a list of things, and then the reading has to pass four clauses.

\begin{definition}[Certified Thiele-complete]\label{def:u34-thiele-complete}
A \emph{machine} has states, moves, a step, a cost for each move and a yes/no reading of each
state, the \emph{record}. A \emph{claim language} for it is fixed first: a type of \emph{claims}
with an exact equality test; a \emph{meaning} for each claim, a statement about a state; for each
claim, a relation ``the thing this claim is about is unchanged'' between two states; and a proof
that a meaning that holds keeps holding while the thing it's about is unchanged. An
\emph{interface} over that language hands over the rest: a \emph{universal base}, meaning
a two-counter configuration shown by each state (its \emph{window}), a class of \emph{live}
states, a move for each two-counter instruction, and a loaded state for each start $(1, a, b)$;
a \emph{checker} that says yes or no to a claim at a state, with a proof that a yes means the
claim's meaning, as the language gives it, holds of that state, with a CHECK of a claim passing at a state only when the checker says yes to it there; a set of \emph{clean} states; a
\emph{ledger}, a whole number for each state; and a \emph{kind} for each move: a base move, CHECK
of a claim, COMMIT of a claim, or CERTIFY. The machine is \emph{certified Thiele-complete over} the language
when some interface over it meets four clauses, and \emph{certified Thiele-complete} when some interface,
with claims, meanings and ``unchanged'' of its own, meets them.
(a) \emph{Universal base.} Each two-counter instruction's move is a base move; from every live
state it acts on the window exactly as the instruction acts on a configuration, and the state
stays live; each loaded state shows $(1, a, b)$, is live and is clean; no base move changes the
record; and a record that's up stays up under every move.
(b) \emph{Earned record.} Clean states have the record down. Every run from a clean state that
ends with the record up contains a CHECK of some claim $c$ that passes, then a COMMIT of $c$
with the thing $c$ is about unchanged at every state in between, then a CERTIFY, with the record
down just before that CERTIFY and up just after it. A passing check means its claim's meaning
holds of the state it checked, and a meaning that holds keeps holding while the thing it's
about is unchanged.
(c) \emph{Exact prices.} Base moves cost 0, each CHECK, COMMIT and CERTIFY costs exactly 1, and
each move adds its cost to the ledger.
(d) \emph{Non-vacuity.} Some claim $c$ has a CHECK, a COMMIT and a CERTIFY such that those three
moves alone, from a loaded state, raise the record exactly when $c$'s meaning holds there; and
$c$ holds at some loaded state and fails at another.
A Thiele machine \emph{has a universal base} when some interface gives it one and nothing more is
asked: each two-counter instruction's move acts on the window exactly as the instruction does from
every live state and keeps it live, and each start $(1, a, b)$ has a live loaded state showing it.
Nothing is asked about the record, the prices or the checker.
\end{definition}

An interface over a language becomes an interface
of the plain kind by taking the language's claims, meanings and ``unchanged'' as its own, and the four clauses are asked of that.

\begin{theorem}[Certified is Thiele-complete plus certification]\label{thm:tc-decomposition}
(1) A machine is certified Thiele-complete through an interface exactly when it is Thiele-complete
through the same interface read without its claims (its base moves as the computing moves, CHECK,
COMMIT and CERTIFY as the marks, the same universal base and ledger), and three more things hold:
its loaded states are clean; clause (b) holds, so clean states have the record down, every run
from a clean state that ends with the record up contains the earned chain, a passing check means
its claim holds, and a claim that holds keeps holding while what it's about is unchanged; and
clause (d) holds,
a claim whose three moves raise the record from a loaded state exactly when it holds there, true
at one loaded state and false at another. (2) Every certified Thiele-complete machine is
Thiele-complete. (3) Not the other way: the two counters with two marks, one that raises the
record when $A$ is 0 and one that always raises it, and a ledger beside them that each move adds its cost to, is Thiele-complete, and certified
Thiele-complete through no interface.
\end{theorem}

So certification is exactly three things on top of the bar: clean starts, an earned record (the
chain, a checker proved sound, claims that stay true while what they're about is unchanged), and a
claim that could have come out either way. Everything Section~\ref{sec:complete} proved holds of
every certified machine. What this section adds is what those three things buy, from a clean start: a floor of three instead of one, a raise that is always a CERTIFY, and no receipt without a check; and, from the loaded states, a check that can fail.

The language comes first for a reason. If the interface could name the meanings, it could take ``the
checker says yes'' as the meaning of every claim, and checker soundness would hold for nothing. In
the language-first form the meanings are written down before any checker is handed over, so ``a
checker proved to mean what it says'' has the content it has on the small machine: what the claim
says is the language's, and the proof is that a yes from the checker makes it true. The theorems
that build certified Thiele-complete machines ask the checker for one thing more, that it say yes exactly when
the claim holds, because clause (d) needs a check that can come out yes. The plain form
stays because the theorems that follow hold for any interface at all, so they hold over every
language too.

Read the four clauses as four plain sentences. (a) It computes anything a Turing machine can,
under Minsky's $2^n$ encoding, done outside the machine, and the computing can't reach the record: the base moves are the two-counter machine, they
leave the record exactly
where it was, and a record that goes up stays up. (b) From a clean start, the record goes up only on
top of the small machine's chain: a check that passed, a commitment
to the same claim with nothing changed underneath, then a certificate. (c) Computing costs nothing on the ledger and each of the three acts costs one unit, so the ledger
counts the acts. The steps in between add nothing to it. (The unit belongs to the ledger; it isn't a unit of heat.) (d) Somewhere there's a claim whose check could have said no. A
machine doesn't qualify over a language whose every claim is ``at least 0'', and neither does
one with nothing to certify.

And look where the checker's honesty sits. In (b), ``a passing check means its claim holds''
is a condition on the interface, proved once, outside every run, exactly where it sat for the small
machine (Lemma~\ref{lem:u34-eval}), and over a claim language it's proved against meanings the
interface didn't write. The definition doesn't ask the machine to vouch for its own
checker. The nearest reasons not to ask are Gödel's second theorem, about a system proving its
own consistency, and the diagonal of Section~\ref{sec:diagonal}, about deciding a property of every
program. Neither forbids a machine checking a proof about a checker, so keeping that proof
outside is a choice.
Nor does it ask the checker to be something a machine could run. Clause (b) wants a yes or no for
every claim at every state, sound for the claim's meaning, and nothing in it says that answer has
to be computable. Every instance in this book has a checker that's an actual program, so anything
those machines certify could be listed by a program that runs them. For example, take a function that answers whether an arbitrary program runs forever.
No program computes it (that's Turing's halting theorem). Hand the definition a checker built on it and it's still met, and the machine certifies claims no program could ever confirm. Nothing
in this book builds one. So ``a certified Thiele-complete machine certifies
exactly the claims a program could list'' is true only with ``and its step and its checker are computable'' added, and the definition doesn't require either.

The same scope goes on Section~\ref{sec:established}'s ``it isn't more computing power.'' The record (the flag
and the ledger) is never read by the step that computes the next core, so
erasing it changes no run, and it makes nothing faster. The checker is part of the machine under
the record. With an oracle checker that machine decides what the oracle decides, since whether a
CHECK traps is there for anyone to see; with a checker that's a program, it does what programs do.
And a checker that's a program still makes CHECK a test at unit cost.

\begin{theorem}[CHECK is a unit-cost test]\label{thm:u34-check-unit}
From any start $(a, b)$, the single move CHECK $\mathsf{Ge}(m)$ $A$ costs 1, and the trap latch is
down after it exactly when $a \ge m$.
\end{theorem}

The two counters alone look at $A$ only through the zero test of DEC, and started one pebble below a threshold or right at it, that test can't come out differently until the pile has been counted down to zero, so they need about as many steps as the threshold for what CHECK does in one. Theorem~\ref{thm:time-tax} gives the precise lower bound and the matching ladder. So the record makes nothing faster,
and CHECK, a primitive test, can.

\emph{Three kinds of exact.} The word exact does three jobs in this book, and each job keeps one
name. \emph{Rise-only pricing} charges at least one unit on every step that raises the reading and
nothing on any other step; its tightest form, one unit per raise and nothing else, is the
\emph{exact schedule}. \emph{Instruction-exact pricing} charges every instruction a fixed listed
price. On this machine that's one unit each for CHECK, COMMIT and CERTIFY and
nothing for the rest, and the definition of certified Thiele-complete calls it \emph{exact prices}, not a toll, because it charges acts that raise nothing. CHECK
and COMMIT raise no flag, so the small machine is instruction-exact without being rise-only, and
no rule is broken by that. What the three units pay for isn't one thing three times. A passing CHECK adds a fact to a table that holds sixteen, so it takes a fresh state and fills one slot of a bounded store; a passing COMMIT overwrites the channel and a passing CERTIFY raises the flag, each sending two states with the same base to one (Theorem~\ref{thm:record-merge-toll}). And each of the three, when it fails, sends the failed state and its twin with the trap already up to one state. So a passing CHECK pays by taking a fresh state, and every one of the three units is also a price on a merge. \emph{Simulation-exact accounting} is a host whose ledger equals its
guest's ledger at every point where the two runs match, as the universal program $U$'s does (Appendix~\ref{sec:universal-host}). A
machine can meet any one of the three and fail the others.

One word needs pinning before the theorems, using those three names.
In clause (c), \emph{exact prices} means instruction-exact pricing: each of the three acts, CHECK
and COMMIT included, costs exactly one, so the ledger counts the moves filed as record moves.
Those are the acts that make a claim, and also a check that fails, a COMMIT of a fact already
committed, and a CERTIFY with the flag already up: each is filed as a record move and costs one. Back in
Section~\ref{sec:counter}, one result is about
rise-only pricing in its tightest form, the \emph{exact schedule}: one unit on the step that
raises the flag and nothing on any other step. Measured against the flag alone, the small machine
isn't rise-only, since CHECK and COMMIT raise no flag and cost one each; measured the way (c)
asks, it is instruction-exact. Neither is wrong. They answer different questions: does the ledger
count the acts, or does it charge the flag and only the flag? Simulation-exact accounting, the
third name, comes up with the universal program $U$ in Appendix~\ref{sec:universal-host}.

\begin{theorem}[What every certified Thiele-complete machine has]\label{thm:u34-tc-consequences}
Let a machine be certified Thiele-complete through an interface. (1) It obeys the toll and has a
universal base; its base runs every two-counter program step
for step from every live state,
and halting from a loaded state matches halting from the configuration it shows. (2) Along any run, the ledger rises by exactly the number of
CHECK, COMMIT and CERTIFY moves in it. (3) A run from a clean state that ends with the record up contains at least three of those moves and raises the ledger by at least 3, and the three moves of clause (d), from a loaded state where its claim holds, raise it by exactly 3. (4) Some CHECK fails
at some loaded state, on a claim that is false there, and some run from some loaded state raises
the record. (5) On a run from a clean state, the move that raises the record is a CERTIFY. (6)
Every move costs 0 or 1, and some move costs 0.
\end{theorem}

\begin{theorem}[The language first, and what a certificate then means]\label{thm:u34-tc-over}
(1) A machine that is certified Thiele-complete over some claim language is certified Thiele-complete, so every part
of Theorem~\ref{thm:u34-tc-consequences} holds of it. (2) If a machine is certified Thiele-complete through
an interface whose claims have an exact equality test, it is certified Thiele-complete over the language
made of that interface's claims, meanings and ``unchanged,'' with clause (b)'s proof that a
meaning keeps holding as the language's own. (3) If a machine is certified Thiele-complete over a claim
language $L$, some interface over $L$ meets the four clauses and has this property: every run
from one of its clean states that ends with the record up contains a CHECK of a claim $c$ of $L$
that passed, then a COMMIT of $c$, and $L$'s meaning of $c$ holds at the state the CHECK
saw and at the state the COMMIT saw.
\end{theorem}

Part (2) says where the strength of the language-first form lives: in the naming. Any reading can
be turned into a language when its claims come with an exact equality test, so for those
``certified Thiele-complete over some language'' asks no more than certified Thiele-complete does. ``Certified Thiele-complete over $L$,'' for an $L$ written down without looking at
the checker, says that what the machine certifies is what $L$ says. Every instance below is
stated that way.

Each item in the proof names the clause it comes from, and that's all the definition gives.
A certified Thiele-complete machine needn't have a fact table, a CHECK that traps when it fails, a COMMIT
that looks anything up, or claims about counters. The definition asks for an earned chain on the
runs that raise the record and leaves the rest to the machine. So the small machine's table
provenance, its trap, and its channel are facts about that machine, and they don't carry over to
every machine that meets the definition. Through a plain interface its claims can be any labels
the interface names; over a claim language they are the language's. Either way its checker can
be an oracle, as the halting-function checker above shows.

The second half of (6) in Theorem~\ref{thm:u34-tc-consequences}, that some move costs 0, is the one
that rules out the clock: a machine that charges every move. It doesn't rule out a program that
makes the ledger tick. On the host of Appendix~\ref{sec:universal-host}, a program that checks one
fact once and then commits it again before every instruction runs any two-counter program with the
flag down and pays one unit a step, plus one for the check: the checked fact is about a counter the
program never writes, so its edition never changes, and every COMMIT finds it.

\begin{observation}[A clock is a Thiele machine and isn't certified Thiele-complete]\label{thm:u34-clock}
Take the two-counter machine, let every move cost 1, and take as the record any yes/no
question about its configuration whatsoever; or take a flag that switches on, and stays on, the
first time any chosen test of the configuration says yes. Either way the machine obeys the toll
and has a two-counter machine running inside it, and it is not certified Thiele-complete. More
generally, no machine whose every move costs at least 1 is certified Thiele-complete.
\end{observation}

(It also follows from Observation~\ref{obs:tc-clock} and Theorem~\ref{thm:tc-decomposition}(2).) The proof is short, and it's short because clause (c) does exactly the job it was written for: it
says computing costs nothing, and a clock charges for computing. The other clauses do work
too, and that's worth checking, because a definition with a clause nothing needs is a
definition hiding something.

\begin{theorem}[Each clause does work]\label{thm:u34-tc-clauses}
(1) Take a two-counter base priced 0 with a ledger beside it, and one extra move TICK that costs 1
and raises a flag. It obeys the toll, and read with TICK as CERTIFY it meets the universal base
clause and the exact prices clause; it is not certified Thiele-complete, and more generally no machine with a
move that puts the record up from every state is. (2) Take a two-counter base priced 0 and a record
that is down at every state, so it never rises. It obeys the toll and meets clauses (a), (b) and (c); it is not
certified Thiele-complete, and more generally no machine whose record never rises is.
\end{theorem}

The paid latch, the TICK machine of (1), is the receipt for nothing from Section~\ref{sec:established}, wearing a universal base:
honest about the price, silent about what was bought. The definition throws it out.

The silent machine of (2) meets three clauses and fails only (d). The paid latch fails (d) as well as
(b), so it doesn't show (b) doing work alone, and nothing above says whether a clock meets (a),
(b) and (d). The sharper check, the usual evidence that a definition isn't written
around one machine, is a machine for each clause that meets the other three and fails it.

\begin{theorem}[The four clauses are independent]\label{thm:tc-independent}
For each clause of Definition~\ref{def:u34-thiele-complete} there is a machine with an interface
that meets the other three clauses and fails that one. Each of the four machines obeys the toll and
has a two-counter machine running inside it, and none of them is certified Thiele-complete through
any interface at all. (a) fails for this chapter's machine with one more move, DROP, priced 0 and read
as a base move, which lowers the flag and trips the trap latch. (b) fails for this chapter's
machine with one more move, FORGE, costing one and read as a CERTIFY, which raises the flag on any
state whose trap is down, with no commitment behind it. (c) fails for this chapter's machine with
every move costing one. (d) fails for the silent machine of Theorem~\ref{thm:u34-tc-clauses}(2).
\end{theorem}

Three of the four are this chapter's
machine with one thing changed, and in each case the one clause it fails is the only thing keeping
it out.
Clause (c) also fixes something it has no business fixing: the size of the unit. Double every price
on the small machine and it leaves the class, because its CHECK now costs two and (c) asks for
one. A definition that a change of currency breaks would be a strange one, so the next theorem states (c) with
the unit left open.

\begin{theorem}[Any unit]\label{thm:u34-tc-unit}
Call a machine \emph{certified Thiele-complete at unit $c$}, for a whole number $c \ge 1$, when some
interface meets clauses (a), (b) and (d) and this form of (c): base moves cost 0, each CHECK,
COMMIT and CERTIFY costs exactly $c$, and each move adds its cost to the ledger. (1) At unit 1
that is certified Thiele-complete. (2) Multiply every price of a certified Thiele-complete machine by $c \ge 1$ and the
machine is certified Thiele-complete at unit $c$; divide every price of a machine certified Thiele-complete at unit
$c$ by $c$ and the machine, with the same states, moves, step and record, is certified Thiele-complete,
with the ledger read as $\lfloor\ell/c\rfloor$.
(3) At every unit some move costs 0, so the clock of Observation~\ref{thm:u34-clock}, with any reading
or any latch as its record, is certified Thiele-complete at no unit, and neither is any machine whose every
move costs at least 1. (4) At unit $c$, a run from a clean state that ends with the record up
contains a CHECK, a COMMIT and a CERTIFY and raises the ledger by at least $3c$; the claim it
stands on held when it was checked and still held when it was committed; and the move that raises
the record is a CERTIFY. (5) The small machine with every price doubled is certified Thiele-complete at unit 2
and isn't certified Thiele-complete.
\end{theorem}

So the class doesn't care what a unit is worth. It cares that computing costs nothing and that the
three acts cost the same. The clock fails at every unit for the reason it failed at one: it
charges for computing. I keep unit 1 in the definition because the small machine is built that
way and every count in this chapter is in its units.

Now to the other direction, because a definition nobody meets is as useless as one everybody meets.

\begin{theorem}[The small machine, the generic one and the sorted one are certified Thiele-complete]\label{thm:u34-tc-instances}
(1) This chapter's machine is certified Thiele-complete, read with INC, DEC and HALT as base moves, CHECK $p$ $c$
and COMMIT $p$ $c$ as record moves on the claim ``$p$ holds of $c$,'' CERTIFY as CERTIFY, the checker as $\mathrm{eval}$ on the value of $c$, the window as (line, $A$, $B$), live states as those with the trap down, $\mathrm{start}(a, b)$ as
the loaded state for $(1, a, b)$, ``unchanged'' as ``same value and same version of $c$,'' clean
states as in Definition~\ref{def:u34-state} and the ledger as $\mu$. (2) So is the same machine over any property
language with an exact equality test and a checker proved to say yes exactly when the claim holds, as long as some
property is true of one number and false of another. (3) In particular the machine whose
property is ``this counter, read as a list, is sorted'' is certified Thiele-complete. (4) Every
certified Thiele-complete machine is a Thiele machine (Definition~\ref{def:certsys}) and obeys the
universal floor; for this chapter's machine, that system agrees step for step and cost for
cost with the one of Theorem~\ref{thm:u34-instance}, and for the generic machine with the
generic machine's own. (5) Each of the three is certified Thiele-complete over its own claim language,
fixed before the checker. For this chapter's machine the language's claims are pairs of a
property and a counter, equal when both parts are, the meaning of ``$p$ holds of $c$'' is the
meaning $\mathrm{holds}$ of Definition~\ref{def:u34-props} applied to $c$'s value, written
without the checker, and ``unchanged'' is ``same value and same version of $c$.'' For the generic machine the
language is built from the properties, their equality test and what they mean, and the checker
isn't part of it; the sorted machine is that case.
\end{theorem}

Appendix~\ref{sec:every-model} carries the same fixed-language form to other bases.
The claims keep their stated meanings; the new interfaces supply the checkers.

There's one thing every machine that meets the definition can't help doing. It has
Section~\ref{sec:windows}'s shape again, only this time it isn't about one machine but about all of
them.

\begin{observation}[Every certified Thiele-complete machine hides its record and its ledger from its window]\label{thm:u34-tc-window}
Let a machine be certified Thiele-complete through an interface, and let its window be the two-counter
configuration its universal base shows. (1) For every state $s$ there's a loaded state and one
base move priced 0 from it that lands on a state with $s$'s window, with the record down and the
ledger unchanged. (2) There's a loaded state and two runs from it that end with the same window,
one with the record up and one with it down; and likewise two runs from one loaded state that end
with the same window and ledgers at least 3 apart. (3) So no function of the window gives the
record at the end of every run from every loaded state, and no function of the window gives the
ledger. (4) The same holds of the Thiele machine built from the machine.
\end{observation}

The observer in that theorem sees one thing, the final window. An observer who sees the window
before and after each step, or the whole trace of windows, or who is also handed the program and
the start, is a different observer, and the collision says nothing about them; given the program
and the start, the run is fixed and so is its record. The theorem's two runs also start from the
same loaded state, so they're reachable witnesses, which is stronger than two arbitrary states
that happen to look alike. And ``no function of the window gives the record'' is all it says. It
doesn't say the window tells you nothing about it.

Notice how little that took, and say it out loud so nobody mistakes it for more. It's the fibre
fact from Section~\ref{sec:windows} plus one observation: a universal base can print any window
at no cost, and that is the point. The window here is everything a Turing machine would call the
machine, the line and the counters. On a certified Thiele-complete machine, that's never enough to tell you
what was established or what it cost, and it can't be, because computing costs nothing and can't touch
the record. A richer window, one that shows the flag, is a different window, and nothing above
says it fails.


\subsection{Parts I and II, run on the example machine}
\label{sec:part3-on-machine}

Parts I and II proved their counts and their collisions for any system with the named pieces,
and they left a promise that the machine of this part would be put through them, and this section is where it's kept. The eight-state
machine of Section~\ref{sec:a2-not-arbitrary} gets run by the small machine, price for price.
Structural entitlement, the count of what narrowing a list of rivals costs, is proved for every
certified Thiele-complete machine, with the small machine as a worked member, a larger search on the host
that has a counter for every number, and a run that shows the count's one awkward hypothesis, the
covering, can't be dropped. And the time a check saves gets an exact exchange rate, for the part
of the machine that has no record. (The no-bare-verifier result is in Section~\ref{sec:windows},
next to the collision it comes from.)

\textbf{The eight-state machine, run by the small machine.} Look at the small machine through a
window that shows two things: its line number read modulo four, which I'll call its \emph{slot} (the
eight-state machine's word from Section~\ref{sec:a2-not-arbitrary}, unrelated to the slot registers of the universal program $U$ in Appendix~\ref{sec:universal-host}),
and its flag. Call a state \emph{live} here, a narrower use of the word than in the definition of
certified Thiele-complete, when its trap latch is down, the fact ``$A \ge 0$'' about
$A$'s current version is in its table, and its channel holds a commitment, which is what CERTIFY
needs in order to run.

\begin{definition}[The eight-state machine, compiled]\label{def:frag-small}
Stamp is the single instruction CERTIFY. Next-slot is the single instruction INC $B$. Jump to
slot $a$, for $a$ one of $0, 1, 2, 3$, is the four instructions COMMIT ($A \ge 0$) $A$; COMMIT
($A \ge 0$) $A$; INC $B$; DEC $B$ $a$. The \emph{setup} is the two instructions CHECK
($A \ge 0$) $A$; COMMIT ($A \ge 0$) $A$.
\end{definition}

The two COMMITs in a jump change nothing but the line and the channel, and each costs one, which is
where the jump's price of two comes from. The INC makes $B$ positive, so the DEC takes one off and
jumps to line $a$ at no cost, and $B$ is back where it was.

\begin{theorem}[The small machine runs the eight-state machine]\label{thm:frag-small}
(1) For every live state $s$ and every move $m$ of the eight-state machine, running $m$'s
instructions from $s$ ends in a live state whose window is the eight-state machine's step of $s$'s
window by $m$, and the ledger has risen by exactly $m$'s charge: one for stamp, two for jump, zero
for next-slot. (2) The same holds for any list of moves: the window after the compiled list is the
eight-state machine's run of the list, the state is live, and the ledger has risen by the sum of
the charges. (3) So a run of compiled moves from a live state that takes the flag from down to up
has raised the ledger by at least one. (4) From any start $(a, b)$, the setup ends in a live state
with window (slot 3, flag down) and the ledger at 2; after it, any compiled list runs as on the
eight-state machine from (slot 3, flag down), with the ledger at 2 plus the sum of the charges; and
if the flag is up at the end, the whole run, setup included, contains the earned chain of
Theorem~\ref{thm:u34-cert-earned}.
\end{theorem}

There are three reservations, and I'll say them now. The setup's claim is true of every value, so its CHECK can't
fail; it's there because the small machine won't CERTIFY without a checked, committed claim, and
not as a test. And the jump's price of two is two paid COMMITs, a choice a program makes that the
machine allows and doesn't force. The window is a projection, too: the small machine's full state
keeps far more than eight states' worth.


\textbf{Structural entitlement, for every certified Thiele-complete machine.} Theorem~\ref{thm:ent-cs} was
about any Thiele machine, and what the run paid went in as a hypothesis. For a certified Thiele-complete machine (Definition~\ref{def:u34-thiele-complete}) what the run pays is known. Call the CHECK, COMMIT and
CERTIFY moves of a run its \emph{record moves}; each costs exactly one, so the ledger rises by the
number of record moves, and clause (b) says what a raised record must contain.

\begin{theorem}[Structural entitlement represented]\label{thm:ent-machine}
Let a machine be certified Thiele-complete through an interface. Take a clean start $s_0$ and a run $t$ from
it; a type of rivals and a type of observations with an exact equality test; an observation
$\mathrm{obs}$ of each rival; a decoder $d$ from runs to
observations; lists $\Omega$, $\Omega'$ of rivals and a decision tree $T$. Suppose
(M1) $\Omega'$ strictly narrows $\Omega$, as in Theorem~\ref{thm:ent-cs}, and (M2) some dropped
rival shows an observation different from every survivor's; (M3) the run ends with the record up
and the decoded run $d(t)$ is shown by some survivor; (M4) $\mathrm{depth}(T)$ is at most the number
of record moves in $t$; (M5) $|\Omega'| > 0$; and (M6) the covering:
$|\Omega| \le L(T) \cdot |\Omega'|$. Then (1) the membership test of $\Omega'$ is
strictly stronger than that of $\Omega$; (2) $t$ contains the earned chain of clause (b), with the
meaning of the claim true at the CHECK and at the COMMIT; (3)
$\lceil \log_2 |\Omega| \rceil \mathbin{\dot{-}} \lceil \log_2 |\Omega'| \rceil$ is at most the number of
record moves of $t$; (4) the number of record moves of $t$ is exactly the ledger's rise; and (5) the
ledger rose by at least 3.
\end{theorem}

It's a checklist, as I said in Section~\ref{sec:entitlement}, and the small print is the
same kind. (M3)'s second half feeds no conclusion: only the record's being up is
used. (M1) is used only for the inclusion; its gap feeds nothing, since (M2) already hands over a dropped rival. And (M5) is used but can go: (M2) puts a rival in the prior, and the covering then forces a survivor. What the machine adds to Theorem~\ref{thm:ent-cs} is conclusions
(2), (4) and (5): the payment is the machine's own three acts, each at exactly one, and a raised
record carries the earned chain.

There are two more, one where the tree comes from the run's own questions and one about what the record adds.

\begin{theorem}[A floor with no supplied tree]\label{thm:ent-questions}
Take a machine and an interface that meets the exact-prices clause (c) of
Definition~\ref{def:u34-thiele-complete}, so each record act costs exactly one. Call the \emph{questions} of a run the claims of its CHECK moves, and say two rival states
\emph{answer alike} when every one of those claims gets the same verdict at both (each verdict taken
at the rival itself, as it starts; the states the run passed through don't enter). If, for a list $\Omega$ of rival
states and a number $m > 0$, the rivals of $\Omega$ that answer alike with any one rival, itself
included, are never more than $m$, then
$\lceil \log_2 |\Omega| \rceil \mathbin{\dot{-}} \lceil \log_2 m \rceil$ is at most the number of CHECK
moves of the run, which is at most its number of record moves, which is the ledger's rise.
\end{theorem}

The tree has gone, and the run's own
questions stand in its place. There are limits, and they're the ones in the statement. The verdicts are taken at the
rivals as they start, and the states the run had reached when it asked don't come into it. A
branch that costs nothing, DEC on zero, can still narrow a list without adding anything to the ledger. So the theorem prices a narrowing that the paid checks
explain, and no more.

The next statement takes the verdicts where the run took them. A rival is a state the run could
have started from; a CHECK the run passed is asked of a rival by running the moves before it from
the rival and checking the same claim at the state that reaches. Nothing is supplied but the prior.

\begin{theorem}[Structural entitlement, with the narrowing made by the run]\label{thm:ent-run-machine}
Let a machine be certified Thiele-complete through an interface. Take a clean start $s_0$, a run $t$ from it
that ends with the record up, and a list $\Omega$ of rival starts that contains $s_0$. Let the
posterior $\Omega'$ be the rivals that pass every CHECK $t$ passed, each asked as above, in the
order $t$ made them, and let $b_i$ and $a_i$ be the numbers of rivals standing just before and just
after the $i$th. Then (1) $\Omega'$ sits inside $\Omega$ and contains $s_0$; (2) every rival of
$\Omega$ that isn't in $\Omega'$ fails, asked that way, a CHECK that $t$ made and passed; (3)
$|\Omega| \cdot \prod_i a_i = |\Omega'| \cdot \prod_i b_i$; (4) if every CHECK $t$ passed keeps at
least a $2^{-k}$ share of the rivals standing before it, $b_i \le 2^k a_i$, then
$\lceil \log_2 |\Omega| \rceil \mathbin{\dot{-}} \lceil \log_2 |\Omega'| \rceil$ is at most $k$ times
the number of CHECKs $t$ passed, that number is at most the record moves of $t$, the record moves
are the ledger's rise, and so the index bits are at most $k$ times the ledger's rise; (5) $t$
contains the earned chain and the ledger rose by at least 3; and (6) with no share assumed, the
index bits are at most $\lceil \log_2 (|\Omega| / |\Omega'|) \rceil$, which is
$\lceil \sum_i \log_2 (b_i / a_i) \rceil$ and at most the sum over the CHECKs $t$ passed of
$\lceil \log_2 (b_i / a_i) \rceil$.
\end{theorem}

Part (3) is the exact form of ``index bits are the sum of what each
test took'': take base-two logarithms and $\log_2 |\Omega| - \log_2 |\Omega'|$ is the sum of
$\log_2 (b_i / a_i)$. Parts (4) and (6) are what survives rounding to whole bits, and (6) rounds once, at the end, so it can be a bit lower for every test than rounding test by test. Rivals $(a, 0)$ for $a = 0, \ldots, 26$ with CHECKs $A \ge 9$, $A \ge 15$, $A \ge 19$ stand at 27, 18, 12 and 8: the index bits are $5 - 3 = 2$, rounding once gives 2, and rounding each test gives 3. Here the run does
the narrowing, every candidate dropped was dropped by a check the run paid for, and what stays a premise is the prior and the share, a bound fixed in advance; part (6) shows what the run pays without one.

\begin{proposition}[One CHECK can take four bits]\label{prop:ent-share-needed}
Take the sixteen starts $(a, 0)$ with $a = 0, \ldots, 15$, and the run CHECK ($A \ge 15$) $A$;
COMMIT ($A \ge 15$) $A$; CERTIFY from $(15, 0)$. It raises the record. Its one passing CHECK keeps
one rival of sixteen, so the posterior is $(15, 0)$ alone and the index bits are 4, while the run
has 3 record moves. So index bits at most the record moves, the bound of
Theorem~\ref{thm:ent-run-machine}(4) read with $k = 1$, fails here.
\end{proposition}

The share here is $2^{-4}$, and with $k = 4$
part (4) reads $4 \le 4$. It's the example of Section~\ref{sec:entitlement} that showed the covering
can't be dropped, seen from the run's side: one test can take many bits, and the bound of part (4) covers the
bits per test times the tests.

\begin{theorem}[What the record adds]\label{thm:ent-observed}
Let a machine be certified Thiele-complete through an interface. (1) From a clean start, the record is up
at the end of a run exactly when the run contains the earned chain, and a run with the record up
raised the ledger by at least 3. (2) Conclusions (1), (3) and (4) of
Theorem~\ref{thm:ent-machine} hold without a clean start and without the record being up, for a
machine that only meets the exact-prices clause (c): the hypotheses (M1), (M2), (M4), (M5) and (M6),
and the decoded run being shown by some survivor, are all it needs.
\end{theorem}

So the record doesn't make narrowing cost more. It adds the earned chain and
a floor of 3, and that's all it adds: a run meeting the same counting and payment hypotheses is still bound by that index-bit
floor even if it never raises the record, and the floor of 3 no longer applies to it.

\emph{Weights, fibres and partitions.} Three more pieces of the same count. First,
suppose the rivals aren't equally many: give each a whole-number weight, a multiplicity, and call
the sum of the weights of a list its \emph{mass}. The count survives with masses in place of
lengths. If the prior's mass is at most $L(T)$ times the posterior's, the rounded index bits of the
two masses differ by at most $\mathrm{depth}(T)$; so for any Thiele machine a run that paid at
least that many charged steps has a cost at least the drop, and for a certified Thiele-complete machine the
representation theorem goes through with masses and the same five conclusions. On such a machine
the ledger rises by the same number from every start, the number of record moves, so the weighted sum of
ledger rises is the prior's mass times the record moves, and the drop in rounded bits
times that mass is at most the weighted sum. Divide by the total mass and you'd have an average
rise; I leave the weights as whole numbers, cross-multiplied, so no probability comes into it.
Second, a way to meet the covering (M6) that you can see. Give each survivor a little crowd of
rivals, its \emph{fibre}, no bigger than the tree has leaves, and let the fibres' sizes add up to
at least the prior. Then $|\Omega|$ is at most the sum of the fibres, which is at most $L(T)$
times the number of survivors, and that's (M6). Third, a plainer way to get fibres. Pick a
\emph{representative} observation of each rival, and say the rivals are \emph{partitioned} by it
when every rival looks like some survivor under it and no observation class, counted inside the
prior, has more than $L(T)$ members. Take each survivor's class as its fibre: every rival is in
the class of some survivor, so the fibres add up to at least the prior, and (M6) holds. The
converse is false: take four rivals that all look alike, two survivors and a tree with two
leaves, and give each survivor two of the four as its fibre. That covers, but the one class has
four members, more than two.

\emph{A worked member on the small machine.} Eight rivals: the starts $(a, b)$ with $a \in
\{0, 1, 2, 3\}$ and $b \in \{0, 1\}$, counter $A$ at $a$ and counter $B$ at $b$. The run is CHECK
(Even) $A$; CHECK ($A \ge 2$) $A$; COMMIT ($A \ge 2$) $A$; CERTIFY, from $(2, 0)$. The program
raises the flag from a rival exactly when $a$ is even and at least 2, so exactly from $(2, 0)$ and
$(2, 1)$, and the posterior is those two. The decoder reads one ``passed'' for each CHECK in the
run's receipts; the tree is the complete tree of depth 2; and the covering comes from fibres
by a representative observation, ``$B$ is 0,'' which isn't constant: each survivor's fibre is the
four rivals that share its $B$.

\begin{proposition}[The eight-start member]\label{prop:ent-small}
For every $a$ and $b$, the run raises the flag from $(a, b)$ exactly when $a$ is even and at least
2. All hypotheses of Theorem~\ref{thm:ent-machine} hold, with eight rivals narrowed to two: the test
of the posterior is strictly stronger; the index bits are $3 - 1 = 2$; the run has 4 record moves
and the ledger rises by exactly 4; and the earned chain is in the run. The run's own questions
(Theorem~\ref{thm:ent-questions}) are its 2 CHECKs, so the floor is met with equality. With the two
CHECKs alone, the record stays down, the ledger rises by 2, and the bound $2 \le 2$ still holds.
\end{proposition}

It's a real narrowing, and not the cups of Section~\ref{sec:entitlement} sitting under it in a costume.
The posterior is exactly the set of
rivals the run certifies, so here a proof ties the lists to the run, and nobody supplies them.
Read with Theorem~\ref{thm:ent-run-machine}, the run's two passing CHECKs, replayed from the eight
rivals, leave four and then two, the same two, so the posterior is the one the run makes; each
keeps half, so $k = 1$, and the index bits, 2, are at most $1 \cdot 2$ and at most the ledger's
rise of 4.

\emph{A search through $n$ bits.} The same shape comes back at a size you can't enumerate, on the machine of
Definition~\ref{def:u34-multi}, the host in Appendix~\ref{sec:universal-host} with a counter
for every number, with the small machine's property language
of Zero, Even and At-least-$n$, read at registers,
in place
of $\mathsf{Slot}$. That machine is certified Thiele-complete too, by the argument of
Theorem~\ref{thm:u34-U-complete} with these three properties in place of $\mathsf{Slot}$. Hide $n = k + m$ bits, one in each of registers 0
to $n - 1$. The program asks $k$ questions, one per bit of the first $k$: CHECK ``at least 1'' of
register $i$ for a 1, CHECK ``is 0'' for a 0. Then it COMMITs the first claim and CERTIFYs. The
prior is every $n$-bit value; the posterior is the $2^m$ values that begin with the answers; the
tree is the complete tree of depth $k$; and a survivor's fibre is the $2^k$ values that agree with
it on the $m$ bits nobody asked about.

\begin{theorem}[A search through $n$ bits]\label{thm:ent-search}
For $1 \le k \le 16$ and any $m$: from a world holding a value $v$ of at least $k$ bits, the program
certifies exactly when $v$ begins with its answers, so the posterior is exactly the set of rivals it certifies from; the
test of the posterior is strictly stronger than the prior's; the earned chain is in the run; the
prior has $2^{k+m}$ members and the posterior $2^m$, so the index bits are exactly $k$; the run has
exactly $k + 2$ record moves and the ledger rises by exactly $k + 2$; and the run's own questions
are $k$ CHECKs, so Theorem~\ref{thm:ent-questions} is met with equality. With the COMMIT and the
CERTIFY left off, the questions alone narrow $2^{k+m}$ values to $2^m$, the record stays down, and
the ledger rises by exactly $k$. And no run from a clean start can execute more than 16 CHECKs before its
first successful CERTIFY: a failed CHECK would trap, and the table holds only
16 passing ones. Thus this program with 17 questions followed by its first
COMMIT and CERTIFY certifies from no start. A trace that certified earlier
can keep its flag up even after later failed instructions.
\end{theorem}

The sixteen belongs to the machine, and the
program has no say in it: it's the same ceiling as the item ``Sixteen checks per run, ever'' in the list of what the machine
doesn't do, near the end of this chapter, and it's why this search never asks more than sixteen questions however many bits
are hidden. Unasked bits can be any number. The lists, the
tree and the fibres are supplied here, as in the representation theorem, and tied to the run only
for this member. The covering still prices only the size of the drop, and its reason goes unpriced.

\emph{The covering can't be dropped.} It's a hypothesis, and you'd hope it's a nuisance you could remove. It isn't, and the small machine shows it. Fix $n \ge 4$, take the $2^n$ starts
$(a, 0)$ with $a$ below $2^n$ as the rivals, and run the three-move chain on the claim ``$A \ge 2^n - 1$'': CHECK,
COMMIT, CERTIFY. It certifies exactly from $A = 2^n - 1$, so the posterior is one start, $n$ index
bits are gone, and the ledger rose by 3. So for $n \ge 4$ the claim ``the index-bit drop is at most
the ledger's rise'' is false for this run, and no tree of depth at most 3 covers the narrowing:
(M6) would need $2^n \le 8 \cdot 1$. One property with a big constant in it carries $n$ bits in
a single yes-or-no question; what the bound prices is a narrowing whose classes are balanced, a
tree with enough leaves. A narrowing of any other shape is outside it.

\textbf{The time a check saves.} The question is whether counter $A$ is at least $m$. There are two
ways to find out. The \emph{unpaid} way uses base moves only, INC and DEC, which add nothing to
the ledger. But a base move reads $A$ one pebble at a time: a DEC takes a pebble off $A$ or finds
it empty. After $n$ moves, a run that started with $A$ at least $n$ has never seen $A$ empty, so
it can't have behaved differently for different such $A$ except by a constant shift in $A$
itself. The \emph{paid} way is one CHECK ($A \ge m$): one move, a ledger of one, and the line it
ends on says which way it went. Certifying the answer takes the whole chain, CHECK, COMMIT and
CERTIFY.

\begin{theorem}[The time tax]\label{thm:time-tax}
(1) For every program of two-counter instructions and every $n < m$, the line number and $B$
after $n$ moves (or where the run stopped, if it stopped sooner) are the same from $A = m - 1$ as
from $A = m$, with the same $B$. So no such program, read off its line number and $B$ after fewer
than $m$ moves, tells whether $A \ge m$ at every start; in particular none that stops within
fewer than $m$ moves does. (2) A program of $m$ DEC/INC pairs (a ladder: pair $i$ is DEC $A$ on
line $2i - 1$, jumping on success to the next pair at line $2i + 1$, and INC $B$ on line $2i$,
where a DEC that finds $A$ empty falls through) decides it when observed after $m$ moves: after $m$
moves its line is $2m + 1$ exactly when $A \ge m$; and it costs 0. (3) Executed as a
three-move trace from $\mathrm{start}(a,b)$, the paid chain raises the flag exactly when
$a \ge m$ and has a ledger of exactly 3, whatever $m$ is. As a stored program it pays 3 on success
and stops after its first failed CHECK with a ledger of 1. (4) From any loaded start $\mathrm{start}(a,b)$, a single CHECK ($A \ge m$) decides it in one move from the line number, with a ledger of exactly 1, whatever $m$ is. (5) So if a
unit on the ledger is worth $\lambda$ moves, the paid route to deciding costs $1 + \lambda$ and the
unpaid route costs at least $m$, and the paid route is cheaper exactly when $1 + \lambda < m$;
the paid route to certifying costs $3 + 3\lambda$.
\end{theorem}

When $3+3\lambda<m$, the paid chain both answers and certifies within a
combined budget smaller than the unpaid lower bound for answering alone.
That compares different guarantees; the unpaid route never certifies.
For a bare answer on both sides, the comparison is $1+\lambda<m$. The exchange rate is sharper than a comparison of two programs I
happened to write would give: the unpaid side is a floor under every program of INC and DEC there is, so
no choice of program comes into it, and the paid side doesn't move with $m$ at all. The bound has limits. It's
for programs that make no record move, INC and DEC with HALT or falling off the end as
stopping. It doesn't cover every certified Thiele-complete machine, because the definition doesn't forbid a
base move that changes a counter by more than one pebble, and such a machine is a different
machine, one the bound doesn't speak for. And the two routes put their answers in different
places, both bare decisions in the line number, with the paid chain also putting a historical record in the flag; this compares two ways to
answer one question, not two certifications. Zero cost isn't zero heat, either:
every move on either route is an event, and the moves are what $\lambda$ prices.


%===========================================================================
\section{From ``Five small models''}
%===========================================================================
The fragment is a three-instruction
memory-bounds language with the policy ``every address is below the limit'': the checker
accepts exactly when the condition holds, and a certificate implies the condition. One example
sits on top of that, a computation you can do in your head: a program that reads address two under a limit of two is rejected.
The fragment isn't the paper's real case study, and nothing connects it to anything else in this book. And it has a limit: a
program that writes garbage to address zero passes if zero
is below the limit. The proof is of the stated policy and nothing else.

And what was the paper's real case study, the one my fragment isn't? It was network packet filters, the
little programs an operating system runs on incoming network traffic to decide which program
gets to see which packet. They run inside the kernel, the most trusted part of the system, and
they have to be fast, so they're exactly the code you'd most like to take from a stranger and
least like to trust. In the 1996 paper that first set out the idea, Necula and Peter Lee wrote
several filters by hand in the assembly language of a DEC Alpha, shipped each with its proof,
and reported a one-time cost of 1 to 3 milliseconds to check the proof, after which the filter
ran with no checks at all. Those milliseconds put a number on ``checking is cheap by design''. What the paper reports changes
nothing about my fragment, which still isn't that case study, and it doesn't add a price: the
milliseconds measure the checker's work in that experiment. They are not
a per-crossing toll theorem, whether or not an implementation saves the result.

\begin{definition}[The bounds fragment]\label{def:u13-pcc}
Instructions are READ $a$, WRITE $a\ v$, and HALT. For a memory limit $L$, the policy allows
READ $a$ and WRITE $a\ v$ when $a < L$, and always allows HALT. The \emph{verification condition}
of a program is that the policy allows every instruction in it. A \emph{certificate} is a list of
per-instruction policy facts, one for each instruction in order. The \emph{checker} tests the
policy on each instruction and returns the conjunction.
\end{definition}

\begin{theorem}[The checker is the verification condition]\label{thm:u13-pcc}
For every limit $L$ and program $P$: the checker accepts $P$ exactly when the verification
condition holds; any certificate for $P$ gives the verification condition; and with $L = 2$
the checker rejects READ 2 followed by HALT.
\end{theorem}

\begin{proof}
Lay the instructions in a row. For each READ or WRITE, compare its address
with $L$; put yes beneath it exactly when the address is smaller. Put yes
beneath HALT without asking about an address. The checker accepts exactly
when every mark in that row is yes. The verification condition asks for
exactly those same per-instruction facts. An empty row has no failing
instruction, and extending the row preserves acceptance precisely when
the new instruction also passes. This covers programs of every finite length.

A certificate supplies one policy fact at each place in the same row.
Reading them in order gives every fact the verification condition asks
for, so it implies that condition too.

With limit two, the first instruction READ 2 asks whether $2<2$. It does
not. The later HALT is allowed, but cannot undo that failed comparison,
so the checker rejects the two-instruction program.
\end{proof}


\subsection{One example, three levels}

Before I lay all five side by side, let me walk one of them up and down the three levels from
the start of the section, because every claim in here lives at exactly one level and it's
easy to let a claim slide up a floor without noticing. I'll take gas, since it's the one place all
three line up. Think of a restaurant's real recipe, my photocopy of one page of it, and a
cookbook chapter on what any recipe has to contain. I can check my photocopy line by line.
That tells you nothing about what the kitchen cooks tonight.

Start with the real specification. Ethereum's execution specification says, in a document that runs real
money, what a storage write costs, what clearing it refunds, and how the refund is capped at
the end of the transaction. None of what I proved is at this level. What the network
actually charges is whatever the clients people run actually compute, and the evidence for
that is the clients and the chain. This book has none of it.

Next comes the transcription. I copied one mechanism out of that document, one slot in one transaction,
gas and refund arithmetic only, at a named commit
(Definition~\ref{def:u13-slot}, Theorem~\ref{thm:u13-sstore}). What the result tells you: \emph{if} the
storage rule is what I copied, then leaving a fresh slot set costs at least twenty thousand,
and set-then-clear nets twenty-three hundred at the counter. What it can't tell you: anything
about the refund cap (not copied, which is why what a real set-then-clear pays came out so different from the counter
earlier), the stack, calls, other slots, or whether any client implements the rule as written.
The boundary of the claim is the edge of the copy.

Last comes the wrapper. The gas schedule of Definition~\ref{def:u13-gas-schedule} isn't Ethereum. It's any
system with steps, a cost per step and a yes/no commitment flag, plus one thing \emph{stipulated}:
a charge predicate, with the premise that a charged step costs at least one and an uncharged
one costs zero. That premise is the whole of what the wrapper adds, and it's added by
hand; no specification supplies it. Given it, the theorem says every run costs exactly its number of commits, never letting a commit go unbilled, never billing one twice and never billing a non-commit, exactly when the charge predicate is the commit predicate at unit price. What that tells you is what exact pricing would have to
look like, in any system. It doesn't say Ethereum meets it, and Ethereum doesn't: real gas
bills computation that commits nothing, which the transcription shows in numbers. So the
wrapper prices a definition, the transcription prices a copied rule, and neither prices the
chain.

The other four sit on the same three levels, and only one more wrapper is spelled out in this
section. Casper's transcription is the accountable-safety proof above; its wrapper is a small
gadget with the stake at risk handed in. A finalize step that risks nothing turns the record
on at no charge, which breaks Section~\ref{sec:toll}'s toll, and one that risks at
least one unit has a floor. It records exposure to loss (no collateral is destroyed in it), has no Casper
checkpoints in it at all, and is priced by stipulation, the same way the gas wrapper is.
Attestation, transparency and proof-carrying code have transcriptions that prove what their verifiers do and no price theorem, because no specification of theirs puts a price on the verdict. Their
wrappers, which I don't spell out here, are Section~\ref{sec:windows}'s verifier argument told in each system's vocabulary,
with a colliding pair behind each one, and that argument is all they carry here.

The book's rule is that every claim gets a way to break it, so I'll say how to knock down the four claims with the most weight on them.
For the gas transcription: a list of stores from a fresh slot that ends nonzero with net charge under 22100. For the gas wrapper: a schedule in which every run costs exactly its number of commits, whose charge predicate still differs from the commit predicate at some step. For the finality gadget: a finalizing trace from unfinalized to finalized whose
charged stake is zero while the gadget claims a floor. For the attestation fragment: a
verifier reading only the quote's fields that decides a runtime flag the quote ignores.


%===========================================================================
\section{From ``The pointer thesis''}
%===========================================================================
\subsection{Where the shape comes from, and how much of it survived}

The shape is from physics, and I want to tell you where exactly, because I borrowed a word
and I didn't get to keep the whole thing.

Quantum mechanics has a famous problem: a system can be in a blend of states, and when you
measure it you find one definite state, and nobody agrees on what ``measuring'' did. In 1981
Wojciech Zurek asked a sharper question: of all the ways a measuring device \emph{could} end up,
why does it end up recording one particular kind of quantity (the position of a pointer on a
dial, say)? Why never some blend? His answer, developed over twenty years and summarised in a
2009 paper called ``Quantum Darwinism,'' is that the \emph{environment} decides. The air, the light,
the table, everything the device touches, is constantly interacting with it, and only
certain states of the device survive that interaction; those are the \emph{pointer states}. More
than that, the environment selects them, and then it \emph{copies} them. The light bouncing off a
dial carries the dial's reading to every corner of the room, redundantly, so that a thousand
observers can each read a different scrap of the environment and all agree on where the
pointer is, without touching the dial. Objectivity, on this view, \emph{is} redundancy: a
property is objective when many independent fragments of the world each carry it.

Zurek isn't the only one behind it. Fernando Brand\~ao, Marco Piani and
Pawe\l{} Horodecki (2015) prove an approximate statement about most small
fragments of a large output: their channels can be approximated by measuring
the input in one common way and preparing fragment states from that result.
The error and the exceptional fraction are part of the theorem. Agreement on
outcomes also needs the fragments to carry enough information about them.

Xiao-Liang Qi and Daniel Ranard (2021) strengthen the bound on the exceptional
region. Outside a region whose size is controlled by the input size, allowed
error and size of the fragments being examined, locally accessible information
is approximately classical. This does not say that every observer gets an
exact record, or any useful information at all. My definitions below ask for
exact Boolean answers from specified observers. That is a simpler chosen
contract, not the quantum theorem with its qualifications removed.

What I borrowed is that shape. Certification is to a machine's ecosystem (the machine plus everything around it that
reads, stores and relies on its records) what the pointer
reading is to the room: the thing that gets copied everywhere. It's a good shape and I think
it's the right intuition. I'm going to be exact about what survived into the definitions below,
because the word ``redundant'' is doing less than it does for Zurek.

Zurek's redundancy is \emph{graded} (each fragment carries almost all the information, with some
small loss) and it requires the fragments to be \emph{disjoint}: different observers reading
different parts of the environment. The definition below is the zero-loss case with \emph{no
disjointness at all}. An ecosystem is a state type, some number of observers, and one bit
per observer; an observer \emph{records} an event when its bit is true exactly when the event
holds; an event \emph{proliferates} when every observer records it. Nothing in that says two
observers draw on separate parts of the state. They could all be reading the same bit, and the definition permits it. So I kept the word and dropped the independence that makes Zurek's version
mean ``objective.'' That leaves a gap between the thesis's prose and that first formal object, and the
prose is where the thesis lives. The definitions under the thesis below close part of the gap:
they define independent carriers, and the definition turns out to force something
(Theorem~\ref{thm:u14-independent-constant}).

The independence I dropped is exactly what the thesis's closest published neighbour keeps.
Jess Riedel, Zurek and Michael Zwolak (2016) took histories of a closed quantum system, the
story a whole sequence of events tells, and asked which of them the world keeps records of in
many separate pieces of the environment. As I read their paper, redundant records in separate pieces make a history
consistent, so ordinary yes-or-no reasoning about it works, and they pick out a small set of
histories from the huge number that are merely consistent: in their words, an objective past.
Their result is ``what proliferates is what counts'' in quantum form, with the disjointness in it, and it
came before this book. Halliwell, in a second 1999 paper, ``Somewhere in the universe,'' had already carried the link between records and
consistent histories from pure starting states to mixed ones, where the records still exist but
are necessarily imperfect.

There's a deeper mismatch in the borrowing, and it runs against the book, so it goes here. In
Zurek's picture the whole thing, system and device and environment together, evolves by quantum
mechanics' rule for closed systems, and that rule never sends two states to one; every step can
be run backwards. So, seen from the whole, the environment's copying is just that, copying: the
dial's reading gets written onto fresh light, and no two states anywhere are merged. A simple copying model sends each new copy to a previously unused carrier.
That is a use of fresh space, like Bennett's blank page. It is not a proof
that every closed quantum evolution visits a new state forever; finite closed
systems can recur.

Looking only at the dial discards the correlations with the light. In an ideal
complete-dephasing channel, distinct input phases can give exactly the same
reduced output. In gradual decoherence they may only become harder to tell
apart; they need not become exactly equal at any finite time. An exact merge
in a reduced description, approximate loss of accessible coherence, and a
heat transfer are three different claims. To obtain a heat floor one still
needs the physical preparation, the relevant entropy change and the thermal
premises. Copying alone supplies none of those conclusions.

That is the difficulty for the borrowing. The thesis points at copies being
kept and used. The quantum account explains particular ways information can
spread and become accessible. It does not attach this book's toll to that act.

A published no-go cuts closer still. Philippe Allard Gu\'erin, Veronika Baumann, Flavio Del
Santo and \v{C}aslav Brukner (2021) gave an argument that, under assumptions that look natural in quantum
mechanics, an observer's own perceptions of her outcome at two different times can't share one
reality: no joint chance for the two, depending linearly on the starting state, reproduces what
unitary quantum theory predicts for each. A record that persists, for the one holding it, is
exactly what gets into trouble once the whole is treated as one-to-one. It isn't a theorem
about my machine, which isn't quantum. It's still the strongest published reason I know to hold
the word ``permanent'' loosely when a physicist is in the room.

Which leaves the question to ask of any borrowed shape: has anybody checked the original? For a
long time quantum Darwinism was a theory without a lab. Then, in 2018 and 2019, three groups
working separately, in Italy, China and Germany, went looking for its fingerprint, as Philip Ball's 2019 report in \emph{Quanta Magazine} tells it. The German
one, at Ulm, with Zurek among the authors, watched the spin of a single defect inside a diamond
and the carbon atoms around it, and saw the redundancy the theory predicts: the defect's state
turned up copied in its surroundings, and it didn't take much of the surroundings before more of
them told you nothing new. The other two did it with photons. Is that the bottom? It isn't, not by a long
way, and the people involved said so. Zurek himself pointed out that the photon experiments make
their copies artificially and that the diamond's environment is small, and others noted that all
three are schematic versions of a real environment and don't rule out rival accounts. So the
shape I borrowed has passed some first tests in a lab, at the scale of a handful of atoms or
photons, with the caveats written on it. None of that touches the thesis here, which is
about machines and never rested on quantum physics. It changes the picture: what I borrowed is a
lead somebody is checking. It isn't settled, and I don't get to lean on it as if it were.


\subsection{What's proved}

This section has definitions, two toys, a theorem about independence, a survey, and two refutations of things that would have been stronger.

The definitions are the ones above, plus ``unique pointer among rivals'': the event proliferates and none
of a listed set of rival events does. Uniqueness is \emph{relative to the list}, necessarily,
since any event that holds on exactly the same states as a proliferating event proliferates
whenever it does: the same fact under another name, so someone has to say which names count
as rivals. And the first criterion is stated
as a schema: for every ecosystem in some class, its metered event is the unique pointer among
its rivals. The class and the rivals are inputs. Nothing about them is proved. A second set of
definitions, after this one, tightens ``proliferates'' to ``copied and relied on'' and drops the
rival list.

The formal object is small enough to hold in one hand.

\begin{definition}[Ecosystem, records, proliferates, pointer]\label{def:u14-ecosystem}
An \emph{ecosystem} is a type of states, a number $N$ of observers, and for each observer $i$ a
Boolean fragment $F_i$ on states. Observer $i$ \emph{records} an event $E$ (a property of states)
when, for every state $s$, $E(s)$ holds exactly when $F_i(s)$ is true. $E$ \emph{proliferates} when
every observer $i < N$ records it. $E$ is the \emph{unique pointer among} a list of rival events
when $E$ proliferates and no rival in the list does. Given a class $C$ of ecosystems, and for
each member a designated metered event and a list of rivals, the \emph{pointer criterion holds}
for $C$ when every member's metered event is the unique pointer among its rivals.
\end{definition}

Nothing in the definition makes two fragments read different parts of the state, and
``metered'' is only a label: the definition has no metering in it.

One bit per observer can't even state the objection a replicated system raises. Every full
node of a blockchain holds the whole shared state, so every fact about it, ``gas used is at least
one'' as much as ``final,'' sits on every node. The tightened definitions give each observer
whatever it holds, an action it takes on what it holds, and a flag for whether it produced the
record.

\begin{definition}[Copied and relied on, independent carriers, the pointer]\label{def:u14-reliance}
An \emph{ecosystem with reliance} is a type of states, a number $N$ of observers, and for each
observer $i$ a view $V_i$ (a function from states to some type of things an observer can hold), an
action $a_i$ (a Boolean function of the view) and a producer flag. Fix a set $R$ of reachable
states. An event $E$ is \emph{copied} on $R$ when for each $i < N$ some Boolean function $d$ of the
view has $E(s)$ exactly when $d(V_i(s))$ is true, for every $s$ in $R$. Observer $i$ \emph{relies on}
$E$ when $a_i(V_i(s))$ true gives $E(s)$ for every $s$ in $R$, and $a_i(V_i(s))$ is true for some $s$ in
$R$. $E$ is \emph{relied on by parties who didn't produce it} when some observer $i < N$ is not a
producer and every such observer relies on $E$; it is \emph{copied and relied on} when it is both.
$E$ is the \emph{pointer} on $R$ when it is copied and relied on, and every event that is copied and
relied on holds at every state of $R$ where $E$ holds. The carriers are \emph{independent} under a family of relations $P_i$ (states that agree on
observer $i$'s own part of the world) when each view gives equal values on states related by its
$P_i$, and for any two observers $i \ne j$ below $N$ and any states $s$, $t$ there is a state $u$
with $P_i(u, s)$ and $P_j(u, t)$: the two parts can be set separately.
\end{definition}

Two pointers on the same $R$ agree on every state of $R$, so the pointer is unique with no list
of rivals. And the definition of independence has a consequence I
didn't put in on purpose.

Put one yes/no card in front of each of two readers. If the cards can be set
separately, all four arrangements are allowed: no/no, no/yes, yes/no and yes/yes.
Ask each reader to recover one shared event using only their own card. Suppose
that event is false at no/no and true at yes/yes. At no/yes, the first reader
sees exactly what they saw in the false case, while the second sees exactly
what they saw in the true case. Their answers conflict. That mixed arrangement
is exactly what the independence condition hands the proof.

Now impose a copying procedure that only leaves no/no or yes/yes. On those
reachable arrangements each reader can report the shared bit correctly.
Separate cards are things that \emph{could} carry a copy; what makes them copies is that
the arrangements you actually reach have the two agreeing. The theorem below is about
recovering the event exactly on \emph{all} arrangements, the mixed ones included, which
the copying never leaves behind.

\begin{theorem}[Independent copies on every state: the event is equivalent at any two states]\label{thm:u14-independent-constant}
In any ecosystem with reliance whose carriers are independent and which has at least two observers,
an event copied on the set of all states holds at one state exactly when it holds at any other,
so it holds at every state or at none.
\end{theorem}

\begin{proof}
Let observers 0 and 1 decide the event by $d_0$ and $d_1$. Given states $s$ and $t$, take $u$ with
observer 0's part from $s$ and observer 1's from $t$. Then the event at $s$ is $d_0$ of observer
0's view, which is the same at $u$, so it equals the event at $u$; and the event at $u$ equals the
event at $t$ the same way through observer 1.
\end{proof}

So
separate parts of the world can each carry an exact copy of a real record only on the states the
world actually reaches, where the copying made them agree, and that's why the definition asks for
copies on a set $R$ of reachable states. Zurek's copies are the same: the light carries the
dial's reading because the dial put it there.

\begin{theorem}[A replicated header: copying doesn't separate, reliance does]\label{thm:u14-ledger}
Let the states be a block header (a finalized bit, a gas-used count) and three nodes' copies of it,
with $R$ the states where every copy equals the header; node $i$'s view is its own copy, its part of
the world is that copy, node 0 is the producer, and every node acts exactly when its copy says
finalized. Then the carriers are independent; ``finalized'' and ``gas used is at least one'' are
both copied on $R$; ``finalized'' is the pointer on $R$; and ``gas used is at least one'' is not
relied on by parties who didn't produce it.
\end{theorem}

\begin{proof}
Independence: to give node $i$ the copy from $s$ and node $j$ the copy from $t$, overwrite node
$j$'s copy in $s$. On $R$ every copy is the header, so each node reads either fact off its copy.
Node 1 acts only where its copy, and so the header, says finalized, and acts at the finalized
header with no gas used; that state shows gas used fails node 1's reliance. Any event relied on by
node 1 holds wherever node 1 acts, which on $R$ is wherever the header is finalized, so
``finalized'' is the strongest.
\end{proof}

It's a sanity check with maps I wrote, and it shows only that the tightened definition can tell the two apart where plain copying can't. Which of the two it picks, the actions pick: let every node act when its copy says gas was used, and gas used is the pointer. When what the other parties act on is copied, that is the pointer, so here the argument lives in the actions, and I wrote those too.


The toy has three observers on a replicated ledger, where certification proliferates and ``work
done'' doesn't. It's an illustration, and that's all it is; it isn't evidence.

\begin{theorem}[The toy]\label{thm:u14-toy}
Let the toy have states (a certified bit, a work count), three observers, and every fragment
equal to the certified bit. Then ``certified'' is the unique pointer among the one rival
``work at least 1''.
\end{theorem}

\begin{proof}
Every fragment is the certified bit, so every observer records ``certified'' by definition.
For the rival, take the state with the bit false and work 1: the event holds and observer 0's
fragment is false, so observer 0 does not record it.
\end{proof}

The survey covers twelve events in one theorem: for each of the five systems in Section~\ref{sec:models}, the
designated event (finalized, committed, attested, logged, checked) proliferates in the
observer map I wrote for that system, and a rival (proposer work, a scratch counter,
measurement noise, submission latency, prover steps) doesn't. The five together are one
statement, an instance of the schema above.
The survey adds a symmetric message-authentication code, which does \emph{not} proliferate, and a digital
signature, which does, for twelve rows in all. The five maps all use the same construction: a state has a designated
Boolean event and a separate nonnegative rival count, three observers each return
the event bit, and the rival asks whether its count is at least one. The bit and
count can vary independently. The signature map likewise gives all three observers
the authentic-signature bit. In the MAC map, observers 0 and 1 receive the authentic-tag
bit and observer 2 always returns false. These maps, not a protocol running
on a network, are the entire tests here. And there's a thirteenth, which I think is the most important one: swap which bit
every observer copies, and the criterion picks the \emph{other} event. The criterion selects
whatever the observer map copies. The map is where the argument lives, and I wrote the maps.
So the five ``confirmations'' are five instances of one construction with the fields renamed.
I'm saying that twice on
purpose, so it can't slip past once.

\begin{theorem}[The survey, and the swap]\label{thm:u14-survey}
(1) In each of the five observer models (proof-of-stake, gas, attestation, transparency,
proof-carrying), the designated event proliferates and the designated rival does not; in the
symmetric-MAC model the authentic-tag event does not proliferate; in the signature model the
authentic-signature event does. (2) Let the states be pairs of bits (first, second), with
three observers whose fragments all equal the second bit. Then ``second is true'' is the
unique pointer among the one rival ``first is true''.
\end{theorem}

\begin{proof}
Each confirmation holds by definition: each observer's fragment is true exactly when the
event holds.
Each refutation exhibits one state where the event holds and some observer's bit is down.
In (2) that state is (true, false).
\end{proof}

Two stronger claims are both refuted, and I'd ask you not to run them together, because they fail for
different reasons.

The first stronger claim is that \emph{every forgery-resistant commitment proliferates}. If that were
true, certification would be the pointer because it's the commitment. It's refuted in three
observer models (deniable authentication, where only the designated verifier's bit tracks the
event; a symmetric MAC, where the key holders share a secret and nobody else can check; an
object capability, an unforgeable token that grants access to whoever holds it) and
confirmed in two (a public log; a digital signature). Build the first
with three observers: give one a copy of the event's bit and the other two a constant false.
The one records it; it doesn't proliferate. The proved half is the map computation. Whether
my maps are \emph{fair} to the real designs is argued in prose, and you should argue back.

\begin{theorem}[One blind observer is enough]\label{thm:u14-blind}
(1) In any ecosystem, if some observer $i < N$ has $F_i(s_0)$ false at a state $s_0$ where
$E$ holds, then $E$ does not proliferate. (2) Let the deniable model have states (authentic,
verifier-accepts), three observers, observer 0's fragment the authentic bit and observers 1
and 2 constantly false. Then observer 0 records ``authentic'', and ``authentic'' does not
proliferate. The symmetric-MAC and object-capability models also fail to proliferate their
events; the public-log and signature models proliferate theirs.
\end{theorem}

\begin{proof}
(1) Recording would require $E(s_0) \iff F_i(s_0)$, and the left side is true while the right
is false. (2) Observer 0's fragment is true exactly when the event holds. Apply (1) with observer 1 and
the state (true, true). The other two refutations apply (1) to a blind observer of their own model (the symmetric-MAC map gives observer 2 a constant false; the object-capability map gives the token holder the bit and the other two constant false). The two confirmations hold by definition: the public-log map gives all three observers the logged bit, and the signature map gives all three the authentic-signature bit.
\end{proof}

The theorem is about the observer maps, and only them. That a third
party in deniable authentication really cannot be convinced is the modelling claim, and the
model has no transcript, key or verifier in it to check that against.

The second stronger claim is that \emph{agreement plus accuracy plus locality imply permanence}. That is,
if all observers agree, each sees the event exactly, and each updates by the same local rule
with no coordinator, then the recorded event can't be undone. The smallest thing
imaginable refutes it: one bit, two observers who both see it, and the rule ``flip.'' All three conditions
hold and true becomes false in one step. Do
it on paper; it takes a minute and it's humbling. What \emph{does} give permanence is durability,
``an observer's true view stays true,'' added as a premise, and then it follows in a few lines. So durability is the premise doing the work, and
durability is Section~\ref{sec:a2-not-arbitrary}'s permanence wearing a different coat. Agreement doesn't buy it. Do the
paper version again with one change, the rule ``stay true once true'' in place of ``flip,'' and
the bit can't fall, because the rule won't let it. The observers agreeing has nothing to do with it. (The
small machine of Section~\ref{sec:smallest} has a reading that's durable in exactly this sense: its flag never comes
down, by its step rule (Theorem~\ref{thm:u34-toll}).) None of this rules out a narrower theorem with real assumptions about the network,
the adversary or durability; it rules out these two strong forms.

\begin{definition}[Observation game]\label{def:u14-game}
A \emph{game} has states, a number $N$ of observers, a Boolean event $e$, a Boolean view $V_i$ for
each observer, and a step map. For $i, j < N$ it has \emph{consensus} when $V_i = V_j$ on every
state; \emph{authenticity} when $V_i(s) = e(s)$ for every state and every $i < N$;
\emph{coordinator-free update} when one function $\ell$ on Booleans has
$V_i(\mathrm{step}(s)) = \ell(V_i(s))$ for every state and every $i < N$; \emph{durable views} when
$V_i(s)$ true gives $V_i(\mathrm{step}(s))$ true for every $i < N$; and a \emph{permanent event}
when $e(s)$ true gives $e(\mathrm{step}(s))$ true.
\end{definition}

\begin{theorem}[What buys permanence]\label{thm:u14-toggle}
(1) It is false that every game with $N > 0$, consensus, authenticity and coordinator-free
update has a permanent event. (2) Every game with $N > 0$, authenticity and durable views has
a permanent event.
\end{theorem}

\begin{proof}
(1) The toggle game: states are Booleans, two observers, event and both views the identity,
step ``flip''. Consensus and authenticity hold trivially, the local rule is ``flip'', and the
state true steps to false. (2) If $e(s)$ is true, observer 0's view is true by authenticity,
stays true by durability, and authenticity turns it back into $e(\mathrm{step}(s))$.
\end{proof}


\subsection{What isn't proved, and must stay that way on the page}

What isn't proved is that certification \emph{in deployed systems} is singled out by being copied and relied on across
\emph{independent} carriers. The criterion defines independence; it has no precise definition of ``deployed'' or ``faithfully model,''
and no cost law connecting being copied and relied on to price. Every correspondence between a mathematical observer
map and a real system is a sentence I wrote. The thesis is a thesis, and anyone who
says the formal object proves it is reading more than is there.

\subsection{What the thesis asks of the examples}

There are three concrete comparisons to make. First, look at the record a deployed
discipline prices. Do parties who did not produce it copy and rely on it? A priced record
that nobody else relies on would contradict the thesis's second half.

Second, compare certification with rival events. If an event not implied by certification
is copied and relied on to the same standard, certification is not the unique pointer.
Multiplication counts or head-turn counts would qualify if they met that same test in the
system being described.

Third, apply the independence condition of Definition~\ref{def:u14-reliance} to a faithful
model of the ecosystem and ask whether certification itself meets the test. The symmetric-MAC
example shows why the carriers and their access matter. The two stronger forms already
refuted above help locate the distinction this thesis needs.

These comparisons supply the evidence the worldly claim asks for. The mathematical
criterion supplies a definite question once the model and its observer maps are chosen.


%===========================================================================
\section{From ``The counter against physics''}
%===========================================================================
\subsection{The calibration is a unit conversion}

Multiply the counter by an energy per unit of cost that you choose,
and you get a theorem: a state with one unit, at the scale kT ln 2 per unit, carries kT ln 2 of
energy, for every k and T. Read that for exactly what it is: a unit conversion, true by definition,
with the scale as an unforced input. Nothing in it picks the scale. The theorem would be just as
true with the scale set to a dollar a unit. What \emph{would} connect units of cost to joules is the premise in Section~\ref{sec:a2-not-arbitrary},
Landauer's principle, which is a physical premise. And even
as a scale, $k_B T \ln 2$ a unit is the floor for one particular squeeze, a halving.
Corollary~\ref{cor:heat-floor}'s floor for a flip that brings $k$ states in beside $m$ already
reading yes, starting
from the even spread on those $m+k$ states, is the larger of $k_B T \ln \frac{m+k}{m}$ and $\frac{2k}{m+k}\,k_B T\ln 2$, and for one newcomer among many yes-states that's far below a halving, about $1.4\,k_B T / m$. Price every unit at $k_B T \ln 2$ and you can overstate the floor
of a single narrow flip by a factor that grows with $m$.

Remember the constant from Boltzmann's grave in Section~\ref{sec:a2-not-arbitrary}, the exchange rate between a count and
heat? It has a newer chapter, and it makes the point sharper than I can. The $k_B$ in
$k_B T \ln 2$ turns a temperature into an energy, and for most of its
history it was a measured number, with error bars, pinned down by careful experiments. Since
20 May 2019 it isn't. That day the international system of units was redefined so that
Boltzmann's constant is exactly $1.380649 \times 10^{-23}$ joules per kelvin, by definition, and
the kelvin is whatever makes that true. So the conversion from temperature to energy has been a
convention since then, in the strictest sense: its numerical value in those units is fixed. Realising the kelvin in a laboratory still
requires measurements, with uncertainty. What's still
physics is everything around it, whether wiping a bit at a temperature really sheds that much
heat, which is what the Lyon bead was measuring. The calibration theorem is the same kind of move
as that redefinition. It fixes a conversion, and the physics, if there is any, is somewhere else.

\begin{definition}[Energy at a scale]\label{def:u24-scale}
For a real number $c$ (joules per unit) and a state $s$, the energy of $s$ at scale $c$ is
$\mu(s) \cdot c$. The \emph{Landauer calibration} at $k_B$ and $T$ is the statement that every
state with $\mu(s) = 1$ has energy $k_B T \ln 2$ at scale $k_B T \ln 2$.
\end{definition}

\begin{theorem}[The calibration is a unit conversion]\label{thm:u24-calibration}
(1) For all real $k_B$ and $T$, the Landauer calibration at $k_B$ and $T$ holds. (2) For any counter, there are two real-valued functions of
the state that give 1 and 2 on every state with
one unit, so the counter alone fixes no energy value.
\end{theorem}

\begin{proof}
(1) The scale says what one ledger unit is worth. At a ledger reading of
one, multiplying by that scale returns the scale itself:
$1\cdot k_B T\ln2=k_B T\ln2$. No condition on the signs of $k_B$ or $T$
is needed for that multiplication.

(2) Choose a scale of one and a scale of two. The same ledger reading of
one is assigned energy one by the first choice and energy two by the
second. The ledger has not changed; only the conversion I supplied has.
Thus that conversion was not determined by the counter.
\end{proof}

The physics is the premise of
Section~\ref{sec:a2-not-arbitrary}, restated here once in the form the next theorems test. In words first: a \emph{dissipation}
hands each instruction a whole number of heat units, each unit worth kT ln 2. The floor premise
says an instruction that takes at least one state where some yes-or-no
property is false to a state where it's true, never sends a true one back to false, and sends
two different states to one, dissipates at least one unit. That's a different statement from Landauer's principle as
Premise~\ref{prem:landauer} has it. That one depends on the spread over states,
$Q \ge k_B T \ln 2 \cdot \Delta H$, and gives no floor for a known state (there the distinction was
absent from the chosen point distribution; that alone says nothing
about a prior preparation bill). This one gives every such
instruction a floor of one whole unit, whatever the spread. To get the second from the first
you'd need more, at the least a spread with a full bit of uncertainty removed at every such
step, and the known-state theorem and the narrow squeeze of Section~\ref{sec:a2-not-arbitrary}
show that isn't automatic. So I give it its own name, \emph{structural dissipation pricing}. The calibration says no instruction dissipates more units than the counter charges
it. One word about that word: ``dissipation'' here means heat handed to the surroundings, the
quantity Section~\ref{sec:a2-not-arbitrary} read off Sagawa as owed even when no entropy is
produced, the heat that pays for where the forgotten bit goes. It is not what Section~\ref{sec:what-a-chair-is}
calls dissipating, which is entropy produced, and it has nothing to do with waste, or with doing
the job badly. There are two premises here, nothing is proved about either,
and the only question is whether they can hold together:

\begin{premise}[Structural dissipation pricing]\label{prem:u24-landauer-pair}
A \emph{dissipation} is a function $\delta$ from instructions to naturals, counted in units of
$k_B T \ln 2$. The \emph{floor}, a hypothesis separate from Premise~\ref{prem:landauer} of
Section~\ref{sec:a2-not-arbitrary}:
every instruction that collapses a yes-or-no property has $\delta(i) \ge 1$, where $i$
\emph{collapses} $P$ when some state with $P$ false is sent to a state with $P$ true, every
state with $P$ true is sent to one with $P$ true, and some two different states are sent to the
same state. The \emph{calibration}: $\delta(i) \le
\mathrm{cost}(i)$ for every instruction. Neither is proved anywhere; the theorems below ask
whether they can hold together.
\end{premise}

The last clause of the floor is there because the first two don't make a step forget anything
on a machine with infinitely many states. INC $A$ sends a state with $A = 0$ to one with $A = 1$
and keeps every positive $A$ positive, so it raises ``$A \ge 1$'' one way, and it is
one-to-one. Any one-to-one map that misses some state does the same to the property ``is in my
image''. Its input can be read back off its output, so it loses no entropy, and Landauer's premise as
this book takes it then puts no heat floor under it at all. It pays in fresh states instead: INC moves
every value of $A$ up one, Hilbert's front desk, and only a counter with no top can keep doing
it. A thermal implementation can dissipate under the model-specific conditions above; the premise only sets no floor under it. A floor read without the last clause would charge a whole unit. On a finite machine the clause asks
nothing extra. A step that sends one false state into the true ones and keeps every true state true
puts one more state into the true set than the set holds, so two of them land on the same state. And on any machine the property comes along for the ride: an instruction that sends two different states $a$ and $b$ to one collapses the property ``is reached from $a$ in one or more steps'', since one of $a$ and $b$ lies outside it and lands in it. So the floor is a price on merging, one unit for every instruction that merges, read as heat.


\subsection{The calorimeter protocol, as a specification}

What follows is how one \emph{would} measure this, and it's here because it makes the boundary exact. Take a two-state
register with energies zero and a gap, a population p in the excited state, and the simplest
possible rate equation for how p changes in one time step (a rate for climbing into the
excited state, a rate for falling out of it, and the length of the step). For the canonical reset
(population one half to zero, with the rates chosen so the equation closes), the heat
delivered to the bath is exactly half the gap. Choose the gap
to be twice kT ln 2 and that heat is exactly kT ln 2.

Read that twice, because the temptation is to read it as a derivation of Landauer's number
and it's the opposite. The gap was \emph{chosen} to make the heat come out as Landauer's number.
The same populations and rates with half that gap give half the Landauer value, and there's
a theorem exhibiting identical population dynamics with different heat. The rate equation fixes the \emph{shape}; the energy
scale is a selected Hamiltonian (the physicist's name for the rule that says what energy each
state has), and selecting it is not deriving it.

The canonical reset is a bookkeeping toy, and its rates show it. A register that trades heat
with a bath at temperature $T$ has rates in \emph{detailed balance}: the rate up over the rate
down is $e^{-\Delta/k_B T}$, which at the gap $2 k_B T \ln 2$ is $1/4$. The canonical rates are
0 and 1, and the factor $e^{-\Delta/k_B T}$ is never 0, so no gap and no temperature put them in
that ratio. With rates that are in it, a
relaxation at a fixed gap moves toward the equilibrium population $1/(1 + e^{\Delta/k_B T})$ and,
with steps short enough not to overshoot, never goes below it. At the
gap $2 k_B T \ln 2$ that population is $1/5$. The
toy has no work term either; nothing in it pays for the squeeze. Part (4) of the theorem below
says what follows from that: at the gap $k_B T \ln 2$ the toy books less heat than Landauer's bound
asks for the same squeeze from one half to zero, so the toy has no second law.

\begin{definition}[The two-state protocol]\label{def:u24-calorimeter}
For a gap $\Delta$, the register's mean energy at excited population $p$ is $\Delta \cdot p$.
One step of the rate equation is
\[
p' = p + \delta t \, \big(k_{01} (1 - p) - k_{10} \, p\big),
\]
and the heat to the bath over a step from $p$ to $p'$ is defined as $\Delta p - \Delta p'$. The
canonical reset takes $p = \tfrac12$, $\delta t = 1$, $k_{01} = 0$, $k_{10} = 1$, ending at
$p' = 0$, and counts one logical event.
\end{definition}

\begin{theorem}[The heat is fixed by the chosen gap]\label{thm:u24-calorimeter}
(1) The canonical numbers satisfy the rate equation. (2) For every $\Delta$, the canonical
reset's heat is $\Delta / 2$. (3) For all $k_B$, $T$, the gap $\Delta = 2 k_B T \ln 2$ gives
heat $k_B T \ln 2$. (4) If $k_B > 0$ and $T > 0$, the heat is below $k_B T \ln 2$ exactly when the gap is below $2 k_B T \ln 2$; at the gap $k_B T \ln 2$ it is half of $k_B T \ln 2$. (5) There are two different gaps for which the same canonical populations
satisfy the rate equation and the heats differ.
\end{theorem}

\begin{proof}
(1) Half the population starts excited. The upward rate is zero, so none
enters the excited state; the downward rate is one for a step of length
one, so the entire excited half leaves. Substitution gives
$\tfrac12+1\cdot(0-\tfrac12)=0$.

(2) The register began with mean energy $\Delta/2$ and ended with zero.
The defined heat is their difference, $\Delta/2$.

(3) Choose $\Delta=2k_B T\ln2$. Half of that gap is $k_B T\ln2$.

(4) Asking whether $\Delta/2<k_B T\ln2$ is the same as multiplying both
sides by two and asking whether $\Delta<2k_B T\ln2$. At the particular
gap $k_B T\ln2$, part (2) gives half that value. With positive $k_B,T$
it is strictly smaller because $\ln2$ is positive.

(5) Take gaps one and two. The population equation in part (1) contains
neither gap, so the same populations solve it in both cases. Their heats
are one half and one respectively. The rate equation alone did not select
between them.
\end{proof}

A reset that pays its way is driven. Raise the gap a little, let the register settle with
rates in detailed balance, raise it again, and so on, starting from a gap of zero, where the
two states are alike and the population is one half. Raising the gap while the population
holds is work done on the register; settling at the new gap is heat to the bath. That's the textbook way to wipe a bit slowly, and here it is with its numbers worked out.

\begin{definition}[The driven reset]\label{def:u24-driven}
Write $kT$ for $k_B T$. At gap $\Delta$ the equilibrium excited population is
$\pi(\Delta) = 1/(1 + e^{\Delta/kT})$ and the register's free energy is
$F(\Delta) = -kT \ln(1 + e^{-\Delta/kT})$. Rates satisfy \emph{detailed balance} at gap $\Delta$
when $k_{01} = k_{10} \, e^{-\Delta/kT}$. A \emph{driven reset} along a schedule of gaps
$g_0, g_1, \dots, g_n$ takes, at step $m$, the gap from $g_m$ to $g_{m+1}$ with the population
held at $\pi(g_m)$, which is work $\pi(g_m)(g_{m+1} - g_m)$, and then lets the register settle to
$\pi(g_{m+1})$, which is heat $g_{m+1}\big(\pi(g_m) - \pi(g_{m+1})\big)$ to the bath. The reset's
work and heat are the sums over its steps. The uniform schedule from 0 to $D$ in $N$ steps has
$g_m = D m / N$.
\end{definition}

\begin{theorem}[A driven reset pays $kT \ln 2$ in work]\label{thm:u24-driven}
(1) With rates in detailed balance at gap $\Delta$, one step of the rate equation takes $p$ to
$\pi(\Delta) + \big(1 - \delta t\,(k_{01} + k_{10})\big)\big(p - \pi(\Delta)\big)$; when
$\delta t\,(k_{01} + k_{10}) = 1$ this discrete update lands on $\pi(\Delta)$ from any $p$. (2) For any schedule, work
minus heat is $g_n \pi(g_n) - g_0 \pi(g_0)$, the change in the register's mean energy. (3) If
$kT > 0$ and the schedule never falls, the work is at least $F(g_n) - F(g_0)$; if also every
step raises the gap by at most $h$, the work is at most $F(g_n) - F(g_0) + h\big(\pi(g_0) -
\pi(g_n)\big)$. (4) If $kT > 0$, $D \ge 0$ and $N \ge 1$, the uniform reset from 0 to $D$ in $N$
steps does work between $kT \ln 2 - kT e^{-D/kT}$ and $kT \ln 2 + D/(2N)$, and leaves an excited
population of at most $e^{-D/kT}$.
\end{theorem}

\begin{proof}
(1) Detailed balance gives
$k_{01}=(k_{01}+k_{10})\pi(\Delta)$. In the rate equation, the upward
contribution minus the downward one is
$k_{01}-(k_{01}+k_{10})p$. Substitute the balance equation and this becomes
$-(k_{01}+k_{10})(p-\pi(\Delta))$. The update therefore subtracts the
fraction $\delta t(k_{01}+k_{10})$ of the distance from equilibrium. What
remains is exactly
\[
 p'=\pi(\Delta)+\bigl(1-\delta t(k_{01}+k_{10})\bigr)
                   \bigl(p-\pi(\Delta)\bigr).
\]
When that fraction is one, none of the distance remains: $p'=\pi(\Delta)$.

(2) Follow one raise from gap $a$ to gap $b$. Its work is
$\pi(a)(b-a)$ and its heat is $b(\pi(a)-\pi(b))$. Subtract the heat from
the work. The two occurrences of $b\pi(a)$ cancel, leaving
$b\pi(b)-a\pi(a)$, the final mean energy minus the initial mean energy.
At the next raise, that final energy becomes the next initial energy.
Adding over the schedule cancels every intermediate energy and leaves
only $g_n\pi(g_n)-g_0\pi(g_0)$.

(3) Differentiating the displayed free energy gives slope $F'(\Delta)=\pi(\Delta)$.
With $kT>0$, this slope decreases as the gap rises. Over $a\le b$, the
increase of $F$ therefore lies between the widths times the two endpoint
slopes:
\[
 \pi(b)(b-a)\le F(b)-F(a)\le\pi(a)(b-a).
\]
The work uses the right-hand quantity, the slope before the raise. It is
at least the free-energy increase. Its excess is at most
$(b-a)(\pi(a)-\pi(b))$. If the raise is at most $h$, that excess is at
most $h(\pi(a)-\pi(b))$.

Add those bounds over the schedule. Intermediate free energies cancel,
and intermediate populations cancel in the excess bound. I obtain work
between $F(g_n)-F(g_0)$ and that amount plus
$h(\pi(g_0)-\pi(g_n))$.

(4) For the uniform schedule, $h=D/N$, $\pi(0)=1/2$, and
\[
 F(D)-F(0)=kT\ln2-kT\ln(1+e^{-D/kT}).
\]
For $x\ge0$, $0\le\ln(1+x)\le x$: it starts at zero and its slope
$1/(1+x)$ is between zero and one. The free-energy difference is thus
at least $kT\ln2-kT e^{-D/kT}$ and at most $kT\ln2$.
The excess from (3) is at most $(D/N)(1/2)$, giving the work bounds.
Finally the denominator of $\pi(D)=1/(1+e^{D/kT})$ is at least
$e^{D/kT}$, so the remaining excited share is at most $e^{-D/kT}$.
\end{proof}

Try one finite step in units where $kT=1$: start with gap zero and half the
registers excited, then raise the gap to $\ln2$. The work is $\tfrac12\ln2$.
After settling, the excited share is $1/3$, so the bath receives
$(\tfrac12-\tfrac13)\ln2=\tfrac16\ln2$ and the registers retain
$\tfrac13\ln2$ more energy. Work equals those two amounts added together.
A third are still excited: one such step has not wiped the bit.

The word ``step'' matters here. Part (1) is a discrete update of the population.
Its multiplier can be zero. Continuous relaxation at fixed positive rates
instead multiplies the distance from equilibrium by
$e^{-(k_{01}+k_{10})t}$, which is positive at every finite time.
Exact settling in the schedule is an idealisation; the algebra does not prove
that a continuous bath completes it in one finite waiting period.

Here nobody picked the gap to make the number come out. Take $D$ large and $N$ much larger
than $D$: the work closes in on $kT \ln 2$ from both sides, and the population left in the
excited state goes to zero. Perfect erasure is the limit, not the result at any finite gap. To reuse the
original symmetric memory, the gap must also return to zero while transitions
are suppressed. Lowering it at the final fixed population recovers
$D\pi(D)$ of work and leaves that population unchanged. This term tends to
zero in the same limit. The reset of the population, restoration of the trap
and time needed for settling are separate parts of the physical protocol.
Detailed balance supplies the thermal premise; the displayed sums supply the
work and heat accounting for the idealised schedule.

What ties it to a machine is that the protocol counts one logical event, by definition, and one is what the small machine's CERTIFY charges, the last of
the three paid acts behind a raised flag (Section~\ref{sec:smallest}). The tie is two numbers that happen to match, the protocol's one by definition and the machine's the least the record-layer premise lets CERTIFY cost, and nothing more. No theorem maps the
protocol's two register states onto any state of the machine. And ``heat to the bath'' here
means the register's mean energy before minus after, which quietly assumes the bath gets
exactly what the register loses; nothing here establishes that. So a device test would have to measure the gap, the
populations, the rates, and the bath heat independently; separately establish
that the device's rates are in detailed balance (that the up and down rates stand in the ratio the
temperature dictates), which the driven reset assumes and the canonical one breaks; and separately establish that the device's
physical reset \emph{is} the machine's certify event. The driven reset is an executable experimental
specification. The canonical reset is a bookkeeping toy. Neither is a measured device, and neither is an unconditional conversion of
units of cost to joules, and anyone who tells you otherwise has skipped the conditions.


\subsection{What this section is}

This section holds a calibration that's a unit conversion; a pair of premises, consistent on
the steps the toll is about and inconsistent on the whole instruction set; a history that makes any step injective by paying with a fresh state at every step; and a
measurement protocol that's a specification with a selected scale. The physics in this book is
in Part~I: Landauer's principle as a named premise with its critics named (Section~\ref{sec:a2-not-arbitrary}), and conditional finite-time dissipation bounds, cited, and proved for one model (Section~\ref{sec:what-a-chair-is}). Here there's
accounting, and I've tried to call every line by its right name. That's the end of the example machine.
What's left, after one suspicion I can't prove, is how you'd break any of this, and where it stops.


%===========================================================================
\section{From ``The claims and their conditions''}
%===========================================================================
\subsection{The diagonal and the worked models}

The diagonal argument supplies the contradiction inside a substrate that can represent
the candidate's flip and has the stated recursion property. Its fixed point makes the
proposed decision disagree with itself. An outside calculation of the natural-number
example does not supply an inside program for that substrate. Section~\ref{sec:diagonal}
and the appendices state the input conventions separately for the natural-number family,
L, the counter programs and the host.

The protocol examples in Section~\ref{sec:models} isolate particular mechanisms: finality,
a charge on commitment, an attestation window, an inclusion proof and a bounds check. Their
statements follow the operations in those fragments. The pointer thesis then asks why
certification is the event to choose. It proposes a connection to records copied and relied
on by independent parties. The counting and window arguments apply to any record meeting
their premises; the thesis supplies the proposed reason for this choice.

The game checks in Appendix~\ref{sec:chsh} give another worked construction. Four fixed
bits yield the classical bound by the displayed sign calculation. The matrix arguments
relate contraction, positive semidefiniteness and the score bound. The integer checker
implements its specified test on tallies with every pair sampled. Its earned chain connects
a raised flag to that test. The pinned-matrix test and the completion test ask different
questions: the appendix's all-ones and PR-box examples make the difference visible.

Those connections are what the book has built: a count, a schedule, an observation and an
earned record, with instructions carrying them through actual runs.

\subsection{Sources and contributions}

Appendix~\ref{sec:whats-mine} sets out the sources and contributions item by item. Here is
the short form: the earlier results I use, the statements proved here, and the constructions
that put them to work.

\emph{Borrowed, with names.} As \hyperref[sec:point]{The point} warned, I usually had the hunch first and
found these names after, and that order doesn't make any of it mine. Everything in
Section~\ref{sec:structure} and the Afterword (every
source is cited where it's used), the category theory (Eilenberg, Mac Lane) included. Minsky's
two-counter machines and the constructions that prove their halting undecidable and compile
computable functions into them (Forster, Larchey-Wendling and coauthors). Turing, Church, Kleene, Rice, Cantor, Gödel, the
diagonal pattern, Lawvere's unification, and later detailed developments:
Michael Norrish's, Mario Carneiro's, and Yannick Forster, Fabian Kunze and Gert Smolka's. Bell, CHSH, Tsirelson, Stephanie Wehner's semidefinite proof of Tsirelson's bound (2006), NPA, Landau,
Masanes, Fine, the elliptope, and Echenim and Mhalla's detailed treatment of Tsirelson's bound. Landauer, Bennett, Szilard,
Maxwell, and their critics. The potential method (Tarjan) and later treatments (Nipkow, Charguéraud and Pottier), Lamport's clocks, Floyd and Hoare. Gettier, Necula, Cook, Hartley,
Shannon, Rissanen, Zurek. The guarantee about every program: the reference monitor (Anderson),
execution monitors (Schneider), and monitors checked at every instruction and proved correct
(Azevedo de Amorim and coauthors). The counting under the permanence results: Poincaré's
recurrence, as Zermelo used it. The seL4 and CompCert discipline. The facts under the permanence results are Landauer's own example and Shannon's
bounds, and ``set a bit to one whatever it held'' is my way of saying his RESTORE TO ONE.

\emph{What I proved, each inside its stated interfaces,} is listed in the claims section (Section~\ref{sec:ledger}) with what each stands on.

\emph{Interpretation, which is mine and not proved.} That structure is what a computation
establishes and the toll is the right thing to price. That certification is the pointer
event. That the record axis ``belongs in the step.'' Those are proposals, with the arguments
I've given, and the proposals are where I'd aim first if I were you.

%===========================================================================


%===========================================================================
\section{Things I had to learn first}
%===========================================================================

I came in with none of the formal fields this work turned out to need. By the time the
argument had taken shape I'd needed enough of half a dozen of
them that I had to put down somewhere what those fields had words for. This appendix is that
place. None of it is the argument; it's what I needed to be able to make it. Everything
below is what I understood after I went and read about it; I'm not claiming to teach you
these subjects, and if you already know a part, skip it.

\subsection{Bits, bytes, and what a CPU sees}

A bit is one digit that's either zero or one, the smallest fact a machine can hold.
Everything else is bits stacked up. Eight of them make a byte, which lands you any number
from 0 to 255, and the first thing that genuinely surprised me is that the same byte is also
a letter: a convention called ASCII simply agrees to call 65 the letter A, 48 the digit 0,
32 a space. Nothing in the byte knows which it is. The meaning is a convention laid on top.

You can see the convention's seams with a pencil. Write 65 in binary: 1000001. Knock off the
first two digits and what's left, 00001, is 1, and A is the first letter. B is 66, which ends in
00010, two. Now lower case: a is 97, which is 1100001, the same pattern with one more bit
switched on, the one worth 32. Every capital and its small letter differ in exactly that bit,
and that's on purpose. The first version of the code, in 1963, had no lower case at all; when lower
case arrived, in the 1967 revision, it was placed so the two cases differ by one bit, which made
matching them, and building keyboards and printers, simpler. Nothing in the byte knows it's a
letter, but somebody made the letters easy for a circuit to tell apart.

Everything in this subsection is about real CPUs. None of it belongs to the abstract model,
which has a kind of state and a step rule and nothing about widths or slots, and the small
machine of Section~\ref{sec:smallest} needs very little of it: its two counters have no width at all, and hold any
whole number.

A word is a chunk of bits the CPU treats as one unit, in 2026 usually 64 bits (0 to about
18 quintillion) on a desktop or server and 32 (0 to about 4 billion) on many small devices. The abstract model doesn't care which width.

A register is a named storage slot inside the CPU that holds one word, the fastest place to
read or write. Memory is a longer list of word-sized slots accessed by index, the address. The
program counter is one specific register that tracks where in the program the machine is;
after every instruction it moves on to the next instruction unless the instruction says jump. On the small
machine it's called the line number. A stack pointer is a register pointing at the top of a
region used for nested calls. A port is a named channel for talking to the outside world.

Modular arithmetic means counting wraps around: counting modulo 16, the numbers run 0 to 15, and 17
wraps to 1.
Bitwise operations work on the bits of a value one position at a time: AND is 1 only when
both bits are, OR when either is, XOR (exclusive or) exactly when they differ. A logical shift moves
bits left or right: for unsigned values it multiplies or divides by a power of two,
discarding overflow on a fixed-width left shift and rounding down on a right shift. Popcount is the
count of one-bits in a value. Hexadecimal is base 16, digits 0 to 9 then A to F, written
with a 0x prefix: 0xFF is 255.

A trap, on a CPU, transfers control to a handler; the interrupted program may later resume.
It can report an error or request a service. On the small machine a trap is one bit that goes up when an instruction
refuses, and then the core stops moving. An ISA, instruction set architecture, is the list of
every instruction a CPU recognises, each identified by an opcode, a numeric tag the CPU
decodes. The small machine's has six. The abstract model has an instruction set and a step rule, but prescribes no binary opcode encoding.

\subsection{The notation the formal parts use}

I don't think in these symbols, and I had to learn to read them. The natural numbers are 0,
1, 2, and so on, written with a double-struck N; the integers add the negatives; the reals
are the whole number line. A function from a set A to a set B is written with an arrow. For
a finite set, vertical bars around it mean its size; for a number, its absolute value. A big
sigma is a sum over every member of a set. A delta before a quantity means the change in it.
An upside-down A means ``for every''; a backwards E means ``there exists.'' A double arrow is
implication; a double-headed one is ``if and only if.'' A small circle between two functions is
composition, apply the right one first. Angle brackets mean one of three things, and the sentence around them says which: a pair of numbers packed
into one, $\langle m, n \rangle$ (Appendix~\ref{sec:universal-host}); a dot product of two vectors; or, around a random quantity,
its
long-run average, which I picked up from the Bell-test papers. The circled times $\otimes$ in the quantum derivation combines two systems:
$A\otimes B$ acts as $A$ on the first and $B$ on the second; Appendix~\ref{sec:chsh-derivation} gives
its entries. The integer check itself uses neither it nor any quantum system. A superscript T on a matrix is its transpose, rows
and columns swapped, and I with a subscript n is the identity matrix.

\subsection{Category theory beyond the sketch}

Section~\ref{sec:structure} sketches objects, arrows, composition, identities, associativity,
and the identity laws. Beyond that: a functor is a structure-preserving map from one category to
another, sending objects to objects and arrows to arrows while respecting composition and
identities; think of it as a translation rule between two worlds of structure. An embedding
is a functor that throws nothing away, and faithful is the precise version for arrows: two distinct arrows between the same two objects never land on one. Conservative is the version for inverses: if the
image of an arrow has an inverse, the arrow already had one, so the functor can't
manufacture an isomorphism (an arrow with an inverse) that wasn't there. The two properties
are different, and neither gives the other. Read the natural numbers under addition as a
category with one object and an arrow for each number, and the same for the integers. The
inclusion of the first into the second is faithful, since distinct naturals stay
distinct
integers. It isn't conservative: 1 has no inverse among the naturals, and its
image has one,
minus 1. The other way round, the functor from the integers
onto the trivial one-arrow group
is conservative, since every arrow of the integers already
has an inverse. It isn't faithful:
it sends every integer to the same arrow. An embedding needs faithfulness specifically. A
monoidal product combines two objects or two arrows into one, with a unit and laws that make
the combining behave. An initial object has exactly one arrow into
every object; the counter's uniqueness in this book is a uniqueness-on-reachable-states
fact, and no category-theoretic initial object is involved. Yoneda's lemma was the piece that
made the categorical engine snap into place for me (Section~\ref{sec:origin}), and nothing in
the formal development uses it.

\subsection{Bell-test vocabulary}

A correlation between two outcomes is the average of their product; with outcomes plus or
minus one it lies between minus one and one, plus one meaning always match, minus one always
differ. A correlator is one such average in a labelled setting, E00 for Alice's setting 0
and Bob's 0. The marginal of a joint distribution is the distribution of one side alone with
the other averaged out; zero-marginal means each side's outcome is plus or minus one equally
often on average. A local hidden-variable model is one where both outcomes are fixed by a
shared random variable agreed before the game plus each side's own setting; Bell's theorem
says no such model produces certain correlations quantum mechanics does. No-signalling means
neither side's marginal depends on the other's setting. Quantum-realisable means compatible
with some quantum state and measurements; the NPA hierarchy gives a ladder of necessary
tests, each a positive-semidefiniteness condition on a moment matrix. For the four correlators
of this game alone, the first rung is already exact; for the full table, each side's own odds
included, Chaturvedi's 2026 preprint argues that no finite rung is, even in this simplest
case. The PR box is the no-signalling correlation that wins the game every
time, stronger than quantum mechanics allows. Cauchy-Schwarz is
the inequality that a dot product is bounded by the product of the lengths. The Tsirelson
bound is two root two, the largest game score quantum mechanics allows; the classical bound
is two and the PR bound is four. (If you go looking for Boris Tsirelson's 1980 paper, look
under Cirel'son, the spelling the journal printed, or you may not find it.) The elliptope is
the set of correlators that can be filled
out into a positive-semidefinite matrix with ones down the diagonal.

\subsection{Other terms I use without rebuilding}

Provenance is the recorded chain of where an object came from; on the small machine, a
commitment's provenance is the passing check that came before it (Section~\ref{sec:smallest}).
Simulation, in computer science, is one system
imitating another's input-output behaviour; a Turing machine can simulate a system with effective finite encodings of its states and moves
and a computable step on those encodings, my small machine included, by representing its state on tape, and the difference
this book is about is what a window on that tape shows. That's real, and it raises a fair
question: if a tape can stand in for the state, what was the
difference? The tape carries a picture of the state, the way a photograph carries a
fire, and a photograph of a fire keeps you exactly as warm as any other photograph. The counter and the toll live in
the step rule. An encoding copies the bytes, and the step rule isn't in them. A Turing machine that bumps a
ledger on its tape whenever its table sets its flag keeps that ledger in its state, and a ledger
kept in the state isn't a price on the move. The toll prices the machine's moves: with one move,
the next step its table picks, every step costs at least one once the flag can go up at all, and
with a move for each row of the table, every row that can set the flag costs at least one wherever
it's taken (Section~\ref{sec:machine-vs-program}). The ledger on the tape only simulates that. What no machine can do is make a simulated rule binding just
by writing it down on tape.

Projection is a function that drops information on purpose. The window on the small machine
keeps three things, the line number and the two counters, which are what Minsky's machine
already has, and that's why projecting lands you back in ordinary computation. It drops
everything else, including the parts doing the new work: the versions, the table of checked
facts, the channel, the trap, the ledger and the flag. Nothing is wrong with it, and it's a perfectly good
function; it just loses you what made the state more than a pair of counters. When what's left runs
as the machine underneath runs, reaches what it reaches and still can't see the record or the
ledger, the projection is a slice (Definition~\ref{def:slice}): every Thiele-complete machine has one
onto two counters, its window, and every computer with a move for each two-counter instruction is a
slice of a certified Thiele-complete machine (Theorem~\ref{thm:slice}). A quotient is the operation of declaring two things the same when
they agree on what you chose to look at and working with the classes instead. A projection
does exactly that. Dropping fields one at a time also glues together states that
the dropped fields kept apart. The fibre of a function over a value is
everything the function sends to that value. An injection sends no two inputs to the same
output; a surjection covers every output. Decidable means an algorithm always halts with a
yes-or-no answer; keep the claim and the algorithm apart, since the algorithm decides the
claim only once there's a proof that its answer matches. Polynomial time means running time
bounded by a power of the input size; P is the problems solvable that way, NP the problems
where a candidate solution can be \emph{checked} that way, and whether the two are the same is
open.

``Kernel'' has two uses here. In mathematics, the kernel of a map is what it sends to the
identity; for a projection the useful version is the relation ``these two states project
to the same thing.'' An operating-system kernel is the core of an operating system, as in
the packet-filter example in Section~\ref{sec:models}.

\subsection{Matrices}

A matrix is positive semidefinite (PSD) when $v^{\mathsf T} M v \ge 0$ for every vector $v$;
the game check is all about one matrix being PSD. The determinant of a square matrix is a
single number that is zero exactly when the matrix squashes some direction flat. For a
two-by-two matrix with rows $(a, b)$ and $(c, d)$ it's $ad - bc$. For a bigger one, expand
along a row: for each entry, delete its row and column, take the determinant of what's left,
multiply by the entry, alternate the signs, and add up; the recursion bottoms out at two by
two. A principal submatrix is what you get by choosing some of the indices and keeping only
those rows and the same columns; its determinant is a principal minor. A symmetric matrix is
PSD exactly when every principal minor is at least zero, every one of them, and the step from
the game check to Tsirelson's bound uses two of the three-by-three ones (Appendix~\ref{sub:chsh-five}). The Schur complement turns a question
about a block matrix into a question about a smaller one. For a symmetric block matrix with
corner block $P$, off-diagonal block $Q$ and other corner $R$, it's $R - Q^{\mathsf T} P^{-1}
Q$, and when $P$ is positive definite the whole matrix is PSD exactly when that complement is
(Lemma~\ref{lem:chsh-schur}). Positive definite matters; invertible isn't enough. Take $P = -1$
and every other entry zero: the complement is zero, which is PSD, and the whole matrix
isn't.

\subsection{Borrowed ideas the book pushes on}

These get their full story in the chapter that uses them, and each gets a one-line version here for
when you've forgotten which one Landauer was.

The pigeonhole principle: put more pigeons than holes and some hole holds two. Section~\ref{sec:a2-not-arbitrary}'s room with too few chairs is that, and so is the proof that, when only finitely many states read yes, a
permanent flag can only be raised by a move that forgets something, two states landing on one. Landauer's principle (Rolf Landauer, 1961): a step that wipes
information, lowering the entropy of what a machine holds by some number of bits, gives off mean heat
at least Boltzmann's constant times the temperature times the natural log of two per bit,
under a physical model relating that logical entropy to the memory's thermodynamic entropy. The bound is a
floor that going slowly doesn't lower, and it isn't all the heat: specified thermal dynamics can require additional entropy production for a finite-time change of distribution. In this book Landauer's floor is always a premise with its name on it. Nothing here derives it (Section~\ref{sec:a2-not-arbitrary}, Section~\ref{sec:landauer}).
Entropy, in Claude Shannon's sense: for a spread of chances over states, roughly the average number
of yes-or-no questions you'd need to pin the state down, so an even spread over $n$ states
has $\log_2 n$ bits. A spread, in this book, is just a list of chances that add up to one.
Maxwell's demon (Section~\ref{sec:narrowing}): a little creature that watches molecules and sorts fast from slow,
seemingly for nothing; Charles Bennett's answer was that the looking can be done without
merging anything, so Landauer's premise sets it no floor (it can dissipate under the stated thermal-model assumptions, less the slower it goes), and the floor lands when the demon wipes its notes to look
again. The potential method (Robert Tarjan, 1985): pay
for a run with a quantity stored in the state, so a step's cost plus the change in the
stored quantity is what you count; Section~\ref{sec:counter} reads the toll as that method with one particular
potential.

The halting problem (Alan Turing, 1936, who essentially proved it without stating it in
this form; Stephen Kleene stated it in 1952, and Martin Davis seems to have been the first to name it, in 1958,
Section~\ref{sec:diagonal}): no program decides, for every program and input,
whether it stops. Kleene's recursion theorem (Stephen Kleene, 1938): for any transformation
of programs that can itself be computed, some program computes the same partial function as its own
transform. Its timing, storage and records need not agree. Rice's theorem (Henry Rice, 1953): any property that depends only on what a
program does, and isn't true of all programs or of none, can't be decided by a program. The
diagonal is the move under all three: build the thing that does the opposite of what the
table says about it (Section~\ref{sec:diagonal}). L is a small lambda calculus, a language made of nothing but
functions applied to functions, which Yannick Forster and Gert Smolka presented and which
Section~\ref{sec:diagonal} uses as a tool. A cryptographic hash is a short fingerprint of a longer piece of data, made so
that finding two inputs with the same fingerprint is supposed to be hard; a hash chain folds
each new record into the fingerprint of the last, and a Merkle tree folds them pairwise into
one root (Section~\ref{sec:toll}, Section~\ref{sec:models}).


%===========================================================================
\section{What's mine, and what isn't}
\label{sec:whats-mine}
%===========================================================================

Every chapter of this book leans on somebody, and I've named them where I lean. This appendix
does a different job, and it's the one where I'd most like to be caught out, and caught out
early. For each thing the book could be accused of claiming as its own, it says what was
already known, what this work actually defines, builds, proves, or connects, and what kind of
contribution that is. Five kinds, and they aren't the same claim: a \emph{new result} (a statement
nobody had); a \emph{new proof} of something known; an \emph{implementation} (something built that runs); a \emph{synthesis} (known
pieces put together under one question); and a \emph{conjecture} or \emph{thesis} (a claim stated and not proved; a thesis ties defined notions to the world, so no proof can settle it). Every kind can be worth doing.
None of them is the others. And if you read down the list and notice that ``new result''
hardly shows up, that's good, and it means you're reading it right.

Two warnings about my own evidence before the list. First, ``I didn't find it'' is a fact about
my searches. The literature is bigger than my searches. I searched with every tool I had, and some tools will
cheerfully invent a citation if you let them, so where it matters I say where I looked, and
you should assume it may well be sitting in a paper I'd need a few more lifetimes of homework
to read. Second, nothing below is inferred from how many theorems there
are or from a good story about who came before. A proof, if it's right, tells you the statement
is true. It doesn't tell you it's new. A thousand correct proofs could be a thousand careful
ways of agreeing with Tarjan.

After the prior-work map come the items themselves, fifteen of them, numbered as their subsections are, in the order the
book reaches them: what keeping takes first, certification last.

\subsection*{The prior-work map}
\phantomsection\addcontentsline{toc}{subsection}{The prior-work map}

\emph{What I do and don't know about prior work.} I came to this with no formal training in the
fields it turned out to need, and I've read what I could reach, with every tool I had. The
prior-work map below is orientation. It doesn't pretend to be scholarship. For each piece of this book it says who did
the thing it leans on, and what, if anything, this development adds. A row that's wrong gets
fixed, and, as the end of Section~\ref{sec:ledger} says, I'd want to hear of one.

There's a law about rows like these, and it's a joke that keeps turning out true. In 1980 the
statistician Stephen Stigler proposed that no scientific discovery is named after the person
who first found it, and then credited the idea to the sociologist Robert K. Merton, who had
been writing about fights over scientific priority since 1957. So Stigler's law is an example
of itself, which is either the best joke in the history of science or a very tidy accident.
Is it a law? No, it's a pattern, and nobody proved it. But this book has already tripped over it
twice, in \hyperref[sec:point]{The point} at the front and in Section~\ref{sec:a2-not-arbitrary}: the pigeonhole principle wears Dirichlet's name, and
Leurechon had it two centuries earlier. And I put my own name on a machine. By Stigler's law, that makes it
somebody else's. The table below is my best attempt at finding out whose.

What I compare each time is what a result is \emph{about}. Landauer's bound is about heat, a
physical quantity under a thermodynamic model. No Free Lunch (an optimisation result, in the
list below the table) is about averages over a whole spread of problems. The toll is a floor on one kind of step, the step that turns the
reading on. Rise-only pricing adds a second rule on top: charge nothing on a step that doesn't turn it on. So the
comparisons are about the shape of each result. They aren't claims about who got somewhere
first, and they don't show that certification is the only event worth pricing.

The right-hand column uses words built in the body of the book (the latch, the fibre, the
universal program, the integer check, the matrix test); the subsections after the table say where
each one is built.

\begingroup\small
\begin{longtable}{@{}>{\raggedright\arraybackslash}p{0.242\linewidth}>{\raggedright\arraybackslash}p{0.344\linewidth}>{\raggedright\arraybackslash}p{0.344\linewidth}@{}}
\toprule
\textbf{Prior work} & \textbf{What this book leans on} & \textbf{What this development adds} \\
\midrule
\endhead
\bottomrule
\endlastfoot
Landauer 1961; Bennett 1973, 1982, 2003; Zermelo 1896; Bocchieri and Loinger 1957; Pin 1992; Brodsky and Pippenger 2002; Ouldridge and Wolpert 2023; Murashita, Funo and Ueda 2014; Li and Vitányi 1996; Demaine, Lynch, Mirano and Tyagi 2016 & the heat of erasure and merging; a unit price per irreversibly erased bit as an axiom of a model of computation, and on the random-access machine; a measurement need not merge, and the demon pays when it wipes its notes; recurrence (a step that loses nothing, on finitely many states, comes back around); machines whose moves all reshuffle states recognise only group languages; recognising a language whose smallest automaton has to merge states produces entropy; absolute-irreversibility bounds require forward and reverse path probabilities & the ledger form of a counting core that's classical (recurrence's finite cousin, folklore for permutation automata) and already priced in entropy for automata, and priced per erased bit and per instruction by Li and Vitányi and by Demaine and coauthors: a certification toll derived from finitely many yes-states, permanence and that rule, with merge-or-revoke, the squeeze and entropy forms, and each premise's counterexample built \\
Maroney 2009; Sagawa and Ueda 2009; Sagawa 2014 & thermodynamic bounds that depend on the actual spread over states; the entropy-loss bound is zero for a point-mass input; further work bounds depend on control, reliability and input distribution & the certificate-specific entropy bound with its support and permanence premises out in the open; a heat floor still needs the Landauer premise, and no new physical law is proved \\
Shannon 1948; Hartley 1928 & entropy and log-of-count bounds & the same squeeze bound in counted states and, from an even spread, in bits, for a permanent flip \\
Tarjan 1985; Lyapunov 1892; Nipkow; Charguéraud and Pottier & the potential method and later treatments & the toll as that method with the flag as potential \\
Katsumata 2014; Orchard, Liepelt and Eades 2019; Danner and Licata & cost carried as a grade that composes with the computation & a run's cost as a writer cost and as a grade fixed before the run; the counter rides in the state, put there by the step rule, and it is their writer exactly, while their type systems and before-the-run checking of grades are not embedded, by design \\
Lamport 1978; Floyd 1967; Turing 1949 & one-way counters give order and termination & the receipt's three properties (it only climbs, it equals the sum of the charges, it's the only such total on the states a run from the start reaches) stated and proved for any machine with a counter and a price list fixed in advance \\
Moore 1956; Kalman 1960; Cousot 1977; Codd 1970; Sampath and coauthors 1995; Saboori and Hadjicostis 2007 & what a window loses; when an observer can work out a hidden event, and when a secret state stays hidden & the fibre fact with a price in place of a decoder; the small machine's named collision, and the window every certified Thiele-complete machine is blind through \\
Smith 2009 (quantitative information flow) & how much an observer learns about a secret from what it sees & none for leakage: the collisions rule out exact decoding through named windows, which isn't zero leakage, and no leakage measure is supplied \\
Gettier 1963; Necula 1997; Cook 1971 & entitlement is not truth; proof-carrying checks & entitlement as object, relation, paid event, with a counted bound \\
Dwork and Naor 1992; Back 2002 (Hashcash, first proposed 1997) & make establishing something cost, in a way anyone can check cheaply & a comparison, with no claim that the two are the same: proof of work checks a found value against a difficulty target, while the toll charges one unit at the step that certifies and verifies nothing about effort \\
Eccles and Jordan 1918 (patent filed) and 1919 (paper); Plotkin 1989; Shapiro and coauthors; Abadi and Lamport 1991; Ahman and coauthors 2018 & the latch; write-once and grow-only state; history variables; monotonic state whose witnessed facts can be recalled later & the latch form over any base, with its two counterexamples and converse \\
Cantor, Gödel, Turing, Kleene, Rice; Lawvere 1969 & the diagonal & the shortcut wall on any substrate; L's recursion theorem by proof \\
Bell, CHSH, Tsirelson, Wehner, NPA, Landau, Masanes, Fine; Echenim and Mhalla 2023 & the game, its walls, and the matrix test & an integer check on counts that implies the pinned matrix is PSD, which implies the bound \\
Buterin and Griffith; Wood; TCG; RFC 9162; Necula & five record disciplines & five fragments read through one question: is the write priced \\
Zurek 2009; Ollivier, Poulin and Zurek 2004; Blume-Kohout and Zurek 2006 & pointer states and redundancy & a thesis borrowing the shape, with the gap to the physics stated \\
Leroy 2009; Klein et al. 2009 & the discipline of saying where the proof stops & the separation of what's proved about the model from what any device running it does, inside its bounds \\
Minsky 1961, 1967; Forster, Larchey-Wendling and coauthors 2019, 2020 & two-counter machines are universal, and their halting is undecidable & a machine of six instruction kinds on two counters where earned commitments, no forging and an exact minimum price of three are theorems, with halting reduced from two-counter halting \\
Forster and Smolka 2017; Forster, Larchey-Wendling and coauthors; Norrish 2011 & L as a model of computation; counter-machine constructions; a compiler from computable functions to counter machines & L's recursion theorem, Rice and halting; a counter compilation putting every computably presented Thiele machine on the universal host \\
Turing 1936; Post 1947; Popek and Goldberg 1974; Hamlen, Morrisett and Schneider 2006 & the universal machine as one fixed table; hosts that keep guests apart; which policies a monitor can enforce & a definition of certified Thiele-complete that every clock fails; one fixed program in the small machine's own instruction kinds that runs every program of it with halting, flag and ledger exact, earning its own flag; and one fixed program, with a paid move added, that runs every computably presented Thiele machine, its ledger exact to within a surcharge of at most two when the guest's reading starts at no, which can't be avoided in general \\
\end{longtable}
\endgroup

Everything in the right column is stated with its premises in the body of the book, and
everything in the middle column is cited where it's used. The table is the short form. The
long form is the rest of this appendix, which takes each row, says what was
already known with sources, and says what kind of addition each one is. Very little of the right column is a new result. Most of
it is things built, new proofs of known things, and known pieces put together under one question.

A few more labels I went and looked up are neighbours. Some of them also sit in
the table above,
and one, the reference monitor, is the move the small machine's step makes. I list them
so you can see why I don't think the toll is one
of them under a new name. Each is a claim inviting correction.

\begin{itemize}
\item Linear logic (Girard 1987) treats an assumption as a resource that gets used up. It
constrains what a program, or a proof, may consume. The toll constrains the model's step
function instead. It's the same instinct at a different layer.
\item Foundational proof-carrying code (Appel 2001) checks safety proofs against the smallest
trusted base it can: a logic plus the machine's semantics. Here no proof rides along with
the program. The step rule writes the receipts as the run goes.
\item Session types (Honda 1993) constrain how programs talk to each other so the conversation
can't go off script. The toll is about what one computation has paid to establish,
which is a different problem.
\item No Free Lunch (Wolpert and Macready 1997): averaged over every problem, whatever an
optimiser gains on some problems it gives back on others. That's an average over a whole
problem space, and the edge is paid for elsewhere. The toll charges one step in one run,
with no refund anywhere.
\item Kolmogorov complexity (Kolmogorov 1965): the description complexity of a string is the length of the
shortest program that prints it, relative to a fixed description language. That's a property of a finished string, with no steps and
no moment of paying. The toll is a law about a step.
\item Irreversibility cost and information distance (Zurek 1989; Bennett, Gács, Li, Vitányi and
Zurek 1998; Li and Vitányi 1996): the fewest bits a reversible machine must take in and erase
irreversibly to turn a string $x$ into a string $y$, about $C(x \mid y) + C(y \mid x)$, where
$C(x \mid y)$ is the length of the shortest program that prints $x$ given $y$; and Kolchinsky and
Wolpert's thermodynamic complexity of a Turing machine's outputs (2020). These tie algorithmic
information to steps and erasures, and they price a transformation between strings, or an output,
in erased bits. The toll prices the rise of a chosen reading.
\item Cook and Levin (Cook 1971; Levin 1973): checking can be cheap while finding looks
expensive, and whether finding is really harder is the open P-versus-NP question, a gap in
how difficulty grows with size. The toll puts a fixed floor under the one step that
commits, however easy or hard the finding was.
\item The thermodynamic cost of choosing (Rovelli 2024): settling a binary choice at temperature
T dissipates at least kT ln 2, with no outside heat bath. That's the closest thing to the toll itself I've found, because it
charges a commitment. An object or an average isn't what it prices. But it's physics, indexed to a
temperature, and its choice is unpredictable, a branching, and stable only for a while, where
the toll prices a deterministic no-to-yes transition. On my side, finitely many yes-states and
permanence are what force the raising move to merge states. The toll charges the commitment in a model of computation, with no
temperature in it. It's the closest neighbour, and it's still a different thing.
\item Persistent copies (Ouldridge and ten Wolde 2017; Ouldridge, Govern and ten Wolde 2017):
making a copy that stays correlated with its original after the two come apart has a
thermodynamic floor. The nearest physics I've found to ``learning borrows, committing pays.'' It's
work and free energy for a correlation, in chemistry; the toll is a count in a model of computation.
\item Thermodynamics of finite automata (Ouldridge and Wolpert 2023): if the smallest automaton
for a language has a letter that sends two states to one, recognising that language on a
device that repeats one local step produces entropy, and if the smallest automaton only
reshuffles, recognition carries no such floor. That's Part~I's counting of states with a physical price on it, and one
of their worked examples is the latch's language turned inside out. It's the closest neighbour
of all to the road the toll is derived on; Section~\ref{sec:a2-not-arbitrary} says exactly
what's left over for me.
\item CerCo, the Certified Complexity project of Amadio and coauthors, builds a verified C compiler that carries the object
code's execution cost back up to annotations on the source. There the cost decorates the
program and survives compilation. Here it lives in the model's step rule.
\item Reference monitors and execution monitors (Anderson 1972; Schneider 2000; Azevedo de
Amorim and coauthors 2015). A small piece every action passes through, enforcing a rule on
every program. Schneider argued that a monitor watching a run can enforce only rules whose
breaking shows up at a finite point, and ``no certificate without an earlier passing check and
commit'' is one of those; the micro-policies check a tag rule at every instruction, proved correct. That's the move the small machine's step makes, so this one isn't a neighbour at all: the small machine's step is a monitor of this kind. What's added is the edition-bound check-then-commit with a paid, permanent record on top,
and the derivation of why the payment can't be zero. Kevin Hamlen, Greg Morrisett and Fred
Schneider (2006) went on to sort which rules monitors, program rewriters and static checkers can
each enforce, and they differ, which is one reason this book doesn't claim that every enforcement
mechanism is a Thiele machine underneath.
\item Trusted components in replicated services (Chun and coauthors 2007; Levin and coauthors
2009; Veronese and coauthors 2013; Decouchant and coauthors 2022; Clement, Junqueira, Kate and
Rodrigues 2012): a small record inside each replica that a replica taken over still can't make
say two things, with a proof of exactly what that buys. The record is theirs. The toll isn't in
their picture.
\item Ghost state (Owicki and Gries 1976; Filliâtre, Gondelman and Paskevich 2014): variables
added for a proof that can be erased without changing what the program computes. The flag and
the counter are ghost in that sense. The small machine's table of checked facts isn't, because
the step reads it and refuses on it, for every program. The slice of
Appendix~\ref{sec:every-model} is that erasure on the lift, while its trap is down, with two things
on top: the erased layer still earns its record, and what's left can't read the record or the
ledger (Theorem~\ref{thm:slice}).
\item Verifiable computation and proof-carrying data (Parno, Howell, Gentry and Raykova 2013;
Chiesa and Tromer 2010): a short proof, checkable by a stranger, that a computation's output is
right. That establishes the claim for the stranger. Producing it may require a witness and substantial computation; it does not identify the original computing device unless that is part of what the proof establishes. Another prover may reproduce it,
the copier included. The record here is about one machine's own history. A proof like that is
one way to carry the local record out of the box.
\item Commitments and entitlements (Brandom 1994): talk as score-keeping, tracking who is
committed to what and who is entitled to it. Brandom's entitlements can be defeated and can be
passed on by testimony. Here entitlement is priced, step by step, and a copy carries it only as
far as the copy is genuine. The words are the same and the machinery is different.
\item Correspondence assertions (Woo and Lam 1993; Gordon and Jeffrey 2003). In protocol
verification, ``this end event must be preceded by a begin event for the same thing'' is a
correspondence assertion, and a type system can prove a protocol keeps it against any
opponent. ``No certificate without an earlier commit, and no commit without an earlier check
of the same claim'' has that shape. Theirs is proved about programs by typing. Mine is enforced
by the step, which is the LCF comparison again, at the level of single instructions.
\item Monotonic parallel composition (\AA{}man Pohjola and Parrow 2016) and persistence in the
pi-calculus (Palamidessi, Saraswat, Valencia and Victor 2006). \AA{}man Pohjola and Parrow
separate models where a neighbouring process can never stop a step from models that can solve
consensus, and test-and-set, an instruction that writes a cell and reports what was there, is
on the consensus side. Palamidessi and coauthors compare messages that are used up with
messages that stay, and the end of their study is a reduction from Minsky machines. The small
machine, put into \AA{}man Pohjola and Parrow's setting with threads pooled on one shared state
(two states that have both moved jam), has a consensus process in their sense, so it lands on
their non-monotonic side; their theorem and that placing are argued in a few lines in
Section~\ref{sec:smallest}. The comparison with Palamidessi and coauthors' persistence stays a
comparison.
\item Constructor theory (Deutsch and Marletto 2015; Marletto 2016; Marletto, Vedral and
coauthors 2022; Violaris and Marletto 2022): physics as which tasks a device that runs in a
cycle can do, with irreversible meaning possible while the inverse isn't. It's the nearest
physics I know to ``committing costs'' stated for a cycling device. It's quantum and
approximate, and Part~I's counting of states is classical, finite and exact
(Section~\ref{sec:a2-not-arbitrary}).
\item Accountability and fault detection (Haeberlen and Kuznetsov 2009; K\"usters, Truderung
and Vogt 2010; Treitz and K\"unnemann 2026): when an outside party can tell who misbehaved.
Their windows are this book's windows, and the trust they leave over, an honest log or a
trusted step, is left over here too (Section~\ref{sec:windows}).
\item Transferability (Easwaran 2009): a deductive proof lets you learn from someone you don't
trust because you can check every step, and a probabilistic proof doesn't. That's the argument
for carrying the object and re-running the check (Section~\ref{sec:windows}).
\end{itemize}

And then the silly question that list raises, which I'm going to chase anyway: why is an
optimisation theorem named after lunch? Wolpert and Macready borrowed an old saying, ``there
ain't no such thing as a free lunch,'' and the saying is older than almost everyone who gets
credit for it. American saloons in the 1800s put out a free lunch for anyone who bought a
drink, and the drink paid for the lunch. A Washington newspaper was already saying in 1909
that somebody has to pay for a free lunch; a newspaper fable by Walter Morrow had the whole
line in 1938; a book by Pierre Dos Utt already had the squeezed form, TANSTAAFL, as its title in
1949; Robert Heinlein's 1966 novel \emph{The Moon Is a Harsh Mistress} made that squeezed form
famous; and the economist Milton Friedman put the saying on the cover of his 1975 book of essays on public policy. I'd have credited
Heinlein or Friedman. As far as anything I could find goes, neither coined it. That's Stigler's law again,
at lunch, and it matters here because No Free Insight is named like a cousin of the saying, and the
family resemblance misleads in one way. The saloon's lunch was paid for somewhere else, in the price of the beer, and that's
Wolpert and Macready's shape: a gain here, handed back over there. The toll isn't that.
On the ledger nothing gets paid somewhere else: the step that raises the flag pays, on the spot. A
blank tape prepared by wiping an uncertain one may have an earlier reset cost under the thermal premises. A supplied blank alone does not prove that payment; the toll is the
step that won't hide its share that way.

\subsection{Every step is a fresh state, a merge, or a loop}

Known: the step law is the first-return argument every textbook on dynamics makes, and under it
sits the pigeonhole. A one-to-one move that raises a reading and keeps it is Hilbert's hotel,
everybody moving up a room so that a room is left empty (Hilbert's lecture course of 1924--25, made
famous by Gamow in 1947), and a collection that can be matched one-to-one with a part of itself
that leaves something out is exactly what Dedekind (1888) called infinite. Recurrence for
reversible finite machines is Poincar\'e's (1890) and Zermelo's (1896) in its finite form. The
noisy steps are the fact that a doubly stochastic chain has no closed class it can enter from
outside, with Gibbs's inequality for the entropy gap; the unitary steps are that a unitary
preserving a finite-dimensional subspace preserves its complement; the bilateral shift is the
standard example of what goes wrong without finite dimension. Heat against time is Bennett's
Brownian computer (1982) with Schnakenberg's formula for the entropy a biased walk hands its surroundings; the bounds for finite-time changes of probability distributions under specified stochastic dynamics belong to finite-time stochastic thermodynamics (Schmiedl and Seifert, 2007; Shiraishi,
Funo and Saito, 2018); and the trade of space for time by recomputing is Bennett's (1989). The measurements are B\'erut and
coauthors' (2012), Jun, Gavrilov and Bechhoefer's (2014), Hong and coauthors' (2016) and
Gaudenzi and coauthors' (2018). Mine: the reading, that the write of a 0 on a cell is a merge and
the toll is the price of something being the case; putting those old facts side by side as one
law with three columns (a fresh state, a merge, a loop that makes nothing new) and no fourth, for
deterministic runs, beside separate preservation theorems for the stated noisy and quantum models; the merge-or-fresh-state dichotomy for a reading that stays
up, with no finiteness in it; and the exact count of states for a bounded memory,
$M \ge 1 + (T+1)(N-1)$, with a building that reaches it (Section~\ref{sec:a2-not-arbitrary},
Section~\ref{sec:what-a-chair-is}); and that a merge seen in a coarse description forces, in any finer description it's exact about,
a fine merge, a fine step out of the set, or infinitely
many fine states (Theorem~\ref{thm:every-scale}). Kind: synthesis and a reading, with the count
of states and the every-scale theorem as small exact lemmas. The mathematics is old. What I'd defend is the arrangement: endless memory,
history-keeping and reversible machines aren't a way out of what keeping takes, they're the fresh-state column,
which dodges only the floor a merge adds and does not settle the implementation's finite-time dissipation; and finiteness
is what takes that column away.

\subsection{The universal floor and the potential method}

Known: the potential method of amortized analysis, a function on states whose drop bounds
the cost of a run, is Tarjan's (1985). Later developments include Nipkow's accounts
(2015 and 2019) and Charguéraud and Pottier's time credits (2019). It has been built
into type systems for automatic resource analysis (from Hofmann and Jost, 2003, and Hoffmann and coauthors, 2011, onward), and
recast in cost-aware type theories (Niu, Sterling, Grodin and Harper, 2022). Mine: choosing
the potential that is one when a designated reading says no and zero when it says yes, so
that ``the crossing costs'' (the toll) and ``every run that crosses paid'' (the floor) are the
potential method's two halves (Section~\ref{sec:counter}), and its count, one unit per crossing
(Observation~\ref{thm:per-crossing}), the same sum read crossing by crossing. Kind: a known technique, used with
one particular potential. The choice is the reading, and it's the point: it makes ``was certified'' something the account must carry. The numbers come with it. A natural-number potential, as defined here, that reads only the reading has to drop by exactly one from no to yes, or its inequality says something other than the toll on a machine of two states and one move; so it's one at no and zero at yes, up to a constant added to both. The theorem is not new mathematics. If it looks clever, that's Tarjan's cleverness; my part was picking which reading to plug in.

Go back to Tarjan's paper and there's a sentence that sits right on the line this book draws.
He opens with the dictionary (to amortize is to put money aside at intervals, as in a sinking
fund) and then gives two views of the method. In the banker's view the computer is
coin-operated, and each operation is handed credits to spend or save; in the physicist's view,
which he credits to Daniel Sleator, every state has a potential and you count its changes.
Then he stops to say that the credits are only an accounting device: the programs that actually
run contain no mention of credits or debits. That's exactly right for what he wanted, a bound
proved about a program from outside. It's also exactly the gap. In his method the coins live in
the proof. The one move I chose was to put them in the step, where no program takes a step
without it being counted. The arithmetic is the same, and only the place is different.

\subsection{The toll from finiteness, permanence, and a price on merges}

Known: that steps which merge states (no single-valued inverse) are what Landauer's principle
charges for is Landauer's (1961), sharpened by Bennett (1973, 1982, 2003; the words ``the erasure of a bit
or the merging of two computation paths'' are from the 2003 abstract); generalizations and the dependence on the input
distribution are Maroney's (2009), Sagawa's (2014), Wolpert's; the critique is Earman and
Norton's and Norton's. Write-once memory is an old topic in coding theory (Rivest and
Shamir, 1982), where the question is reuse, and what it takes is counted in cells, which in this book's words is paying in fresh states; heat and prices never come up, and Chris Heegard (1985) worked out the capacity of permanent memory with noise, again a question of how much fits.
Statisticians priced permanence on finite memory in their own currency: in Hellman and Cover's
analysis of hypothesis testing with finite memory (1970), a machine that settles its answer by
falling into states it never leaves does worse than one that doesn't, so a decision held for
good costs memory states. Heat never enters their analysis. And the counting under the theorem is older than
any of that, twice over. Theorem~\ref{thm:perm-flip-merges} is the finite, discrete cousin of the
recurrence argument Zermelo, using Poincaré's theorem of 1890, turned on Boltzmann in 1896, here
aimed at a reading. Zermelo aimed it at entropy, and Bocchieri and Loinger gave the quantum recurrence in
1957. In automata form it is classical: a finite machine whose moves are all one-to-one
reshuffles recognises only group languages, and the inputs that raise a permanent reading form a
language that is never one of those unless it is empty (Pin's 1992 paper on reversible automata
is the reference I'd check; automata theorists treat the step as folklore), and going quantum
doesn't change it (Brodsky and Pippenger, 2002: quantum finite automata that measure once, at
the end, accept with bounded error exactly the group languages). The closest of all is recent.
Ouldridge and Wolpert (2023) give a proof that if the smallest automaton for a language has a step
that merges states, every automaton for it does, and argue that recognising it on a device that repeats
one local step carries an entropy floor, while a language whose smallest automaton only
reshuffles carries no such floor; one of their worked examples is the strings with no $b$, the complement
of the latch's language. That's the core with a physical price already on it. And the
logarithm in the heat floor matches the absolute-irreversibility bound of Murashita, Funo and
Ueda (2014), which Manzano, Karde\c{s}, Rold\'an and Wolpert (2024) carried over to computation.
The logarithms coincide for the stated counts. This does not supply a physical reverse
process: absolute irreversibility needs forward and reverse path probabilities that the
logical map alone does not specify.
Constructor theory states a cousin of the same core for devices that run in a cycle: Marletto,
Vedral and coauthors (2022) and Violaris and Marletto (2022) built an irreversible task out of
the quantum homogeniser and argued that erasure may dissipate more than Landauer's bound.
Section~\ref{sec:a2-not-arbitrary} tells all of it. The price isn't mine: Li and Vitányi priced each erased bit (1996), and Demaine and his coauthors
priced named instructions that destroy bits (2016). Mine is what a flat price of at least one on
every merging move does to a reading that stays: it turns ``the flip merges'' into ``the
flip costs'' on any machine with finitely many states reading yes, whether or not anybody's
physics is right, and gives the toll (Observation~\ref{thm:toll-derived}); what comes with that form, namely the characterisation of
the moves every merge-pricing schedule must charge (the merging moves, so the core again in
pricing words), the merge-or-revoke dichotomy for a single move, the squeeze and entropy forms
for a permanent flip, and a built counterexample for each premise
(Section~\ref{sec:a2-not-arbitrary}). The counterexample for finiteness is the endless strip, the machine that keeps every history, and
it says what the premise is for: the strip raises its reading with no merge by taking a fresh state
at every step, so a price list that charges merges and nothing else lets it certify at a cost of
zero. That's a price list that charges nothing for the fresh states, not a machine that got out of
paying, and its physical cost remains implementation-dependent. Kind: the core is classical, and somebody else's: Zermelo's in
physics, folklore in automata theory, with Pin's paper as the place to check it. The pricing
form isn't wholly mine either: priced in entropy, for automata, it's Ouldridge and Wolpert's.
They derive the price from the physics. An assumed price is older than mine too: Li and Vitányi
(1996) took a unit charge per irreversibly erased bit, with reversible steps priced at zero, as
axioms of a model of computation, and Demaine and his coauthors (2016) put that charge on named
instructions of the random-access machine. What's mine is the toll derived from a price on merging
moves for a reading that stays, its single-move and counted forms, and the counterexamples.
The physics is a named premise and is Landauer's. Searches: the core
turned up in the places just named, and its physical price in Ouldridge and Wolpert; I didn't
find the derived toll stated anywhere I looked, and nobody's search covers a
whole field, mine least of all.

\subsection{Learning borrows, committing pays}

Known: that a measurement can be made without merging anything, so that Landauer's premise
sets it no floor, and that the floor lands when the demon wipes its notes to look again, is
Bennett's resolution of Maxwell's demon (1982), with the quantitative trade-off by Sagawa and
Ueda (2009); that the entropy can be split between measuring and erasing in more than one way is
Fahn's (1996), and that the work an erasure takes depends on what the eraser already knows is del Rio
and colleagues' (2011). The closest physics I found to the line between learning and committing
is Ouldridge and ten Wolde's fundamental costs of persistent copies (2017; with Govern, 2017),
copies that stay correlated with their originals after the two come apart; theirs is a work and free-energy bound for a
correlation in chemistry, and Section~\ref{sec:narrowing} says how it differs. The sharp
chemical form, that a persistent copy grown at equilibrium carries no information about its
template, is Poulton, ten Wolde and Ouldridge's (2019). Mine: the four-state demon, where copying
a hidden bit onto a display known to be blank merges nothing and wiping the display merges two
pairs of states; the second copy that hands the display back blank while the original is still
there; the three-state ring and the two-state impossibility (Section~\ref{sec:narrowing}); and
the certainty bound (Theorem~\ref{thm:certainty-bill}), which splits the uncertainty about the start
into the part a run destroyed, paid for in units of cost, and the part the end state still holds. That split
is Shannon's chain rule (1948), and bounds of this kind on the actual spread are Maroney's (2009) and
Sagawa's (2014); mine is reading it on the ledger, about what an observer saw. Kind:
finite models and one bookkeeping identity illustrating a known physical point inside my ledger,
and not a new physical claim.
Bennett did the hard part in 1982. My four dots put the line where you can work it through
yourself: the copy owes no floor and can cost nothing, but it still takes something. Its physical cost depends on the dynamics and protocol. And it uses up a blank
page, which is a fresh state; while the original is still there it can hand the page back, so learning
borrows. Resetting the uncertain page without using its correlations has the
stated floor. A supplied blank alone does not prove a previous payment. Committing makes the record permanent, so there's nothing to hand back, and with
finitely many states it has to merge.

\subsection{The counter as a receipt}

Known: one-way counters that give order and termination are old: Turing's 1949 note on
checking a large routine, Floyd's ranking functions (1967), Lamport's clocks (1978).
Carrying a run's cost beside its result is the writer cost of Danner, Licata and Ramyaa (2015)
and of Danner and Licata (2022), and the
graded effect semantics of Katsumata (2014), with a working-language version in Orchard,
Liepelt and Eades (2019). The receipt side of time credits is Mével, Jourdan and Pottier's
(2019): in their extension of the Iris logic every step uses up a credit and
produces a receipt, receipts prove lower bounds, and enough of them prove that a counter can't
overflow in any feasible time, which is the closest published neighbour of a total that only
climbs. Mine: putting three properties side by side (the total only
climbs, it equals the sum of the charges on every path including the failures, and it is the only such total, on the states a run from the start reaches, for a given price list), the argument for any system with a counter
and a price list, and the sum proved from the small machine's own step (Section~\ref{sec:counter},
Section~\ref{sec:smallest}). Kind: synthesis; the abstract statements are proved for any system
with a counter and a price list, and on the small machine the sum is proved too. Against the graded
frameworks the answer is split: the counter is their writer exactly, and their type systems, bounding relations and
before-the-run checking of grades are not embedded, because the grade here is written by the step
during the run (Section~\ref{sec:counter}). For the well-graded, loop-free programs defined in Section~\ref{sec:counter}, on runs
that do not trap, the counter's rise equals the fixed grade. The loop built there costs
its input value, so no fixed whole number is its exact grade on every input; an
input-dependent grade is exact. The arguments are given
there. Cost-aware type theories such as calf allow grades that depend on
the inputs.

\subsection{Windows that can't price what they hide}

Known: a function of an observation is constant on the observation's fibres; that's
elementary, and the systems version of it is Moore's distinguishing experiments (1956),
Kalman's observability (1960), and the completeness question of abstract interpretation
(Cousot and Cousot, 1977; Giacobazzi, Ranzato and Scozzari, 2000); in control theory it is
diagnosability and opacity of discrete-event systems (Sampath and coauthors, 1995; Saboori and
Hadjicostis, 2007), where a collision is a hidden event no observer can diagnose without delay;
in economics it is the
hard-evidence view of mechanism design (Green and Laffont, 1986; Bull and Watson, 2007), where
what a mechanism can tell apart is limited by the evidence on offer; and its sharp form for
fingerprints is Naor and Rothblum's lower bound for checking a stored file (2005); in distributed systems it
is Haeberlen and Kuznetsov's (2009), that only omission and commission faults can be detected from outside,
and in philosophy it is Easwaran's transferability (2009). Mine: stating the pricing
version (no window price is exact across a collision), exhibiting a colliding pair for the
small machine's window, and the theorem that every certified Thiele-complete machine has two runs from one
clean start that end looking the same through the window its base shows, one certified and one
not (Section~\ref{sec:windows}, Section~\ref{sec:smallest}). Kind:
an elementary observation about what a window hides, plus specific facts about the machine and the definition.

\subsection{Structural entitlement and its bound}

Known: decision-tree counting bounds (a tree that tells $N$ cases apart has depth at least
$\log_2 N$); certifying algorithms that hand over checkable certificates (McConnell,
Mehlhorn, Näher and Schweitzer, 2011); proof-carrying code (Necula, 1997); the gap between
justified true belief and knowledge (Gettier, 1963). Mine: the definition of entitlement as an
object, a checked relation, and a paid event on the record, and the count that, given lists, a
tree and a covering, a run that paid for the tree's depth pays at least the index bits it
commits to (Section~\ref{sec:entitlement}). Kind: synthesis; the counting idea is standard, the
assembly and its vocabulary are this work's, the pricing policy is assumed (nothing derives it), and the
count is proved in its cost form over any Thiele machine, and for the record moves of
every certified Thiele-complete machine on a run from a clean start, with the earned chain in the run, a floor that takes the run's own
checks in place of a tree, and a counterexample showing the covering can't be dropped
(Section~\ref{sec:part3-on-machine}); and in a form where the run makes the posterior from a
prior, every drop is made by a check the run passed, and the index bits are at most $k$ times the
cost when each passing check keeps a $2^{-k}$ share, a hypothesis the same counterexample shows
is needed (Theorems~\ref{thm:ent-run} and~\ref{thm:ent-run-machine}).

\subsection{The record axis is a latch}

Known: monotone, write-once state is old: the latch (Eccles and Jordan, patent filed 1918, paper 1919), sticky bits
(Plotkin, 1989), grow-only state in replicated data types (Shapiro and coauthors), history
variables and stuttering in refinement (Abadi and Lamport, 1991), append-only authenticated sets
for transparency logs, which fit revocation as a once-only event (Tomescu and coauthors, 2019),
and, in program logic, monotonic state, where a property witnessed once on state that only moves
forward can be recalled at any later point, with a soundness proof and a working home in
F$^\star$ (Ahman, Fournet, Hri\c{t}cu, Maillard, Rastogi and Swamy, 2018): the type-theory name
for a record that stays up. Past the finite case, inductive and quasi-inductive definitions
(Richter and Aczel, 1974; Burgess, 1986) and the flag trick of infinite time Turing machines
(Hamkins and Lewis, 2000) are the transfinite versions of a record that only grows and a
toggle, and Danos and Krivine's (2004) reversible calculus needs an action declared
irreversible in order to commit.
Mine: the exact statement
that a record driven by the computation that never switches off is, over any base, a latch
on one event, with its converse and a counterexample for each of the two conditions, and the
argument that base-driven records agree up to the price list while the stronger uniqueness fails
(Section~\ref{sec:latch}). Kind: a small structural lemma, I think not stated before in this
form (I found no prior statement in the searches I ran), built from old notions. The proofs
show why the record must latch and why its price list remains a separate choice.

\subsection{The diagonal and L}

Known: Kleene's recursion theorem and Rice's theorem; Lawvere's fixed-point form (1969) and
Yanofsky's retelling (2003); later detailed developments of recursion theorems and Rice
(Norrish, 2011; Carneiro; Forster, Kunze, Smolka and coauthors, including L itself, 2017, and
their library of undecidability results). Mine: the diagonal stated once over any substrate that
supplies a recursion theorem; and Rice and halting proved again for L from its reduction rules
(Section~\ref{sec:diagonal}); and Rice's and Kleene's theorems proved for the small machine's
host, with Rice's on every input of the two-counter machine itself, Kleene's there with inputs
written as powers of two, and Kleene's refuted for plain inputs (Appendix~\ref{sec:small-diagonal}).
Kind: new proofs of known theorems of computability, for particular machines. The theorems hold
for any model with universality and s-m-n, and the work is in building those for each machine;
the refutation on plain inputs is a small fact about two counters, resting on Schroeppel's.

\subsection{An example machine that meets the definition}

Known: two-counter machines and the undecidability of their halting are Minsky's (1961, 1967),
developed further by Forster, Larchey-Wendling and coauthors (2019, 2020). The guarantee about every program is older than any of this, and it
isn't mine either. It's the reference monitor (Anderson, 1972): a check that sits in the path of
every operation, so no program gets around it. Schneider (2000) worked out what a monitor that
watches a run can enforce: only safety properties, the kind a bad run breaks at some finite
moment, and ``no CERTIFY without an earlier earned COMMIT'' is one of those. Azevedo de Amorim
and coauthors (2015) verified monitors that check tags at every instruction. The shape
of ``earned,'' a certificate nothing can make except through the rules, is Milner's LCF, where a
theorem is a value only the logic's rules can build (Gordon, Milner and Wadsworth, 1979; Gordon,
2000); the small machine is that architecture at the level of single instructions. ``No certify
without an earlier commit'' is also a correspondence assertion in the sense of Woo and Lam
(1993), which Gordon and Jeffrey (2003) check with a type system, and a store of commitments
that the rules of a dialogue keep, each speaker's updated by the rules as the moves are made, is
Hamblin's (1970, chapter 8, on formal dialectic); the difference from all three is
again that here the step enforces it. A check that
blesses a value and a write that takes the blessing away is close kin to dynamic taint tracking,
whose labels go back to Denning's lattice model of information flow (1976), with the blessing
itself, controlled endorsement, given a type system and a security theorem by Cecchetti, Myers
and Arden (2017), and what a monitor
can enforce with bounded memory is Talhi, Tawbi and Debbabi's question (2008). A small monotone
record inside each replica of a replicated service is A2M (Chun and coauthors, 2007), TrInc
(Levin and coauthors, 2009), MinBFT (Veronese and coauthors, 2013) and Damysus (Decouchant and
coauthors, 2022), and Clement, Junqueira, Kate and Rodrigues (2012) worked out, with proofs, what such a record
buys and what it doesn't. The signature route to the same goal, which needs no trusted step,
is IA-CCF (Shamis and coauthors, 2021): replicas sign and log their protocol messages in the
ledger, so blame lands on individual members however many of them misbehave. And ``check once, issue a capability the system trusts, rely on it
later'' is built and deployed in Garg and Pfenning's proof-carrying file system (2010), whose
logic lets a policy's consequences change with time and state, and in Sirer and coauthors'
logical attestation for the Nexus operating system (2011), unforgeable statements about programs
used as credentials; those are the closest systems I've found to check, commit, certify. Mine: one worked example, a
machine of six instruction kinds on two counters, proved to meet the definition so that the definition is not empty, with a version
for each counter, a bounded fact table and a channel; its theorems that, from a clean start, a committed fact was checked
first about the same edition, that nothing is forged, that nothing is evicted,
that with the flag up from a clean start the channel names a fact an earlier passing COMMIT committed, and a later commit can replace it, as a two-commit run shows, and that a run from a clean start to a raised flag adds at least three to
the ledger, with three reached by a run whose check could have failed (the step doesn't
police trivial checks, so one that can't fail costs three too); the same theorems for any
property language with an exact equality
test whose checker is proved to say yes exactly when the claim holds, a sorted list read out of a
counter included; and the step-for-step match with
the two-counter machine, so halting is undecidable for it (Section~\ref{sec:smallest}).
Kind: an instance of monitor enforcement with two things on top
that I didn't find put together where I looked: the check is bound to the edition of the thing
it checked, and a toll is paid for the commitment. The first of those alone is old: checking at
commit time that nothing you read has changed since is optimistic concurrency control (Kung and
Robinson, 1981), and each counter's version here is that scheme's version number. What I didn't find is that
check put together with a toll. The universality and the
undecidability are borrowed, and the other properties are short consequences of the step
rule. Whether the design is new I can't say, and I make no claim beyond the proofs.

\subsection{Thiele-complete and certified Thiele-complete, defined so a clock fails}

Known: Turing completeness, and the observation every theorist makes in four seconds, that a
toll alone is cheap: a machine that charges for every step pays any toll. Enforcement mechanisms
have been compared and sorted before, and they aren't one notion seen several ways: James Morris
(1973) separated access limitation from authentication in programming languages, and Kevin Hamlen, Greg
Morrisett and Fred Schneider (2006)
classified which policies execution monitors, program
rewriters and static analysis can each
enforce, and they differ; Wai Wong (1995), on a neighbouring track, recorded
and checked HOL proofs so a theorem could carry its derivation. Mine: a definition of
Thiele-complete with four clauses that say nothing about what the record is about (computing that runs every program, costs nothing and can't touch the record, while every other move costs one; a record that stays; blank starts;
and some run whose record tells two starts apart), with the proofs that a clock and a silent
machine fail it, that each clause is needed, that every universal base becomes one by adding
marks, and that every machine meeting it hides its record and ledger from its own window
(Section~\ref{sec:complete}); and a definition of
certified Thiele-complete with four clauses (a universal base that can't touch the record and a record that stays up once it's
up, an earned
record, exact prices, nothing for the base and one for each of the three record acts, and some claim whose check can fail), the proofs that every
clock-style machine and a paid latch and a silent machine fail it while the small machine, the
same machine over any property language with an exact equality test, a checker proved to say yes exactly when the property holds and a property that holds of one number and fails of another, and the universal program's host meet it,
that the four clauses are independent, each failed by some machine with a two-counter machine
inside that meets the other three (Theorem~\ref{thm:tc-independent}),
and the theorem that every machine meeting it hides its record and ledger from its own window
(Section~\ref{sec:smallest}), with the theorem that it is exactly Thiele-complete with clean
starts, the earned chain and a claim that could have failed added
(Theorem~\ref{thm:tc-decomposition}). Kind: definitions with short consequences. The window theorem is
the fibre fact of item 6 plus one observation, that a universal base prints any window at no cost on the ledger.
I'm not claiming the definition is the only reasonable one, or that LCF kernels, reference
monitors and proof-carrying code are all certified Thiele-complete machines in disguise; Hamlen, Morrisett
and Schneider's sorting says such mechanisms genuinely differ, and the book doesn't have a
theorem showing the definition holds up across such differences.

\subsection{Five record disciplines, one question}

Known: each system and its guarantee belongs to its designers: Casper (Buterin and Griffith,
2017), Ethereum's storage gas and refunds (EIP-2200, EIP-3529), extend-only TPM registers,
Certificate Transparency (RFC 9162), proof-carrying code. Mine: asking all five whether the
write is priced, with a small worked fragment of each (Section~\ref{sec:models}). Kind:
synthesis and an organizing reading; the fragments are not verifications of those systems.

\subsection{The pointer thesis}

Known: Zurek's quantum Darwinism (2009), with Ollivier, Poulin and Zurek (2004) and
Blume-Kohout and Zurek (2006): preferred states are the ones whose records proliferate
redundantly in the environment. Mine: borrowing that shape for the question of which event
deserves a record, with definitions (copied and relied on by parties who didn't produce it, independent
carriers), two toys, and refutations of two stronger forms
(Section~\ref{sec:pointer}). Kind: thesis. No claim about quantum physics.

\subsection{The universal host for the example}

Known: Turing's universal machine (1936), which runs any machine given as data, and the
universal register machines that run two-counter programs. Keeping a guest and the code running
beside it out of each other's state is what a virtual machine monitor promises (Popek and
Goldberg, 1974). Compiling computable functions into counter machines is Forster, Larchey-Wendling and coauthors' work. Mine: one fixed program, $U$, written in the small machine's own six kinds of
instruction on a host with a register for every whole number, with no guest step built in, that
runs every program of the small machine from every clean start $\mathrm{start}(x, y)$, halts exactly when its guest does, raises
its own flag by its own check, commit and certificate if and only if the guest's rises somewhere in its run, and
keeps a ledger equal to the guest's at every matching point; its host machine is certified Thiele-complete,
and whether $U$ halts on a loaded guest is undecidable (Appendix~\ref{sec:universal-host}). Two
small facts force its shape, and they're stated there: reading a counter above 0 into the program's control flow on this machine writes it, and a finite program can't keep the guest's facts on a copy of the guest's counter, by the
pigeonhole. Beside it, a simpler host with a built-in guest step runs guests by construction and
keeps a mirror of the guest's flag, with the mirror, isolation and earned-chain theorems.
Code that shares unprotected memory with a record can overwrite it, and ordinary computers
isolate memory and enforce rules with monitors too, so what I claim is narrower: the simulation
and the record guarantee proved for this one fixed program $U$. Kind: an implementation, with its
proofs, of a design whose shape is Turing's universal
machine, with the record and the toll carried through. The universality is borrowed from the
small machine and from Minsky. And the same program with one paid move added, $U_P$, runs every
computably presented Thiele machine, with its halting matched, its reading carried as a latch on
the host's own earned flag whose claim is that reading, and its ledger to within a surcharge of
at most two when the reading starts at no (exactly three when it starts at yes), which no program on the host, run from a clean start, can avoid when the guest certifies with
its whole ledger below three (Appendix~\ref{sec:presented-universal});
the counter-compilation construction is given in Appendix~\ref{sec:counter-constructions}. Kind for that part:
an implementation with its proofs again. That a universal machine runs any computable machine is
Turing's; what's carried through here is the record, earned again, and the prices, exact up to
the surcharge, and I didn't find that put together where I looked. In the words of
Definition~\ref{def:universal-thiele}, $U$ is a universal Thiele machine for the small machine's
programs with no surcharge, and $U_P$ is one for every computably presented Thiele machine with a
surcharge of at most two when the guest's reading starts at no and three when it starts at yes.

\subsection{The game check}

Known: that positive semidefiniteness of a correlation matrix gives the Tsirelson bound is
Tsirelson's (1980) and Stephanie Wehner's semidefinite proof (2006), from a paper on Tsirelson bounds for generalized CHSH inequalities; the moment-matrix hierarchy, the
Schur-complement analysis of the level-one matrix of this scenario and the level $1+AB$ are Navascués, Pironio
and Acín's (2007, 2008), with both four-body cells unconstrained where my matrix ties them; that every point with a PSD completion is the correlator table of a quantum strategy, built from anticommuting observables in a maximally entangled state, is Tsirelson's too (1980); the arcsine characterization is Landau's (1988) and Masanes's (2003); Echenim and Mhalla (2023) give a further detailed derivation of the CHSH and Tsirelson bounds. Rational sum-of-squares certificates are Helfried Peyrl and Pablo A. Parrilo's
(2008), from a paper on computing sum-of-squares decompositions with rational coefficients. Mine: a decision procedure in integer arithmetic on outcome counts; a proved
chain from ``the counts pass'' to ``the pinned level-one matrix is positive semidefinite'' to
``$S^2 \le 8$''; explicit constructions of the representation in finite dimension, on two four-level systems with real amplitudes and two two-level ones with complex, and a proof of the arcsine characterization; proofs of the contraction equivalence by explicit test vectors; and the
level-$1+AB$ cascade (on the slice $\langle \{A_0, A_1\} \otimes \{B_0, B_1\} \rangle = 0$) and an elliptope gate in whole numbers, with a soundness proof
(Appendix~\ref{sec:chsh} and its derivation, Appendix~\ref{sec:chsh-derivation}). Kind: an integer procedure connecting
counts to these known bounds, with known theorems (Tsirelson's representation, the arcsine
characterization, the equivalence of contraction and PSD) proved again. I found no earlier chain
that takes counts as input where I looked. The same check earned on the small machine's host, a raised flag standing on
it, is built and proved (Theorems~\ref{thm:chsh-exact} and~\ref{thm:chsh-earned}); the idea isn't new.

\subsection{What that adds up to}

Read straight down the list. The mathematics underneath this
book is mostly known, and some of it is very old. The question organizing this book is what it costs for a distinction to come into being and
stay, written into the step of a model of computation, with certification, a record that a claim
was checked, as its sharpest case, and what's mine in the answer isn't one result. It's four, and I name them the same way wherever they
come up. The toll derived from finiteness, permanence and a price on merges (item 3; the
counting is classical, the finite cousin of Zermelo's recurrence objection and folklore about
permutation automata, and Ouldridge and Wolpert already priced it in entropy for automata, so
what's mine is the toll a price on merging moves gives for a reading that stays, after Li and
Vitányi's price per erased bit and Demaine and his coauthors' price on instructions, and the price is
an abstract premise motivated by Landauer's reset). The pricing form of the window fact (item 6; the fact itself is elementary). The
latch form of a record that never comes down (item 8; latches and history variables are old).
And a record that's priced, and on the small machine earned, under the reference monitor's
every-program move (items 10, 11 and 14; the move is Anderson's, Schneider's and the micro-policy
authors', and what's added is the check bound to its edition, the toll on top, a definition of
Thiele-complete that a clock fails, over any base and any event, a certified refinement the small
machine meets, and universal programs that carry the check and the toll to every
small-machine program and to every computably presented Thiele machine, earning their own flag).
Under all four sits item 1, which is old mathematics read one way: every step of every machine run by one rule takes a fresh state, merges, or goes round a loop, a merge in a coarse description is, in every
finer description it is exact about, a fine merge, a fine step out of the set, or endlessly many fresh states, specified thermal models add their own dissipation bounds, and
finiteness is what takes the fresh states away.
The first is a ledger
form of a core that's classical and already priced in physics; the second and third are small exact statements I didn't find
stated where I looked; the fourth is built, with its proofs on the page, and its idea isn't new. The general statements are proved for every machine that meets the definition; the small machine and the programs that run on it are the worked instance that shows the definition can be met. Around it is my best
attempt at saying which is which,
which you should check as hard as anything else in here. If someone shows me
item 3's ledger form, the pricing form in item 6, or item 8
stated earlier, I'll move it up the list.


%===========================================================================
\section{From ``The universal machines''}
%===========================================================================
\subsection{A simpler host, with the guest step built in}
\label{sec:simpler-host}

There's a much simpler host than $U$, and I keep it in the book because it's easy to read and
it shows the copy rule in one line. Its universality is by construction: it has a move that runs
one guest step by calling the small machine's own step. It isn't a universal Thiele machine in the
sense of Definition~\ref{def:universal-thiele}, and the reasons are exact. Its interpreter loop
$U_0$ never halts: it keeps returning to its first line even after the guest stops.
The host itself does have an OWN HALT instruction, which another program can use.
The loop also keeps charging one per pass, even when the guest has stopped or
its next step costs nothing, so no fixed surcharge bounds it. And it isn't a register host, since its guest step
does the computing. A host that fixed the first two and kept the built-in step would still fail the
last (the argument after the definition). $U$ is the universal Thiele machine of this book. Read it as a picture of what $U$ does the long way, and as a host that
starts from any guest state at all, clean or not.

It carries a guest inside its own state: the guest's program as data, the guest's counters, its
table, its ledger, its flag. The host has one extra move, run one guest step. In that one move it
looks at the guest's flag, takes the step, looks again, and if the flag went from down to up, the
host's own copy of that flag goes up in the same move, and the host pays for it.

\begin{definition}[The host]\label{def:u34-host}
A \emph{guest} is a program $P$ of the small machine together with a state $g$ of it
(Definition~\ref{def:u34-state}). A \emph{host state} $h$ has six parts: its own core
$\mathrm{core}(h)$, its ledger $\mu(h)$, its own flag $\mathrm{own}(h)$, the guest program
$P(h)$ held as a list of instructions, the guest state $g(h)$, and the \emph{mirror} $m(h)$, a
yes/no copy of the guest's flag. The host has two kinds of move. $\mathrm{OWN}\ i$ runs the
instruction $i$ on the host's own core, costs what the small machine charges for $i$,
and raises $\mathrm{own}$ exactly when that machine's flag would rise. $\mathrm{GSTEP}\ b$
costs $b$. If the host's core is trapped it does nothing else. Otherwise it moves the host to
its next line and, if the guest's next instruction costs at most $b$, takes one program step of the guest, which leaves a stopped guest where it is; in the same move the mirror becomes $m(h) \lor (\lnot\mathrm{cert}(g(h)) \land
\mathrm{cert}(g'))$, where $g'$ is the new guest state. The host's reading is
$\mathrm{own}(h) \lor m(h)$. Its stored program is $U_0 = [\mathrm{GSTEP}\ 1;\ \mathrm{OWN}\
(\text{INC } A);\ \mathrm{OWN}\ (\text{DEC } A\ 1)]$: one guest step, then two moves that bring the host back to its first line. It \emph{loads} a program $P$ from $(x, y)$ with $g(h) = \mathrm{start}(x, y)$ and the mirror and its own flag down. \end{definition}

Every instruction of the small machine costs at most one, so $\mathrm{GSTEP}\ 1$ runs every
guest instruction there is. Why charge the budget $b$? Why not the cost of whatever
instruction got fetched? Because charging the budget keeps the price a function of the move
alone, which is what the definition of Section~\ref{sec:toll} asks for: you read the price of a
host move off the move, the way you read the small machine's prices off its table, without
peeking inside the guest. What that costs the host is easy to see. A host that steps a guest
instruction priced 0, an INC say, still pays one for the move.

So the host's ledger is closer to a clock than to a receipt for its guests: one unit per guest
step, whatever the step was. It never falls behind the guest's own ledger (part (3) of
Theorem~\ref{thm:u34-host-toll} below), and it matches that ledger only when every guest step
happens to cost one. Could a cleverer host have simulation-exact accounting, a ledger equal to
its guest's? Not with this one fixed-price move: $\mathrm{GSTEP}\ 1$ is the same
move whether the guest's next instruction is an INC priced 0 or a paid CHECK, so one
price can't match both. Two moves can: one that costs nothing, takes the guest's step when that step costs nothing and otherwise jumps, and one that costs one and takes it. A loop that tries the first before the second keeps the host's ledger equal to the guest's at every pass. That host still has its guest step built in. $U$ gets the match with
ordinary instructions: it pays only on the record moves it runs itself, so its
ledger equals its guest's (Theorem~\ref{thm:u34-U-flag}).

And there's one thing to say about the definition before the theorems, so you know which of them carry
weight. $\mathrm{GSTEP}$ runs the guest's instruction by calling the small machine's own step on
the guest's state: the guest's step is a primitive of the host. So ``the host runs any guest''
(Theorem~\ref{thm:u34-universal}) is true by construction. The host was built to do exactly
that, and the proof follows the three moves around the loop. The theorems with something in them come after it.
From a loaded host, or any state where the two agree, the mirror always equals the guest's flag (Observation~\ref{thm:u34-mirror}). Nothing raises the
mirror unless the guest itself passes a CERTIFY inside a guest step the host paid for
(Theorem~\ref{thm:u34-no-free-host}). No $\mathrm{OWN}$ move, in any host program at all,
touches the guest's state or the mirror (the second sentence of that same theorem), and no host
move of either kind rewrites the guest's program. And
the mirror is tied to the guest's earned chain: after any host run from a loaded
host, if the mirror is up, the guest's own run from its start contains a
passing CHECK, then a COMMIT of the same claim, then a CERTIFY, since the mirror is the guest's
flag and the guest started clean (Theorem~\ref{thm:u34-cert-earned}). So this
host runs any small-machine program by construction and keeps the guest's record by theorem, as a copy of the
guest's flag. The mirroring program $U_0$ does not earn the host's own flag; a different program
could do so with the host's OWN CHECK, COMMIT and CERTIFY. $U$ does both the long way: it runs the guest
with its own instructions, and it earns its own flag.

\begin{theorem}[The host runs any guest]\label{thm:u34-universal}
(1) If the host's core is not trapped, then after $n$ moves $\mathrm{GSTEP}\ 1$ the guest state
is the guest program's own run of $n$ steps from $g(h)$, the guest program is unchanged, and
the host is still not trapped. (2) If the host's core is at line 1 and not trapped, then $3n$
steps of the program $U_0$ leave the guest exactly where $n$ steps of its own program put it.
(3) For every two-counter program $M$ and inputs $a$, $b$: $M$ halts from $(1, a, b)$ exactly
when the host program made of $\mathrm{OWN}$ moves of the compiled $M$ halts, from a host
whose core starts at $(a, b)$, whatever guest that host carries.
\end{theorem}

\begin{observation}[The copy is the record]\label{thm:u34-mirror}
Say a host state $h$ \emph{agrees} when $m(h) = \mathrm{cert}(g(h))$. Every host move takes an
agreeing state to an agreeing state, so every host trace and every stored host program
does too, and a freshly loaded host agrees. In particular, if $h$ agrees and is not trapped,
the mirror after $n$ moves $\mathrm{GSTEP}\ 1$ is the guest's flag after $n$ steps of its own
program.
\end{observation}

\begin{theorem}[The host's step collects the guest's toll]\label{thm:u34-host-toll}
(1) If $h$ agrees, is not trapped, and the guest step that budget $b$ allows takes the guest's
flag from down to up, then the mirror is down before $\mathrm{GSTEP}\ b$ and up after it,
the host's ledger rose by exactly $b$, and $b \ge 1$. (2) In the run of (1) of
Theorem~\ref{thm:u34-universal}, from a host that agrees, if the guest's flag goes up between its steps $n$ and $n + 1$,
the mirror goes up between the host's moves $n$ and $n + 1$, and that move costs 1. (3) For every
host trace $t$, $\mu(g(\mathrm{run}(t, h))) + \mu(h) \le \mu(\mathrm{run}(t, h)) +
\mu(g(h))$: the host's ledger grows at least as fast as the guest's.
\end{theorem}

\begin{theorem}[No unpaid host certification]\label{thm:u34-no-free-host}
If the mirror is down at $h$ and up after a host move $i$, then $i$ is $\mathrm{GSTEP}\ b$ with
$b \ge 1$, the host was not trapped, the guest's next instruction was CERTIFY, and the guest's
flag went from down to up in that move. A trace made only of $\mathrm{OWN}$ moves leaves the
mirror and the guest state exactly where they were. And the host's own flag rises only by an
$\mathrm{OWN}$ CERTIFY that passes on the host's core.
\end{theorem}

\begin{theorem}[The host is a Thiele machine]\label{thm:u34-host-instance}
For every host state and move, if the host's reading goes from no to yes, the move costs at
least 1; the same holds for the mirror alone and for the own flag alone. So the host read
through $\mathrm{own} \lor m$, the host read through $m$ alone, and the guest class are
Thiele machines (Definition~\ref{def:certsys}), and any host trace that raises the
host's reading costs at least 1.
\end{theorem}

Two statements have been proved here. The designated guest-step runs follow the
guest computation. Separately, every permitted host instruction preserves agreement
with a genuine guest run, so arbitrary host instruction sequences from a loaded
start cannot fabricate the guest's mirror. The second statement uses the full host
instruction set. It does not follow merely from the first, or from charging for a
raise. The combined reading own-or-mirror remains a different observation: the host
may establish its own record while the guest's record stays down.

\begin{remark}[Two guests, run]\label{rem:u34-host-runs}
Load the guest CHECK \emph{is zero} $A$; COMMIT \emph{is zero} $A$; CERTIFY; HALT with both
counters at 0. After two moves $\mathrm{GSTEP}\ 1$ the mirror is down; after three it is up,
the guest's flag is up, the host's own flag is still down, and the host has paid 3. Under the
program $U_0$ the mirror is down at host step 6 and up at step 7, with 3 paid. The guest
CERTIFY; CERTIFY traps on its first step: after five moves $\mathrm{GSTEP}\ 1$ the mirror is
down, the guest is trapped, and the host has paid 5 for nothing.
\end{remark}

That copy can't be faked. The proof walks every host move there is, and the copy only rises
in a guest step where the guest itself ran CERTIFY and passed. A host program that never
touches the guest leaves the copy exactly where it was. From a loaded host the copy always equals
the guest's
flag, and the host's ledger always grows at least as fast as the guest's, and a three-instruction
loop runs any small-machine program one guest step per pass. Here the mirror, the crossing test
and the charge sit in the host's step rule, so Theorems~\ref{thm:u34-host-toll}
and~\ref{thm:u34-no-free-host} hold for every sequence of host moves, including ones that never
run a guest. The fixed program $U$ supplies simulation, an exact ledger at
matching points and an earned flag using ordinary instructions. Its guest
simulation theorem concerns that program from its loaded starts. The host's
native earned-chain rule covers other programs too, but does not make every
such program simulate the same guest.

\subsection{The diagonal, on this machine's programs}
\label{sec:small-diagonal}

Section~\ref{sec:diagonal} left the diagonal standing on any substrate, on a little family built to fit
it, and on L, handled on its own terms rather than as a substrate. It also said what a machine has to
supply to fill in the same shape, two things: programs that can be handed to programs as numbers, and a
recursion theorem for the transformations it can carry out. This machine's universal program reads a whole guest program out
of one register, with a numbering of programs that goes there and back
(Definition~\ref{def:u34-U}), and halting is undecidable here in the reduction sense of Theorem~\ref{thm:u34-undecidable}. What's left are the two classical results that turn a
universal machine into a place where the diagonal bites: Rice's theorem and Kleene's recursion
theorem. They come with a wrinkle I should put first, because it decides which machine they're
about.

The wrinkle is the one from Section~\ref{sec:smallest}, at the end of the subsection on what two counters compute. The small machine takes its inputs straight
into two counters, and Schroeppel's memo gives a proof that two plain counters can't even compute $2^n$ from $n$. Rice's
reduction has to keep a program's input safe while a copy of a halting instance runs. Kleene's
construction has to run an evaluator while keeping its input. Two counters have nowhere to keep it.
So the two theorems land where they're true and where I could prove them. The host machine of
Definition~\ref{def:u34-multi}, the one with a counter for every number that $U$ runs on, is a fine
programming system: its programs are numbers, an input sits in register 1, and there's room to
spare. On it both theorems hold in full. On the two-counter machine itself the story splits, and
I'll tell it at the end of the subsection: Rice's theorem holds on every input, by a trick for
keeping the input that I'll show you; Kleene's holds when inputs are written as powers of two, the
way Minsky packed them; and on plain inputs Kleene's is false, and that's a theorem too.

The host machine's programs are numbered in the same style as $U$ numbers a guest's. Write $\langle m, n \rangle$
for $(2n + 1) \cdot 2^m$. INC $r$ is $\langle 0, r \rangle$, DEC $r$ $j$ is
$\langle 1, \langle r, j \rangle \rangle$, HALT is $\langle 2, 0 \rangle$, CHECK $r$ is
$\langle 3, r \rangle$, COMMIT $r$ is $\langle 4, r \rangle$, CERTIFY is $\langle 5, 0 \rangle$, and a
program is the list code: the empty program is 0, and an instruction $x$ followed by a program $t$ is
$\langle x, \mathrm{code}(t) \rangle$. An instruction number that is none of these codes reads as
HALT. Then every number decodes to some program, since a positive number is
$\langle x, n \rangle$ with $n$ smaller than it, and decoding a program's code gives the program
back. When I put one program after another, the later one's jumps move down by the length of what
comes before it, and a jump to a line where it has no instruction, line 0 included, goes to the line
just after it, so a stop inside it hands on to whatever follows. Say a host program \emph{outputs} $y$ on input $x$ when, started with $x$
in register 1 and every other register 0, with a clean table, it halts with $y$ in register 0.

Rice comes first, because it's the shorter road. Henry Rice, whom you met in
Section~\ref{sec:windows}, proved in his 1951 doctoral thesis, and published in 1953, that no program decides any
property of what programs \emph{do} unless every program has
it or none does. Properties of how a program is written aren't covered. The road from halting is a single detour. Take a property of what programs do,
a program $y$ that has it, and a program $n$ that doesn't. Given an instance of two-counter halting,
a Minsky program $H$ with starting counters $a$ and $b$, build a host program that first loads $a$
and $b$ into two fresh registers (not register 1, and named by neither $y$ nor $n$), then runs a copy of $H$ moved onto those registers, then empties
both, then runs $y$. If $H$ halts on $(a, b)$ the built program behaves like $y$ and has the
property. If $H$ runs forever the built program never stops on any input, and it behaves like the
program that loops, which lacks the property (if the looping program has it instead, swap the roles
of $y$ and $n$). Everything the built program does before $y$ is INC and DEC, which add nothing
to the ledger and touch no fact, channel, trap latch or flag. So a decider for the property would be a decider for
halting, and by Theorem~\ref{thm:u34-undecidable} that would let a program list every single-tape
Turing machine that doesn't halt. Two words carry the weight, and they're the
ones to check in the statement: what ``behaves the same'' sees, and where the program starts.

\begin{definition}[Behaving the same, on the host]\label{def:host-same}
A host program \emph{ends} in a state when it halts there, from the start with $x$ in register 1 and
the rest clean; a trap counts as stopping. Two host programs \emph{behave the same} when for every
input $x$, any state one ends in has a state the other ends in that agrees with it on every
register's value, the fact table, the channel, the trap latch, the ledger and the flag, and the same
the other way round. The line number and the versions aren't compared. A property of host programs
is \emph{extensional} when behaving the same carries it from one program to the other.
\end{definition}

\begin{theorem}[Rice's theorem on the host machine]\label{thm:host-rice}
Let $\Pi$ be an extensional property of host programs that some program has and some program lacks.
Then $\Pi$ is undecidable in the reduction sense of Theorem~\ref{thm:u34-undecidable}.
\end{theorem}

A property that depends only on the partial function a program computes, whether it stops
and what it leaves in register 0, is extensional, so it is covered. That case also follows
from the recursion theorem below.

On the small machine itself the same road works from the clean start, and only from there. Say two
two-counter programs \emph{behave the same from the clean start} when, whatever state one halts in
from counters $(0, 0)$, the other halts in a state with the same counters, trap latch, ledger and
flag, and the same the other way round. The fact table and the channel aren't compared, because the
prefix moves the versions of $A$ and $B$ by fixed amounts and the facts, the channel's included, carry versions. There is no input to keep,
since the program starts at $(0, 0)$, and $H$'s instance is loaded straight into $A$ and $B$.

\begin{theorem}[Rice's theorem for the small machine, from the clean start]\label{thm:guest-rice}
Let $\Pi$ be a property of two-counter programs that's carried from a program to any program that
behaves the same from the clean start, held by some program and lacked by some other. Then $\Pi$ is
undecidable in the reduction sense of Theorem~\ref{thm:u34-undecidable}. In particular ``halts from the clean start'' is undecidable.
\end{theorem}

That theorem covers the clean start only. Every input is harder, because the standard reduction has to keep the
input safe while the halting instance runs, and two counters have nowhere obvious to keep it. The
way through is at the end of the subsection.

Kleene's recursion theorem is the deeper one, and it's a construction, which is
Section~\ref{sec:diagonal}'s point about Kleene's little gem (Moschovakis's phrase): the content
is the building, and the contradiction in Section~\ref{sec:diagonal} was just the last line. It
says: for every transformation $F$ of programs that some program can carry out on numbers, there's a
program $p$ that computes the same thing as $F(p)$. The shape of the building is the old one. From
the program that carries out $F$, build an \emph{evaluator} that reads a number, treats it as a
program's code, specialises that program to its own code, runs $F$'s program on the result, and runs
whatever program comes back. The specialising step is what logicians call s-m-n: fix one input of a
program and get a new program that needs one input fewer. Then the witness is the evaluator
specialised to its own code. Run on any input, it builds its own code, hands it to $F$'s program,
gets the code of $F(p)$, and runs that.

Each piece on the host machine comes next, with its exact premise.

\begin{theorem}[s-m-n, and a universal program for host programs]\label{thm:host-smn}
(1) For a number $c$ and a host program $V$, let
$\mathrm{spec}(c, V)$ be $c$ copies of INC 2 followed by $V$ with its jumps moved to follow them.
Then for all $x, y$, $\mathrm{spec}(c, V)$ outputs $y$ on input $x$ if and only if $V$, started with
$x$ in register 1 and $c$ in register 2, halts with $y$ in register 0. (2) There is a host program
$W$ made only of INC, DEC and HALT such that, for all $c, x, y$, $W$ started with $x$ in register 1
and $c$ in register 2 halts with $y$ in register 0 exactly when the host program numbered $c$
outputs $y$ on input $x$.
\end{theorem}

One thing to be exact
about: $U$ is not that universal program. $U$ runs two-counter guests, and a two-counter guest
given its input plainly can't run the evaluator. With inputs written as powers of two one can, as
Theorem~\ref{thm:tc-kleene} near the end of this subsection shows. $W$ runs host programs, and it can, because it's a host program on a machine
with registers.

\begin{theorem}[Kleene's recursion theorem on the host machine]\label{thm:host-kleene}
(1) For every host program $T$ there is a host program $e$ such that for all $x, y$: $e$ outputs $y$
on input $x$ if and only if, for some $d$, $T$ outputs $d$ on the number of $e$, and the program
numbered $d$ outputs $y$ on $x$. (2) So for every transformation $F$ of host programs that some host
program carries out on numbers (on the number of $p$ it outputs the number of $F(p)$, for every $p$),
there is a host program $e$ such that for all $x, y$, $e$ outputs $y$ on $x$ if and only if $F(e)$
does. (3) The same on numbers: for every function $f$ on numbers, defined everywhere, that some host program computes,
some number $e$ has the program numbered $e$ and the program numbered $f(e)$ output the same on
every input.
\end{theorem}

That fixed point is of the \emph{function}: $e$ and $F(e)$ stop on the same inputs and
leave the same number in register 0. Their ledgers, flags and fact tables can differ: $e$ is all INC
and DEC, so whatever $F(e)$ pays, $e$ only simulates. I'd like a fixed point that agrees on the record too, and the next two theorems say how
much of that can be had. The ledger and the flag can be had, but agreement on every register can't.

\begin{theorem}[The recursion theorem, up to the record]\label{thm:host-kleene-record}
Let $F$ be a transformation of host programs that some host program $T$ carries out on numbers.
Then there is a host program $e$ such that, for every input $x$, any state $e$ ends in has a state
$F(e)$ ends in that agrees with it on register 0, the trap latch, the ledger, the flag, the number of
facts in the table, and whether the channel is empty, and the same the other way round. So $e$ stops
on $x$ exactly when $F(e)$ does, and the ledgers are equal exactly, with no surcharge. $T$ may itself
check, commit and certify.
\end{theorem}

Look only at the function computed and you get the
usual fixed-point statement back. There are two limits, and the second one is a theorem. First, $e$ pays
after it has computed everything. It is not the way $U$ runs a guest, mirroring every move: $e$ pays
the counts, CHECKs, then COMMITs, then the CERTIFY, whatever order $F(e)$ used, so the order of the
record's events isn't preserved, and for a program $F(e)$ that never stops, $e$ never stops and pays
nothing, though $F(e)$ may pay without bound. The relation says nothing about states in the middle of a
run. Second, no fixed point agrees on every register, or on what the facts name.

\begin{theorem}[No fixed point of full behaviour]\label{thm:host-no-exact}
There is a transformation $F$ of host programs, carried out by a host program, such that no host
program $e$ behaves the same as $F(e)$ in the sense of Definition~\ref{def:host-same}.
\end{theorem}

A variant of $F$ that also
CHECKs register $r$ and commits to it gives $F(e)$ a fact and a commitment about $r$ that $e$ can't
have, so no fixed point agrees on the fact table alone, or on the channel alone, or on the versions of
the registers. What survives is what Theorem~\ref{thm:host-kleene-record} states, plus any register
named in advance, or any finite set of them: before it pays, the evaluator reads that register's final
value off the simulated run of $F(e)$ and writes it into its own copy of the register, keeping its
scratch work and its two payment registers elsewhere. The register $F$ above writes is chosen by
looking at $e$, and that's the one no fixed point can follow.

And with the recursion theorem proved, the wall of Section~\ref{sec:diagonal} goes up on this
machine without assuming that theorem as an extra premise. Call a property of host programs \emph{functional} when it
depends only on the partial function a program computes.

\begin{theorem}[The diagonal on the host machine]\label{thm:host-wall}
Let $\Pi$ be a functional property of host programs that some program has and some program lacks.
(1) Read ``$f$ is representable'' as ``some host program carries $f$ out on numbers.'' Then no
function $d$ from programs to Booleans whose flip is representable is right about $\Pi$ on every
program: this is the argument of Theorem~\ref{thm:u12-no-inside-decider}, reproved for host programs, with its recursion premise a theorem here.
(2) No host program decides $\Pi$ outright: there is no program $D$ that, on the number of every
program, halts with a number in register 0 that's 1 exactly when the program has $\Pi$.
\end{theorem}

Rice and Kleene are the ingredients the diagonal of Section~\ref{sec:diagonal} asks a substrate for,
and the wall they put up doesn't care about the record: the base does the computing and the flag
watches. A program's ledger and flag are among what ``behaves the same'' sees on the host,
so a nontrivial property of those is undecidable too, in the reduction sense of Theorem~\ref{thm:host-rice}, and the fixed point is up to the record, as
far as the theorems above say and no further.

Now the two-counter machine itself, which I promised at the start of the subsection: Minsky's two
bins of pebbles and nothing else. Rice comes first. Say a two-counter program \emph{behaves the same} as
another on a set $X$ of inputs when, on every input $n$ in $X$, started with $n$ pebbles in bin A and
none in bin B, one stops exactly when the other does, and then they agree on both bins, the trap
latch, the ledger, the flag, and on which property and which counter each recorded fact and the
channel are about. The versions are left out: they say which write a fact is about, and the trick
below moves them.

The trick is the whole problem. Rice's reduction has to keep the program's input safe while a copy of
a halting instance $H$ runs, and there's no third bin to keep it in. My first guess was bin B, and
that dies at once, because $H$ needs both bins. The second guess is Minsky's own packing: keep several
numbers in one bin as powers of different primes, $H$'s two counters as the exponents of 2 and 3, and
the input as the exponent of 5. That dies more quietly. To put the input there, the program would
have to turn $n$ pebbles into $5^n$, an exponential, and Schroeppel's memo, from
Section~\ref{sec:smallest}, gives a proof that two plain counters can't even build $2^n$ from $n$. I had no way to build $5^n$.

What works is to stop keeping the input beside the work. Turn $n$ pebbles into $6n + 5$, which is only
adding, and then multiply in $H$'s counters as powers of 2 and 3. $H$ runs by halving the pile and
splitting it in thirds and building it back up, and the $6n + 5$ rides along in the middle of every
one of those piles. It's odd, and counted off in threes it leaves two over, so no halving or thirding
ever reaches into it.

\begin{theorem}[Rice's theorem for the two-counter machine, on any inputs]\label{thm:tc-rice}
For any set $X$ of inputs, a property of two-counter programs that's carried from a program to any
program that behaves the same on $X$ is undecidable in the reduction sense of Theorem~\ref{thm:u34-undecidable}, unless every program has
it or none does, and then it's decidable. $X$ can be every input, the powers of two, or $\{0\}$.
\end{theorem}

So $6n + 5$ is the place to keep the input that two counters seemed not to have. It isn't beside the
work but inside it, in a factor the work never touches. I like that one more than I can justify: the
input hides in the leftover that every halving and thirding walks straight past.

Kleene needs more than keeping an input. It needs a program that runs a program. So write a number
$x$ as $2^x$ pebbles and read a result as the exponent of 2, the way Minsky did. With inputs written
that way, there's one two-counter program, call it $U_2$, that, started on $2^x \cdot 3^e$, stops with $2^y$
exactly when the program numbered $e$ gives $y$ on $x$; specialising a program to a number is exact;
and the recursion theorem holds. ($U_2$ is a different program from the universal host program
$U$, which runs on the machine with a register for every number.)

\begin{theorem}[Kleene for the two-counter machine, inputs written as powers of two]\label{thm:tc-kleene}
(1) For every transformation $F$ of programs that some program carries out on program numbers, there's
a program $e$ that computes the same function as $F(e)$, with inputs and outputs written as powers of
two. (2) The fixed point is of the function only: there's such an $F$ for which no $e$ behaves the same
as $F(e)$ once the properties and counters of the recorded facts count.
\end{theorem}

The counterexample in (2) is short. Number the programs the way $U$ numbers its guests
(Appendix~\ref{sec:universal-blocks}), so a program's number is above the code of every property
it names. Let $F(p)$ be the one-line program that checks that bin A holds at least the number of
$p$ plus one. A program can only record facts about properties it names, so no $e$ can record the
fact $F(e)$ records, and from a large enough input $F(e)$ does record it.

Even with Schroeppel in hand, my first guess about Kleene was that the powers of two were a
convenience, and that plain inputs would only take more work. They don't take more work: with plain inputs, Kleene's theorem is false. A
transformation is \emph{carried out} if some program, started on $2^n$ where $n$ is the number of a
program $p$, stops with $2^m$ where $m$ is the number of $F(p)$.

\begin{theorem}[No plain recursion on two counters]\label{thm:no-plain-recursion}
(1) No program of $n$ lines stops with $c \cdot x$ pebbles in bin A, started on $x$ pebbles in bin A
and none in bin B, for every $x$, where $c$ is the least prime above $N = 16(n+1)(4n+1)^{16}$. (2) So there's a
carried-out transformation $F$ such that no program $e$ computes the same plain function as $F(e)$.
\end{theorem}

Nothing in that proof used more than this: two bins that move one pebble at a time, finitely many controls, a next step fixed by the control and the bins, and bins that, above some threshold, are read only through their parity. Any
machine like that has the same obstruction, and the same computable transformation $F$ with no
program $e$ equivalent to $F(e)$ on plain inputs.

I went looking for Schroeppel's own memo and found a scan of it. As I read it, his route is close in spirit: he trades a two-counter machine that
starts with one counter empty for a machine with one register that can only add,
multiply and
divide by fixed numbers and take one away, each of those built out of one of the program's loops, and that's the
stages here under another name. I haven't followed his proof line by line, so I compare the two
at that level and no further; the argument this book stands on is the one above. The $c$ it uses is astronomically bigger than anything a real
multiplier needs, and finding the best bound is a question I leave alone.

So on two bins, Rice holds on every input; Kleene holds with inputs written as powers of two, for the
function, and fails for the facts; and with plain inputs Kleene fails outright. The last one is the
answer I didn't see coming: the powers of two weren't a convenience, and on two bins they're the
difference between a recursion theorem and none.

That failure has a setting computability theorists know. The mathematician Hartley Rogers, in his
1967 textbook, calls a numbering of the computable functions \emph{acceptable} when every other
computable numbering of them can be translated into it by a computable map, and in every acceptable numbering
Kleene's theorem holds. Read on plain inputs, the two-bin programs number something smaller than
all the computable functions, since none of them computes $2^n$ from $n$ (Schroeppel's theorem; compare Theorem~\ref{thm:no-plain-recursion}), so they fall outside Rogers's
class, and his route to the theorem isn't open there. The computer scientists James Royer and John
Case (1994) studied that territory: programming systems for classes of functions smaller than all
the computable ones, which can't be acceptable, and what can still be programmed and proved in
them. And the logician and puzzle
writer Raymond Smullyan's two books on recursion theory and self-reference (1993, 1994) take the
fixed-point theorems apart axiomatically, so you can see which property of a numbering each one
leans on. The plain-input failure is a small instance on the map they drew.

\subsection{The universal machine's guest interface}

I'll say what the two universal programs cover plainly, a sentence at a time. The stated class of guests of $U$ is the small
machine's programs, every one of them, from every clean start $\mathrm{start}(x, y)$.
The whole-run theorems don't cover a guest handed facts in its table, a commitment in its channel
or a raised flag at the start; the one-step simulation, by itself, works from any related pair. It's the same
kind of class Turing's universal machine has: a fixed instruction set, and a finite program held
as data, here packed into one number the way Turing packed his into a description number. Each
guest keeps its computation, its halting, its answer, its record and its prices, exactly: the
host halts when it halts, its flag goes up when the guest's would, on a chain of its own, and its
ledger is the guest's.

The guests of $U_P$ are every Thiele machine whose driver, step, prices and reading are computed
by programs, with the reading decided on every state. Each keeps its computation and its halting;
its reading becomes the host's flag as a latch, raised on an earned chain whose claim is the
guest's own reading; and its prices become the host's ledger, exactly, except for a surcharge at
the move where the reading first says yes: at most two when the reading starts at no, exactly
three when it starts at yes. The surcharge is zero when that move costs three or more, and it can't always be avoided: when the
guest certifies with a ledger below three, no program on that host, run from a clean start, matches its ledger exactly (Theorems~\ref{thm:u34-presented}
and~\ref{thm:u34-presented-surcharge}). Picture a clerk who'll stamp anybody's receipt. He copies
it onto his own forms, bills at your rates to within two when your reading starts at no and to within three when it starts at yes, and the stamp at the bottom is one he
earned by checking your own reading, not one he copied.

What's outside the class is outside the theorems. A machine whose reading no program decides
(an oracle checker, say, which the definition of certified Thiele-complete doesn't rule out) doesn't come in. Nor
does a machine whose next move is nobody's computation. The two universal programs have different start conditions: $U$'s whole-run result starts
its guest clean; $U_P$ starts its guest where the presentation says, even if that state
already has a yes reading. The host starts clean and earns a claim about that reading.
It does not reconstruct or authenticate the guest's earlier history.

And insisting on no built-in step is the point, and a plain interpreter shows why. A plain interpreter on an ordinary
computer runs guests too. If its guest record is freely writable data, its
meaning depends on the interpreter following the intended update rule.
An ordinary computer can also protect that data with a monitor or another
enforced interface. Here the rule is explicit in the host's instruction
semantics: $U$ and $U_P$ are fixed interpreters, and nothing in them raises the
host's flag except the host's own CERTIFY, and the host's step won't let that pass without a
COMMIT, or the COMMIT without a CHECK on a slot nobody has touched. The guest's record isn't data
the host updates. It's a record the host earns again, under the same rule, on its own registers.


%===========================================================================

\section{Every model}
\label{sec:every-model}
%===========================================================================

After the smallest honest machine (Section~\ref{sec:smallest})
and the universal programs of the previous appendix, there was a worry I couldn't put down. Everything in them was built on
Minsky's two counters. The definition of certified Thiele-complete names a two-counter window in clause (a).
The universal program runs two-counter guests. So is all of this a fact about counters? Pick a
different machine to start from, a Turing machine with its tape, a register machine, Alonzo
Church's lambda calculus, and does the record still go on the same way, or did the definition
depend on my example?

I chased it with a picture: a box with a panel bolted to its front. The box computes, with levers
for moves and a window showing what clause (a) asks to see, two numbers and a line number. The panel is the example machine's earned layer: version tallies that click
when the box writes, a rack for facts, the flag, and the three buttons CHECK, COMMIT and CERTIFY at a cost of
one each. The worry is whether that panel only bolts onto the one box I built.

So this appendix goes both ways. First the panel goes onto boxes I didn't build, and I say
exactly what a box has to supply for it to fit. Then other boxes, panel and all, get run
inside mine. If the record is really about the record and not about pebbles, both trips
should work, and the conditions on the way should say what has to survive the move.

\subsection{From any base up}

Start with the direction I care about most: the panel onto a box I didn't build. The box can be any
machine at all that has a universal base in the sense of clause (a), and in the picture that means two
things: it shows a two-counter window, and for each two-counter instruction it has a lever that acts
on the window as that instruction does. What else is inside, what its states look like, I don't get
to know.

My first guess was to keep whatever the box came with. Plenty of machines have a meter of their own,
or a light on the front that means something's done, so leave those where they are, bolt the panel on
beside them, and see what happens. It can't work, and the definition says why. If a lever can cost
anything, computing adds to the ledger, and certified Thiele-complete asks that it add nothing, which
is a rule of the price list, not a claim about heat. If a lever can switch on the box's own light, a record went up with
no CHECK behind it. Certified Thiele-complete rules out both. So the meter and the light come off, along with
everything else the box already had: its costs, its record, all of it, except the states, the levers
and what each lever does to the state.

Then bolt on a panel cut down from the small machine's, and nothing else:
one version number for the whole box that every lever
raises, which is the tally, a table of facts, a channel, the flag, and three new moves, CHECK,
COMMIT and CERTIFY, each costing one, with every lever costing nothing. The three act as on the
small machine, with the tally in place of a counter's version: CHECK files its claim with the tally
when the checker says yes and the table has room, COMMIT puts the fact in the channel when the table
holds its claim at the current tally, and CERTIFY raises the flag when the channel holds a fact.
Call that the \emph{lift} of the machine. A lever moves the box only while the trap is down. The trap works as on the small
machine: a CHECK, COMMIT or CERTIFY that fails sets it, nothing takes it down, and once it's up
every move, a lever included, changes nothing but the ledger.

\begin{theorem}[Every universal base lifts]\label{thm:lift}
(1) The lift of any box with a universal base, over a language of claims about the box's state (``unchanged'' meaning the box's state is the same), with an exact equality test and a checker proved to say yes exactly when the claim holds, is certified Thiele-complete, as long as the claims its CHECK can
ask include one that's true at one loaded state and false at another, and the fact table has room
for at least one fact. (2) The claims ``this register is at least $n$'' and ``the program is at line $n$,''
read off the window, always include such a claim, so with those claims every universal base lifts.
(3) Conversely, every certified Thiele-complete machine has an interface meeting the four clauses under which
its base moves, taken alone, form a machine with a universal base. (4) The
panel leaves the box alone: from a state whose trap is down, a run of lever pulls in the lift moves
the box exactly as the box's own run does, and the trap stays down.
\end{theorem}

\begin{proof}
(1) The four clauses, one at a time. (a) The lift keeps the box's states, levers and step and adds
the tally, the table, the channel, the flag, the trap and the ledger. With the trap down, a lever
moves the box's part of the state as the box's step does, raises the tally by one, and touches
nothing else. So from a live state, which is a state whose box part is live and whose trap is down, the window moves as the instruction does and the state
stays live; no lever changes the flag; the loaded states are the box's loaded states with an empty
panel, so they're clean; and the flag, once up, is lowered by nothing. (b) Start with the panel
empty. The only move that raises the flag is CERTIFY, and CERTIFY traps unless the channel holds a
fact; the channel is filled only by a COMMIT, which traps unless the table holds its claim at the
current tally; and the table is filled only by a CHECK that passed, which files the claim with the
tally it was checked at. So in front of the raise there's a passing CHECK of some claim $c$, then a
COMMIT of $c$ at the same tally, then the CERTIFY. The trap is down from the CHECK to the COMMIT,
since the COMMIT passed and nothing takes the trap down, and with the trap down the tally rises
with every lever pull, so an unchanged tally means no lever was pulled in between, so the box's
state is unchanged, and that's
what ``the thing $c$ is about is unchanged'' says in the lift's language, whose claims are about the
box's state. A passing CHECK means the checker said yes, and the checker was proved to mean what it
says. (c) is the price list: levers at a cost of zero, the three buttons at one each, the ledger the running
sum. (d) Take the claim that's true at one loaded state and false at another. From the first, CHECK
passes and files it, the table having room, COMMIT finds it at the same tally, and CERTIFY raises
the flag. From the second the CHECK traps, and a trapped machine raises nothing. (2) ``$A$ is at
least 1'' is true at the loaded state showing $(1, 1, 0)$ and false at the one showing $(1, 0,
0)$. A window claim is a register and a number, or a line number, so two claims can be tested
equal, and its checker reads the window, which is its meaning, so the checker is sound by
definition. (3) Take an interface the machine is certified Thiele-complete with and throw away every move
whose kind isn't base. What's left has the same states, the base moves as levers, and the
interface's window, live states and loaded states, and clause (a) says of exactly these that the
window follows every instruction from every live state, that the state stays live, and that each
loaded state shows $(1, a, b)$ and is live. That's a universal base. (4) A lever pull with the trap
down changes the box's part by the box's step and the tally by one, and nothing else; the trap
goes up only on a CHECK that fails, a COMMIT without its fact, or a CERTIFY with the channel empty,
never on a lever. So along a run of lever pulls the box's part follows the box's own run, pull for
pull, and the trap stays down.
\end{proof}

Of the checker's exactness, (1) uses one half whole and the other half once: (b) needs the checker
to say yes only when the claim holds, and (d) needs it to say yes on the one claim at the one loaded
state where that claim is true. The theorem asks for both halves whole, the way every theorem
about the small machine does, so one checker serves all of them.

Clause (b), the small machine's proof with the bins swapped out for whatever the box has,
is the part I'd braced for, and it turned out to be nothing. The chain CHECK, COMMIT, CERTIFY
never looked at anything but the tally and the claim. The levers changed only the tally on the panel, making old facts stale; they never inserted
a fact, made a commitment or raised the flag. Neither
side of the proof ever needed to know it was looking at pebbles.

So the earned layer isn't something two counters do. It sits on top of any box with a universal base, and it never needed to know which one. I'd half expected to be rewriting the proof once per model.

\subsection{What never lifts}

A theorem that says the panel fits every box is one I don't trust until I've watched it fail to fit
something. So I went looking for boxes it won't turn into a certified Thiele-complete machine, and the first one I tried
was the one I was surest would work.

It was a Turing machine, with one fixed table, and each lever does one row of the table. That looked like the
obvious base to me, and it isn't one at all. Clause (a) asks for a lever for every two-counter
instruction, and there are infinitely many two-counter instructions, since a jump can go to any line.
A fixed table has finitely many rows, so it has finitely many levers, and it runs out of levers long
before it runs out of instructions. The first two items below are that failure in two forms, and
the others fail for other reasons, and each is proved.

\begin{theorem}[What doesn't lift]\label{thm:lift-fails}
(1) A box with only finitely many levers has no universal base, and no box with a universal base is
finitely branching, meaning that from every state all its levers together lead to only finitely many
different states. (2) The lift of a finitely branching box is never certified Thiele-complete, whatever claims
and fact table it gets. (3) Neither is the lift of a box whose step ignores the state it's in, where
each lever sends every state to the same place. (4) On a record axis where every place is at or below
every other, there is no certified Thiele-complete machine at all. (5) With a checker that says yes exactly when the claim holds, the lift's own non-vacuity clause holds exactly when the fact table has room and some claim is true at one loaded state and false at another.
\end{theorem}

\begin{proof}
(1) Take a universal base and its loaded state showing $(1, 1, 0)$. For each $j$ the lever for
DEC $A$ $j$ must take it, in one pull, to a state showing $(j, 0, 0)$: $A$ was above 0, so one
comes off and the line jumps to $j$. Different $j$, different window, different state. So one state
has infinitely many different successors, which is what not finitely branching means; and since one
lever lands in one place from one state, the levers for the DEC $A$ $j$ are all different, so there
are infinitely many of them. (2) In the lift the moves that cost nothing are the levers, since each
button costs one, and from any state the levers together lead to finitely many states: with the trap
down, the box's finitely many successors, each with the tally raised by one, and with it up, the state
itself. (The buttons can lead to infinitely
many, one CHECK for each claim, but each costs one.) If the lift were certified Thiele-complete through some
interface, that interface's base moves would cost nothing by clause (c), so they'd all be levers, and
by Theorem~\ref{thm:lift}(3) they'd form a universal base on their own. By (1) a universal base has
a state from which its moves reach infinitely many states, and levers of a finitely branching box
can't. (3) Suppose each lever lands in one place whatever the
state. The move an interface gives an instruction is a base move, at a cost of zero by clause (c), while each
button costs one, so it's a lever. Pull the lever for INC $A$ twice from a loaded state showing
$(1, 0, 0)$, and pull the lever for INC $B$ and then the lever for INC $A$ from the same state.
Both runs end with the lever for INC $A$, so both end with the box in the one place that lever
lands, with the tally two higher and the panel untouched: the same state of the lift. But the window
would have to show $(3, 2, 0)$ after the first run and $(3, 1, 1)$ after the second, and a window
is a function of the state. The same two runs refute the box on its own. (4) On a record axis,
raising the record means leaving the down-set of where it stood. If every place is at or below
every other, every down-set is the whole axis and no move ever leaves one; but clause (d) asks for
three moves that raise the record from some loaded state. (5) If the table has no room, every
CHECK traps and no CERTIFY ever fires, so (d) fails. If no claim is true at one loaded state and
false at another, (d)'s second half has no witness. And when both hold, (d) was proved in (1).
\end{proof}

So my first test case failed, and at first I read item (1) as saying a Turing machine can't be a
base at all. It says something narrower. A Turing machine \emph{with one fixed table} can't be one. A classic model comes in the other way round, with one lever for each program that carries out one two-counter instruction, run to its end. Pull it and the box runs that program from start to stop in a single move. That's how
the register machine below comes in, and it's the honest reading of ``a Turing machine is a base'' too.
The levers are the programs, not the rows.

That still leaves a worry about clause (a) itself. It asks for one lever per two-counter
instruction, which is a fact about how the two-counter machine is presented, and a box that does
each instruction in several pulls looks shut out for a reason that has nothing to do with
computing. It isn't shut out, and the fix is to let the pulls be grouped. Call a box a \emph{base
by simulation} when it has a window reading a two-counter configuration off each state, through
any encoding at all, a class of live states, a finite list of its own levers for each two-counter
instruction, and a live loaded state showing $(1, a, b)$ for each start, such that from every live
state the list acts on the window exactly as the instruction acts on a configuration and the
state stays live. The \emph{grouped box} has the same states, and its levers are finite lists of
the box's levers, a list pulled in order.

\begin{theorem}[The one-lever presentation loses nothing]\label{thm:lift-macro}
(1) Every box with a universal base is a base by simulation. (2) Every base by simulation is a
universal base of its grouped box, and a run of the grouped box is the box's run along the joined
lists. (3) So for every base by simulation, the lift of its grouped box, with the window claims and
a fact table with room for at least one fact, is certified Thiele-complete. (4) The box with five levers, add
one to $A$ and go to the next line, the same for $B$, add one to a jump target (leaving the window
alone), and for each of $A$ and $B$ ``if it's above 0,
take one off and jump to the target, otherwise go to the next line, and put the target back to 0
either way,'' is a base by simulation, has no universal base, and its grouped box lifts to a
certified Thiele-complete machine.
\end{theorem}

\begin{proof}
(1) Each list has one lever. (2) The grouped box's step is the box's run along the list, so the
simulation clause of a base by simulation is clause (a) for the grouped box. (3) is (2) with
Theorem~\ref{thm:lift}. (4) INC $r$ is one pull; DEC $r$ $j$ is $j$ pulls of the target lever,
then the decrement lever for $r$, from a live state, where live means the target is 0; the
decrement lever puts the target back to 0, so the state stays live. The box has five
levers, so item (1) of Theorem~\ref{thm:lift-fails} says it has no universal base.
\end{proof}

So clause (a) can be read by simulation, and nothing is lost: a box that carries out each two-counter instruction with one fixed finite list of its own moves, through any window, is a universal base the moment its moves are
grouped, and it lifts. Item (1) of Theorem~\ref{thm:lift-fails} is about the presentation, and the
five-lever box is the proof that a finite lever set is not what fails. What this doesn't give is a
Turing machine's rule table running a tape: the five-lever box is a register box, and the next
subsection is still how Turing machines come in.

Item (3) is a box with no memory. Pull a lever and it lands wherever that lever always lands,
whatever happened before. So a window on it can't tell INC $A$ followed by INC $A$ from INC $B$
followed by INC $A$, and a box that can't tell those apart can't keep two counters in its window.

Item (4) is the trouble with bare preorders from Section~\ref{sec:direction}, where the cups stopped telling you the place, showing up one level higher. If
every place on the axis is at or below every other, the flag has no up to go to. No step ever turns a cup over, nothing ever has to be paid for, and the earned chain has nothing to earn.

\subsection{Computation is the slice}

The lift says the panel goes onto any box with a universal base. Turn it round and it says something
about the box: the box is what you see of the lift when you stop looking at the panel. That's where
ordinary computation sits in this book, so I want it said exactly, with a name.

\begin{definition}[Universal-base slice]\label{def:slice}
Take a machine with a ledger, a whole number for each state, and a loaded state for each start, and
take a box with a universal base. A \emph{universal-base slice}, called a slice
in this appendix, is different from the global sliced interface of
Definition~\ref{def:passive-sliced-thiele}. It requires a universal base, a working
domain and reachable hidden witnesses; it does not require all non-computing moves
to fix the view. It is a map from the
machine's states to the box's states, a move of the machine for each lever of the box, called its
\emph{computing move}, and a class of the machine's states, called \emph{working}, such that:
(i) each loaded state is working, and the map sends it to the box's loaded state for the same start;
(ii) from a working state, a computing move lands on a working state, and the map sends where it
lands to where the box's lever takes the map's image; (iii) from some loaded state there are two runs
whose ends the map sends to one box state, with the record up at one end and down at the other, and
two runs whose ends it sends to one box state, with different ledgers.
\end{definition}

When there's a slice, I'll say the box is a slice of the machine. The first two clauses say the slice
runs as the box runs. Start with a working state and pull one computing lever.
Clause (ii) says that its picture moves exactly as the box moves, and that the
full state is working again. That last part lets you apply the same rule to the
next lever, then the next. A whole run therefore follows the box's run, move
for move. Clause (i) puts the two at matching starts, so they finish at matching
pictures too. The third clause
says it can't see what was kept or what it cost: a function of the box state that gave the record at
the end of every run from every loaded state would have to give both answers at the one box state
the clause names, and the same goes for the ledger. Nothing else is hidden in the definition. Read
the box's window through the map and take as live the working states whose image is live, and the
first two clauses make the computing moves for the box's instruction levers a universal base of the
machine. So a machine has a slice onto some box exactly when it has a universal base whose window
meets the third clause, and then the window itself is a slice onto the two-counter box, whose states
are configurations and whose levers are the instructions. The third clause does work. A
Thiele-complete machine's computing moves, taken alone, are a box, and the map that forgets nothing
meets the first two clauses onto it and fails the third, because the record and the ledger are
functions of the state; and the window of the first clock of Observation~\ref{obs:tc-clock}, whose
record is a question about its counters, fails it too.

\begin{theorem}[Computation is the slice]\label{thm:slice}
(1) Take the lift of a box with a universal base, as in Theorem~\ref{thm:lift}(1). Forgetting the
panel, which keeps the box's part of each state and drops the tally, the table, the channel, the
trap, the ledger and the flag, is a slice of the lift onto the box, with the levers as its computing
moves and the states whose trap is down as its working states. Beyond that, CHECK, COMMIT and
CERTIFY never move the box's part, from any state; from any box state with an empty panel, the
lever runs are sent, under forgetting, to exactly the box's own runs from that state, lever list for
lever list; and from a loaded state where the claim of Theorem~\ref{thm:lift}(1) is true, the run of
no moves and the run CHECK, COMMIT, CERTIFY of that claim end on one box state with the flag down
and up and ledgers exactly 3 apart, CERTIFY three times ends there too with the flag down and the
same ledger as the second run, and CERTIFY $d$ times, for any $d$, ends there with the flag down and
a ledger $d$ above the first run's. (2) With no room in the fact table, the lift of any box, with any
claims, is certified Thiele-complete through no interface at all. (3) Let a machine be
Thiele-complete through an interface; every certified Thiele-complete machine is
(Theorem~\ref{thm:tc-decomposition}). Its computing moves are exactly its moves that cost nothing,
so they're the same through every interface it's Thiele-complete through; taken alone, with the
interface's window, live states and loaded states, they form a machine with a universal base; and
the window is a slice of the machine onto the two-counter box, with the instructions' moves as
computing moves and the live states as working. (4) Let a machine be run as one rule: from each
state, every move leads to the same next state. It has no universal base, no base by simulation and
no universal-base slice onto any box. It obeys the toll exactly when every move costs at least one or no state
reading no steps to one reading yes, so it's a Thiele machine only as a clock or with a reading that
never rises.
\end{theorem}

\begin{proof}
(1) The lift's loaded states are the box's loaded states with an empty panel and the trap down, so
they're working, and forgetting sends each to the box's loaded state for its start. From a state with
the trap down a lever moves the box's part as the box's step does and leaves the trap down
(Theorem~\ref{thm:lift}(4)), which is the second clause. Start with any box state
and an empty panel. The first pull matches the box's move and keeps the trap down;
the next pull therefore starts under the same conditions. Repeat through the
word and every pull matches. A CHECK, COMMIT or CERTIFY that
passes changes only the table, the channel, the flag and the ledger; one that fails sets the trap and
adds to the ledger; and with the trap up every move changes nothing but the ledger. None of them
touches the box's part. Now take the loaded state where the claim is true. The run of no moves ends
there with the flag down and the ledger at 0. On the other run the CHECK passes and files the claim,
the table having room, the COMMIT finds it at the same tally, and the CERTIFY raises the flag: three
moves at one each, none of them moving the box, so the same box state, the flag up and the ledger
at 3. CERTIFY from a clean state finds the channel empty and traps, and every later move changes only
the ledger, so CERTIFY $d$ times ends at the same box state with the flag down and the ledger at $d$,
which at $d = 3$ is the chain's ledger. So the third clause holds, for the flag and for the ledger,
and no function of the box state and the ledger together gives the flag. (2) Take any interface and
the three moves its clause (d) names. Each costs one by clause (c), so each is a button. Take a loaded
state from which they raise the flag; it's clean by (a), so its flag is down by (b). If its trap is
up, nothing passes and nothing rises. If its channel holds a fact, CERTIFY alone raises the flag from
it: a run of one move from a clean state that ends with the flag up, where (b) asks for a CHECK and a
COMMIT before the raise. If its channel is empty, some COMMIT has to pass inside the three moves,
before the CERTIFY that raises the flag. Until then no CHECK has passed, because a CHECK passes only
with room in the table, and no button moves the tally, so that COMMIT found its fact in the loaded
state's own table at the loaded state's own tally; that COMMIT and a CERTIFY, run from the loaded
state, then raise the flag in two moves, against (b) again. So the three moves raise the flag from no
loaded state, and clause (d) fails. (3) By (C), computing moves cost 0 and every other move costs 1,
so the computing moves are the moves that cost 0, whichever interface the machine is read through.
Throw the rest away and what's left has the same states, the interface's window, live states and
loaded states, and a move for each instruction that acts on the window as the instruction does from
every live state and keeps it live, with each loaded state live and showing $(1, a, b)$: a universal
base, as in Theorem~\ref{thm:lift}(3). The same facts are the first two clauses of a slice onto the
two-counter box, and the third is Theorem~\ref{thm:tc-window}(2): two runs from one loaded state
end on one window, one with the record up and one with it down, and their ledgers differ. (4) From
each state all the moves lead to one state, so the machine is finitely branching, and by
Theorem~\ref{thm:lift-fails}(1) it has no universal base. A slice would give it one, as the remark
after Definition~\ref{def:slice} says. A base by simulation fails too. A list of $p$ moves does $p$
steps of the rule, so two lists run one after the other end in the same place in either order. From
the loaded state showing $(1, 0, 0)$, INC $A$ and then DEC $A$ to line 5 shows $(2, 1, 0)$ and then
$(5, 0, 0)$, while DEC $A$ to line 5 and then INC $A$ shows $(2, 0, 0)$ and then $(3, 1, 0)$, and one
state can't show both. For the toll: if a state reading no steps to one reading yes, every move
takes that step, so every move costs at least one, every step does, and the machine is a clock
(Observation~\ref{obs:tc-clock}); if none does, the reading never rises and the toll asks nothing.
\end{proof}

Not every part of that carries the same weight, and I'd rather say which before somebody else does.
Some of it holds by construction. That the slice runs as the box runs, and that CHECK, COMMIT and
CERTIFY leave it alone, are how the lift was built, and the blindness in (1) follows from the
second once the chain raises the flag: the chain and the run of no moves end on one box state
because no button moves the box. In (3), the window meeting the first two clauses of a slice is the
definition of Thiele-complete read out, and so is the fact that the moves that cost nothing are the
same through every interface. Three parts could have failed. That the lift is certified at all is
the earned chain of Theorem~\ref{thm:lift}, clause (b). That the window of every Thiele-complete
machine can't see what was kept or what it cost is Theorem~\ref{thm:tc-window}, and it could have
gone the other way: the first clock's window sees its record. That a table with no room certifies
through no interface is (2). And (4), that a machine run as one rule has no universal-base slice and is a Thiele
machine only as a clock or with a reading that never rises, is short, and its second half rests on
one choice in Definition~\ref{def:certsys}: the toll puts the price on the move, the same from
every state. With a price that could read the state, a machine run as one rule could charge exactly
its raises.

What ``computation'' means in the theorem's name is a box's levers acting on its states. A model that
comes with a move for each two-counter instruction is covered as it stands, and one with a fixed list
of its own moves for each is covered once its moves are grouped (Theorem~\ref{thm:lift-macro}). A model run as one rule, a Turing machine
stepping by its one table, is covered by (4) only: it has no universal-base slice, and it cannot supply the appendix's universal-base target. Its ordinary
transition rule still lifts to the global sliced machine of
Definition~\ref{def:passive-sliced-thiele}: put the lamp and tally beside it.
The fixed-list obstruction concerns its instruction interface. A Turing machine read the other way, with a lever for each program that carries out one
two-counter instruction, as in the previous subsection, is a box like any other. Nothing in the
theorem uses the equivalence of the seven models below (Theorem~\ref{thm:lift-models}(1)), which is
other people's work; that equivalence is what says a box with a universal base computes whatever the
other models compute.

A machine run as one rule can still run programs, and its program can keep a flag. Whether that flag
is a record the machine keeps or a bit it only simulates can be said exactly too, once the program
is part of the state.

\begin{observation}[A record or a simulated bit]\label{obs:program-in-state}
Take a machine whose moves are operations on configurations, with a reading of the configurations,
and put the program in the state: a state is a program, which picks an operation for each
configuration, together with a configuration, and one step applies to the configuration the
operation the program picks there and keeps the program. Take every program. Then the reading never
comes down under any operation from any configuration exactly when, read on these states, it never
comes down from any state.
\end{observation}

\begin{proof}
If no operation brings the reading down anywhere, no step does, since each step applies an
operation. If an operation brings it down from some configuration, the program that picks that
operation at every configuration is one of the programs, and the step from the state made of that
program and that configuration brings it down.
\end{proof}

So whether a cell is a record or a simulated bit is a fact about the one rule that steps program and
configuration together, and it takes the whole set of states to see it. A certification cell that
some program can clear, like the one a universal machine keeps as plain data on its tape
(Section~\ref{sec:machine-vs-program}), is a simulated bit, though on a run whose program never
clears it the cell goes up and stays, as a record would.

\subsubsection*{Why the live-state clause stays in the slice}

Go back to the blown fuse. A free lever turns the live machine's gears and leaves the
trapped machine still. Hide the difference between those states and one camera picture
would have to prescribe both outcomes. Theorem~\ref{thm:native-static-shadow} gives the
exact consequence: hide every native CHECK globally, require autonomous free moves,
and every projected free move becomes identity.

Keep the camera promise on live states instead. The computing moves act there and keep
the machine live. A failed panel check leaves that domain. That is the boundary in the
slice above. It lets the computational view move while the native failure rule keeps
its power to stop the run.
 Restricting the promise to reachable states would not do this:
append a failed CHECK to a run and the trapped state is reachable too. It is the live
boundary, which that check crosses, that makes the difference.

The certificate and ledger remain passive beside the core. The trap is inside the core,
where the gears read it. Add the publication file and you can keep which claims were
certified, with the native lamp as its nonempty reading. The fuse still works as before.
The file records what was published; the live-state clause says where the computing
promise holds. Both boundaries belong in the description of the machine.
 Theorem~\ref{thm:native-archive-lamp} proves that the lamp reports
whether the publication file is empty. Theorem~\ref{thm:native-history-fibre} counts
how many different files can sit above the same full native state.

\subsection{The same thing, in seven models}

Now the other direction, the one the worry was really about. Put boxes with different insides in a
row: one with a tape, one rewriting lambda terms, one with bins. Is a panel riding on one of them the
same as a panel riding on another? The honest way to ask is whether the boxes agree on what they
compute.

My first idea was to wire the boxes to each other myself, two at a time, with a relation that
says one box's run keeps pace with another's, but that's the wrong tool. Section~\ref{sec:latch}'s
weak base equivalence is about runs keeping pace, and what I need here is coarser: do the boxes
compute the same relations on numbers? For that the work is already done, and it isn't mine. The seven models
are the weak call-by-value lambda calculus L, Turing machines, two Minsky instruction
conventions, binary stack machines, Conway's FRACTRAN, and the $\mu$-recursive functions,
and that they all compute the same things is Minsky's, Shepherdson and Sturgis's, and Forster,
Larchey-Wendling and their coauthors' work. The translations, with how each one writes its
inputs and reads its outputs, are theirs (Appendix~\ref{sec:counter-constructions}).

\begin{theorem}[The base doesn't matter]\label{thm:lift-models}
(1) A relation on numbers that's computable in any one of those seven models is computable in all
seven. (2) The box that keeps a whole two-counter configuration in one number, with one lever for
each instruction, has a universal base, its step is computable in all seven models, and its lift,
with a table of sixteen facts, is certified Thiele-complete. (3) A register machine of the kind Cook and
Reckhow described, where each two-counter instruction compiles to at most five of its own instructions, has a universal base too, and its lift is certified Thiele-complete. (4) Any box on numbers
with a universal base lifts to a certified Thiele-complete machine, and its step, computable in one model, is
computable in all seven.
\end{theorem}

\begin{proof}
(1) Compose the translations (Appendix~\ref{sec:counter-constructions}).
(2) Encode a configuration as $\langle\ell,\langle a,b\rangle\rangle$.
Decoding, performing one INC or DEC, and encoding the result are the total finite
arithmetic operations constructed there. So the window follows
each instruction exactly; loading is the encoding of $(1,a,b)$.
(3) Let RAM registers 0, 1 and 2 hold the simulated line, $A$ and $B$, and keep
register 3 equal to 1. To increment counter register $r$, add register 3 to $r$
and then to register 0. To decrement it with target $j$, use these five lines, numbered 0 to 4: if $r>0$ jump to 3; add register 3 to register 0; jump to 5; subtract register 3 from $r$; set register 0 to $j$. At line 5 the finite block has ended. A base move restarts a block at its first line,
runs it for its length, and keeps the resulting registers; an ended block stays put.
The zero and positive cases give exactly the required window, and neither changes
register 3. Loading puts $(1,a,b,1)$ in registers 0 to 3.
(4) Apply the lift theorem to that universal base and apply (1) to its step function.
\end{proof}

So the answer to the worry is that the record doesn't care. A box with a universal base, described in any of the models people have
used to say what computing is, takes the same panel, becomes the same
kind of machine, and pays the same price, one each, for the same three acts.

And the trip goes back down. So far the panel went
onto other boxes. Now put another box, panel
and all, inside mine. Appendix~\ref{sec:presented-universal} proved that the fixed program $U_P$ on the small
machine's host runs every computably presented Thiele machine. In terms of the seven models, that
reads: if the four relations that describe a Thiele machine (its driver, which picks the next move or says
stop, its step, its cost and its reading, on codes) are computed in any one of the seven by programs that answer on every input, then the machine is computably presented,
and $U_P$ runs it, halting when it halts, with its flag, and with a surcharge of at most two when its reading starts at no.

So a Thiele machine described in lambda calculus runs, under one fixed program, on the small
machine's host, which gives model independence in both directions: up, by lifting any universal base, and
down, by running anything presentable on one fixed program. When I saw both arrows close I stopped
worrying that the definition depended on my example. A definition that did wouldn't survive being carried both
ways.

\subsection{Why two counters}

If the box doesn't matter, why did I build on two bins? I wanted the smallest box that holds
up, and the obvious smaller one has a single bin. Nothing earlier in this book says one bin is too
few. So before I'd let myself say two is the fewest, I tried one.

Watch what a program on one bin can see. While the bin isn't empty, it can't tell how full it is: a
stretch of lines that never finds the bin empty does exactly the same thing from a fuller bin, with
the extra pebbles riding along. So the only news a one-bin program ever gets is the bin running dry.
Between two of those moments it's just walking its own lines, and among more moments than there are
lines, two are on the same line, which is the pigeonhole again. Put those together and a one-bin
program either stops within a count you can write down in advance, or never stops at all.

\begin{theorem}[One counter is too few]\label{thm:one-counter}
(1) A program $P$ of INC, DEC and HALT instructions on one bin, started with $n$ pebbles in it, halts exactly when it has stopped
within $k(n + k + 1)$ steps, where $k$ is the number of lines in $P$. So whether it halts is
decidable. (2) Halting for the book's two-counter machine is undecidable. (3) So no computable
translation takes each two-counter program and start to a one-counter program and start that
halts exactly when the original does.
\end{theorem}

\begin{proof}
(1) If $k=0$, the program has already stopped. Otherwise suppose it is still running
after $T = k(n+k+1)$ steps, so the $T+1$ configurations at times $0, 1, \ldots, T$ are all live.
If the counter has stayed at most $n+k$ up to time $T$, those live configurations have only
$k(n+k+1) = T$ possibilities: $k$ lines and $n+k+1$ counter values. There are $T+1$ of them, so
one repeats, and the deterministic run loops forever.

If instead the counter reached $n+k+1$ by time $T$, take the first time it did, and look back to
the last time before that when it was at most $n$. In between it stayed above $n$, so above 0. It
moves at most one per step, so on the way up it arrived at each of the $k+1$ levels
$n+1,\ldots,n+k+1$; mark the first arrival at each. There are $k+1$ arrivals and $k$ lines, so two
arrivals have the same line $\ell$, at levels $a<b$. Every DEC in the stretch between them found
the counter positive. Start that same stretch at level $b$: every counter value is larger by
$b-a$, so still positive, every DEC goes the same way, and the run comes back to line $\ell$ at
level $2b-a$. Repeat forever. Either way a run still live at time $T$ never stops, so a run that
stops does so within $T$ steps, and running $P$ for $T$ steps decides it. (2) The tape simulation,
prime packing and halting diagonal (Appendix~\ref{sec:counter-constructions}), carried over to the
book's instruction set line by line. (3) Such a translation, followed by the test in (1), would
decide two-counter halting, against (2).
\end{proof}

So two counters are the fewest in exactly this sense: one counter can't compute everything, because its halting can be decided, and two, handed their inputs packed, can. I like where the answer came from. The pigeonhole that left no spare chair in the closed room is the same one that says why the box has two bins.

\subsection{The lift and its instruction interface}

The lift gives each box with a universal base an earned panel. In the other direction,
every certified Thiele-complete machine supplies a universal base through its own base
moves. The obstruction results also identify two kinds of base that never lift. Together
these describe the coverage proved here. Two further classification questions concern
bases outside those classes and the possibility of certification through an interface
other than the lift's own.

The tape-to-counter construction (Appendix~\ref{sec:counter-constructions}) connects
bases through what they compute. The lift then supplies the record operations. Its tally
is one edition number for the whole box, whereas the small machine keeps an edition per
counter. That coarser tally is enough for the lift's arguments. To run the lifted machine
inside the presented-machine interpreter, supply the four computability relations for
its chosen presentation.

The construction lets you follow the box and the panel separately: the base provides its
computing moves, the panel earns its record, and the live projection recovers the box's
working run.


%===========================================================================

\section{Running it}
\label{sec:running}
%===========================================================================

Everything so far is definitions and proofs about them. Does any of it run? This appendix is
me finding out. I built the small machine twice, once straight out of the definitions and once
by hand from the same description, ran both, and compared the whole state after every step.
Then I ran the universal programs on guests small enough to finish, and real programs through
the part of the compiler that executes.

One rule holds for the whole appendix: the runs test the builds, and the proofs carry the theorems. A
run shows what a build did on the inputs it was given, and no theorem in the book leans on one: every theorem in it stands on its argument, on the page.

There are two counts to keep apart as you read. The ledger counts what the price list charges. The
number of interpreter steps counts the work this particular build does to carry the run out.
One guest step can take millions of host steps and add one to the ledger.

\begin{center}
\small
\begin{tabular}{@{}p{.27\linewidth}p{.63\linewidth}@{}}
Stage & What the next stage must preserve \\ \hline
Source orders & Variable values and completed procedure calls \\
Assignments and jumps & The same branches, loops and final answer \\
Counter blocks & Values at block boundaries; spare counters restored \\
Host instructions & The same result, with the opposite DEC convention translated \\
Packed two-counter guest & The result as an exponent in the packed counter \\
Fixed interpreter $U_P$ & Guest counters at matching points; decode the exponent to read the answer \\ \hline
\end{tabular}
\end{center}
This is a path through the compiler, not a claim that one step at each stage
takes the same time. The source compiler inserts no certification chain. Running the executable uses physical hardware and energy. These runs did not measure
heat per host step, and the ledger is not a heat measurement. The times only
tell you which runs are practical on this build. The table shows what ran and what didn't.

\begin{center}
\begin{tabular}{@{}p{0.42\linewidth}p{0.5\linewidth}@{}}
\textbf{What} & \textbf{Status} \\ \hline
The small machine & run: the executable versions agree on every tested step \\
The host machines, with and without PAY & run, compared the same way \\
$U$ and $U_P$ on tiny guests & run, against a separate guest implementation \\
$U$ on a guest that certifies & proved and counted, not run \\ $U_P$ on a guest that certifies & proved, not run \\
Source programs through the compiler & run as far as the host program \\
The two-counter guest and $U_P$ running it & proved, not run \\
The presented machine carried by $U_P$ & proved, not run \\
\end{tabular}
\end{center}

\begin{theorem}[The compiler is correct]\label{thm:running-compiler}
For every source program with no recursion, and every input: the program computes $y$ exactly when
the runner (the bounded evaluator of the compiled host program) halts with $y$ for some amount of fuel; exactly when the compiled host program
halts with $y$; exactly when the two-counter guest halts with $y$; and exactly when $U_P$ running that
guest halts with $y$, with its first counter at 0, its ledger at 0 and its flag down. So a terminating compiled program delivers its answer and certifies nothing. To certify a
claim about that answer, a target program must add a CHECK, a COMMIT and a CERTIFY;
this compiler does not insert that chain. The runner also
takes exactly as many steps as the host program does, with the rest of its fuel left over, and a run
that doesn't halt has used all its fuel.
\end{theorem}


%===========================================================================
\section{The CHSH check}
\label{sec:chsh}
%===========================================================================

The example machine checks whether a number is zero, even, or at least something. That's a
fine way to show the chain, and a modest thing to certify. This appendix builds a check worth earning
from zero, and then shows the one thing the small machine's host does with it, which is smaller and more
exact than ``it checks for quantum behaviour,'' a sentence that's tempting and wrong.
In brief, so you know where this is going: the check never computes the famous
number and never writes down the famous bound. It tests seven facts about integers in eight
counts. If they pass, a certain five-by-five table has a property called positive
semidefinite, and that property forces the famous number under the famous bound. The wall
is a consequence of the check, and it never goes into it. Everything else is explaining those sentences.

Why put it in the book? Because it pulls four things apart that the small examples keep
bunched together: the counts, the relation the counts satisfy, the act of checking that
relation, and the record kept of the act. Writing a tally down costs nothing on this
machine. The checker decides a real condition about a matrix, using nothing but integer
arithmetic. The step ties the passing check to this edition of the tally, and the earned
chain costs three, the same three it costs for ``$A$ is zero.'' How deep the condition
is and what the three acts cost are different numbers, and here you can see them apart.

And the tally on its own carries no history of where it came from. A tally somebody made up
passes the same arithmetic as one that was played. What's claimed here is about the numbers
and what the machine does with them, and the guarantee paragraph below says exactly how far
that goes.

\subsection{The game}

There are two players, Alice and Bob, in separate rooms. A referee flips two independent fair coins, independent also of the players' plan,
and hands Alice a question x, zero or one, and Bob a question y, zero or one. Alice answers a, zero or one;
Bob answers b. No talking during the game. They may agree on a plan beforehand, as detailed
as they like. They win a round when their answers
\emph{differ exactly when both questions were one}. In words: if either question is zero, match; if both are one, differ.

Try to win every time. A plan is a list: what Alice says to x equals zero, what she says to
x equals one, and the same two for Bob, four bits in all. Now write the four requirements the
rounds impose: a0 equals b0 (both asked zero); a0 equals b1; a1 equals b0; a1 differs from
b1. Chain the first three: a0 equals b0 equals a1, and a0 equals b1, so b0 equals b1. But
the fourth says a1 differs from b1, and a1 equals b0 equals b1, a contradiction, so one of the
four must break. So any plan loses at least one of the four question pairs, and if the
referee asks all four equally often, the best a plan can do is three out of four, seventy-five percent. Do it the long way too: there are sixteen plans (four bits), and if you write
them all out and count wins you'll find eight win three rounds and eight win one, and none
wins four (the broken requirements always come in odd numbers, which you can see by adding
up the four equations and keeping only whether the totals come out even or odd, with ``equal'' as 0
and ``differ'' as 1: each of the four answers appears twice on the left, so the left sides add to an
even number, and the right sides add to one). Sharing a
coin before the game to pick a plan at random doesn't help; a mixture of numbers no larger
than three-quarters is no larger than three-quarters. Arthur Fine, a philosopher of physics at the
University of Washington, much of whose work is on the physicist Albert Einstein and the quantum theory, proved
in 1982 that the
whole ``pre-agreed answers, no influence between rooms'' story is \emph{exactly} the mixtures of
those sixteen lists, one scoreboard for all four questions at once; the sixteen are the
whole classical story. That one I don't have to take on his word: the form of his theorem that
speaks of the four averages is proved in this book, and the end of this appendix says where.

\begin{center}
\small
\begin{tabular}{@{}ccccc@{}}
Alice's question & Bob's question & Both answer 0 & Required relation & Result \\ \hline
0 & 0 & equal & equal & win \\
0 & 1 & equal & equal & win \\
1 & 0 & equal & equal & win \\
1 & 1 & equal & different & lose \\ \hline
\end{tabular}
\end{center}
Changing the four planned answers can move the loss to another question
pair. It cannot remove every loss: the first three equality requirements
already force the fourth equality.

Now the number itself: recode answers zero and one as plus one and minus one. For each question pair,
let E be the average, over many rounds, of Alice's answer times Bob's. The CHSH quantity is
S equals E00 plus E01 plus E10 minus E11. For one fixed plan each win adds one to S and each
loss subtracts one, so S is wins minus losses: a three-win plan has S equal to two, a one-win
plan has minus two. With equal question frequencies, the win rate is one half plus S over
eight. So ``seventy-five percent'' and ``S at most two'' are the same fact, and the second is the
one that got the name.

The ceiling is easy to reach and costs nothing to write down. Tally one round per question
pair with both players answering zero every time: after the recoding every answer is plus one,
so all four correlators are one. That's the plan that answers zero to everything,
which matches on every pair and so wins the
three pairs that ask for a match, $S = 1 + 1 + 1 - 1 = 2$. On the small machine's host a tally like
that is a few INCs on counters, and INC costs nothing. Nothing in it says the rounds were played; anybody can
tally anything at no charge on the ledger. The other side, that no finite mixture of plans beats two, is
Theorem~\ref{thm:chsh-fine} in Appendix~\ref{sec:chsh-derivation}, one direction of Fine's result. The
full equivalence, in the form that speaks of the four averages, appears in the last subsection of this appendix.

\subsection{Where it came from, and why anyone cared}

So who started this, and why would anybody care whether two rooms can win a coin game? The names
come in a line, and the line didn't mean much to me until I saw that each of them moved the same
thing a notch: from an argument about what's real, which you can only agree with or not, toward a
pile of counts a referee can judge. Each one rests on a different piece of the game above.

The trail starts in 1935, with Einstein, by then at the Institute for
Advanced
Study in Princeton, and two younger colleagues there, Boris Podolsky and Nathan Rosen, asking whether quantum mechanics describes
everything that's real about a pair of particles. Their argument combines certain prediction with a no-disturbance premise:
if measuring one particle lets you predict the other's result without
disturbing it, their criterion treats that result as an element of reality.
They argue that the quantum description leaves something out. What they rest on is the plan in each player's pocket, a pre-agreed list of
answers, the plan above. And the first thing worth checking is whether the paper even says what
Einstein meant, and by his own account it doesn't quite. Podolsky did the writing, and that June Einstein
wrote to Erwin Schrödinger, the Austrian physicist whose 1926 wave equation is at the heart of quantum mechanics, that the essential thing had been ``smothered by formalism.'' The person who
lays that letter out is Arthur Fine, the same Fine whose theorem closed the sixteen plans above, in
his encyclopedia entry on the argument. So Fine stands at both ends of the line: he checked the
founding paper against what its author meant, and he proved exactly what the plans can do. Even the
founding paper is a lead to check, not a verdict.

For thirty years that was philosophy. In 1964 John Bell turned it into arithmetic: if each particle
carries pre-agreed answers and nothing travels between the rooms, certain averages can't exceed a
fixed number, and quantum theory predicts they do. What that rests on took me a few reads. For this game the assumptions can be said without that shorthand: the
questions are independent of the shared preparation, and each room produces
its answer using its own question and that preparation, with any local
randomness included. Even stochastic response rules under those conditions
are mixtures of the sixteen plans. Giving up predetermined answers alone
therefore does not evade the bound. What fails is this local explanation
under the stated independence condition. And where does a
paper like that get published? In a brand-new journal called \emph{Physics Physique Fizika}, which
charged authors no page fees and in fact paid them by the word, ten cents a word, and which folded in
1968 after four volumes, or
three. I went looking for that count, and the sources don't agree on it: four is the
number the American Physical Society's archive of the journal gives, and Iulia Georgescu, writing
about the journal in \emph{Physics World} in 2026, says three. She's also where the ten cents comes
from. As Andrew Whitaker tells it in his biography of Bell, the journal gave out no free copies,
so Bell spent his payment buying copies of his own paper to mail to other physicists. A result
nobody else holds a copy of hasn't left the room yet, and Bell paid to send his across. In 1969 four
physicists, John Clauser, then a graduate student at Columbia, Michael Horne and Abner Shimony at
Boston University, and Richard Holt at Harvard, rewrote it as the four-average form above, S
at most two, which could be tried in a lab with imperfect equipment. It's called the CHSH inequality, after their initials. That notch rests on what a real lab can
deliver. Shimony was in philosophy and in physics at Boston University both, which is the whole line
in one man.

Then the obvious next question: did anybody actually try it? Stuart Freedman, a graduate student at
Berkeley, and Clauser, by then a postdoc there too, tried it in 1972. The physicist Alain Aspect's group at the Institut d'Optique in Orsay, with
two
colleagues there, Jean Dalibard and Gérard Roger, did it in 1982 with the measurement settings switched while the
photons were in flight. Their switches changed the analyzer faster than a
light signal could cross between the stations. They operated at fixed,
incommensurate frequencies, however, rather than choosing fresh independent
random settings on every trial. Rapid switching and unpredictable questions
are separate requirements; this experiment did not close every loophole. Before that, in 1975, Aspect
went to see Bell at CERN to tell him what he planned, and by Aspect's own account Bell heard him out
in silence and then asked, first thing, whether he had a permanent position. Aspect did. Bell's
reason, as Aspect tells it, was that most physicists would take the experiment for a waste of time,
or take him for a crackpot. I read that as the toll in its oldest form: a check that could have come
out the other way costs somebody something to run, and Bell wanted to know who'd pay. And in 2015
three groups of experimental physicists (Bas Hensen's in Delft, with electrons 1.3 kilometres
apart; Marissa Giustina's in Vienna; Lynden Shalm's at NIST, the National Institute of Standards and
Technology, in Boulder) closed the two main escape routes at once. In 2022 Clauser, Aspect and the Austrian
physicist Anton Zeilinger shared the Nobel Prize ``for experiments with entangled photons,
establishing the violation of Bell inequalities and pioneering quantum information science.'' The
2015 tests came out of that line of work, and Zeilinger was one of the authors of the Vienna one. So
in under fifty years the subject went from a question you needed a permanent job to risk to a Nobel.
The lab number sits above two. A remaining escape, that the referee's coins were secretly correlated
with the particles, is a philosophical one and I'm not going to pretend it's closed.

Now put this book's check on the same line, because that's where my position gives. It sits past the
far end, at the counts. Those fifty years went into making the counts worth judging, rooms really
apart, questions really late, escape routes really shut, and my check starts after all of it. It's
handed eight counts and can't tell a lab's tally from one somebody wrote down at no charge. So the check is a referee of counts, and everything that made these
counts worth refereeing happened upstream of it.

What quantum play does: if the players share an entangled pair and choose what to measure
from their question, the best they can do is S equals two times the square root of two,
about 2.83, a win rate of about eighty-five percent. Boris Tsirelson, a Soviet mathematician who later taught at Tel Aviv University, in a 1980 paper called ``Quantum generalizations of Bell's inequality,'' proved that that's
a \emph{ceiling}: quantum correlations violate Bell, but only up to there. I prove that ceiling
below, and show the state that reaches it, with real amplitudes and with complex ones. The bound holds in arbitrary
dimension, including infinite dimension; the construction attaining it needs
sufficient room, supplied by the two-qubit example below.
The first thing proved is the algebra.

Then there's the wall above that, which is the one I find most instructive. In 1994 the physicists
Sandu Popescu and Daniel Rohrlich asked: what if you only demand that nothing can be \emph{signalled}, that Alice's
answers on their own have the same odds whatever Bob is asked, and the other way round?
Then you can write down a ``box'' (the PR box, after their initials) that wins \emph{every} round (its answers differ exactly when
both questions are one, by fiat, each answer alone a fair coin), S equals four, and nothing
about it lets either room send a message. So no-signalling on its own doesn't explain the quantum
ceiling; the PR box meets that condition and beats 2.83. (A box isn't a physical theory, and
nobody is claiming one that meets everything relativity asks.) The order is: two for plans, 2.83 for
quantum, four for no-signalling.

And the box has the same kind of history as everything else in this book: somebody got there
first. Popescu and Rohrlich's initials stuck, but the big review of the field by Nicolas Brunner
and four coauthors credits boxes like it to Peter Rastall, whose 1985 paper on locality and Bell's theorem appeared in \emph{Foundations of Physics},
and, separately, to Leonid Khalfin, a physicist at the Steklov Institute's Leningrad branch, and Tsirelson himself, both in 1985, nine years before. Then the question that really matters: if
relativity isn't what stops things at 2.83, what is? People have tried candidates. In 2005 the
computer scientist Wim van Dam, who works on the theory of quantum computing, argued, in a paper on
the implausible consequences of such boxes, that with enough PR boxes, Bob could work out any yes-or-no function of both
their inputs after Alice sends him a single bit, which would make the whole question of how much
two parties need to talk collapse, and that's a strong hint the boxes are too good to be real.
In 2009 Marcin Pawłowski, a physicist at the University of Gdańsk, and five coauthors proposed a principle they called information
causality: Bob can't get more information about Alice's data than the bits she actually sends
him. They argue that ordinary and quantum physics both obey it, and that any theory that beats 2.83
breaks it. How much of the rest of the quantum
boundary it
reaches is something I haven't checked. Which principle fixes the
whole quantum boundary, if any one does,
this book won't answer. The
check below is a fence at the middle step.

\subsection{Five sentences that are easy to say by accident}

Each one comes below with what I can stand behind
beside it.

``Faster than light.'' The correlations alone cannot carry a chosen message
between the rooms: each room's marginal probabilities do not depend on the
other question. That operational statement does not settle every account of
what happens between measurements. ``Quantum means random.'' The state and
the chosen measurements determine the outcome probabilities and expected
score; a particular outcome can be certain or random. ``A Bell violation
proves nonlocality.'' A statistically supported violation rejects the local
response model under the setting-independence and experimental assumptions.
Calling it a rejection of ``realism'' does not identify those assumptions.

``A tally above two
proves nonlocality.'' No: S is a \emph{population} average, and one finite tally can land above two
by luck. Four rounds, one per question pair, with coin-flip answers, reach S equals four one
time in sixteen. That gap between a population bound and a finite sample is the problem
Yanbao Zhang, Scott Glancy and Emanuel Knill, at NIST in Boulder, work through for real Bell tests (Physical Review A, 2011 and 2013),
and the book doesn't solve it. The way they solve it is worth knowing, because it has the shape
of this book's record. Turn each round into a small bet that no local plan can expect to win,
even a plan that remembers every earlier round, multiply the bets as the rounds come in, and
the French mathematician Jean Ville's inequality, from his 1939 thesis, one of the works that
started the theory of martingales, says the chance the product ever climbs past one over $\delta$
is at most $\delta$. So a flag can go up for good the moment the product crosses, with its error
stated up front. It needs one thing the small machine doesn't have: settings drawn underneath the
program, where the program can't choose them. And how much can one round tell you? Wim van Dam, Richard Gill, a mathematical statistician who later held
the chair in that subject at Leiden, and Peter Gr\"unwald (2005) measured exactly that, in bits of evidence against every local plan. Of the
two-party proofs they compare, CHSH comes out the strongest, ``much stronger'' than the others in
their words, at about 0.046 of a bit per round, a few hundredths, so a convincing record takes
hundreds of rounds, and four won't do. About a finite tally, this book proves only the narrow thing, that a tally above two, with every question pair asked at least once, is not explained by \emph{one fixed plan replayed every round}, and that narrow thing is what
the sentence has to stay. The bound for a finite local model is a separate theorem, and it says
nothing about one tally. And ``the PSD test proves the data is
quantum.'' (PSD is short for positive semidefinite, the property of the five-by-five table from
the opening of this appendix; the next subsection builds both.) It doesn't, and the direction matters. Passing
implies the five-by-five table, pinned in a way the next subsection explains, is PSD, and that
implies S squared is at most eight; both are proved. The
correlators of a PSD table like that are the correlators of some quantum strategy, by
Tsirelson's representation theorem, which is proved here: the strategy is built on two systems of four
levels each, or two levels each with complex amplitudes. The exact edge of the quantum set
for this game was found by Lawrence Landau, whose 1969 Berkeley thesis was on the quantum theory of local observables, and Lluís Masanes, then at the University of Barcelona (next subsection), and that's proved here too. Failing says nothing about whether the play was quantum: the check is complete for
the pinned matrix, the quantum set is wider than that, and quantum behaviour can fail it. Every one
of the sixteen plans fails it, and I'll show you why. What it guarantees is in one paragraph, once
the check is on the host machine.

\subsection{The five-by-five, and what positive semidefinite means}

The mathematics the check rests on is one idea: a table of numbers
\emph{could be a table of dot products}. Take a two-by-two symmetric table with a and c on the diagonal and b
off it. It's positive semidefinite (PSD) when a and c are at least zero and a times c minus
b squared is at least zero. Check two by hand: ones on the diagonal and one half off, the
determinant is three quarters, passes; ones on the diagonal and 1.2 off, the determinant is
negative, fails. The reading: PSD means the table could be the dot products of some arrows,
and 1.2 is too big a dot product for two arrows of length one. Everything else is that idea,
generalised to bigger tables.

Now the table. Rows and columns labelled I, A0, A1, B0, B1 (the identity, Alice's two
questions, Bob's two). Diagonal all ones. The I row and column zero off the diagonal (that's ``zero
marginals'': the table takes each answer alone to average to zero, a fair coin; the counts
record no marginals and the check doesn't test this). A0 with A1 zero, B0 with B1 zero (no round asks
Alice both questions, so no tally holds that average, and the table sets it to zero). Those
zeros are choices. They pin
the table to one slice, and I'll come back to what they cost.
And the four Alice-Bob cells hold the four correlators E, which leaves a five-by-five with four free
numbers in it. It splits into the one in the corner and a four-by-four block made of an
identity and the two-by-two grid E of correlators, and PSD of that block is the same as
saying \emph{E shrinks every vector} (E is a contraction), which is three inequalities: one minus
E00 squared minus E10 squared at least zero; one minus E01 squared minus E11 squared at least
zero; and their product at least the square of (E00 E01 plus E10 E11).

Run it. All four correlators 0.6 except E11 at minus 0.6, S equals 2.4: the two diagonal
terms are 0.28 and 0.28, the cross term zero, passes; above two, below 2.83. All four at one
over root two (E11 negative): diagonals zero, cross term zero, passes with no slack, S equals
two root two exactly, the 2.83; nudge any entry's size up and a diagonal goes negative. The all-ones plan (one, one,
one, one): first diagonal is one minus one minus one, negative, \emph{fails}, though S equals two
and it's a perfectly legal classical play. The PR box (one, one, one, minus one): fails the
same way. The check is a fence at the middle step, and a fence that keeps out every one of the sixteen
plans: each has all four correlators at plus or minus one, so the first diagonal is always
one minus one minus one. Of the classical plays, only diluted mixtures of plans get through.

That point, one over root two in size in every entry, is four real numbers, and no tally of rounds
ever lands on it. A tally's correlators are ratios of whole numbers, and the point needs each of
them to be one over root two, which would make root two a fraction, and you killed that fraction
with a pencil in Section~\ref{sec:structure}. That point and its mirror image, every sign flipped, are the only places
the bound is reached (Remark~\ref{rem:chsh-equality} in the appendix). So a passing tally always has
S squared strictly below eight. It can get as close as you like: the (17, 3) tally in the five
examples below has S squared equal to 7.84, and tallies of the same shape with a fraction nearer one
over root two get nearer still. Two root two is the least bound the check can promise, and no tally
the check accepts ever reaches it. Hold on to that corner; it comes back.

Picture what quantum play can reach as a field, and every test as a fence. In 2007 Miguel Navascués,
Stefano Pironio and Antonio Acín, in a paper called ``Bounding the Set of Quantum Correlations'' (NPA, after their initials),
came at the edge from outside, resting on necessary conditions. In their terms a table like this one
is the first rung of a ladder: bigger tables of expected products, each PSD, each a necessary
condition for quantum behaviour when its entries are left free, each rung tighter, and (their 2008
paper) converging to the quantum set (strictly, the set from commuting operators, which can differ
from the tensor-product set in general scenarios). Each rung is a fence with the whole field inside
it. The check's table is the first rung with the marginals and the two within-party entries pinned
to zero. Pinning is a stronger demand than the rung makes, so passing puts you inside the first
rung, and a quantum point can still fail. My fence cuts across the field. Lawrence Landau in 1988
and Lluís Masanes in 2003 came at the edge itself, for this two-questions-each case only. I'd have
drawn a flat wall at 2.83. The exact boundary they found is a curved surface (the arc-sine form)
touching 2.83 at one point. It's proved here: a table has a PSD completion exactly when its
entries lie between minus one and one and their arcsines pass four inequalities, so those are exactly the quantum tables.

So the obvious next question: if every rung is tighter than the last, does some rung finally land
exactly on the quantum set and stop? My first guess was yes, at least here. A 2026 preprint by Anubhav
Chaturvedi, a quantum-information theorist at the University of Gdańsk, says no: no \emph{finite}
rung of the ladder captures the whole quantum set even in this simplest case. And for the four correlators alone, a 2008 paper by the same Navascués, Pironio and Acín works out that the first rung already lands on Landau and Masanes's curved edge. That looked like a
contradiction until I saw what each one fences. Landau and Masanes draw the edge for the four
correlators alone. Chaturvedi's is the full behaviours, each room's own odds included. So there's
the corner again from the other side: tallies never reach it, and, if his proof holds, rungs never close on the field. I
wanted that to be one fact, but it's two. Mine is root two again, about which points a count can hit.
His is about the fences. His result backs, from
outside, the caution this book keeps about the ladder, that
no finite rung may close on the whole
quantum set. It says
nothing about my own check, though. My check's caution, that a quantum point can fail it, comes from
the pinning, and his paper never pins anything.

\subsection{The check, in integers}

There are eight counts: for each of the four question pairs, one counting rounds where the answers
matched and one where they differed. A round is one bump to one count; nothing
about the round
is tested, and no marginals are recorded. Then the check reads the
eight and tests seven integer facts. For each pair, let N be same plus different and D be same
minus different (so E is D over N). The seven: the four N are strictly positive (every pair was
sampled; a pair with no rounds is rejected outright, and never read as a correlator of zero); A, which is N00
squared N10 squared minus D00 squared N10 squared minus D10 squared N00 squared, is at least
zero; B, the same with the other column, is at least zero; and with C equal to D00 D01 N10
N11 plus D10 D11 N00 N01, C squared is at most A times B. Those are the three contraction
inequalities with the denominators cleared, and there's no division in them, no reals, no square root and no S.
Try it: counts (same, different) of (4, 1), (4, 1), (4, 1), (1, 4) give E of 0.6, 0.6, 0.6,
minus 0.6; N is five and D is three (minus three for the last); A is 175, B is 175, C is zero;
passes. These counts are the first example above, the S equals 2.4 one, written in integers.

\begin{center}
\small
\begin{tabular}{@{}ccccc@{}}
Question pair & same & different & total & difference \\ \hline
00 & 4 & 1 & 5 & 3 \\
01 & 4 & 1 & 5 & 3 \\
10 & 4 & 1 & 5 & 3 \\
11 & 1 & 4 & 5 & $-3$ \\ \hline
\end{tabular}
\end{center}
First check the four totals: all are positive. Each column's remaining
allowance is $625-225-225=175$. The cross term is $225-225=0$.
Thus the final comparison is $0\le175^2$, and the check passes. These
steps verify the arithmetic; they supply no evidence that rounds were played.

The theorem chain, in words: a passed check gives PSD of the five-by-five built from the
counts; contraction and PSD are equivalent, both directions; and PSD forces S squared at most
eight. The names are in the statements below. The underlying algebra is one identity: four times
the sum of the four correlators squared, minus S squared, is a sum of six squares, so if the row sums of squares are at most one, S squared is
at most eight, and the point one over root two reaches it. So Tsirelson's bound comes out as \emph{algebra}, with no Hilbert space (the
mathematical home of quantum states) in it.

The chain goes below as three statements. The derivations are in Appendix~\ref{sec:chsh-derivation},
and the statements say what they prove.

\begin{definition}[Correlators from the counts; the pinned matrix]\label{def:u23-pinned}
For eight counts and a question pair $xy$, let
$E_{xy} = (\mathrm{same}_{xy} - \mathrm{diff}_{xy}) / (\mathrm{same}_{xy} +
\mathrm{diff}_{xy})$, read as 0 when both are 0. For reals $E_{xy}$,
$M(E)$ is the symmetric five-by-five matrix on $I, A_0, A_1, B_0, B_1$ with ones on the
diagonal, $E_{xy}$ in the cells $(A_x, B_y)$ and $(B_y, A_x)$, and zero everywhere else
(Definition~\ref{def:chsh-five}). $M$ is \emph{PSD} when $v^{\mathsf T} M v \ge 0$ for every real
vector $v$, and $S(E) = E_{00} + E_{01} + E_{10} - E_{11}$.
\end{definition}

\begin{theorem}[The integer check implies the pinned matrix is PSD]\label{thm:u23-notrap-psd}
If eight counts pass the seven integer facts, then $M(E)$ for their correlators is symmetric
and PSD.
\end{theorem}

\begin{proof}
The four positivity facts make every denominator positive, so each $E_{xy}$ is $D_{xy}/N_{xy}$,
and dividing the other three facts by positive squares turns them into the three contraction
inequalities (Lemma~\ref{lem:chsh-clear}). Contraction is equivalent to PSD of $M(E)$ by a Schur
complement on the four-by-four block (Corollary~\ref{cor:chsh-iff}).
\end{proof}

In the appendix this is Theorem~\ref{thm:chsh-check-sound}, and Corollary~\ref{cor:chsh-iff} gives the equivalence of contraction and PSD in both directions.

\begin{theorem}[PSD forces the Tsirelson bound]\label{thm:u23-psd-tsirelson}
For all real $E_{00}, E_{01}, E_{10}, E_{11}$: if $M(E)$ is symmetric and PSD, then
$S(E)^2 \le 8$.
\end{theorem}

\begin{proof}
Call the four entries $a,b,c,d$, read across the rows. PSD means that
no choice of five weights makes the matrix's weighted total negative.
Choose weights $(0,1,0,-a,-b)$ in the order $(I,A_0,A_1,B_0,B_1)$.
The diagonal contributes $1+a^2+b^2$, and the two cross terms contribute
$-2a^2-2b^2$. The total is $1-a^2-b^2$, so $a^2+b^2\le1$.
Weights $(0,0,1,-c,-d)$ give $c^2+d^2\le1$. All four squares therefore
add to at most two.

Now compare four times that sum with the squared score. Their difference is
\[
\begin{aligned}
4(a^2+b^2+c^2+d^2)-(a+b+c-d)^2
={}&(a-b)^2+(a-c)^2+(a+d)^2\\
 &+(b-c)^2+(b+d)^2+(c+d)^2.
\end{aligned}
\]
To check it, each squared letter occurs three times on either expanded
side; the cross terms are $-2ab,-2ac,+2ad,-2bc,+2bd,+2cd$ on both.
The difference cannot be negative because it is a sum of squares.
Thus the squared score is at most four times two, which is eight.
\end{proof}

This is Theorem~\ref{thm:chsh-psd-tsirelson} in the appendix.

\begin{corollary}[A passed check means the bound]\label{cor:u23-notrap-bound}
If eight counts pass the seven integer facts, their correlators satisfy $S^2 \le 8$. In fact
$S^2 < 8$, and passing tallies come as close to 8 as you like.
\end{corollary}

\begin{proof}
Theorem~\ref{thm:u23-notrap-psd} makes the pinned matrix of these four correlators PSD, and
Theorem~\ref{thm:chsh-psd-tsirelson} in the appendix, which holds of every such matrix, gives
$S^2\le8$. The score is a ratio of whole numbers, and no such ratio
squares to 8, so the bound is strict. The tallies whose four correlators are
$n/q, n/q, n/q, -n/q$ with $2n^2 \le q^2$ pass, and their scores $4n/q$ come as near
$2\sqrt 2$ as you like.
\end{proof}

\begin{remark}[Sound, and incomplete]\label{rem:u23-sound-not-complete}
The pinned matrix is stricter than the quantum question it is usually read against. PSD of
$M(E)$ gives the correlators a PSD completion, since the zeros are one choice of the two free
entries. The converse fails.
The all-ones plan has a PSD completion, and every
deterministic plan is a quantum strategy; yet its
minor on rows $A_0, B_0, B_1$ is $1 - 1 - 1 = -1$, so $M(E)$ is not PSD there, and a tally of
one matching round per pair gives $A = 1 - 1 - 1 = -1$ in the integer check, so the check
fails. The second failure is a named theorem on the machine: the tally of one matching round
per pair is refused forever (Proposition~\ref{prop:chsh-demos}). Every deterministic plan fails
the same way. The tally of one round per pair fails the check for each of the sixteen plans, and so does
any tally of one plan replayed over any number of rounds (a pair left unsampled fails the
positivity facts outright), because no deterministic plan's pinned matrix is PSD. So a pass proves the bound, and a failure proves only that the tally's pinned matrix isn't
PSD or a pair went unsampled; it says nothing about whether the play was classical or quantum. The tally from the start of this appendix, written down at no cost, is
exactly such a tally: it reaches $S = 2$ and the check throws it out.
\end{remark}

None of that algebra is new, and I don't want it to read as if it were. Mnacho Echenim
and Mehdi Mhalla, computer scientists at Grenoble, give a detailed treatment of the CHSH
and Tsirelson bounds (2023). Navascués, Pironio and Acín already analyse this very matrix by Schur complement and discuss
the level above it. What I think is mine here is smaller: an integer decision procedure on
the counts, and a proved link from ``the check passed'' to PSD to the bound.

The picture under the check is older still. Tsirelson showed, in a longer paper of 1987 after his
1980 note, that the correlations two separated quantum parties can reach are exactly
the tables of inner products of unit vectors, one vector for each question on each side. With the
marginals pinned to zero the pinned matrix is that picture, and by its Schur complement it is PSD
exactly when the two-by-two table of correlators, read as a matrix, has operator norm at most one.
The physicists Reinhard Werner and Michael Wolf's survey of Bell inequalities (2001) sets out that
correlation-matrix view for two-party inequalities in general.

\subsection{The check, earned on the host machine}

Now put it on the machine. A check on the small machine is a property: a checker, which is a
little program that says yes or no, a meaning, which is a statement, and a proof that the two
agree. The version of the small machine with its property language left open
(Section~\ref{sec:smallest}, Theorem~\ref{thm:u34-tc-instances}) takes any property language with an exact equality
test whose properties come
with those three things, and for every one of them a flag raised from
a clean start stands on a passing CHECK of
one claim, a later COMMIT of the same claim with its counter untouched in between, and a later CERTIFY, the three acts costing one each.
For the game, pack the eight counts into one counter, as a list, the way the
sorted-list machine packs a list into one number (Section~\ref{sec:smallest}). Let the checker be
the seven integer facts read off the unpacked counts. And let the meaning be what the facts are
for: every pair was sampled, and the pinned five-by-five is positive semidefinite. Then the proof
the generic machine asks for, that the checker says yes exactly when the meaning holds, is the
chain above with the denominators cleared in both directions, and
a flag raised on that machine from a clean start
means: these counts, at this edition of their counter, passed the check, somebody committed to
that with nothing changed underneath, and the correlators they give satisfy $S^2 \le 8$.
The theorems below state that exactly, and they're stronger than the chain above in one way I want to say once
and plainly. The earlier theorems said that a passed check makes the pinned matrix PSD. The check is
also complete for it: it passes \emph{if and only if} every pair was sampled and the pinned matrix
is PSD, because clearing the denominators works in both directions. The statement is over the machine of Definition~\ref{def:u34-multi}, the
small machine with a counter for every number,
with PCHSH in place of its one property, $\mathsf{Slot}$. A CHECK reads one counter, which is why the eight
counts have to travel in one.

\begin{definition}[The packed tally, and the property PCHSH]\label{def:chsh-packed}
Read a whole number as a list in the list code (0 is the empty list, and $2^x(2y+1)$ is $x$ followed
by the list $y$), and take as its \emph{tally} the first eight entries, in the order
$\mathrm{same}_{00}, \mathrm{diff}_{00}, \mathrm{same}_{01}, \mathrm{diff}_{01}, \mathrm{same}_{10},
\mathrm{diff}_{10}, \mathrm{same}_{11}, \mathrm{diff}_{11}$, a missing entry reading 0. Every number
is some tally. The property PCHSH has the checker ``the tally passes the seven integer facts'' and
the meaning ``every question pair was sampled, and the pinned matrix $M(E)$ of the tally's
correlators is symmetric and PSD.''
\end{definition}

\begin{theorem}[The check is exact]\label{thm:chsh-exact}
(1) The seven integer facts hold of a tally if and only if every pair was sampled and $M(E)$ is
symmetric and PSD. (2) For every whole number $v$, the checker of PCHSH says yes at $v$ if and only
if the meaning holds of $v$'s tally. (3) Reading the number that packs a tally gives that tally
back. (4) If the facts hold, the score of the tally satisfies $S^2 \le 8$ and $|S| \le 2\sqrt 2$;
in fact $S^2 < 8$ and $|S| < 2\sqrt 2$, and passing tallies come as close to 8 as you like.
\end{theorem}

\begin{proof}
(1) The forward half is Theorem~\ref{thm:u23-notrap-psd}. For the backward half, PSD gives the three
contraction inequalities (Corollary~\ref{cor:chsh-iff}); with every $N$ positive, multiplying them
through by the squares of the denominators gives $A \ge 0$, $B \ge 0$ and $C^2 \le AB$, which is
Lemma~\ref{lem:chsh-clear} read the other way. (2) is (1) with the packing, since every number is a
tally. (3) is the code and its decoder being inverse. (4) is (1) forward, then
Theorem~\ref{thm:u23-psd-tsirelson}; the strict form and the closeness are
Corollary~\ref{cor:u23-notrap-bound}.
\end{proof}

\begin{theorem}[A raised flag stands on the check]\label{thm:chsh-earned}
On the machine above, let $\tau$ be
a trace from a clean start $s_0$
that ends with the flag up. Then,
for some counter $c$, \[ \tau \;=\; \mathit{pre} \cdot \text{CHECK PCHSH } c \cdot \mathit{mid}_1 \cdot \text{COMMIT PCHSH } c
\cdot \mathit{mid}_2 \cdot \text{CERTIFY} \cdot \mathit{post},
\]
and with $v$ the value of $c$ after $\mathit{pre}$: the tally of $v$ passes the seven integer facts;
every pair was sampled and the pinned matrix is symmetric and PSD; its score satisfies $S^2 \le 8$
and $|S| \le 2\sqrt 2$, and in fact $S^2 < 8$, though passing tallies come as close to 8 as you
like; and $c$ still held $v$ when the COMMIT ran. For any stored program run from a fresh start (counters set, every version 0, line 1, the table and channel empty, the trap and the flag down, the ledger at 0)
that raises the flag, the trace it executed contains a CHECK PCHSH $c$ that passed, and the tally
$c$ held when it ran passes the seven integer facts, has every pair sampled and the pinned matrix
symmetric and PSD, and has $S^2 \le 8$ and $|S| \le 2\sqrt 2$. Conversely,
from a fresh start in which a counter
$c$ holds the packing of a tally $t$, the three instructions CHECK PCHSH $c$,
COMMIT PCHSH $c$, CERTIFY raise the flag if and only if $t$ passes the check, and then the ledger
reads exactly 3; if $t$ fails, the trap goes up at the CHECK and no number of further steps raises
the flag.
\end{theorem}

\begin{proof}
A raised flag from a clean start has the earned chain (Theorem~\ref{thm:u34-cert-earned}, which
holds for this machine as for the small one): a passing CHECK, then a COMMIT of the same claim with
the counter's edition unchanged, then a CERTIFY. A passing CHECK of PCHSH at $c$ means the
checker says yes at $c$'s value, which is the meaning by Theorem~\ref{thm:chsh-exact}(2), and the
bound is (4) of the same theorem. An unchanged edition means an unchanged value. For the converse,
the CHECK passes exactly when the tally passes; if it does, COMMIT and CERTIFY find their
facts and the three acts cost three; if it fails, the trap latch goes up, and no later move can raise the flag. A stored program stops at the trap;
an explicitly supplied instruction trace may still incur its defined charges
(as in Theorem~\ref{thm:u34-tight}).
\end{proof}

So the one thing the
machine does with the check, as this appendix's opening said, is that a raised flag can stand on it. Five
tallies, each a theorem with its counts written out, show the check isn't vacuous and where its
edges are.

\begin{proposition}[Five tallies]\label{prop:chsh-demos}
Counts are given as (same, diff) for the pairs 00, 01, 10, 11. (4,1), (4,1), (4,1), (1,4) has
$S = 12/5$: the check passes, the flag rises, the ledger reads 3. (17,3), (17,3), (17,3), (3,17)
has $S = 14/5 = 2.8$, just under $2\sqrt 2$, with $S^2 = 7.84$: it passes, the ledger reads 3.
(9,1), (9,1), (9,1), (1,9) has $S = 16/5$ with $S^2 > 8$: refused forever, whatever follows.
(1,0), (1,0), (1,0), (0,1), the PR box, has $S = 4$: refused forever. (1,0), (1,0), (1,0), (1,0),
the all-ones classical plan, has $S = 2$: refused forever.
\end{proposition}

\begin{proof}
A pair's correlator is $(\mathrm{same} - \mathrm{diff}) / (\mathrm{same} + \mathrm{diff})$, and $S$
adds the first three and subtracts the fourth. $(4, 1)$ gives $3/5$ and $(1, 4)$ gives $-3/5$, so
$S = 12/5$; the seven facts are the worked tally after Corollary~\ref{cor:chsh-check-bound}, with
$A = B = 175$ and $C = 0$, so the check passes, and Theorem~\ref{thm:chsh-earned} gives the flag
and the ledger at 3. $(17, 3)$ gives $7/10$ and $(3, 17)$ gives $-7/10$, so $S = 14/5$ and
$S^2 = 196/25 = 7.84$; with $N = 20$ in every pair and $D = 14, 14, 14, -14$,
$A = B = 400 \cdot 400 - 196 \cdot 400 - 196 \cdot 400 = 3200$ and
$C = 14 \cdot 14 \cdot 400 - 14 \cdot 14 \cdot 400 = 0$, so the seven facts hold and the same
theorem applies. $(9, 1)$ gives $4/5$, so $S = 16/5$ and $S^2 = 256/25 > 8$; the PR box gives
$1, 1, 1, -1$ and $S = 4$; both fail the check, because a tally that passes has $S^2 \le 8$
(Corollary~\ref{cor:chsh-check-bound}), and a failed check traps at the CHECK, after which no step
raises the flag (Theorem~\ref{thm:chsh-earned}). The all-ones plan gives $N = D = 1$ in every
pair, $S = 2$, and $A = 1 - 1 - 1 < 0$, the worked tally's second computation, so it fails the
check and is refused the same way.
\end{proof}

The last one is the
check being stricter than the bound, as a named theorem: a perfectly legal classical play with
$S = 2$ is turned away.

What the theorem doesn't say stays what it was. The tally is \emph{placed}. Nobody plays it on the
machine: the
theorem starts from a counter that already holds the packing, and no instruction plays a round, so a
program can put any code in a counter at no cost. There are no marginals. The pinned matrix is stricter than the quantum set, so failing says nothing about
whether the play was quantum, and there is no margin for sampling noise at
Tsirelson's point.

Writing a tally costs nothing. Each count is a counter, a round is an INC on one of them, and INC costs
nothing, so a pile of trials costs nothing to write down (what it spends is fresh states, in a counter that
never runs out of numbers); the packing into one counter is a program's
business, and in the theorem the packed counter is simply there at the start. What costs is the
claim about the pile. On the ledger, statistics cost nothing; the claim about them is what costs.

So what the check guarantees fits in one paragraph. If it passes, then every question
pair was sampled, the four ratios from the counts satisfy three polynomial inequalities, so the
pinned five-by-five is PSD, so S squared is at most eight. It does \emph{not} guarantee that
the tally came from honest play (a program can INC
whatever counts it likes
at no cost, including four rounds that put S at four). It does not guarantee any
marginal is zero (no marginals are recorded and none are tested). It does not put the players'
state behind the counts. What it does give is a quantum strategy with the same correlators: a
passing tally's correlators have a PSD completion, and
every point with one is the correlator table of a strategy on two four-level systems.
And it fails the tally of every single plan, honest classical play included. It has no margin
for sampling noise either. Honest play at Tsirelson's point aims at correlators of plus or minus
one over root two, and there every pair of squared correlators the check adds up comes to exactly
one, the edge of what it accepts. So an honest tally aimed at that point sits right on the line,
and ordinary luck over a finite number of rounds pushes it over most of the time. I simulated it,
drawing each round's answers at Tsirelson's correlators: the check fails 92 percent of the time at
10 rounds per question pair, 87 percent at 100, and between 85.4 and 85.9 percent from 1,000 to 100,000. In those simulated trials, fair
coins passed every time from 100 rounds per pair on, and failed about 2.5 percent of the time at 10.
I saw that in the trials, and it's no guarantee about any one tally.
The limit has a closed form. At Tsirelson's point the correlator matrix, read as a two-by-two,
turns arrows without stretching them (it's a mirror), and every one of its entries wobbles by the same amount, independently. To first
order the largest singular value of the noisy matrix is one plus the larger eigenvalue of
$(G + G^{\mathsf T})/2$, where $G$ is the wobble turned back by that mirror, a two-by-two of
independent Gaussian entries of equal size,
and the check passes only when that eigenvalue is at most zero. Write that symmetric matrix's
eigenvalues as $u \pm \sqrt{v^2 + w^2}$, with $u, v, w$ independent and of equal spread (scale
them to standard normals). Both are at most zero exactly when $u \le -\sqrt{v^2 + w^2}$. Now
$v^2 + w^2$ has the exponential law with mean 2, so the chance is
$\tfrac12\,\mathbb E\big[1 - e^{-u^2/2}\big] = \tfrac12\big(1 - 1/\sqrt2\big) = (2 - \sqrt 2)/4$.
So honest play at Tsirelson's point fails with probability
$(2 + \sqrt 2)/4 \approx 0.854$ in the long run. The arithmetic is on the page; the first-order
step is not proved here, and the simulation is what backs it.

The check can also allow a tolerance. For $\varepsilon = p/q$ it tests the same seven
facts in integers with the 1 in the contraction inequalities replaced by $(1 + \varepsilon)^2$. At
$\varepsilon = 0$ it is the check above. A pass means the two-by-two correlator matrix has
operator norm at most $1 + \varepsilon$, and $S^2 \le 8(1 + \varepsilon)^2$: divide every correlator
by $1 + \varepsilon$ and the tolerance check becomes the exact one, so the chain above applies to
the shrunken table. A tally of 100 rounds
per pair with correlators $0.72, 0.70, 0.72, -0.72$ fails the exact check and passes at
$\varepsilon = 1/10$. In the same kind of simulation, with
$\varepsilon = 3/\sqrt n$ for $n$ rounds per pair, honest Tsirelson play failed it less than once in
ten thousand trials at every size from 100 to 100,000 rounds per pair: under once in a hundred
thousand at 100, five at 1,000, eight at 100,000, close to the long-run rate the same first-order picture gives, about
nine in a hundred thousand. In that Gaussian calculation the exact expression
is $1-\Phi(6)+e^{-9}\Phi(3\sqrt2)/\sqrt2$, where $\Phi$ is the standard normal
cumulative probability. Dropping its tiny tail corrections gives
$e^{-9}/\sqrt2$. This is the asymptotic calculation under independent sampling,
not a finite-sample confidence guarantee. What the tolerance gives up shows in the bound: a pass caps $S$ at $2\sqrt 2\,(1 + \varepsilon)$,
a little outside Tsirelson's number. The exact check remains the default. With $p$ and $q$
fixed for a machine instance, the tolerance check also runs as CHECK on the same versioned
registers: its CHECK, COMMIT, CERTIFY chain certifies exactly when the integer tolerance
test passes, costs three on success, and traps permanently on failure. From a clean
start, every raised flag has that earned chain behind it, with the same register value at
CHECK and COMMIT, and establishes the relaxed bound above: it's Theorem~\ref{thm:chsh-earned}
again, with the tolerance test as the checker and ``every pair sampled and operator norm at most
$1 + \varepsilon$'' as the meaning, the scaling just described being the proof that they agree.  The check certifies a \emph{shape} of the tally. What the shape means is a question for whoever
supplied the rounds.

And notice which way the fence faces. The most practical use anyone has found for a Bell test,
certified randomness, needs the opposite bound. Stefano Pironio and colleagues (2010) argued, and backed
with an experiment, that a Bell violation, with settings the devices can't choose, certifies that the outputs are random,
and a group of quantum cryptographers, Rotem Arnon-Friedman, Fr\'ed\'eric Dupuis, Omar
Fawzi, Renato Renner and Thomas Vidick (2018), made the accounting tight with what they call entropy accumulation. Under those assumptions, a predetermined local output strategy passes a
properly calibrated test only within its stated false-acceptance probability, so it's the cleanest case I know, outside this book, of a check that can
fail, separating output that was established from output that was printed. That needs $S$
bounded from below, with an error probability. This check bounds $S$ from above and fails every
single plan, so it's a fence that keeps out the super-quantum and the fixed plan. As a
nonlocality test it's no use. The deployed version exists, too: Gautam Kavuri and colleagues, physicists at NIST and
the University of Colorado (2025), run a
public randomness beacon whose outputs each trace back through a Bell test to a public record,
the nearest working system I know to a certified claim written down for good. And records cut
both ways there. Jonathan Barrett, Roger Colbeck and Adrian Kent, working on quantum cryptography (2013), pointed out that a device
which keeps its own memory of earlier rounds can leak it through outputs it makes public later,
so reusing a device can undo a guarantee the user thought was settled. A permanent record in the
wrong box is an attack.

\subsection{A rational stand-in for an irrational number}

One useful rational bound: 5657 over 2000, which is 2.8285, and two root two is
2.82842... It's the four-decimal round-up, chosen so that rational arithmetic can test ``above the
Tsirelson bound'' without square roots, and it's safe because 5657 squared is 32,001,649 and
eight times 2000 squared is 32,000,000, so the fraction sits above root eight by about seven
hundred-thousandths. It's used in a few rational
theorems. It is \emph{not} used in the integer check, and the real-number theorems use root
eight.

\subsection{Two more tests, as mathematics}

The second rung: a nine-by-nine table on I, the four questions, and the four Alice-Bob
\emph{products}, with the correlators and five extra free numbers for three- and four-body
entries. Its two four-body entries are tied, one the negative of the other. That tie is a
constraint, $\langle \{A_0, A_1\} \otimes \{B_0, B_1\} \rangle = 0$: the zero entries of the first
rung don't imply it, and a classical strategy breaks it. So the second rung is level 1+AB on that slice, and
that's what I mean by level 1+AB everywhere in this book. Its PSD is again equivalent to a contraction-like condition.
The check for it proves something stronger, that the matrix is positive definite, by a cascade of
two Schur steps down to a four-by-four that Sylvester's criterion closes.
Sylvester's name on a test about matrices fits in a way I enjoy, because the word
\emph{matrix} is his. In 1850, in the \emph{Philosophical Magazine}, the English
mathematician James Joseph Sylvester used it for an ``oblong arrangement of
terms'' out of which you cut square blocks and take their determinants, and \emph{matrix} was
already an old word for the place something grows from. The test does exactly that:
it cuts the top-left square blocks out of the big grid, one row deep, two, three, four (the
leading minors), and asks each one's determinant to be positive. Was he the first to state the
test? The name doesn't settle that; names stick to people for all kinds of reasons, and I'd
treat this one as a lead.
There's also a fact that's easy to list without noticing what it means: with all five extra numbers at
zero, the condition certifies a strict \emph{subset of the classical set} (the ball where the four
correlators' squares sum to at most one), forcing S at most two; it can't certify anything above the
classical line. To get a point above two you have to supply the extras
(S equals 2.4 with one extra at minus one
half). Appendix~\ref{sec:chsh-derivation} has the algebra.

The two sizes are one test. A single positive-semidefinite predicate takes the size
as a parameter, and the five-by-five and the nine-by-nine are it at two sizes.

\begin{definition}[The size-generic PSD test]\label{def:u23-genrealizable}
A \emph{claim} is a top index $n$ and a matrix $M$ indexed by naturals. It is \emph{realizable} when
$M_{ij} = M_{ji}$ for all $i, j \le n$ and $\sum_{i \le n} \sum_{j \le n} v_i M_{ij} v_j \ge
0$ for every real vector $v$. The CHSH claim of $E$ is top index 4 with the matrix $M(E)$ of
Definition~\ref{def:u23-pinned}; the level-1+AB claim uses top index 8 and the nine-by-nine
matrix.
\end{definition}

\begin{corollary}[The level-1 test is the size-five case]\label{cor:u23-general}
For all real $E_{00}, E_{01}, E_{10}, E_{11}$: the three contraction inequalities hold if and
only if the CHSH claim of $E$ is realizable. The same holds at size nine for the level-1+AB
contraction condition and its claim.
\end{corollary}

\begin{proof}
Top index 4 means the five indices $0,1,2,3,4$. The double sum visits
each of their twenty-five cells once, multiplies it by its row and column
weights, and adds the results. That is precisely the weighted-total test
for the five-by-five matrix. The completed-square calculation in
Corollary~\ref{cor:chsh-iff} equates that test with the three contraction
inequalities. Top index 8 similarly visits the eighty-one cells of the
nine-by-nine matrix. Theorem~\ref{thm:chsh-nine-iff} writes its total as
three squares plus the residual and proves that its test is exactly the
residual condition. The same definition therefore asks the same question
at either size.
\end{proof}

The third thing is the \emph{elliptope}, the set of correlators for which \emph{some} choice of the two
missing entries (A0 with A1, B0 with B1) makes the five-by-five PSD. Every one of the
sixteen plans is inside it, mixtures stay inside, every point inside obeys $S^2 \le 8$, Tsirelson's point is inside and reaches it, the PR box is out, and it's strictly bigger than the classical set. There's a certificate checker for it with a soundness proof, in two
branches: a strict Sylvester check with the completion supplied, and a rational
factorisation certificate of the completed matrix, checked in integers. What's proved is
that a pass puts the correlators inside the set and that no choice of inputs gets the PR box
through, and, by
computation, that some named points pass, the all-ones plan and a point that the arc-sine
boundary puts exactly on the quantum edge among them.
And every point with a rational PSD completion has such a certificate: peel the completed matrix a row at a time, as in Section~\ref{sub:four-vectors}, but keep each pivot instead of taking its square root; a zero pivot forces its row to zero, and every number that comes out is rational. The certificate checker is sound and, on rational completions, complete, as mathematics; none of the small machine's properties is the elliptope check.

\subsection{The quantum set}

Three descriptions of the same set have been running through this appendix: the tables some quantum
strategy can produce, the tables whose five-by-five matrix has a PSD completion, and the tables
inside Landau and Masanes's curved edge. They are the same set, for strategies
on finitely many levels, and the next subsection takes away the ``finitely''. Each piece is taken in turn below, because the constructions are the part
worth seeing.

Quantum tables are inside the completion set. A strategy's correlator is the inner product of two
unit vectors: Alice's observable applied to the shared state, and Bob's applied to it. Four unit vectors always have a PSD table of inner products, so every
quantum table has a completion, and Tsirelson's bound follows from that: $S^2 \le 8$ for every
strategy, and its level-one moment matrix passes the
PSD test with the marginals in. Quantum mechanics uses complex amplitudes,
and the complex case goes through by a plain trick: a complex vector of length $n$ is a real vector
of length $2n$, its real parts followed by its imaginary parts, and the real part of a complex inner
product is the ordinary inner product of those longer real vectors. So the real argument runs
unchanged. The two-qubit
strategy written out in Appendix~\ref{sec:chsh-derivation} reaches $2\sqrt 2$ in both readings.

The completion set is inside the quantum tables. This direction is Tsirelson's representation, and the
construction is his. Take a table with a PSD completion. A PSD matrix is a table of inner
products of vectors, always: peel off one row at a time, and a zero on the diagonal forces its
whole row to zero. So the completion hands over four unit vectors,
$a_0, a_1$ for Alice and $b_0, b_1$ for Bob, with $a_x \cdot b_y$ equal to the table's
entries, and they fit in three dimensions: slide one of the two entries the table leaves open to
the end of the range that keeps the completion PSD, and there the completion is singular, which
puts the four vectors in a space of three dimensions. Section~\ref{sub:four-vectors} has the picture behind the
factorization and the measurements. Now comes the trick. There are three real four-by-four matrices, built from the two-by-two ones that swap
two entries or flip the sign of one, that each square to the identity and anticommute in pairs:
multiply any two of them in either order and the answers differ only in sign. Use a vector's three components to mix
those three matrices, and you get an observable that squares to the identity, a legitimate
measurement. Give Alice and Bob a maximally entangled pair of four-level systems, and the
correlator of two such observables comes out as exactly the inner product of their vectors. So the strategy built from the four vectors produces the table, for both real and complex amplitudes. Four levels on each side is always enough, and with complex amplitudes two are.

The curved edge is the same set again. Each correlator is the cosine of an angle, the angle
between one of Alice's vectors and one of Bob's. Walk the loop $a_0, b_0, a_1, b_1$ and back to
$a_0$: the angle across any one leg can't exceed the sum of the other three, because angles on a
sphere obey the triangle inequality. Flip one of
Alice's vectors to its opposite and walk the loop again, then do the same with her other one, and
between them you get the other four inequalities, with $2\pi$ in them. Written in
arcsines, those eight inequalities are Landau and Masanes's four. Going the other way, any four
angles that pass them can be met by actual vectors: choose an angle between Bob's two that meets both pairs of constraints,
build each of Alice's from the two angles it has to make with them, and check that the determinant that says such a triangle exists comes
out nonnegative. So a table has a PSD completion exactly when its arcsines pass the four
inequalities, and those are exactly the quantum tables. The point $(3/5, 4/5, 4/5, -3/5)$ sits on the edge, because
$\arcsin\tfrac35 + \arcsin\tfrac45 = \tfrac{\pi}{2}$ by the 3-4-5 triangle, and so does Tsirelson's point. The interval construction, including both endpoint cases,
is pictured in Section~\ref{sub:arcsine-construction}.

\subsection{Any number of levels, and one converse that fails}

Nothing here cares about dimension. Quantum mechanics also allows systems with
infinitely many levels, and the arguments don't need finitely many. Whatever the dimension, a shared
state $v$ gives an expectation on the operators the players measure with:
$x \mapsto \mathrm{Re}\,\langle v, x\,v\rangle$. That one line is the whole step. The
proofs use only what that expectation always has: you can add, multiply, scale and take
adjoints, the product spreads over sums, and the expectation of $x^*x$ is never negative. Each
measurement is its own adjoint with the expectation of its square at most one, which covers
unsharp measurements too, and Alice's commute with Bob's. That last demand is weaker than putting
them on two separate systems, so the result covers the wider, commuting model as well. On that
footing Tsirelson's sum of squares goes through, and the CHSH value is at most $2\sqrt2$. The same nonnegative form, with each diagonal entry raised to one,
completes the correlator table to a PSD matrix, so the four arcsine
inequalities hold, and with a unit the full first-level moment matrix,
marginals and all, is PSD. The other
direction needs no help from infinity: every point inside the arcsine edge is reached on two
four-level systems.

Also, the finite-mixture theorem
proves only one direction, that a finite local model satisfies the bound, and its converse as
stated is false: a table with $S = 0$ that no local model produces.
Fine's own theorem, in the form that speaks of correlations, is proved separately as an
equivalence: a table of four correlators comes from a local model exactly when each is between
$-1$ and $1$ and all eight CHSH inequalities hold. Section~\ref{sub:finite-local-construction}
pictures the finite distribution. And one more thing about tallies, which the guarantee paragraph above already used: a tally loaded at no charge
is not thereby classical (the four rounds with S equal to four). The correlators (one, one, one, minus one) are the PR box, which no classical model has as its averages, though coin-flip play lands on that four-round tally one time in sixteen; the classical vertex this
appendix used, which fails the check like every plan, is (one, one, one, one).

That covers the game, the number, the two walls, and the check. The derivations come in full in the
next appendix, Appendix~\ref{sec:chsh-derivation}; what all this costs against physics was
Section~\ref{sec:landauer}.


%===========================================================================

\section{The CHSH derivation, in full}
\label{sec:chsh-derivation}
%===========================================================================

Appendix~\ref{sec:chsh} tells the story; this appendix does the algebra. Each item gives
the statement with its hypotheses and the steps you need to check it. Nothing here says
that a matrix is a physical state. Where quantum strategies come in, I say exactly how they
stand to the matrices.

\subsection{Setting and notation}
\label{sub:chsh-notation}
Start with four piles of tally sheets, one for each pair of settings Alice
and Bob can be given. On each round their answers either agree or disagree.
Give agreement a score of $+1$ and disagreement a score of $-1$. In each pile,
add the scores and divide by the number of rounds. That pile's average is its
\emph{correlator}. It says how strongly the two answers agree for that pair of
settings.

\begin{definition}[Settings, outcomes, correlators]\label{def:chsh-setting}
Alice receives a setting $x \in \{0,1\}$ and Bob a setting $y \in \{0,1\}$. Each answers a
bit, recoded as an outcome $a, b \in \{+1, -1\}$ (bit 0 becomes $+1$, bit 1 becomes $-1$).
The correlator $E_{xy}$ is the average of the product $ab$ over the rounds with settings
$(x, y)$. The CHSH quantity is
\[
S = E_{00} + E_{01} + E_{10} - E_{11}.
\]
I collect the four correlators into the $2 \times 2$ matrix
\[
E = \begin{pmatrix} E_{00} & E_{01} \\ E_{10} & E_{11} \end{pmatrix},
\]
whose rows are indexed by Alice's setting and whose columns by Bob's. The rows of $E$ are
$(E_{00}, E_{01})$ and $(E_{10}, E_{11})$; the columns are $(E_{00}, E_{10})$ and
$(E_{01}, E_{11})$.
\end{definition}

\begin{definition}[Counters]\label{def:chsh-counters}
There are eight counts: for each setting pair, $\mathrm{same}_{xy}$ counts rounds
whose two answer bits matched and $\mathrm{diff}_{xy}$ counts rounds whose bits differed. A
round adds one to exactly one of the eight. Write
\[
N_{xy} = \mathrm{same}_{xy} + \mathrm{diff}_{xy}, \qquad
D_{xy} = \mathrm{same}_{xy} - \mathrm{diff}_{xy}, \qquad
E_{xy} = \frac{D_{xy}}{N_{xy}}.
\]
\end{definition}

A matched round contributes $ab = +1$ and a differing round $ab = -1$, so $D_{xy}$ is the sum
of the products and $D_{xy}/N_{xy}$ is their average. Since $|D_{xy}| \le N_{xy}$, every
correlator from counters lies in $[-1,1]$. When a bucket is empty, its correlator is defined
as 0. The score is $E_{00}+E_{01}+E_{10}-E_{11}$.

\subsection{The classical bound}
\label{sub:chsh-classical}
Write a fixed plan on a card: Alice's answer to setting zero, her answer to
setting one, and Bob's two answers. Each answer is a sign, plus or minus.
Before multiplying anything, look at Bob's pair. Either the signs agree or
they differ. If they agree, adding them gives plus or minus two and subtracting
them gives zero. If they differ, the sum is zero and the difference is plus
or minus two. One of those two quantities vanishes; the other cannot exceed
two in size. The next calculation uses exactly that split.

\begin{definition}[Finite factorizable model]\label{def:chsh-factorizable}
Four real correlators $E_{00}, E_{01}, E_{10}, E_{11}$ are factorizable when there are a
natural number $\lambda_{\max}$, weights $p(\lambda)$ and response signs $A(x, \lambda)$,
$B(y, \lambda)$ for $\lambda = 0, \ldots, \lambda_{\max}$ such that:
$0 \le p(\lambda) \le 1$;
$\sum_{\lambda = 0}^{\lambda_{\max}} p(\lambda) = 1$; every $A(x, \lambda)$ and every
$B(y, \lambda)$ is $-1$ or $+1$; and for all $x, y$,
$E_{xy} = \sum_{\lambda=0}^{\lambda_{\max}} p(\lambda) A(x, \lambda) B(y, \lambda)$.
In words: a hidden card $\lambda$ is drawn with probability $p(\lambda)$, and each player reads
an answer off the card using only their own setting. Only the correlators are modelled.
\end{definition}

\begin{theorem}[The classical bound for a finite local model]\label{thm:chsh-fine}
If $E$ is factorizable, then $-2 \le E_{00} + E_{01} + E_{10} - E_{11} \le 2$.
\end{theorem}

\begin{definition}[A tally consistent with one local plan]\label{def:chsh-wclc}
For bits $a_0, a_1, b_0, b_1 \in \{0, 1\}$ and a counter record $w$, $w$ is locally consistent with
$(a_0, a_1, b_0, b_1)$ when, for each setting pair $(x, y)$: if $a_x = b_y$ then
$\mathrm{diff}_{xy} = 0$, otherwise $\mathrm{same}_{xy} = 0$; and every bucket is sampled,
$N_{00}, N_{01}, N_{10}, N_{11} > 0$.
\end{definition}

\begin{theorem}[A tally above 2 is not one replayed strategy]\label{thm:chsh-tally}
Let $w$ be a counter record whose rational score $S(w)$ is greater than 2.
Then for all bits $a_0, a_1, b_0, b_1$, $w$ is not locally consistent with
$(a_0, a_1, b_0, b_1)$. Every locally consistent record has $S(w) = 2$ or $S(w) = -2$; the first statement is that fact turned around.
\end{theorem}

\subsection{Tsirelson's bound from algebra}
\label{sub:chsh-tsirelson-algebra}
Now put the cards aside and look just at four numbers in the square. I want
to bound their three-plus-one-minus score from their squared sizes. The useful
trick is to write the proposed allowance minus the squared score as a pile
of squares. A square cannot be negative. If the difference really is that
pile, the score cannot outrun the allowance.

\begin{lemma}[The six-square identity]\label{lem:chsh-sos}
For all real $a, b, c, d$,
\[
\begin{aligned}
&4(a^2 + b^2 + c^2 + d^2) - (a + b + c - d)^2 \\
&\quad = (a-b)^2 + (a-c)^2 + (a+d)^2 + (b-c)^2 + (b+d)^2 + (c+d)^2.
\end{aligned}
\]

\end{lemma}

\begin{theorem}[Row-bound form of Tsirelson's bound]\label{thm:chsh-row-tsirelson}
For all real $E_{00}, E_{01}, E_{10}, E_{11}$ with $E_{00}^2 + E_{01}^2 \le 1$ and
$E_{10}^2 + E_{11}^2 \le 1$, $S^2 \le 8$. The absolute form $|S| \le \sqrt 8 = 2\sqrt 2$ follows.
\end{theorem}

\begin{theorem}[Tightness]\label{thm:chsh-tight}
The point $E_{00} = E_{01} = E_{10} = 1/\sqrt 2$, $E_{11} = -1/\sqrt 2$ satisfies both row
bounds and has $S = \sqrt 8$.
\end{theorem}

\begin{remark}[Where equality can happen]\label{rem:chsh-equality}
Equality $S^2 = 8$ under the row bounds forces
equality in both steps of the proof above: the six squares must all vanish, which gives
$a = b = c = -d$, and the rows must be tight, which gives $2a^2 = 1$. So the point of
Theorem~\ref{thm:chsh-tight} and its negative are the only places the bound is reached. A tally's correlators are
ratios of whole numbers and $2a^2 = 1$ has no such solution, so no tally reaches the bound; the integer
check's passing tallies stay strictly below it and come as near as you like. (Take rational
$a, b$ with $a^2 + b^2 = 1$, from a Pythagorean triple, and the correlators $(a, b, b, -a)$: they pass
the check of Section~\ref{sub:chsh-integer} with no slack, and $S = 2(a + b)$, which gets as close to
$2\sqrt 2$ as you like as $a/b$ approaches 1.)
\end{remark}

\subsection{The zero-marginal \texorpdfstring{$5\times5$}{5 by 5} moment matrix}
\label{sub:chsh-five}
The four averages fill only the Alice--Bob part of a larger table. Give that
table five labels: a constant reading 1, Alice's two possible readings, and
Bob's two. Each cell records an average product of its row and column
readings. A reading multiplied by itself is 1, so the diagonal is all ones.
The four averages measured fill the cross-player cells and their reflected
partners.

\begin{definition}[The pinned level-1 matrix]\label{def:chsh-five}
Take rows and columns $I,A_0,A_1,B_0,B_1$. Put 1 on the diagonal and the four correlators in
the Alice--Bob cells, and set the four one-party entries (each player's own average answer) and
the two within-party entries ($A_0$ against $A_1$, $B_0$ against $B_1$) to zero. The resulting
matrix is:
\end{definition}

\[
M(E) \;=\;
\begin{array}{c@{\quad}c}
 & \begin{array}{ccccc} I & A_0 & A_1 & B_0 & B_1 \end{array} \\
\begin{array}{c} I \\ A_0 \\ A_1 \\ B_0 \\ B_1 \end{array} &
\begin{pmatrix}
1 & 0 & 0 & 0 & 0 \\
0 & 1 & 0 & E_{00} & E_{01} \\
0 & 0 & 1 & E_{10} & E_{11} \\
0 & E_{00} & E_{10} & 1 & 0 \\
0 & E_{01} & E_{11} & 0 & 1
\end{pmatrix}
\end{array}
\]

\begin{lemma}[Schur complement, the standard form]\label{lem:chsh-schur}
Let $P$ be symmetric positive definite. The symmetric block matrix
$\begin{pmatrix} P & Q \\ Q^{\mathsf T} & R \end{pmatrix}$ is PSD if and only if
$R - Q^{\mathsf T} P^{-1} Q$ is PSD.
\end{lemma}

\begin{definition}[Column contractive]\label{def:chsh-cc}
Write
$p = 1 - E_{00}^2 - E_{10}^2$, $q = 1 - E_{01}^2 - E_{11}^2$ and
$s = E_{00} E_{01} + E_{10} E_{11}$. The correlators are column contractive when
\[
p \ge 0, \qquad q \ge 0, \qquad p\,q - s^2 \ge 0.
\]
\end{definition}

\begin{corollary}[The equivalence]\label{cor:chsh-iff}
For all real $E_{00},E_{01},E_{10},E_{11}$, the symmetric matrix $M(E)$ is positive
semidefinite if and only if the correlators are column contractive.
\end{corollary}

\begin{theorem}[PSD implies Tsirelson]\label{thm:chsh-psd-tsirelson}
If the pinned matrix $M(E)$ is symmetric and positive semidefinite at
$E_{00},E_{01},E_{10},E_{11}$, then $S^2\le8$ and $|S|\le\sqrt8$.
\end{theorem}

\subsection{The integer check}
\label{sub:chsh-integer}
The machine has whole-number tallies. The three conditions above were written
using averages, which are fractions. You can test the same conditions without
ever dividing: first make sure each denominator is positive, then multiply
it away. Multiplying by a positive number cannot turn a passing inequality
into a failing one or a failing one into a passing one.

\begin{lemma}[The seven facts are the three conditions with denominators cleared]\label{lem:chsh-clear}
Assume $N_{00}, N_{01}, N_{10}, N_{11} > 0$ and put $E_{xy} = D_{xy}/N_{xy}$. Then, with
$p, q, s$ as in Definition~\ref{def:chsh-cc},
\[
\begin{gathered}
p = \frac{A}{N_{00}^2 N_{10}^2}, \qquad q = \frac{B}{N_{01}^2 N_{11}^2}, \qquad
s = \frac{C}{N_{00} N_{01} N_{10} N_{11}}, \\
pq - s^2 = \frac{AB - C^2}{(N_{00} N_{01} N_{10} N_{11})^2}.
\end{gathered}
\]
Because every denominator is positive, $0 \le A \iff p \ge 0$, $0 \le B \iff q \ge 0$ and
$C^2 \le AB \iff s^2 \le pq$.
\end{lemma}

\begin{theorem}[Soundness of the integer check]\label{thm:chsh-check-sound}
If a tally $w$ passes the seven-integer check, then its four correlators
$E_{xy}=D_{xy}/N_{xy}$ are column contractive. Consequently its pinned matrix is positive
semidefinite.
\end{theorem}

\begin{corollary}[A passed check implies the Tsirelson bound]\label{cor:chsh-check-bound}
If a tally $w$ passes the seven-integer check, the correlators of $w$ satisfy
$S^2 \le 8$. It is the composition of Theorem~\ref{thm:chsh-check-sound},
Corollary~\ref{cor:chsh-iff} and Theorem~\ref{thm:chsh-psd-tsirelson}.\proofclosed \end{corollary} Two of the seven facts carry the bound by themselves. With every pair sampled, $0\le A$ and $0\le B$ are $p\ge0$ and $q\ge0$, so the four squared correlators add to at most 2 and Lemma~\ref{lem:chsh-sos} gives $S^2\le8$. The fact $C\cdot C\le A\cdot B$ is what positive semidefiniteness adds, and the bound doesn't use it.
For the form this takes on the small machine, where a raised flag has to stand on a passing
check, see Appendix~\ref{sec:chsh}.

\subsection{Level 1+AB}
\label{sub:chsh-q1ab}
The first table had one place for each player's answer. Enlarge it by giving
each Alice--Bob product a place too. There are four such products, so the
five places become nine. A cell still asks for the average product of its
row and column labels. Some of those products reduce to entries I already
have; five more numbers are needed for the particular filling below.
Call them $\gamma_1$ through $\gamma_5$. The longer labels say which
products they stand for.

\begin{definition}[The level-1+AB matrix]\label{def:chsh-nine}
The $9 \times 9$ real matrix on rows and
columns $I, A_0, A_1, B_0, B_1, A_0B_0, A_0B_1, A_1B_0, A_1B_1$, with the following entries:
\end{definition}

\[
\setlength{\arraycolsep}{3.2pt}
\begin{array}{c|ccccccccc}
 & I & A_0 & A_1 & B_0 & B_1 & A_0B_0 & A_0B_1 & A_1B_0 & A_1B_1 \\ \hline
I      & 1 & 0 & 0 & 0 & 0 & E_{00} & E_{01} & E_{10} & E_{11} \\
A_0    & 0 & 1 & 0 & E_{00} & E_{01} & 0 & 0 & \gamma_1 & \gamma_2 \\
A_1    & 0 & 0 & 1 & E_{10} & E_{11} & \gamma_1 & \gamma_2 & 0 & 0 \\
B_0    & 0 & E_{00} & E_{10} & 1 & 0 & 0 & \gamma_3 & 0 & \gamma_4 \\
B_1    & 0 & E_{01} & E_{11} & 0 & 1 & \gamma_3 & 0 & \gamma_4 & 0 \\
A_0B_0 & E_{00} & 0 & \gamma_1 & 0 & \gamma_3 & 1 & 0 & 0 & \gamma_5 \\
A_0B_1 & E_{01} & 0 & \gamma_2 & \gamma_3 & 0 & 0 & 1 & -\gamma_5 & 0 \\
A_1B_0 & E_{10} & \gamma_1 & 0 & 0 & \gamma_4 & 0 & -\gamma_5 & 1 & 0 \\
A_1B_1 & E_{11} & \gamma_2 & 0 & \gamma_4 & 0 & \gamma_5 & 0 & 0 & 1
\end{array}
\]

Write a vector as $(v_I, v_{A_0}, v_{A_1}, v_{B_0}, v_{B_1}, v_{00}, v_{01}, v_{10}, v_{11})$,
where $v_{xy}$ is the coordinate of $A_xB_y$. Put
\[
\begin{aligned}
L_1 &= E_{00} v_{00} + E_{01} v_{01} + E_{10} v_{10} + E_{11} v_{11}, \\
L_2 &= E_{00} v_{B_0} + E_{01} v_{B_1} + \gamma_1 v_{10} + \gamma_2 v_{11}, \\
L_3 &= E_{10} v_{B_0} + E_{11} v_{B_1} + \gamma_1 v_{00} + \gamma_2 v_{01}.
\end{aligned}
\]

\begin{definition}[The residual]\label{def:chsh-residual}
\[
\begin{aligned}
R ={}& v_{B_0}^2 + v_{B_1}^2 + v_{00}^2 + v_{01}^2 + v_{10}^2 + v_{11}^2 \\
&+ 2 v_{B_0} (\gamma_3 v_{01} + \gamma_4 v_{11}) + 2 v_{B_1} (\gamma_3 v_{00} + \gamma_4 v_{10}) \\
&+ 2 \gamma_5 (v_{00} v_{11} - v_{01} v_{10}) - L_1^2 - L_2^2 - L_3^2.
\end{aligned}
\]
The condition called column contractivity at level 1+AB is $R \ge 0$ for all six real inputs.
\end{definition}

\begin{theorem}[PSD at level 1+AB]\label{thm:chsh-nine-iff}
For all real $E_{xy}$ and $\gamma_1, \ldots, \gamma_5$,
the matrix of Definition~\ref{def:chsh-nine} is positive semidefinite if and only if
the residual $R$ is nonnegative for every choice of its six real inputs.
\end{theorem}

\subsection{The elliptope completion}
\label{sub:chsh-elliptope}
Go back to the two within-player cells I filled with zero. Leave them blank
this time. Keep all four measured cross-player averages fixed, and ask whether
there is any way to fill those two blanks that makes the table pass. A failure
with zeros no longer settles the question; another filling may work.

\begin{definition}[Completion and membership]\label{def:chsh-elliptope}
Complete the matrix of Section~\ref{sub:chsh-five} by setting the
within-party entries to free reals: $x$ in the $(A_0, A_1)$ cells and $y$ in the
$(B_0, B_1)$ cells, marginals still zero. The correlators are elliptope realizable
 when some $x, y$ make the completed matrix symmetric and PSD.
\end{definition}

\[
M(E; x, y) \;=\;
\begin{pmatrix}
1 & 0 & 0 & 0 & 0 \\
0 & 1 & x & E_{00} & E_{01} \\
0 & x & 1 & E_{10} & E_{11} \\
0 & E_{00} & E_{10} & 1 & y \\
0 & E_{01} & E_{11} & y & 1
\end{pmatrix}
\]

Its quadratic form at $v = (v_0, \ldots, v_4)$ is
$\sum_k v_k^2 + 2x v_1 v_2 + 2y v_3 v_4 + 2E_{00} v_1 v_3 + 2E_{01} v_1 v_4 + 2E_{10} v_2 v_3
+ 2E_{11} v_2 v_4$. The pinned matrix is the case $x = y = 0$,
so PSD of $M(E)$ implies membership.

\begin{definition}[The two branches]\label{def:chsh-gate}
The positive-minor branch: with the completion supplied as buckets and all
entries scaled to integers by a positive common factor $K$, the four leading minors of
the $4 \times 4$ block on $A_0, A_1, B_0, B_1$ are strictly positive. The factorization branch: with the completion and a rational certificate supplied (six subdiagonal entries over
a common positive denominator, four pivots $d_k \ge 0$ over another), the ten entry equations
of $M = L D L^{\mathsf T}$ hold in cross-multiplied integer arithmetic. The full gate is their disjunction.
\end{definition}

\begin{theorem}[Soundness of the gate]\label{thm:chsh-gate-sound}
If the full gate returns yes on a counter record, any
completion buckets and any certificate, then the four correlators of the record are
elliptope realizable. On the PR box
tally the gate returns no for every choice of the remaining arguments.
\end{theorem}

\subsection{The four-vector construction}
\label{sub:four-vectors}
Put two equal-length pencils down with their ends together, and draw a line along
one of them. Keep that pencil still and turn the other. Use the square corner of
a sheet of paper to bring the moving pencil's tip straight across to the line,
and mark where it lands. The mark moves from a full pencil-length forward,
through the joined ends when the pencils make a square corner, to a full
pencil-length backward. Call those readings 1, 0 and $-1$. That is the signed
shadow I want. Between those positions the reading is the cosine of the angle.
A dot product of unit vectors measures that shadow.

\subsection{The arcsine boundary, in both directions}
\label{sub:arcsine-construction}
Keep Bob's two directions hinged at one end. Alice's first direction has to make
a specified angle with each. How far can Bob's hinge open while that is possible?
Hold Alice's direction upright and sweep each of Bob's around it at its required
angle. Point them towards the same side to close the hinge as far as it will go;
point them towards opposite sides to open it as far as it will go. There is a
whole interval of openings between those limits. Alice's second direction asks
for another interval.

\subsection{The finite local model, in both directions}
\label{sub:finite-local-construction}
Make sixteen cards, each with four boxes labelled $A_0,A_1,B_0,B_1$, and put every
possible arrangement of plus and minus signs on them, one arrangement per card.
Draw a card before choosing which setting to ask each player about. Alice reads
her chosen box; Bob reads his. Neither needs the other's setting. How often each
card is drawn is the only choice left in this model.

\subsection{Scope of the constructions}
\label{sub:chsh-not-proved}

The constructions above pin down the correlators. What follows is where they stop: they
aren't a completeness claim for the integer check, they say nothing about full behaviours
with marginals, and they say nothing about data from a lab.

The arcsine boundary. For two settings and two outcomes per side, Tsirelson, Landau (1988) and
Masanes (2003) give theorems characterizing the correlators reachable by quantum measurements:
$|\arcsin E_{00} + \arcsin E_{01} + \arcsin E_{10} - \arcsin E_{11}| \le \pi$ and its sign
variants. The characterization holds as an if and only if with the
completion set and, through Tsirelson's representation, with the quantum
set. The point $(3/5, 4/5, 4/5, -3/5)$ is on the boundary:
$\arcsin \tfrac35 + \arcsin \tfrac45 = \tfrac{\pi}{2}$ by the 3-4-5 triangle, so the sum is $\pi$; Tsirelson's point is on it too.
The necessary direction holds in any dimension, for any expectation on an algebra of
observables in which Alice's commute with Bob's.

Equality with the quantum set. That the completion set of Section~\ref{sub:chsh-elliptope} is exactly the set of
quantum correlators is Tsirelson's representation theorem, and it is proved in finite dimension:
for both real and complex amplitudes, with two four-level systems enough for every point, and two
two-level ones with complex amplitudes. That PSD of the level-1 moment matrix is necessary for
quantum behaviour is proved for strategies in finite dimension, with real amplitudes
and with complex ones, and in any dimension, marginals included, for any expectation on
an algebra of observables.

Achievability. Theorem~\ref{thm:chsh-tight} is a point of $\mathbb{R}^4$. A state and
measurements that reach it are explicit: the two-qubit state
$(|00\rangle + |11\rangle)/\sqrt 2$ with $A_0 = Z$, $A_1 = X$, $B_0 = (Z + X)/\sqrt 2$ and
$B_1 = (Z - X)/\sqrt 2$ has correlators $1/\sqrt 2$, $1/\sqrt 2$, $1/\sqrt 2$, $-1/\sqrt 2$ and
$S = 2\sqrt 2$, and every strategy in finite dimension has
$S^2 \le 8$, with real amplitudes and with complex ones. Both are arithmetic about the model.
How close a real laboratory comes, and what that costs in loopholes closed, is the story told in
Appendix~\ref{sec:chsh}.

Completeness of the checks. The 1+AB cascade and the elliptope gate are proved sound. The cascade has no proved converse; the gate's factorization branch accepts every rational PSD completion, by the row-peeling of Section~\ref{sub:four-vectors} with each pivot kept. The level-1 integer check is the exception: on the small machine the same
seven integer facts hold exactly when every pair was sampled and the pinned matrix is positive
semidefinite (Theorem~\ref{thm:chsh-exact}), so that check is
complete for what it tests. The 1+AB cascade is strict and misses boundary points.
The two 1+AB checks don't sit inside one another. The full check accepts the $12/5$ tally above, which the zero-moment check refuses, and it refuses the tally with $E_{00}=1$ and the other three correlators 0, where $H_{11}=0$, which the zero-moment check accepts.

Finite levels. For full behaviours, marginals included, the NPA levels are relaxations.
Chaturvedi (arXiv 2607.14569, 2026) argues, in a preprint, that no finite NPA level characterizes the complete
quantum set even in this simplest Bell scenario. I haven't gone through that proof, and nothing
here rests on it.

Statistics and physics. Nothing here says how many rounds make a tally informative, and
nothing says the counters came from honest play; a program can INC whatever counts it likes.

\subsection{Prior work}
\label{sub:chsh-prior}

Tsirelson's 1980 paper gives the bound $2\sqrt 2$ for quantum correlations, and his theorem
says that quantum two-outcome correlators are inner products of real unit vectors. Landau (1988)
and Masanes (2003, arXiv quant-ph/0309137) each give a condition that, by their proofs,
characterizes the quantum correlators exactly for two settings and two outcomes. Wehner (2006,
arXiv quant-ph/0510076) gets Tsirelson bounds, CHSH among them, from a semidefinite program on
the $4 \times 4$ Gram matrix of those vectors. Navascués,
Pironio and Acín (2007, arXiv quant-ph/0607119; 2008, arXiv 0803.4290) built the hierarchy of
moment-matrix tests; their 2008 paper applies a Schur complement to the level-1 matrix of this
scenario, recovers the arcsine condition, and discusses the 1+AB level. I went back to each of
those papers for these points, and that is what I found in them. The results are theirs. Two of
them are proved again here, in finite dimension: Tsirelson's representation
and the arcsine characterization; the rest of
their proofs are not redone. Echenim and Mhalla give a detailed derivation of Tsirelson's
bound (arXiv 2306.12535), and exact rational
sum-of-squares certificates go back at least to Peyrl and Parrilo (2008). So the
moment-matrix route to $2\sqrt 2$ is standard. What is specific here is narrower: an integer
certificate computed from outcome counts, and the argument from ``the counts pass'' to PSD
of the pinned matrix and on to $S^2 \le 8$. I did not find that chain elsewhere, and I make no
claim beyond it.
On the host machine, the same chain is earned: a raised flag from a clean start
stands on a passing check of the packed counts, and the check passes exactly when the pinned matrix
is positive semidefinite and every pair was sampled (Appendix~\ref{sec:chsh},
Theorems~\ref{thm:chsh-exact} and~\ref{thm:chsh-earned}). The converse direction of the integer check,
integers from PSD, is proved there too, and there too I make no claim beyond the chain.


%===========================================================================
\section{Methodology note}
\label{sec:methodology}

Everything this book claims as its own stands or falls on the arguments written in it; what it
borrows is cited to the people who proved it, and the physics enters as a named premise. If a
proof on the page is wrong, the result is wrong, and nothing else I could tell you would rescue it.

Language models were my hands, as I said in Section~\ref{sec:origin}: they did
symbol-pushing, drafted code and helped work through derivations. I have to answer for
what stayed. Start with the definitions and follow how each step leads to the conclusion.
If I counted a finite construction, its cases belong on the page. If I claimed something
for every run, the argument has to explain why the next step preserves it. A successful
example cannot carry that claim, and a name attached to a proof cannot fill a missing step.

The same demand applies to the picture around the mathematics. A definition has to say
what I meant, and a model of a device has to fit that device. Those questions belong beside the premises and the argument, where you can examine the fit.


%===========================================================================
\section{Rows of the claim ledger for the passages above}
%===========================================================================

The book's claim ledger listed these claims, each with the kind of claim it is and where it lived.

\begingroup\small
\begin{longtable}{@{}>{\raggedright\arraybackslash}p{0.242\linewidth}>{\raggedright\arraybackslash}p{0.344\linewidth}>{\raggedright\arraybackslash}p{0.344\linewidth}@{}}
\toprule
\textbf{Claim} & \textbf{Kind} & \textbf{Home} \\
\midrule
\endhead
\bottomrule
\endlastfoot
The small machine runs the eight-state machine through a window, move for move, paying exactly its charges & proved (from live states; the setup's claim can't fail) & Section~\ref{sec:part3-on-machine}, Theorem~\ref{thm:frag-small} \\
The small machine as a whole is not merge-priced & proved & Section~\ref{sec:part3-on-machine}, Theorem~\ref{thm:frag-merge-split} \\
A decision tree has at most two to its depth leaves & definition; proved & Section~\ref{sec:entitlement}, Definition~\ref{def:u08-dtree}, Lemma~\ref{lem:u08-tree-width} \\
Paid narrowing, counted & proved with the total of the charges, which equals the counter's rise by conservation; both are proved (premises: a covering tree, paid steps) & Section~\ref{sec:entitlement}, Theorem~\ref{thm:u08-count} \\
Structural entitlement for every certified Thiele-complete machine: stronger test, earned chain, index bits at most the record moves, ledger rise at least 3 & proved (premises: a clean start, and M1 to M6) & Section~\ref{sec:part3-on-machine}, Theorem~\ref{thm:ent-machine} \\
A floor with no supplied tree, from the run's own CHECKs & proved (answers taken at the rivals as they start) & Section~\ref{sec:part3-on-machine}, Theorem~\ref{thm:ent-questions} \\
Structural entitlement with the narrowing made by the run, over any Thiele machine whose test steps are named: every drop made by a test the run passed, the exact product count, index bits at most $k$ times the cost & proved (premises: test steps charged at least one, a prior containing the start, a $2^{-k}$ share kept by each test) & Section~\ref{sec:entitlement}, Theorem~\ref{thm:ent-run} \\
The same for every certified Thiele-complete machine: index bits at most $k$ times the passing CHECKs, at most $k$ times the ledger's rise; the rounded sum with no share assumed; the share can't be dropped & proved (premises: a clean start, the record up at the end, a prior containing the start, the share for part (4)) & Section~\ref{sec:part3-on-machine}, Theorem~\ref{thm:ent-run-machine}, Proposition~\ref{prop:ent-share-needed} \\
What the record adds to the price bound from a clean start: the earned chain and a floor of 3 & proved & Section~\ref{sec:part3-on-machine}, Theorem~\ref{thm:ent-observed} \\
The weighted count and the partition form & proved; the weights are whole numbers, and no probabilities enter & Section~\ref{sec:part3-on-machine} \\
The eight-start member, and a search through $n$ bits with up to 16 questions & proved; the 16 is the machine's fact table & Section~\ref{sec:part3-on-machine}, Proposition~\ref{prop:ent-small}, Theorem~\ref{thm:ent-search} \\
The covering hypothesis can't be dropped & proved, by a counterexample on the small machine & Section~\ref{sec:part3-on-machine} \\
The time tax: no program of zero-cost two-counter moves decides, in fewer than $m$ moves, $A \ge m$ from its final line and $B$, one paid move decides it, and three paid moves certify it & proved, for programs of INC and DEC & Section~\ref{sec:part3-on-machine}, Theorem~\ref{thm:time-tax} \\
The lists, tree and fibres are read off the run & not claimed, except in the worked members & Section~\ref{sec:part3-on-machine} \\
The latch on a Turing machine, a register machine, L & definition; proved (and the relation is an equivalence) & Section~\ref{sec:latch}, Definition~\ref{def:u11-weak-equiv}, Theorem~\ref{thm:u11-four-bases} \\
Thiele-complete means a universal base whose moves cost nothing and can't touch the record, every other move costing exactly one, a record that stays up, blank starts, and some run whose record ends up from one start and down from another; nothing in it says what the record is about & definition & Section~\ref{sec:complete}, Definition~\ref{def:thiele-complete} \\
Every Thiele-complete machine obeys the toll, raises its record only by a mark costing exactly one, counts its marks on its ledger, pays at least one on any run that raises the record, and runs every two-counter program & proved, for every machine and every interface meeting the definition & Section~\ref{sec:complete}, Theorem~\ref{thm:tc-has} \\
A clock is not Thiele-complete, nor is a machine whose record never rises, and each of the four clauses is needed & proved & Section~\ref{sec:complete}, Observation~\ref{obs:tc-clock}, Theorem~\ref{thm:tc-clauses} \\
Every Thiele-complete machine hides its record and its ledger from its window, with a gap of one unit, and one is reached & proved & Section~\ref{sec:complete}, Theorem~\ref{thm:tc-window} \\
Every universal base becomes Thiele-complete by adding marks, with no claims and no checker & proved & Section~\ref{sec:complete}, Theorem~\ref{thm:tc-lift} \\
Certified Thiele-complete is exactly Thiele-complete with clean starts, the earned chain and a claim that can fail; the inclusion is strict & proved & Section~\ref{sec:certified}, Theorem~\ref{thm:tc-decomposition} \\
Halving pricing bounds weighted conditional entropy; an infinite landing pile admits no finite halving price & proved; average loss still depends on the distribution & Section~\ref{sec:shadow-exact}, Theorems~\ref{thm:dir-bits-loss} and~\ref{thm:dir-infinite} \\
The nat family: each candidate decider misses program two in the instance built from it & proved & Section~\ref{sec:diagonal}, Theorem~\ref{thm:u12-nat-family} \\
Recursion, Rice and halting for L & definition; proved & Section~\ref{sec:diagonal}, Definition~\ref{def:u12-L}, Theorem~\ref{thm:u12-L} \\
Platform and quote; inclusion check; bounds fragment & definition & Section~\ref{sec:models}, Definition~\ref{def:u13-quote}, Definition~\ref{def:u13-inclusion}, Definition~\ref{def:u13-pcc} \\
The proof-carrying checker is the verification condition & proved & Section~\ref{sec:models}, Theorem~\ref{thm:u13-pcc} \\
Ecosystem, proliferation, pointer; copied and relied on, independent carriers; observation game & definition & Section~\ref{sec:pointer}, Definition~\ref{def:u14-ecosystem}, Definition~\ref{def:u14-reliance}, Definition~\ref{def:u14-game} \\
Independent copies on every state, with two observers or more: the event is equivalent at any two states & proved & Section~\ref{sec:pointer}, Theorem~\ref{thm:u14-independent-constant} \\
In the specified replicated-header toy, gas used and finality are copied, and the chosen action maps rely only on finality & proved for those maps & Section~\ref{sec:pointer}, Theorem~\ref{thm:u14-ledger} \\
In the toy, ``certified'' is the unique pointer & proved & Section~\ref{sec:pointer}, Theorem~\ref{thm:u14-toy} \\
The survey; the swap shows the verdict follows the observer maps & proved & Section~\ref{sec:pointer}, Theorem~\ref{thm:u14-survey} \\
One blind observer stops proliferation & proved & Section~\ref{sec:pointer}, Theorem~\ref{thm:u14-blind} \\
Agreement doesn't buy permanence; durable views that report the event do & proved & Section~\ref{sec:pointer}, Theorem~\ref{thm:u14-toggle} \\
No stored fact is forged on a clean-start trace & proved on the small machine, under the stated premises & Section~\ref{sec:smallest}, Theorem~\ref{thm:u34-no-forging} \\
Bounded table keeps facts and traps when full & proved on the small machine, under the stated premises & Section~\ref{sec:smallest}, Theorem~\ref{thm:u34-bounded} \\
Certified Thiele-complete means a universal base that costs nothing and can't touch the record, a record that stays up once it's up, an earned record with a checker proved to mean its claims, exact prices of one per check, commit and certificate, and a claim whose check can fail & definition & Section~\ref{sec:smallest}, Definition~\ref{def:u34-thiele-complete} \\
A machine certified Thiele-complete over a claim language fixed first is certified Thiele-complete, and from a clean state its raised record stands on a claim whose meaning held when it was checked and when it was committed & proved & Section~\ref{sec:smallest}, Theorem~\ref{thm:u34-tc-over} \\
Every certified Thiele-complete machine obeys the toll, runs every two-counter program, counts its record moves on its ledger, pays at least three on any run from a clean state that ends with the record up, has a check that fails and a run that certifies, and from a clean state raises its record only by a CERTIFY & proved, for every machine and every interface meeting the definition & Section~\ref{sec:smallest}, Theorem~\ref{thm:u34-tc-consequences} \\
A clock (every move charged, any reading or any latch as the record) obeys the toll with a two-counter machine running inside it and is not certified Thiele-complete; nor is any machine whose every move costs something & proved, directly from the rules & Section~\ref{sec:smallest}, Observation~\ref{thm:u34-clock} \\
Each clause does work: a paid latch meets the universal base and exact prices clauses and isn't certified Thiele-complete, and a silent machine meets three clauses and fails non-vacuity & proved & Section~\ref{sec:smallest}, Theorem~\ref{thm:u34-tc-clauses} \\
The four clauses of certified Thiele-complete are independent: for each, a Thiele machine with a two-counter machine inside meets the other three, fails it, and is certified Thiele-complete through no interface & proved & Section~\ref{sec:smallest}, Theorem~\ref{thm:tc-independent} \\
Certified Thiele-complete at any unit: with record moves at one fixed price $c \ge 1$ the class is the unit class with its prices multiplied by $c$, and the clock fails at every unit & proved & Section~\ref{sec:smallest}, Theorem~\ref{thm:u34-tc-unit} \\
The small machine, the same machine over any property language with an exact equality test, a checker proved to say yes exactly when the property holds and some property true of one number and false of another, and the sorted-list machine are certified Thiele-complete & proved on the small machine and on the generic machine & Section~\ref{sec:smallest}, Theorem~\ref{thm:u34-tc-instances} \\
Every certified Thiele-complete machine hides its record and its ledger from its base window & proved, for every machine and every interface meeting the definition, directly from the rules & Section~\ref{sec:smallest}, Observation~\ref{thm:u34-tc-window} \\
The calibration is a unit conversion & definition; proved & Section~\ref{sec:landauer}, Definition~\ref{def:u24-scale}, Theorem~\ref{thm:u24-calibration} \\
Structural dissipation pricing (a whole-unit floor on collapsing instructions, separate from Landauer's principle); dissipation at most cost & assumed (premise) & Section~\ref{sec:landauer}, Premise~\ref{prem:u24-landauer-pair} \\
A collapse needs the merge clause: INC $A$ raises ``$A \ge 1$'' one way, one-to-one and costing nothing; on a finite machine a one-way step always merges & proved  & Section~\ref{sec:landauer}, Premise~\ref{prem:u24-landauer-pair} \\
The heat is fixed by the chosen gap & definition; proved & Section~\ref{sec:landauer}, Definition~\ref{def:u24-calorimeter}, Theorem~\ref{thm:u24-calorimeter} \\
A driven reset with rates in detailed balance, its gap raised from 0 to $D$ in $N$ equal steps, pays between $kT \ln 2 - kT e^{-D/kT}$ and $kT \ln 2 + D/(2N)$ in work & definition; proved & Section~\ref{sec:landauer}, Definition~\ref{def:u24-driven}, Theorem~\ref{thm:u24-driven} \\
The simpler host runs any small-machine program as a guest, and halts with any two-counter program & proved on the small machine, by construction: its guest step runs the guest's own step as a primitive; the guest program is held as a list & Appendix~\ref{sec:simpler-host}, Theorem~\ref{thm:u34-universal} \\
From a loaded host, or any state where the two agree, the simpler host's mirror always equals the guest's flag & proved on the small machine, directly from the rules & Appendix~\ref{sec:simpler-host}, Observation~\ref{thm:u34-mirror} \\
From a host that agrees, the simpler host's step collects the guest's toll; on every trace its ledger covers the guest's & proved on the small machine & Appendix~\ref{sec:simpler-host}, Theorem~\ref{thm:u34-host-toll} \\
The simpler host's mirror rises only in a guest step that passed CERTIFY & proved on the small machine, for every sequence of host moves & Appendix~\ref{sec:simpler-host}, Theorem~\ref{thm:u34-no-free-host} \\
The simpler host's own instructions never touch the guest & proved on the small machine & Appendix~\ref{sec:simpler-host}, Theorem~\ref{thm:u34-no-free-host} \\
From a loaded host, the simpler host's raised mirror means the guest's own run checked, committed and certified the same claim & proved on the small machine  & Appendix~\ref{sec:simpler-host} \\
The simpler host is a Thiele machine & proved on the small machine & Appendix~\ref{sec:simpler-host}, Theorem~\ref{thm:u34-host-instance} \\
The simpler host's ledger equals the guest's & not claimed; a guest step with budget $b$ costs $b$ whatever the guest instruction cost, so its ledger only covers the guest's; $U$'s is exact & Appendix~\ref{sec:simpler-host}, Theorem~\ref{thm:u34-host-toll} \\
Rice's theorem on the host machine, for properties that depend on everything a program ends with (registers, table, channel, trap, ledger, flag) & proved; undecidable in the reduction sense of Theorem~\ref{thm:u34-undecidable} & Appendix~\ref{sec:small-diagonal}, Theorem~\ref{thm:host-rice} \\
Rice's theorem for two-counter programs, from the clean start & proved; the fact table isn't compared & Appendix~\ref{sec:small-diagonal}, Theorem~\ref{thm:guest-rice} \\
s-m-n, and a universal host program for host programs & proved, with no premise on register 2 & Appendix~\ref{sec:small-diagonal}, Theorem~\ref{thm:host-smn} \\
Kleene's recursion theorem on the host machine & proved; the fixed point is of the function & Appendix~\ref{sec:small-diagonal}, Theorem~\ref{thm:host-kleene} \\
Kleene's recursion theorem up to the record on the host machine: stopping, register 0, trap latch, ledger exactly, flag, number of facts, whether the channel is empty & proved; the host program simulates the program $T$ that carries out the transformation, marks and all, and pays for the whole record in one go once it has computed & Appendix~\ref{sec:small-diagonal}, Theorem~\ref{thm:host-kleene-record} \\
A fixed point of full behaviour (every register, what the facts name) on the host machine & refuted for registers, and for the fact table, the channel and the versions  & Appendix~\ref{sec:small-diagonal}, Theorem~\ref{thm:host-no-exact} \\
No inside decider with a representable flip, and no host decider, for functional properties of host programs & proved (the recursion premise is a theorem there) & Appendix~\ref{sec:small-diagonal}, Theorem~\ref{thm:host-wall} \\
Rice's theorem for two-counter programs on every input (stop, counters, trap latch, ledger, flag, and what each fact and the channel are about) & proved & Appendix~\ref{sec:small-diagonal}, Theorem~\ref{thm:tc-rice} \\
Kleene's recursion theorem for two-counter programs, inputs written as powers of two, for the function & proved; refuted once the recorded facts count & Appendix~\ref{sec:small-diagonal}, Theorem~\ref{thm:tc-kleene} \\
Kleene's recursion theorem for two-counter programs on plain inputs & refuted: a carried-out transformation has no fixed point & Appendix~\ref{sec:small-diagonal}, Theorem~\ref{thm:no-plain-recursion} \\
Every universal base lifts with an exact checker, a nontrivial claim and room for a fact; finite simulations lift after grouping & proved under the stated interface conditions & Appendix~\ref{sec:every-model}, Theorems~\ref{thm:lift} and~\ref{thm:lift-macro} \\
Finite branching, state-independent moves and an indiscriminate record order give the stated obstructions to lifting & proved & Appendix~\ref{sec:every-model}, Theorem~\ref{thm:lift-fails} \\
Every box with a universal base is a slice of its lift with the window claims and room for a fact: forgetting the panel runs as the box runs while the trap is down, reaches exactly what the box reaches, isn't moved by CHECK, COMMIT or CERTIFY, and gives neither the flag nor the ledger, the chain's ledger exactly 3 above the empty run's; with no room in the table no interface certifies the lift; a Thiele-complete machine's computing moves are its moves that cost nothing, through every interface, and its window is a slice onto two counters; a machine run as one rule has no universal base, no base by simulation and no universal-base slice, and is a Thiele machine only as a clock or with a reading that never rises & proved; that the slice runs as the box runs and that the record moves leave it alone, directly from the rules, and the converse's universal base and step-for-step clauses directly from the interface & Appendix~\ref{sec:every-model}, Definition~\ref{def:slice}, Theorem~\ref{thm:slice} \\
With the program in the state, a reading is kept by every operation exactly when it never comes down from any state; a cell some program clears is a simulated bit & proved, directly from the rules & Appendix~\ref{sec:every-model}, Observation~\ref{obs:program-in-state} \\
The seven models compute the same numerical relations under their encodings; one counter has decidable halting & proved & Appendix~\ref{sec:every-model}, Theorems~\ref{thm:lift-models} and~\ref{thm:one-counter} \\
The source compiler preserves halting and numerical output through the host, two-counter guest and universal program & proved; final two stages not executed & Appendix~\ref{sec:running}, Theorem~\ref{thm:running-compiler} \\
Executable implementations agree on the tested whole-state transitions & tested, not a guarantee about untested inputs or physical hardware & Appendix~\ref{sec:running} \\
The integer check makes the pinned matrix PSD & definition; proved & Appendix~\ref{sec:chsh}, Definition~\ref{def:u23-pinned}, Theorem~\ref{thm:u23-notrap-psd} \\
PSD forces $S^2 \le 8$ & proved (algebra only) & Appendix~\ref{sec:chsh}, Theorem~\ref{thm:u23-psd-tsirelson} \\
A passed check means $S^2 < 8$, and no smaller bound holds & proved: $S^2 < 8$, with passing tallies as close to 8 as desired & Appendix~\ref{sec:chsh}, Corollary~\ref{cor:u23-notrap-bound} \\
A PSD pinned matrix has a PSD completion & proved & Appendix~\ref{sec:chsh}, Remark~\ref{rem:u23-sound-not-complete} \\
Every deterministic plan fails the check & proved; no such plan's pinned matrix is PSD  & Appendix~\ref{sec:chsh}, Remark~\ref{rem:u23-sound-not-complete} \\
The level-1 test is the size-five case & definition; proved & Appendix~\ref{sec:chsh}, Definition~\ref{def:u23-genrealizable}, Corollary~\ref{cor:u23-general} \\
The 1+AB matrix ties its four-body cells, the constraint $\langle \{A_0, A_1\} \otimes \{B_0, B_1\} \rangle = 0$; the level-1 pinnings don't imply it, and a classical strategy breaks it & proved  & Appendix~\ref{sec:chsh}, Appendix~\ref{sec:chsh-derivation} \\
A tally above two, with every pair asked, isn't one plan replayed & proved & Appendix~\ref{sec:chsh-derivation}, Theorem~\ref{thm:chsh-tally} \\
A tally above two proves nonlocality; a pass proves the data quantum & not claimed & Appendix~\ref{sec:chsh} \\
A finite local model meets the bound & proved; Fine's theorem in correlation form is proved as an iff, and the converse of the single-bound statement is false  & Appendix~\ref{sec:chsh-derivation}, Theorem~\ref{thm:chsh-fine}, Section~\ref{sub:finite-local-construction} \\
The elliptope gate is sound; the PR box fails it & proved & Appendix~\ref{sec:chsh-derivation}, Theorem~\ref{thm:chsh-gate-sound} \\
The arc-sine boundary: a table has a PSD completion, and so is quantum, exactly when its arcsines pass the four inequalities {\footnotesize(if and only if)} & proved & Appendix~\ref{sec:chsh-derivation} \\
Every PSD completion comes from a quantum strategy, and the completion set is exactly the set of quantum correlator tables in finite dimension, real or complex & proved & Appendix~\ref{sec:chsh-derivation} \\
A two-qubit strategy reaches $2\sqrt{2}$; every strategy in finite dimension, with real or complex amplitudes, has $S^2 \le 8$ and a PSD level-1 moment matrix; every deterministic plan is such a strategy & proved & Appendix~\ref{sec:chsh-derivation} \\
The integer route to the bound is new mathematics & not claimed; the route is NPA's & Appendix~\ref{sec:chsh} \\
The integer check passes exactly when every pair was sampled and the pinned matrix is PSD & proved & Appendix~\ref{sec:chsh}, Theorem~\ref{thm:chsh-exact} \\
Honest play at Tsirelson's point fails the exact check about 85 percent of the time & simulated, with a first-order limit of $(2 + \sqrt 2)/4$; no theorem & Appendix~\ref{sec:chsh} \\
The check with a tolerance $\varepsilon$ is the exact check at $\varepsilon = 0$; a pass means operator norm at most $1 + \varepsilon$ and $S^2 \le 8(1 + \varepsilon)^2$ & proved, including the earned CHECK  & Appendix~\ref{sec:chsh} \\
A flag raised on the host machine from a clean start stands on a passing CHECK of the packed counts, a COMMIT and a CERTIFY; the chain run from a fresh start certifies for 3 exactly when the check passes & proved & Appendix~\ref{sec:chsh}, Theorem~\ref{thm:chsh-earned} \\
Five tallies: 12/5 and 14/5 certify; 16/5, the PR box and the all-ones plan are refused forever & proved & Appendix~\ref{sec:chsh}, Proposition~\ref{prop:chsh-demos} \\
The tally was played, or was honest, or puts a bound on nature & not claimed; it is loaded into a counter at no charge on the ledger & Appendix~\ref{sec:chsh} \\
Structural dissipation pricing & every collapsing instruction dissipates at least one unit, and dissipation is at most cost & Section~\ref{sec:landauer}, Premise~\ref{prem:u24-landauer-pair} \\
Covering tree, paid steps & a decision tree whose leaves cover the narrowing, and at least as many steps charged at least one as the tree is deep & entitlement, Theorems~\ref{thm:u08-count}, \ref{thm:ent-cs} and~\ref{thm:ent-machine} \\
Share kept by a test & each passed test keeps at least a $2^{-k}$ share of the candidates, for a $k$ fixed before the run; every run meets it for $k$ its own biggest test & entitlement made by the run, Theorems~\ref{thm:ent-run} and~\ref{thm:ent-run-machine} \\
Lift interface & a claim language with an exact equality test, a checker proved to say yes exactly when the claim holds, a claim true at one loaded state and false at another, room for one fact & the lift, Theorem~\ref{thm:lift} \\
\end{longtable}
\endgroup

\section{Further reading}

The works below were read for the book and are not named in its finished text or in this volume. They are kept here with their annotations, so that nothing read for the book is lost.

\subsection{Structure, information and knowledge}



\noindent Hans Albert. Wikipedia. \href{https://en.wikipedia.org/wiki/Hans_Albert}{Article}. German philosopher (1921--2023), with Popper a leading critical rationalist; Traktat über kritische Vernunft (1968).
\par\medskip

\noindent A.~J. Ayer (1956). The Problem of Knowledge. Macmillan, London, p.~34. The other definition Gettier targets.
\par\medskip

\noindent Max Black. Wikipedia. \href{https://en.wikipedia.org/wiki/Max_Black}{Article}. The Russian-born British-American philosopher, at Cornell from 1946, and his 1952 paper on two spheres.
\par\medskip

\noindent Roberto Bonola (1912). Non-Euclidean Geometry: A Critical and Historical Study of Its Development. Open Court; Dover reprint, 1955. The attempts on the parallel postulate from Proclus through the Islamic geometers and Saccheri, and Farkas Bolyai's letter to his son.
\par\medskip

\noindent G. K. Chesterton (1929). The Thing. The fence passage is quoted in the American Chesterton Society's \href{https://www.chesterton.org/lecture-57/}{Lecture 57: The Thing}.
\par\medskip

\noindent G\"osta Grahne (1998). Updates and counterfactuals. Journal of Logic and Computation 8(1), pp. 87--117. \href{https://doi.org/10.1093/logcom/8.1.87}{DOI: 10.1093/logcom/8.1.87}. Revision of a static world versus update of a changing one.
\par\medskip

\noindent Jonathan Jenkins Ichikawa and Matthias Steup (2001, substantive revision 2026). The Analysis of Knowledge. Stanford Encyclopedia of Philosophy. \href{https://plato.stanford.edu/entries/knowledge-analysis/}{Entry}. Dharmottara's mirage and Peter of Mantua's running Plato.
\par\medskip

\noindent Rosalie Iemhoff (2008, substantive revision 2019). Intuitionism in the Philosophy of Mathematics. Stanford Encyclopedia of Philosophy. \href{https://plato.stanford.edu/entries/intuitionism/}{Entry}. Kreisel's creating-subject axioms CS1--CS3 (1967), CS2 being that the creating subject never forgets what it has proven, and Troelstra's strengthening (1969).
\par\medskip

\noindent Petrus Mantuanus (d. 1400). Logica. Medieval manuscripts in Oxford libraries, Bodleian Library. \href{https://medieval.bodleian.ox.ac.uk/catalog/work_3607}{Catalogue entry}. Peter of Mantua as a logician of the late fourteenth century.
\par\medskip

\noindent Per Martin-Löf (1996). On the Meanings of the Logical Constants and the Justifications of the Logical Laws. Nordic Journal of Philosophical Logic 1(1), pp. 11--60. Lectures given in Siena in 1983. The passage on proof as an act is on p.~28. \href{https://archive-pml.github.io/martin-lof/pdfs/Meanings-of-the-Logical-Constants-1983.pdf}{Lecture text}.
\par\medskip

\noindent Meta Platforms (2021). Introducing Meta: A Social Technology Company. Company announcement, 28 October 2021. \href{https://about.fb.com/news/2021/10/facebook-company-is-now-meta/}{Announcement}.
\par\medskip

\noindent Benjamin Morison (2014, revised 2026). Sextus Empiricus. Stanford Encyclopedia of Philosophy. \href{https://plato.stanford.edu/entries/sextus-empiricus/}{Entry}. The source for his dates, his medicine and his skepticism.
\par\medskip

\noindent J. J. O'Connor and E. F. Robertson (n.d.). Theaetetus of Athens. MacTutor History of Mathematics. \href{https://mathshistory.st-andrews.ac.uk/Biographies/Theaetetus/}{Entry}.
\par\medskip

\noindent J. J. O'Connor and E. F. Robertson (n.d.). Paul Richard Halmos. MacTutor History of Mathematics. \href{https://mathshistory.st-andrews.ac.uk/Biographies/Halmos/}{Biography}.
\par\medskip

\noindent Benoît Rittaud and Albrecht Heeffer (2014). The Pigeonhole Principle, Two Centuries before Dirichlet. The Mathematical Intelligencer 36(2), pp. 27--29. Jean Leurechon's use of the principle in 1622.
\par\medskip

\noindent G\"oran Sundholm (1998). Proofs as Acts and Proofs as Objects: Some questions for Dag Prawitz. Theoria 64(2--3), pp. 187--216. \href{https://doi.org/10.1111/j.1755-2567.1998.tb00187.x}{DOI: 10.1111/j.1755-2567.1998.tb00187.x}.
\par\medskip

\noindent G\"oran Sundholm (2014). Constructive Recursive Functions, Church's Thesis, and Brouwer's Theory of the Creating Subject: Afterthoughts on a Parisian Joint Session. In Constructivity and Computability in Historical and Philosophical Perspective, Logic, Epistemology, and the Unity of Science, Springer, pp. 1--35. \href{https://doi.org/10.1007/978-94-017-9217-2_1}{DOI: 10.1007/978-94-017-9217-2\_1}.
\par\medskip

\noindent Timothy Yenter and Ezio Vailati (2003, revised 2024). Samuel Clarke. Stanford Encyclopedia of Philosophy. \href{https://plato.stanford.edu/entries/clarke/}{Entry}. Clarke as the British philosopher and theologian who defended Newton's natural philosophy in his letters with Leibniz.
\par\medskip

\subsection{Cost, records and thermodynamics}



\noindent J\'anos Acz\'el, Bruno Forte and C. T. Ng (1974). Why the Shannon and Hartley entropies are ``natural.'' Advances in Applied Probability 6(1), pp. 131--146. \href{https://doi.org/10.2307/1426210}{DOI: 10.2307/1426210}. Ng's 1971 Waterloo doctorate under Acz\'el: \href{https://mathgenealogy.org/id.php?id=102040}{Mathematics Genealogy Project}.
\par\medskip

\noindent Roberto M. Amadio, Nicolas Ayache, François Bobot, Jaap P. Boender, Brian Campbell, Ilias Garnier, Antoine Madet, James McKinna, Dominic P. Mulligan, Mauro Piccolo, Randy Pollack, Yann Régis-Gianas, Claudio Sacerdoti Coen, Ian Stark and Paolo Tranquilli (2014). Certified Complexity (CerCo). Foundational and Practical Aspects of Resource Analysis, FOPARA 2013, LNCS 8552, pp. 1--18. \href{https://doi.org/10.1007/978-3-319-12466-7_1}{DOI: 10.1007/978-3-319-12466-7\_1}.
\par\medskip

\noindent Walter J. Boyne. Wikipedia entry: United States Air Force officer, aviation historian, director of the National Air and Space Museum (acting from 1982, director 1983--1986). \href{https://en.wikipedia.org/wiki/Walter_J._Boyne}{Entry}.
\par\medskip

\noindent Léon Brillouin (1951). Maxwell's Demon Cannot Operate: Information and Entropy. I. Journal of Applied Physics 22(3), pp. 334--337. \href{https://doi.org/10.1063/1.1699951}{DOI: 10.1063/1.1699951}.
\par\medskip

\noindent Arthur Charguéraud and François Pottier (2019). Verifying the Correctness and Amortized Complexity of a Union-Find Implementation in Separation Logic with Time Credits. Journal of Automated Reasoning 62, pp. 331--365; first published online in 2017. \href{https://doi.org/10.1007/s10817-017-9431-7}{DOI: 10.1007/s10817-017-9431-7}.
\par\medskip

\noindent Gavin E. Crooks (1999). Entropy production fluctuation theorem and the nonequilibrium work relation for free energy differences. Physical Review E 60(3), pp. 2721--2726. \href{https://doi.org/10.1103/PhysRevE.60.2721}{DOI: 10.1103/PhysRevE.60.2721}. A path's chance against its reverse's is $e$ to the dissipated work.
\par\medskip

\noindent Joe DeBlasio (2025). Discussion: removal of one-6962-log requirement. Chromium ct-policy mailing list, 17 November 2025. \href{https://groups.google.com/a/chromium.org/g/ct-policy/c/amrnTiDFdwM/m/mCyPDc4TAwAJ}{Post}.
\par\medskip

\noindent Jérémie Decouchant, David Kozhaya, Vincent Rahli and Jiangshan Yu (2022). Damysus: streamlined BFT consensus leveraging trusted components. Proceedings of the Seventeenth European Conference on Computer Systems (EuroSys '22), pp. 1--16. \href{https://doi.org/10.1145/3492321.3519568}{DOI: 10.1145/3492321.3519568}.
\par\medskip

\noindent Encyclopaedia Britannica. Donation of Constantine. \href{https://www.britannica.com/topic/Donation-of-Constantine}{Web page}. Composed in the eighth century; the consensus view puts it in the 750s or 760s.
\par\medskip

\noindent Encyclopaedia Britannica. Leo Szilard. \href{https://www.britannica.com/biography/Leo-Szilard}{Web page}. Born in Budapest in 1898; Berlin doctorate 1922 and the Institute of Theoretical Physics there.
\par\medskip

\noindent Ethereum Foundation (n.d.). The Merge. ethereum.org roadmap page, on the transition of 15 September 2022. \href{https://ethereum.org/roadmap/merge/}{Page}.
\par\medskip

\noindent Ethereum Foundation (n.d.). Timeline of all Ethereum forks (Tangerine Whistle, EIP-150, and the Berlin upgrade, block 12,244,000). ethereum.org. \href{https://ethereum.org/ethereum-forks/}{Page}.
\par\medskip

\noindent Ethereum Foundation. Devcon editions (Devcon 0, Berlin 2014; Devcon 1, London 2015; Devcon 2, Shanghai 2016). devcon.org. \href{https://devcon.org/past-events}{Page}.
\par\medskip

\noindent Ethereum Foundation. Ethereum execution specifications, Prague fork, src/ethereum/forks/prague/vm/instructions/storage.py and vm/gas.py, at commit ecb68f21a7d28db101c769878184f6facb319da1. \href{https://github.com/ethereum/execution-specs}{Repository}.
\par\medskip

\noindent Federal Aviation Administration (2008). Revisions to Cockpit Voice Recorder and Digital Flight Data Recorder Regulations. Final rule, Federal Register 73(46), 7 March 2008, pp. 12542--12573. The two-hour cockpit voice recorder.
\par\medskip

\noindent Federal Aviation Administration (2026). 25-Hour Cockpit Voice Recorder (CVR) Requirement, New Aircraft Production. Final rule, Federal Register 91(21), 2 February 2026, pp. 4447--4462. \href{https://www.govinfo.gov/content/pkg/FR-2026-02-02/html/2026-02110.htm}{Federal Register text}. The thirty-minute recorders of 1966, the two-hour rule, and the 25-hour standard of ICAO (2016) and EASA (from 2021).
\par\medskip

\noindent Richard FitzNigel (c. 1180). Dialogus de Scaccario (The Dialogue of the Exchequer). Text and translation ed. Charles Johnson, Nelson's Medieval Classics, 1950. The widths of the cuts on Exchequer tallies.
\par\medskip

\noindent Fox-IT (2012). Black Tulip: Report of the investigation into the DigiNotar Certificate Authority breach. Version 1.0. \href{https://roselabs.nl/files/audit_reports/Fox-IT_-_DigiNotar.pdf}{Report}.
\par\medskip

\noindent Tom Furtwangler (2017). The vaccine vial monitor: the world's smartest sticker. PATH, 6 November 2017. \href{https://www.path.org/articles/vaccine-vial-monitor-worlds-smartest-sticker/}{Article}. The stickers first sold in 1996, for oral polio vaccine.
\par\medskip

\noindent General Conference on Weights and Measures (2018). Resolution 1 of the 26th meeting: On the revision of the International System of Units (SI), effective 20 May 2019.
\par\medskip

\noindent C. H. Jenkinson (1910). Early wooden tallies relating to Surrey. Surrey Archaeological Collections 23, pp. 203--207. \href{https://archaeologydataservice.ac.uk/catalogue/adsdata/arch-379-1/dissemination/pdf/vol_23/surreyac023_203-207_jenkinson.pdf}{PDF}. Hazel; the payer kept the stock, the Exchequer the foil; abolished by statute in 1826; the burning that started the fire.
\par\medskip

\noindent Joint Committee for Guides in Metrology (2012). International vocabulary of metrology: basic and general concepts and associated terms (VIM), 3rd edition, JCGM 200:2012, entry 2.41, ``metrological traceability.'' The definition is on p.~29; ``does not ensure \ldots{} that there is an absence of mistakes'' is Note~5 to 2.41, p.~30. \href{https://www.bipm.org/documents/20126/2071204/JCGM_200_2012.pdf}{BIPM text}.
\par\medskip

\noindent Mikhail Kuter, Marina Gurskaya and Albert Gurskii (2025). Genoa 1340, Double Entry, and `Gold Coasting'. De Computis, Revista Española de Historia de la Contabilidad 22(1), pp. 1--33. \href{https://doi.org/10.26784/issn.1886-1881.22.1.11927}{DOI: 10.26784/issn.1886-1881.22.1.11927}.
\par\medskip

\noindent Klaus-J\"orn Lange, Pierre McKenzie and Alain Tapp (2000). Reversible Space Equals Deterministic Space. Journal of Computer and System Sciences 60(2), pp. 354--367. \href{https://doi.org/10.1006/jcss.1999.1672}{DOI: 10.1006/jcss.1999.1672}.
\par\medskip

\noindent Ben Laurie, Adam Langley and Emilia Kasper (2013). Certificate Transparency. RFC 6962. \href{https://www.rfc-editor.org/rfc/rfc6962.html}{RFC 6962}; \href{https://doi.org/10.17487/RFC6962}{DOI: 10.17487/RFC6962}.
\par\medskip

\noindent Ben Laurie, Eran Messeri and Rob Stradling (2021). Certificate Transparency Version 2.0. RFC 9162. \href{https://www.rfc-editor.org/rfc/rfc9162.html}{RFC 9162}; \href{https://doi.org/10.17487/RFC9162}{DOI: 10.17487/RFC9162}.
\par\medskip

\noindent A. Liapounoff (1907). Problème général de la stabilité du mouvement, traduit du russe par Édouard Davaux. Annales de la Faculté des Sciences de Toulouse, 2e série, 9, pp. 203--474. \href{http://www.numdam.org/item?id=AFST_1907_2_9__203_0}{Numdam}. The French translation the 1992 English one was made from.
\par\medskip

\noindent Sinisa Matetic, Mansoor Ahmed, Kari Kostiainen, Aritra Dhar, David Sommer, Arthur Gervais, Ari Juels and Srdjan Capkun (2017). ROTE: rollback protection for trusted execution. Proceedings of the 26th USENIX Security Symposium, pp. 1289--1306. \href{https://www.usenix.org/conference/usenixsecurity17/technical-sessions/presentation/matetic}{USENIX page}.
\par\medskip

\noindent Clyde Lee Miller (2009; revised 2021). Cusanus, Nicolaus [Nicolas of Cusa]. Stanford Encyclopedia of Philosophy. \href{https://plato.stanford.edu/entries/cusanus/}{SEP entry}.
\par\medskip

\noindent Y\^uto Murashita, Ken Funo and Masahito Ueda (2014). Nonequilibrium equalities in absolutely irreversible processes. Physical Review E 90, 042110. \href{https://doi.org/10.1103/PhysRevE.90.042110}{DOI: 10.1103/PhysRevE.90.042110}. The bound on average entropy production by $-\ln(1-\lambda)$, with free expansion as the first example.
\par\medskip

\noindent Satoshi Nakamoto (2008). Bitcoin: A Peer-to-Peer Electronic Cash System. \href{https://bitcoin.org/bitcoin.pdf}{Paper}.
\par\medskip

\noindent Devon O'Brien (2018). Certificate Transparency Enforcement in Google Chrome. Chromium ct-policy mailing list, 6 February 2018. \href{https://groups.google.com/a/chromium.org/g/ct-policy/c/wHILiYf31DE/m/iMFmpMEkAQAJ}{Announcement}.
\par\medskip

\noindent Garson O'Toole (2016). There Ain't No Such Thing as a Free Lunch, TANSTAAFL. Quote Investigator, 27 August 2016. \href{https://quoteinvestigator.com/2016/08/27/free-lunch/}{Article}.
\par\medskip

\noindent Pavel Pakshin, Elena Panteley and Natalia Pakshina (2017). The Legacy of A. M. Lyapunov to Systems and Control Theory: 110 years of the Toulouse edition. IFAC World Congress 2017, open-track proposal. Davaux's translation, reviewed and corrected by Lyapunov.
\par\medskip

\noindent Lídia del Rio, Johan Åberg, Renato Renner, Oscar Dahlsten and Vlatko Vedral (2011). The thermodynamic meaning of negative entropy. Nature 474, pp. 61--63. \href{https://doi.org/10.1038/nature10123}{DOI: 10.1038/nature10123}; \href{https://arxiv.org/abs/1009.1630}{arXiv: 1009.1630}.
\par\medskip

\noindent J\"urgen Schnakenberg (1976). Network theory of microscopic and macroscopic behavior of master equation systems. Reviews of Modern Physics 48, pp. 571--585. \href{https://doi.org/10.1103/RevModPhys.48.571}{DOI: 10.1103/RevModPhys.48.571}. The entropy a Markov chain hands its surroundings, as the net rate across each pair times the log of the forward-to-backward rate ratio, which at a steady state is the whole entropy production; the biased walk is its smallest case.
\par\medskip

\noindent Securities and Exchange Commission (2022). Electronic Recordkeeping Requirements for Broker-Dealers, Security-Based Swap Dealers, and Major Security-Based Swap Participants. Release No. 34-96034, 12 October 2022. \href{https://www.sec.gov/rules/final/2022/34-96034.pdf}{Final rule}. Adds an audit-trail alternative to the write-once requirement of Rule 17a-4(f) in place since 1997.
\par\medskip

\noindent Caroline Shenton (2012). The Day Parliament Burned Down. Oxford University Press. The tallies burned in the furnaces under the House of Lords, 16 October 1834.
\par\medskip

\noindent Michael Stock, Pauline Barat, Richard S. Davis, Alain Picard and Martin J. T. Milton (2015). Calibration campaign against the international prototype of the kilogram in anticipation of the redefinition of the kilogram part I: comparison of the international prototype with its official copies. Metrologia 52(2), pp. 310--316. \href{https://doi.org/10.1088/0026-1394/52/2/310}{DOI: 10.1088/0026-1394/52/2/310}.
\par\medskip

\noindent Wei Tang (2019). EIP-2200: Structured Definitions for Net Gas Metering. Ethereum Improvement Proposals. \href{https://eips.ethereum.org/EIPS/eip-2200}{eips.ethereum.org/EIPS/eip-2200}.
\par\medskip

\noindent This Day in Aviation (n.d.). 30 October 1935: Boeing Model 299 crash at Wright Field. \href{https://www.thisdayinaviation.com/30-october-1935/}{Page}.
\par\medskip

\noindent Kent Weare (2016). Ethereum Security Alert Issued, Ethereum Foundation Responds with ``From Shanghai, With Love''. InfoQ, 19 September 2016. \href{https://www.infoq.com/news/2016/09/Ethereum-DOS-Attack}{Article}.
\par\medskip

\noindent World Health Organization (2011). PQS performance specification: Vaccine Vial Monitor. WHO/PQS/E06/IN05.2.
\par\medskip

\subsection{Computation, self-reference and checking}



\noindent Keith Adams and Ole Agesen (2006). A Comparison of Software and Hardware Techniques for x86 Virtualization. Proceedings of the 12th International Conference on Architectural Support for Programming Languages and Operating Systems (ASPLOS XII), pp. 2--13. \href{https://doi.org/10.1145/1168857.1168860}{DOI: 10.1145/1168857.1168860}. Binary translation of instructions such as POPF, which on x86 runs in user mode without trapping. Intel's VT-x and AMD-V in 2005 and 2006: Wikipedia, \href{https://en.wikipedia.org/wiki/X86_virtualization}{x86 virtualization}.
\par\medskip

\noindent American Standards Association (1963). American Standard Code for Information Interchange, ASA X3.4-1963.
\par\medskip

\noindent Henk Barendregt. Wikipedia. \href{https://en.wikipedia.org/wiki/Henk_Barendregt}{Article}. A Dutch logician known for his work in lambda calculus and type theory.
\par\medskip

\noindent Gerd Behrmann, Ansgar Fehnker, Thomas Hune, Kim Larsen, Paul Pettersson, Judi Romijn and Frits Vaandrager (2001). Minimum-cost reachability for priced timed automata. Hybrid Systems: Computation and Control (HSCC 2001), Lecture Notes in Computer Science 2034, pp. 147--161.
\par\medskip

\noindent N. G. de Bruijn (1972). Lambda calculus notation with nameless dummies, a tool for automatic formula manipulation, with application to the Church--Rosser theorem. Indagationes Mathematicae (Proceedings) 75(5), pp. 381--392. \href{https://doi.org/10.1016/1385-7258(72)90034-0}{DOI: 10.1016/1385-7258(72)90034-0}. Variables written as numbers that count the binders back, now called de Bruijn indices.
\par\medskip

\noindent Stanley Burris and Marcel Jackson (2010, substantive revision 2026). George Boole. Stanford Encyclopedia of Philosophy. \href{https://plato.stanford.edu/entries/boole/}{Entry}.
\par\medskip

\noindent Stanley Burris and Javier Legris (2009, substantive revision 2021). The Algebra of Logic Tradition. Stanford Encyclopedia of Philosophy. \href{https://plato.stanford.edu/entries/algebra-logic-tradition/}{Entry}.
\par\medskip

\noindent G. K. Chesterton. Wikipedia. \href{https://en.wikipedia.org/wiki/G._K._Chesterton}{Article}. An English writer, referred to as the ``prince of paradox''.
\par\medskip

\noindent Donald W. Davies (2004). Corrections to Turing's Universal Computing Machine. In B. Jack Copeland (ed.), The Essential Turing, pp. 103--124. Oxford University Press. Davies at the National Physical Laboratory from 1947 and the word \emph{packet}: Wikipedia, \href{https://en.wikipedia.org/wiki/Donald_Davies}{Donald Davies}.
\par\medskip

\noindent Charles Dickens (1843). A Christmas Carol. Chapman and Hall, London. Stave Three, the game of Yes and No.
\par\medskip

\noindent William Henry Eccles and Frank Wilfred Jordan (1920). Improvements in ionic relays. British patent GB148582A, filed 21 June 1918, published 5 August 1920. \href{https://patents.google.com/patent/GB148582A/en}{Patent text}. The 1918 date in the text is the filing date. The circuit as the flip-flop, the basis for electronic memory: History of Information, \href{https://www.historyofinformation.com/detail.php?id=3606}{entry}; and Wikipedia, \href{https://en.wikipedia.org/wiki/Flip-flop_(electronics)}{Flip-flop (electronics)}, on the Eccles--Jordan trigger circuit as the first electronic latch, later the basic storage element of digital circuitry.
\par\medskip

\noindent László Fejes Tóth (1953). Lagerungen in der Ebene, auf der Kugel und im Raum. Berlin: Springer.
\par\medskip

\noindent László Fejes Tóth (1964). Regular Figures. Oxford: Pergamon Press.
\par\medskip

\noindent Jean-Christophe Filliâtre, Léon Gondelman and Andrei Paskevich (2014). The spirit of ghost code. Computer Aided Verification (CAV 2014), Lecture Notes in Computer Science 8559, pp. 1--16. \href{https://doi.org/10.1007/978-3-319-08867-9_1}{DOI: 10.1007/978-3-319-08867-9\_1}.
\par\medskip

\noindent Nicholas Griffin, ed. (2003). The Cambridge Companion to Bertrand Russell. Cambridge University Press. G. G. Berry as a junior librarian at the Bodleian, p. 63.
\par\medskip

\noindent Jean-Pierre Marquis (1996, substantive revision 2019). Category Theory. Stanford Encyclopedia of Philosophy. \href{https://plato.stanford.edu/entries/category-theory/}{Entry}. The borrowed words ``category'' and ``functor''.
\par\medskip

\noindent Ross M. McConnell, Kurt Mehlhorn, Stefan Näher and Pascal Schweitzer (2011). Certifying algorithms. Computer Science Review 5(2), pp. 119--161. \href{https://doi.org/10.1016/j.cosrev.2010.09.009}{DOI: 10.1016/j.cosrev.2010.09.009}.
\par\medskip

\noindent Jason Megill. The Lucas-Penrose Argument about Gödel's Theorem. Internet Encyclopedia of Philosophy. \href{https://iep.utm.edu/lp-argue/}{Article}.
\par\medskip

\noindent Glen Mével, Jacques-Henri Jourdan and François Pottier (2019). Time credits and time receipts in Iris. ESOP 2019, LNCS 11423, pp. 3--29. \href{https://doi.org/10.1007/978-3-030-17184-1_1}{DOI: 10.1007/978-3-030-17184-1\_1}.
\par\medskip

\noindent Jeff Miller. Earliest Uses of Symbols of Set Theory and Logic. MacTutor History of Mathematics. \href{https://mathshistory.st-andrews.ac.uk/Miller/mathsym/set/}{Page}. Peano's epsilon for \emph{est}, 1889.
\par\medskip

\noindent J. J. O'Connor and E. F. Robertson (n.d.). The four colour theorem. MacTutor History of Mathematics. \href{https://mathshistory.st-andrews.ac.uk/HistTopics/The_four_colour_theorem/}{Entry}. (De Morgan's letter to Hamilton of 23 October 1852 is quoted from here.)
\par\medskip

\noindent Panu Raatikainen (2013, revised 2025). Gödel's Incompleteness Theorems. Stanford Encyclopedia of Philosophy. \href{https://plato.stanford.edu/entries/goedel-incompleteness/}{Entry}. Source for the cautions about what the Gödel sentence says.
\par\medskip

\noindent Erich Reck (2008, substantive revision 2025). Dedekind's Contributions to the Foundations of Mathematics. Stanford Encyclopedia of Philosophy. \href{https://plato.stanford.edu/entries/dedekind-foundations/}{Entry}. Peano acknowledged Dedekind's priority for the axioms.
\par\medskip

\noindent Neil Robertson, Daniel Sanders, Paul Seymour and Robin Thomas (1997). The four-colour theorem. Journal of Combinatorial Theory, Series B 70(1), pp. 2--44.
\par\medskip

\noindent Jürgen Schmidhuber (2003). Gödel machines: self-referential universal problem solvers making provably optimal self-improvements. Technical Report IDSIA-19-03. \href{https://arxiv.org/abs/cs/0309048}{arXiv: cs/0309048}.
\par\medskip

\noindent Marcel-Paul Schützenberger (1961). On the definition of a family of automata. Information and Control 4(2--3), pp. 245--270.
\par\medskip

\noindent Dana Scott (1962). A system of functional abstraction. Lectures delivered at the University of California, Berkeley, 1962--63; unpublished. The origin of the Scott encoding, as given in Mogensen--Scott encoding, Wikipedia, \href{https://en.wikipedia.org/wiki/Mogensen\%E2\%80\%93Scott_encoding}{Article}.
\par\medskip

\noindent Dana Scott. Carnegie Mellon University's School of Computer Science, \href{https://www.cs.cmu.edu/~scsfacts/turing.html}{faculty Turing Award recipients}: the 1976 award jointly with Michael Rabin for their work on nondeterministic machines. His doctorate at Princeton under Alonzo Church, domain theory and models of the lambda calculus: Wikipedia, \href{https://en.wikipedia.org/wiki/Dana_Scott}{Dana Scott}.
\par\medskip

\noindent Emin G\"un Sirer, Willem de Bruijn, Patrick Reynolds, Alan Shieh, Kevin Walsh, Dan Williams and Fred B. Schneider (2011). Logical attestation: an authorization architecture for trustworthy computing. SOSP 2011, pp. 249--264. \href{https://doi.org/10.1145/2043556.2043580}{DOI: 10.1145/2043556.2043580}.
\par\medskip

\noindent Ken Thompson. Heidelberg Laureate Forum, \href{https://www.heidelberg-laureate-forum.org/laureate/ken-thompson/}{laureate profile}: the 1983 Turing Award with Dennis Ritchie for operating systems theory and the implementation of Unix. Thompson's design and implementation of the original Unix at Bell Labs: Wikipedia, \href{https://en.wikipedia.org/wiki/Ken_Thompson}{Ken Thompson}.
\par\medskip

\noindent Alfred North Whitehead and Bertrand Russell (1910). Principia Mathematica, vol. I. Cambridge University Press. Introduction, chapter II.
\par\medskip

\noindent Robin Wilson (2002). Four Colours Suffice? Gresham College lecture, 21 October 2002. \href{https://www.gresham.ac.uk/node/11353}{Lecture}. De Morgan's letter of 23 October 1852 to the Irish mathematician William Rowan Hamilton.
\par\medskip

\subsection{Correlations and environmental records}



\noindent M.-C. Chen et al. (2019). Emergence of classical objectivity of quantum Darwinism in a photonic quantum simulator. Science Bulletin 64, p. 580.
\par\medskip

\noindent M. A. Ciampini, G. Pinna, P. Mataloni and M. Paternostro (2018). Experimental signature of quantum Darwinism in photonic cluster states. Physical Review A 98, 020101.
\par\medskip

\noindent Mauro Dorato and Jan Faye (2026). Niels Bohr as a Physicist of Principle. \href{https://arxiv.org/abs/2609.04297}{arXiv: 2609.04297}. Quotes Bohr's Unity of Knowledge (1954) on irreversible amplification and permanent marks on a photographic plate, citing p. 73.
\par\medskip

\noindent Lluís Masanes (2003). Necessary and sufficient condition for quantum-generated correlations. \href{https://arxiv.org/abs/quant-ph/0309137v1}{arXiv: quant-ph/0309137v1}. The paper gives his affiliation as the Universitat de Barcelona.
\par\medskip

\noindent Miguel Navascués, Stefano Pironio and Antonio Acín (2007). Bounding the Set of Quantum Correlations. Physical Review Letters 98, 010401. \href{https://doi.org/10.1103/PhysRevLett.98.010401}{DOI: 10.1103/PhysRevLett.98.010401}; \href{https://arxiv.org/abs/quant-ph/0607119v2}{arXiv: quant-ph/0607119v2}.
\par\medskip

\noindent Miguel Navascués, Stefano Pironio and Antonio Acín (2008). A convergent hierarchy of semidefinite programs characterizing the set of quantum correlations. New Journal of Physics 10, 073013. \href{https://arxiv.org/abs/0803.4290v1}{arXiv: 0803.4290v1}.
\par\medskip

\noindent Marcin Pawłowski, Tomasz Paterek, Dagomir Kaszlikowski, Valerio Scarani, Andreas Winter and Marek Żukowski (2009). Information causality as a physical principle. Nature 461, pp. 1101--1104. Pawłowski's affiliation, the University of Gdańsk, is on the paper: \href{https://arxiv.org/abs/0905.2292}{arXiv: 0905.2292}.
\par\medskip

\noindent Erwin Schrödinger. Wikipedia. \href{https://en.wikipedia.org/wiki/Erwin_Schr\%C3\%B6dinger}{Article}. Austrian physicist; the wave equation of 1926.
\par\medskip

\noindent B. S. Tsirel'son (1987). Quantum analogues of the Bell inequalities. The case of two spatially separated domains. Journal of Soviet Mathematics 36(4), pp. 557--570. \href{https://doi.org/10.1007/BF01663472}{DOI: 10.1007/BF01663472}.
\par\medskip

\noindent T. Unden, D. Louzon, M. Zwolak, W. H. Zurek and F. Jelezko (2019). Revealing the emergence of classicality using nitrogen-vacancy centers. Physical Review Letters 123, 140402.
\par\medskip

\cleardoublepage
\thispagestyle{empty}
\vspace*{154pt}
\begingroup\setlength{\parindent}{0pt}\setlength{\parskip}{7pt}
\noindent\BookLabel{Colophon}

\vspace{14pt}
This book is set in Charter, drawn by Matthew Carter in 1987 for the laser printers of its day,
10.5 on 14 point, on a page 189 by 246 millimetres. It was built from its source with pdfTeX.
The cells are drawn in TikZ.

Its own count, taken by the build: \ClaimTotal\ numbered statements, of which \PaidTotal\ are
closed by a proof on the page, and the rest stated and not proved: \OpenClaimList. One cell for
each, in the order they come, filled where the proof closed:

\vspace{7pt}
\noindent\LedgerGrid
\endgroup
\clearpage
\thispagestyle{empty}
\vspace*{0pt}
\noindent\BookLabel{\phantom{Devon Thiele}}

\vspace{14pt}
\PageCellFilled
\par\medskip

\end{document}
